% concatenated from main_arxiv.tex, appendix.tex
\documentclass[11pt]{article}
\usepackage[utf8]{inputenc}


\usepackage{natbib}
\usepackage{wustyle}

\usepackage[colorlinks,citecolor=blue,urlcolor=blue,linkcolor=blue,linktocpage=true,pagebackref=true]{hyperref}
\renewcommand*{\backref}[1]{}
\renewcommand*{\backrefalt}[4]{(\small%
    \ifcase #1 not cited%
          \or cited on page~#2%
          \else cited on pages #2%
    \fi%
    )}
  
\usepackage{multirow}
\usepackage{makecell}
\usepackage{fontawesome5} 
\usepackage[most]{tcolorbox}

% Highlighted theorem box with a very light-gray fill, thin dark frame, and slightly rounded corners.
\newtcolorbox{highlightedtheorem}{%
  enhanced,
  breakable,
  colback=black!2,
  colframe=black!72,
  boxrule=0.65pt,
  arc=2pt,
  outer arc=2pt,
  left=3pt,
  right=3pt,
  top=3pt,
  bottom=3pt,
  before skip=10pt,
  after skip=10pt
}

\usepackage[none]{hyphenat}
\emergencystretch=3em

\newcommand{\wzl}[1]{\todo[color=red!40, size=footnotesize]{wzl: #1}}
\newcommand{\lyh}[1]{\todo[color=blue!40, size=footnotesize]{lyh: #1}}
\newcommand{\zsb}[1]{\todo[color=green!40, size=footnotesize]{zsb: #1}}

\definecolor{sgdblue}{HTML}{0072B2}
\definecolor{polyakorange}{HTML}{D55E00}
\definecolor{nestgreen}{HTML}{009E73}

\newcommand{\sgd}{\mathrm{sgd}}
\newcommand{\polyak}{\mathrm{polyak}}
\newcommand{\nesterov}{\mathrm{nesterov}}
\newcommand{\crit}{\mathrm{crit}}
\newcommand{\la}{\mathrm{la}}
\newcommand{\peak}{\mathrm{peak}}
% \newcommand{\large}{\mathrm{large}}
% \newcommand{\eff}{\mathrm{eff}}

\newcommand{\blfootnote}[1]{%
  \begingroup
  \renewcommand\thefootnote{}%
  \footnote{#1}%
  \addtocounter{footnote}{-1}%
  \endgroup
}


% \title{From Polyak to Nesterov: Landscape-Aware Momentum Methods May Improve Stability and Data Scaling}

\title{Momentum in large-batch training: Polyak enlarges the  critical batch size, Nesterov improves data efficiency}


% \title{\bf Polyak versus Nesterov in Large-Batch Training:  Stability and Scaling Laws}

\author{
    Jia-Nan Wang$^{1,*}$\quad
    Zixun Huang$^{3,*}$\quad
    Kairui Li$^{1,*}$\quad
    Lei Wu$^{1,2,\dagger}$
      \\[.5em]
  $^1$Peking University\quad
  % $^2$Center for Machine Learning Research, Peking University
  % \\
  $^2$AI for Science Institute, Beijing\\[-.1em]
  $^3$University of Pennsylvania\\[.5em]
  \texttt{jiananwang25@stu.pku.edu.cn},\quad  \texttt{zixunh@wharton.upenn.edu}\\[-.1em] 
  \texttt{kaisms@stu.pku.edu.cn},\quad  \texttt{leiwu@math.pku.edu.cn}\\ 
}
\date{\vspace*{-2em}}

\begin{document}


\maketitle

\blfootnote{$^*$Equal contribution. }
\blfootnote{$^\dagger$Corresponding author.}
% \blfootnote{$^1$Code}
% \blfootnote{$1$Code is available at \url{https://github.com/alexhuang13/Momentum_scaling}.
% \quad Code is available at \faGithub\ \url{https://github.com/alexhuang13/Momentum_scaling}}
% \blfootnote{$^*$ Equal contribution.} 
% \blfootnote{$^\dagger$ Corresponding author: \texttt{leiwu@math.pku.edu.cn}}
% \blfootnote{\faGithub\ \url{https://github.com/alexhuang13/Momentum_scaling}}
\begin{abstract}
We study when and how momentum improves large-batch training in the one-pass regime, using power-law kernel regression as a tractable setting. We first characterize risk stability through the \emph{critical learning rate}, defined as the largest learning rate for stable training, and obtain \(\eta_{\sgd}^{\crit}\eqsim 1\), \(\eta_{\polyak}^{\crit}\eqsim\min\{1,B(1-\rho)\}\), and \(\eta_{\nesterov}^{\crit}\eqsim\min\{1,B^\beta(1-\rho)\}\), where \(B\) is the batch size, \(\rho\) is the momentum factor, and \(\beta>1\) is the capacity exponent. Within this admissible region, we derive scaling laws for the full risk dynamics, capturing the progression from an early transient, through power-law decay, to a noise floor. We then minimize the final-step risk over the admissible learning rates and momentum factors under a fixed data budget, yielding a three-regime batch-size phase diagram that reveals how the role of momentum changes with batch size. Notably, Polyak enlarges the critical batch size, the largest batch size preserving the best small-batch data-scaling exponent, thereby enabling greater parallelism without sacrificing data efficiency. In contrast, Nesterov achieves better data efficiency in the large batch regime because its look-ahead mechanism suppresses noise accumulation. Numerical experiments validate the predicted stability boundaries, risk dynamics, and batch-size phase diagram.
\end{abstract}


\vspace{.5em}
\begin{figure}[H]
\centering
\includegraphics[width=0.9\textwidth]{figures/intro_figure.pdf}
% \vspace{-.5em}
\caption{\textbf{Batch-size dependence of momentum benefits:} \textcolor{polyakorange}{\textbf{Polyak}} enlarges the critical batch size, while \textcolor{nestgreen}{\textbf{Nesterov}} improves data efficiency. \textbf{(a)} Theoretical batch-size phase diagram from Theorem~\ref{thm:batch-size-scaling}, showing the optimal data-scaling exponent as a function of the batch exponent \(b\), defined by the scaling \(B=D^{b+o(1)}\). 
\textbf{(b)} Best-tuned final excess risk versus batch size $B$ at fixed data budget \(D=2^{14}\) and \((s,\beta)=(0.3,4)\). The simulation recovers the same qualitative three-regime structure predicted by the theory. Experimental setup and additional results are provided in Section~\ref{sec:experiment}.}
\label{fig:intro-result}
\end{figure}


% \caption{\textbf{Batch-size dependence of momentum benefits: Polyak enlarges the critical batch size, while Nesterov improves data efficiency.} 

\doparttoc
\faketableofcontents
\part{}

\vspace*{-2.5em}
\section{Introduction}

Momentum has become a central mechanism in  optimizer design for modern large-scale training, underlying popular optimizers such as Adam~\citep{kingma2015adam} and Muon~\citep{jordan2024muon,liu2025muon}. Classical theory views momentum primarily as an acceleration mechanism,   quantified through  reduced iteration complexity
\citep{nesterov1983method,su2016differential,wibisono2016variational,
wilson2021lyapunov,gitman2019understanding,velikanov2023view,
wang2024marginal,liu2019acceleratingsgdmomentumoverparameterized,
gupta2024nesterov}. Yet iteration complexity alone does not fully capture the role of momentum in modern large-scale training, where the optimization regime itself changes as training resources, such as data and hardware memory, scale up.

Modern large-scale training, particularly language-model pretraining, often operates in a \emph{one-pass training} regime~\citep{brown2020language}. In this regime, the batch size \(B\) controls not only hardware parallelism but also the optimization horizon: under a fixed data budget \(D\), the number of sequential updates is
\[
T=\frac{D}{B}.
\]
Larger batches improve hardware utilization and training throughput, but necessarily reduce the number of optimization steps available per data pass. This creates a fundamental tradeoff between hardware efficiency and data efficiency. Indeed, recent language-model studies show that the relative advantages of different optimizers can change substantially with batch size, with some optimizers exhibiting markedly stronger data efficiency in the large-batch regime~\citep{marek2026small,shah2025practical,semenov2025benchmarking}. These observations motivate a systematic characterization of the role of momentum across batch sizes. We therefore ask:
\begin{highlightedtheorem}
\begin{center}
\textbf{Question:} {\it Can momentum enable larger batches without sacrificing data efficiency, or\\
even improve data efficiency itself?}
\end{center}
\end{highlightedtheorem}



We address this question for two canonical momentum methods, Polyak~\citep{polyak1964some} and Nesterov~\citep{nesterov1983method}, using power-law kernel (PLK) regression as a tractable learning problem. Its infinite-dimensional power-law spectrum spans a broad range of learning timescales, making the effects of momentum across spectral directions analytically accessible. The problem is characterized by two exponents: the capacity exponent $\beta>1$, which controls the decay of the feature spectrum, and the source exponent $s>0$, which controls how the target-function energy is distributed across the spectrum. Related power-law settings have been widely used to understand scaling laws in large-scale training across a variety of settings~\citep{lin2024scaling,bahri2024explaining,bordelon2025feature,li2026functional,wang2026fast,li2026optimal}. Building on the functional scaling law (FSL) framework of \citet{li2026functional}, we extend the analysis to momentum methods.


Our analysis of one-pass training follows the hierarchy 
\[ 
\mathcal A_B \;\longrightarrow\; \mathcal E(k,\eta,\rho,B) \;\longrightarrow\; \min_{(\eta,\rho)\in\mathcal A_B} \mathcal E\!\left(T,\eta,\rho,B\right),\quad T=\frac{D}{B}. 
\] 
Here, $\mathcal A_B$ denotes the admissible set of learning rates $\eta$ and momentum factors $\rho$ at batch size $B$, and $\mathcal E(k,\eta,\rho,B)$ denotes the expected excess risk. The three stages correspond, respectively, to \textbf{1)} identifying the admissible hyperparameter region through stability analysis, \textbf{2)} deriving scaling laws that characterize the risk dynamics, and \textbf{3)} optimizing the final-step risk subject to both the admissibility constraint $(\eta,\rho)\in\mathcal A_B$ and the one-pass constraint $T=D/B$.

\vspace{-.7em}
\paragraph{Our contributions.}
Our main contributions are summarized as follows.
\begin{itemize}
  \vspace{-.3em}
\item \textbf{Risk stability.}
We characterize risk stability through the \emph{critical learning rate} \(\eta^{\crit}\), the largest learning rate for which the stochastic dynamics remain stable (Section~\ref{sec:stability}). In the large-batch, high-momentum regime, we obtain the sharp scalings
\begin{equation}\label{eqn: stability-0}
\eta_{\sgd}^{\crit}\eqsim 1,\qquad
\eta_{\polyak}^{\crit}\eqsim \min\{1,B(1-\rho)\},\qquad
\eta_{\nesterov}^{\crit}\eqsim \min\{1,B^\beta(1-\rho)\}.
\end{equation}
Thus, stronger momentum requires a correspondingly larger batch to maintain stability. This compensation scales as \(B\) for Polyak but as \(B^\beta\) for Nesterov. Since \(\beta>1\), Nesterov permits substantially higher momentum in large-batch training.

\item \textbf{Scaling laws for risk dynamics.}
Within the risk-stable region, we derive sharp scaling laws for the risk dynamics of Polyak and Nesterov (Section~\ref{sec:scaling law}). In the high-momentum regime, the risk transitions from an early exponential transient to power-law decay and eventually saturates at a noise floor. Writing \(\eta_\rho:=\eta/(1-\rho)\), after the transient the dominant terms take the simple form
\[
\mathcal E_{\sgd}(k)\eqsim(k\eta)^{-s}+\frac{\sigma^2\eta}{B},\qquad \mathcal E_{\polyak}(k)\eqsim(k\eta_\rho)^{-s}+\frac{\sigma^2\eta_\rho}{B},
\]
while Nesterov retains the same accelerated signal term but reduces the noise floor to
\[
\mathcal E_{\nesterov}(k)\eqsim(k\eta_\rho)^{-s}+\frac{\sigma^2}{B}\min\{\eta_\rho,\eta_\rho^{1/\beta}\}.
\]
Thus, Polyak preserves the SGD signal--noise tradeoff while operating at the enlarged effective learning rate \(\eta_\rho\), whereas Nesterov further suppresses noise accumulation at large \(\eta_\rho\) through the curvature-adaptive damping induced by its look-ahead mechanism.

\item \textbf{Data scaling and batch-size phase diagram.}
Let \(B=D^{b+o(1)}\). Optimizing the remaining hyperparameters yields three characteristic batch exponents \(b_1<b_2<b_3\) and a three-regime phase diagram (Section~\ref{sec:data-scaling}). For \(b\leq b_1\), SGD, Polyak, and Nesterov achieve the same data-scaling exponent. For \(b_1<b<b_3\), Nesterov outperforms Polyak, which in turn outperforms SGD, with Nesterov attaining its best data-scaling rate at \(b=b_2\). For \(b\geq b_3\), Polyak and Nesterov again share the same exponent and both outperform SGD. Here \(b_1\) and \(b_3\) mark the critical batch exponents of SGD and Polyak, respectively, showing that Polyak preserves the best small-batch scaling over a substantially wider batch-size range, while Nesterov further improves the achievable data-scaling rate within the large batch regime. Figure~\ref{fig:intro-result} provides a complete view of how the relative advantage of the three methods evolves with batch size.
\end{itemize}
Finally, Section~\ref{sec:experiment} empirically validates these theoretical predictions.

\paragraph{Why momentum helps large-batch training.}
The batch-size dependence of momentum can be understood through the \textbf{per-sample effective learning rate} (ELR),
\[
\eta_{\rm sample}:=\frac{\eta_{\rm eff}}{B},\qquad
\eta_{\rm eff}=
\begin{cases}
\eta, & \textnormal{SGD},\\
\eta_\rho:=\eta/(1-\rho), & \textnormal{momentum}.
\end{cases}
\]
Ignoring the early transient and using \(T=D/B\), SGD and Polyak obey the same data-level signal--noise tradeoff,
\[
\mathcal E(D)\eqsim(D\eta_{\rm sample})^{-s}+\sigma^2\eta_{\rm sample}.
\]
The difference is the admissible range of the per-sample ELR. Stability condition~\eqref{eqn: stability-0} gives \(\eta_{\rm sample}\lesssim1/B\) for SGD, whereas Polyak allows \(\eta_{\rm sample}\lesssim1\). Thus, as the batch size grows, the admissible per-sample ELR of SGD shrinks as \(1/B\), while Polyak can keep it at order one by increasing \(\eta_\rho\) proportionally with \(B\). Polyak therefore extends the SGD signal--noise tradeoff to substantially larger batch sizes without changing its optimum, explaining both the enlarged critical batch size and the unchanged best data-scaling exponent.

Nesterov changes the picture more fundamentally. Its look-ahead mechanism suppresses noise accumulation at large effective learning rates, improving the signal--noise tradeoff itself and thereby achieving better data efficiency in the large-batch regime. At ultra-large batch sizes, however, the limited number of sequential updates \(T=D/B\) becomes the dominant bottleneck, and the two momentum methods recover the same signal-limited scaling.

% \vspace{-.5em}
\section{Related work}

% \vspace{-.5em}
\paragraph{Theory of momentum.}
Classical analyses explain momentum through accelerated convergence, Lyapunov functions, and continuous-time analysis~\citep{nesterov1983method,su2016differential,wibisono2016variational,wilson2021lyapunov,shi2022understanding,kovachki2021continuous}. Under stochastic training, subsequent work has studied convergence, stability, and asymptotic behavior~\citep{vidyasagar2025convergence,liu2020improved,yuan2016influence,gitman2019understanding,velikanov2023view,zhang2024optimalityacceleratedsgdhighdimensional}. However, existing theory does not characterize how the role of momentum changes with batch size under a fixed data budget. We instead treat batch size as a scaling variable and compare Polyak and Nesterov across the resulting training regimes.

\vspace{-.5em}
\paragraph{Scaling-law theory.}
Scaling laws characterize how learning performance scales with resources such as data, model, and compute~\citep{hestness2017deep,kaplan2020scaling,hoffmann2022training,grattafiori2024llama,deepseekai2026deepseekv4}. A growing theoretical literature has derived such laws in linear, kernel, and related tractable models~\citep{sharma2022scaling,hutter2021learning,maloney2022solvable,wei2022more,jain2024scaling,michaud2023quantization,nam2024exactly,atanasov2026scaling,bahri2024explaining,dohmatob2024tale,bordelon2024dynamical,lin2024scaling,paquette2024phases,bordelon2025feature,yan2026larger,lin2025improved,sous2026asymmetric,li2026functional,zhang2025critical,wang2026fast,li2026optimal,ding2025scalinglawstochasticgradient}. Functional scaling laws further characterize risk dynamics through signal learning and noise accumulation~\citep{li2026functional}. We extend this framework to momentum methods. Closely related, \citet{ferbach2026dimension} study momentum scaling laws in noiseless power-law random features with fixed batch size, showing that Polyak preserves the SGD scaling exponents while dimension- and time-adapted Nesterov-type methods can improve them. We instead consider constant-order label noise and scale the batch size under a fixed data budget, obtaining a batch-size phase diagram for canonical Polyak and Nesterov momentum.

\vspace{-.5em}
\paragraph{Stability of gradient-based methods.}
Stability characterizes the admissible hyperparameter region and has been studied extensively for gradient-based methods, with connections to implicit bias, generalization, and optimization geometry~\citep{zhu2019anisotropic,wu2018sgd,nar2018step,ma2021linear,wu2022alignment,wu2023implicit,chemnitz2025characterizing,mulayoff2021implicit,nacson2023implicit,qiao2024stable,kaur2023maximum,wang2023theoretical}. For momentum methods, its joint dependence on batch size and momentum remains less understood. We derive sharp risk stability boundaries for Polyak and Nesterov, revealing a much stronger batch-size compensation for high momentum in Nesterov.

\vspace{-.5em}
\paragraph{Concurrent work.}
\citet{morwani2026compute}, which appeared while this manuscript was in preparation, studies batch-size tradeoffs for stochastic momentum in finite-dimensional linear regression and derives lower bounds on contraction rates. These bounds suggest that heavy-ball may retain data efficiency over a larger batch-size range, but matching upper bounds are left open. In contrast, we derive sharp scaling laws for the full risk dynamics and, after optimizing the learning rate and momentum factor, obtain the data-scaling exponent as an explicit function of batch size. This yields a complete batch-size phase diagram, including a large-batch regime in which Nesterov outperforms Polyak.

% \vspace{-.7em}
\section{Preliminaries}
\label{sec:preliminaries}
\vspace{-.3em}
{\bf Notation.} In this paper, we use $\eqsim$ to denote equivalence up to a constant factor, and $\lesssim$ (resp.~$\gtrsim$) indicates inequality up to a constant factor. For two non-negative functions $f,g$, we write $f(x)\lesssim g(x)$ (resp. $f(x)\gtrsim g(x)$) if there exists a constant $C>0$ such that $f(x)\leq C g(x)$ (resp. $g(x)\leq C f(x)$) for all $x$ under consideration. We write $f(x)\eqsim g(x)$ if both $f(x)\lesssim g(x)$ and $f(x)\gtrsim g(x)$ hold. 
Throughout the paper, \(\langle\cdot,\cdot\rangle\) and
\(\lVert\cdot\lVert_2\) denote the standard inner product and the associated \(2\)-norm, respectively.


\subsection{Power-law kernel regression}
\label{sec:plkr}
We consider a power-law kernel (PLK) regression  on a Hilbert space \(\mathcal H\) with the input \(\bx\sim\cD\), and the observed label is \(y=\langle\bphi(\bx), \btheta^*\rangle + \epsilon\), where \(\epsilon \sim \cN(0,\sigma^2)\) is independent of \(\bx\). Here we assume the feature function $\bphi(\bx)$ to be Gaussian: $\bphi(\bx) \sim \cN(0,\bH)$, where $\bH = \EE_{\bx\sim \cD}[\bphi(\bx) \otimes\bphi(\bx)]$ denotes the feature covariance operator. Let \(\{\lambda_j\}_{j\geq1}\) denote the eigenvalues of \(\bH\), arranged in decreasing order.
We learn the target function using the linear student model $f(\bx;\btheta) = \langle\bphi(\bx), \btheta\rangle$ by minimizing the population risk:
$\cR(\btheta) = \frac{1}{2} \EE_{\bx,y}\left[(f(\bx;\btheta)-y)^2\right].$
A direct calculation gives $\cR(\btheta) = \cE(\btheta) + \frac{\sigma^2}{2}$, where $\cE(\btheta) = \frac{1}{2}\langle\btheta-\btheta^*,\bH (\btheta-\btheta^*)\rangle$ is the  excess risk.

Since SGD, Polyak momentum, and Nesterov momentum are orthogonally equivariant, we work in the eigenbasis of the covariance operator without loss of generality; see Appendix~\ref{app:subsec:rotation invariant}.

\begin{assumption}[Diagonal covariance]
\label{ass:diagonal-covariance}
The covariance operator is diagonal:
\begin{equation}
\label{eqn:diagonal-hessian}
\bH=\diag(\lambda_1,\lambda_2,\ldots),
\qquad
\lambda_1\geq\lambda_2\geq\cdots>0.
\end{equation}
\end{assumption}

Throughout, we focus on the label-noise-dominated setting with a constant-order noise level.
\begin{assumption}[Constant-order label noise]
Assume \(\sigma\eqsim1\).
\end{assumption}

\begin{assumption}[Source and capacity conditions] \label{ass:power-law}
    There exist  $\beta>1$ and $s>0$ such that 
    \[
        \lambda_j \eqsim j^{-\beta},\quad \lambda_j |\theta^*_j|^2\eqsim j^{-(1+s\beta)}.
    \]
\end{assumption}
Here \(\beta\) is the \textbf{capacity exponent}, which controls the decay rate of the feature spectrum, and \(s\) is the \textbf{source exponent}, which describes how the target energy decays across spectral directions. 
We also refer to \(s\) as the \textbf{relative difficulty} of the task: smaller \(s\) means that more target energy lies in low-curvature, small-eigenvalue directions, making the problem harder to learn. We remark that similar assumptions have been widely used in the analysis of kernel methods~\citep{caponnetto2005fast,caponnetto2007optimal,spigler2020asymptotic,bordelon2020spectrum,maloney2022solvable}.  
Our work builds upon and extends this line of research. 

\subsection{Optimization methods}
\label{sec:err recursion}
At each step, we draw a fresh mini-batch $\cB_k =\{(\bx_{k,i}, y_{k,i})\}_{i=1}^{B}$ and denote the  mini-batch gradient by $\bg(\btheta;\cB_k) = \nabla_{\btheta}(\tfrac{1}{2B}\sum_{(\bx,y)\in\cB_k}(f(\bx;\btheta)-y)^2)$. SGD updates as follows
\begin{align}
\textnormal{SGD:}&\qquad
\btheta_{k+1}
=
\btheta_k-\eta \bg(\btheta_k;\cB_k),
\label{eq:sgd-update}
\end{align}
where $\eta$ is the learning rate. Polyak momentum, also called the heavy-ball method, is given by 
\begin{align}
\textnormal{Polyak:}&\qquad
\btheta_{k+1}
=
\btheta_k
-\eta\bg(\btheta_k;\cB_k)+\rho(\btheta_k-\btheta_{k-1}),
\label{eq:hb-update}
\end{align}
where $\rho\in [0,1)$ is the momentum factor. Nesterov momentum differs from Polyak momentum by evaluating the gradient at a look-ahead point:
\begin{align}
\notag \textnormal{Nesterov:}&\qquad\btheta_k^{\la}=\btheta_k+\rho(\btheta_k-\btheta_{k-1})\qquad \text{(look ahead)}\\
&\qquad \btheta_{k+1}=\btheta_k -\eta\bg(\btheta_k^{\la};\cB_k)+\rho(\btheta_k-\btheta_{k-1}),
\label{eq:nag-update}
\end{align}
where 
\(\btheta_k^{\la}\) is the look-ahead point obtained by extrapolating the current iterate along the previous update direction. This look-ahead mechanism distinguishes Nesterov from Polyak.


For clarity, we decompose  \(
\bg(\btheta;\cB_k)
=
\nabla\cR(\btheta)+\bxi_k(\btheta),
\)
where \(\nabla\cR(\btheta)\) denotes the population gradient and \(\bxi_k(\btheta)\) represents  the stochastic gradient noise satisfying
\[
\EE[\bxi_k(\btheta)\mid\btheta]=\bm 0,
\qquad
\EE\!\left[
\bxi_k(\btheta)\otimes\bxi_k(\btheta)
\,\middle|\,\btheta
\right]
=
\frac{1}{B}\bSigma(\btheta),
\]
where \(\bSigma(\btheta)\) denotes the gradient-noise covariance of batch size \(1\). We next recall a noise--curvature alignment property that will be used
throughout our analysis.

\begin{proposition}[Noise--curvature alignment]
\label{prop:noise-geometry}
For any unit vector \(\bnu\in\mathcal H\),
\begin{equation}
\label{eqn:noise-geometry}
\mathbb E\!\left[
\left|
\left\langle\bxi_k(\btheta_k),\bnu\right\rangle
\right|^2
\,\middle|\,
\btheta_k
\right]
\eqsim
\frac{1}{B}
\bigl(\cE(\btheta_k)+\sigma^2\bigr)
\left\langle\bnu,\bH\bnu\right\rangle.
\end{equation}
\end{proposition}

Here, \(\langle\bnu,\bH\bnu\rangle\) is the curvature along \(\bnu\). The two terms in Proposition~\ref{prop:noise-geometry} have distinct origins and roles. The term proportional to \(\cE(\btheta_k)\) arises from mini-batch sampling, depends on the current iterate, and vanishes at the optimum, acting as \emph{multiplicative noise}. The term proportional to \(\sigma^2\) arises from label noise and persists at the optimum, acting as \emph{additive noise}. Both are curvature-aligned but affect the dynamics differently: multiplicative noise feeds back into the dynamics and determines stability, whereas additive noise determines the asymptotic noise floor. This noise--curvature alignment has been observed in neural networks~\citep{zhu2019anisotropic,wu2022alignment} and established rigorously in toy models~\citep{wang2023theoretical}. We use this  constitutive relation throughout our analysis; see Appendix~\ref{app:hessian alignment} for a proof.
% \paragraph{Error recursions.} 
% Let
% \(
% \be_k=\btheta_k-\btheta^*
% \)
% and \(
% \be_k^{\la}
% =
% (1+\rho)\be_k-\rho\be_{k-1}.
% \)
% Our analysis shall relies on  error recursions of three optimizers:
% \begin{align}
% \textnormal{SGD:}&\qquad
% \be_{k+1}
% =
% (\bI-\eta\bH)\be_k
% -\eta\bxi_k(\btheta_k),
% \label{eq:sgd-error-update}
% \\
% \textnormal{Polyak:}&\qquad
% \be_{k+1}
% =
% \bigl((1+\rho)\bI-\eta\bH\bigr)\be_k
% -\rho\be_{k-1}
% -\eta\bxi_k(\btheta_k),
% \label{eq:hb-error-update}
% \\
% \textnormal{Nesterov:}&\qquad
% \be_{k+1}
% =
% (\bI-\eta\bH)\be_k^{\la}
% -\eta\bxi_k(\btheta_k^{\la}).
% \label{eq:nag-error-update}
% \end{align}

\subsection{Regime-dependent damping in momentum methods}
\label{sec:damping-regimes}
% This is the discrete analogue of a damped mass--spring system: the three
% terms represent inertia, friction, and the restoring force, respectively.

\begin{wrapfigure}{r}{0.33\textwidth}
\centering
\vspace{-2em}
\hspace{1em}\includegraphics[width=\linewidth]{figures/momentum_phase_diagram.pdf}
\vspace{-2em}
\end{wrapfigure}

Consider the full-batch Polyak on a one-dimensional quadratic function
$
f(x)=\lambda x^2/2,
$
where \(\lambda>0\) denotes the curvature. Since $\nabla f(x)=\lambda x$, Polyak updates:
\begin{equation}\label{eqn: polyak-1}
x_{k+1} =  x_k-\eta\lambda x_k+\rho(x_k-x_{k-1}).
\end{equation}
The corresponding characteristic polynomial is
$
r^2-(1+\rho-\eta\lambda)r+\rho=0.
$
The relevant transition from real to complex roots occurs at $\eta\lambda= (1-\sqrt{\rho})^2$. 

For fixed \(\eta\) and \(\rho\), different curvature directions can
exhibit two qualitatively different damping behaviors.
\begin{itemize}
    \item \textbf{Overdamped regime}: $\lambda \eta  < (1-\sqrt{\rho})^2$. In this regime, the characteristic roots $r_{\pm}$  are real with the dominant root
satisfying
$
r_+
\approx
1-\frac{\eta\lambda}{1-\rho}
$
when $1-\rho=o(1)$.
Hence, for sufficiently large $k$, it holds that
\[
x_k =c_{+} r_{+}^k + c_{-} r_{-}^k \approx c_{+} r_{+}^k 
\approx
c_{+}\exp\left(
-\frac{\eta\lambda}{1-\rho}k
\right),
\]
where $c_{\pm}$ depend on the initialization. Hence, its update
behaves like gradient descent with the {\it momentum-rescaled learning rate} 
$
\eta_{\rho} = \eta/(1-\rho).
$

\item \textbf{Underdamped regime}: $\lambda\eta > (1-\sqrt{\rho})^2$. In this regime, the characteristic roots form a complex-conjugate pair with modulus
\(\sqrt{\rho}\). The iterate therefore has the form
$
x_k
=
\rho^{k/2}
\bigl[c_1\cos(k\omega)+c_2\sin(k\omega)\bigr],
$
where \(c_1,c_2\), and \(\omega\) depend on the initialization and curvature. Thus, the iterate exhibits oscillatory convergence to zero, with the amplitude decaying as \(\rho^{k/2}\). 
\end{itemize}

% \begin{figure}
% \centering
% \includegraphics[width=0.4\linewidth]{figures/momentum_phase_diagram.pdf}
% \caption{xx}
% \end{figure}

The key difference between the two regimes lies in how curvature affects
the dynamics. In the overdamped regime, the decay rate
\(\eta\lambda/(1-\rho)\) is proportional to the curvature, so flatter
directions converge more slowly. In the underdamped regime, the
oscillation frequency depends on the curvature, but the amplitude
envelope \(\rho^{k/2}\) does not.

\paragraph{Nesterov look-ahead induces curvature-adaptive damping.} On the same quadratic objective, Nesterov updates using the gradient at the look-ahead point $x_k^{\mathrm{la}}= x_k+\rho(x_k-x_{k-1})$:
\begin{align}\label{eqn: curvature-adaptivity}
\notag x_{k+1} &= x_k - \eta \lambda [x_k+\rho(x_k-x_{k-1})] + \rho (x_k - x_{k-1})\\
&=x_k - \eta\lambda x_k + \rho(1-\eta\lambda) (x_k -x_{k-1}),
\end{align}
which is equivalent to Polyak with the momentum factor $\rho(1-\eta\lambda)$. 
For an underdamped direction, it accelerates convergence:
$
\rho^{k/2}
\to 
[\rho(1-\eta\lambda)]^{k/2},
$ 
which damps sharper directions more rapidly.
By contrast, in sufficiently flat directions where
\(\eta\lambda\ll1\), we have
\(\rho(1-\eta\lambda)\approx\rho\), and Nesterov retains
approximately the same  dynamics as Polyak.


\begin{remark}[Spectral decomposition into damping regimes]
For a fixed pair $(\eta,\rho)$, the damping threshold
$
\lambda_{\mathrm{cut}}
:=
\frac{(1-\sqrt{\rho})^2}{\eta}
$
partitions the spectrum into two parts. Directions with
$\lambda_j>\lambda_{\mathrm{cut}}$ are underdamped, whereas directions
with $\lambda_j<\lambda_{\mathrm{cut}}$ are overdamped. Hence, in PLK
regression, the risk dynamics can be analyzed by treating the
underdamped spectral head and overdamped spectral tail separately and
then combining their contributions.
\end{remark}

\section{Risk stability and critical learning rates}
\label{sec:stability}

In this section, we characterize the maximal learning rates that ensure stable training and study their dependence on the momentum factor and batch size. We measure stability directly through the population risk.

\begin{definition}[Risk stability]
Consider a stochastic optimization algorithm for minimizing a population risk $\mathcal R(\cdot)$, producing model parameters $\{\theta_k\}_{k\ge 0}$. We say that the algorithm is \emph{risk stable} if, for every admissible initialization with finite population risk,
\[
\sup_{k\ge 0}\mathbb E[\mathcal R(\theta_k)]<\infty.
\]
\end{definition}
% Any auxiliary states or previous iterates required by the algorithm are treated as part of the initialization.

\begin{definition}[Critical learning rate]
Consider a stochastic optimization algorithm $\mathcal A$ parameterized by a learning rate $\eta$ and additional hyperparameters $\bh$. For fixed $\bh$,  let
\[
\eta_{\mathcal A}^{\crit}(\bh):=\sup\left\{\eta>0:\mathcal A\text{ is risk stable}\right\}.
\]
\end{definition}
Hence, the critical learning rate is the largest learning rate compatible with risk stability. When the dependence on $\bh$ is clear from context, we simply write $\eta_{\mathcal A}^{\crit}$. In our setting, $\bh=B$ for SGD and $\bh=(B,\rho)$ for Polyak and Nesterov. The following theorem characterizes how the critical learning rate scales with batch size and momentum; the proof is deferred to Appendix~\ref{app:sec:stability}.
\begin{highlightedtheorem}
\begin{theorem}[Critical learning-rate scaling]\label{thm:stability}
\normalfont
For sufficiently large \(B\) and sufficiently small \(1-\rho\), the critical learning rates satisfy
\begin{align}
\eta_{\sgd}^{\crit} &\eqsim 1,\\
\eta_{\polyak}^{\crit} &\eqsim \min\{1,\,B(1-\rho)\},\\
\eta_{\nesterov}^{\crit} &\eqsim \min\{1,\,B^\beta(1-\rho)\},
\end{align}
where the implicit constants are independent of \(B\) and \(\rho\).
\end{theorem}
\end{highlightedtheorem}



\paragraph{Joint dependence on momentum and batch size.}
Theorem~\ref{thm:stability} reveals a direct coupling among the learning
rate, momentum factor, and batch size in determining training stability.
Increasing $\rho$ toward one reduces the critical learning rate, whereas
increasing the batch size enlarges it. In the high-momentum regime,
$
\eta_{\polyak}^{\crit}\eqsim B(1-\rho),
\eta_{\nesterov}^{\crit}\eqsim B^\beta(1-\rho),
$
before reaching their order-one ceilings. Thus, the batch-size effect on
stability is substantially stronger for Nesterov, with a $B^\beta$
dependence rather than the $B$ dependence of Polyak. This difference originates from the curvature adaptivity induced by Nesterov's look-ahead mechanism.
\paragraph{Multiplicative-noise feedback governs stability through the renewal kernel.} For all three methods, Appendix~\ref{app:sec:stability} shows that the risk dynamics admit, up to constant factors, a renewal equation
$
\ell_k\eqsim h_k+\sum_{j=0}^{k-1}K_{k-1-j}\bigl(\ell_j+\sigma^2\bigr),
$
where \(\ell_k\) is the risk moment, \(h_k\) is the free response, and the renewal kernel \(K\) describes how stochastic perturbations propagate through the optimizer dynamics. Additive and multiplicative noise play different roles: additive label noise acts only as an external forcing, whereas multiplicative noise feeds the current risk back through the renewal kernel, producing the feedback term \(K*\ell\). Risk stability is therefore determined by the strength of this feedback, characterized by \(\|K\|_1<1\); see Appendix~\ref{app:subsec:label-noise-stability}. Since \(K\) is independent of the label-noise scale \(\sigma\), the stability condition in Theorem~\ref{thm:stability} is also independent of \(\sigma\). 
% This distinction changes the conclusion on acceleration. Under state-independent noise models, Polyak momentum can take \(\rho\to1\) without an accompanying stability penalty and thereby achieve improved statistical rates~\citep{zhang2024optimality}. With mini-batch multiplicative noise, however, the feedback imposes \(\eta\lesssim B(1-\rho)\), so at fixed batch size the effective learning rate \(\eta_\rho=\eta/(1-\rho)\) remains bounded; the faster signal dynamics therefore do not translate into an improved data-scaling exponent.



\section{Scaling laws for risk dynamics across damping regimes}
\label{sec:scaling law}

In this section, we derive sharp scaling laws for the full risk dynamics of Polyak and Nesterov, explicitly characterizing their joint dependence on the hyperparameters \((\eta,\rho,B)\). The analysis follows how increasing \(\rho\) drives transitions in damping and stability: from a fully overdamped regime with SGD-like dynamics, to a mixed-damping regime, and eventually to instability.

As discussed in Section~\ref{sec:damping-regimes}, for a single eigendirection with curvature \(\lambda\), the overdamped--underdamped transition occurs exactly at
$
(1-\sqrt{\rho})^2=\eta\lambda.
$
For the scaling analysis below, where \(\rho\to1\), we use the equivalent high-momentum form
$
\eta\lambda\eqsim(1-\rho)^2.
$
This motivates the curvature cutoff
\[
\lambda_{\mathrm{cut}}
\eqsim
\frac{(1-\rho)^2}{\eta},
\qquad
\II_{\mathrm{over}}
=
\left\{
j:\lambda_j\lesssim\lambda_{\mathrm{cut}}
\right\},
\qquad
\II_{\mathrm{under}}
=
\left\{
j:\lambda_j\gtrsim\lambda_{\mathrm{cut}}
\right\}.
\]
Hence, the sets \(\II_{\mathrm{over}}\) and \(\II_{\mathrm{under}}\) consist of flat overdamped directions and sharp underdamped directions, respectively. Since the PLK spectrum contains arbitrarily small eigenvalues, overdamped directions always remain present.

As $\rho$ increases, the first underdamped directions appear when \(\lambda_{\mathrm{cut}}\lesssim \lambda_1\), corresponding to
\[
1-\rho\eqsim\sqrt{\eta}.
\]
Combining this damping threshold with the risk stability conditions from Section~\ref{sec:stability} yields three regimes:
\[
\begin{array}{lll}
\textnormal{fully overdamped:}
& 1-\rho\gtrsim\sqrt{\eta},\\[1mm]
\textnormal{mixed damping:}
& \eta/B^q\lesssim1-\rho\lesssim\sqrt{\eta},\\[1mm]
\textnormal{unstable:}
& 1-\rho\lesssim\eta/B^q,
\end{array}
\qquad
q=
\begin{cases}
1, & \textnormal{Polyak},\\
\beta, & \textnormal{Nesterov}.
\end{cases}
\]
Thus, the damping threshold \(1-\rho\eqsim\sqrt{\eta}\) and the risk stability threshold \(1-\rho\eqsim\eta/B^q\) delimit the three regimes.

% \paragraph{Scaling law of SGD.}
For reference, the {\bf SGD scaling law} for PLK regression is
\begin{equation}
\label{eqn:scaling-law-sgd}
\mathbb E[\cE(\btheta_k)]\eqsim(\eta k)^{-s}+\frac{\sigma^2}{B}\eta,\qquad k\eta\gtrsim1.
\end{equation}
This follows from \citet[Theorem 5.2]{li2026functional}; see Appendix~\ref{app:sgd_scaling_law} for details. The first term, \((\eta k)^{-s}\), describes signal learning, with the decay rate determined by the source exponent \(s\): larger \(s\) corresponds to an easier target and faster learning. The second term, \(\sigma^2\eta/B\), captures stochastic noise accumulation and determines the asymptotic noise floor.



% \begin{figure}[!h]
% \centering
% \begin{tikzpicture}[x=1.2cm,y=0.8cm]

% % axis
% \draw[->,thick] (0,0) -- (10.5,0) node[right] {$\rho$};

% % regions
% \fill[blue!20] (0,-0.35) rectangle (3.5,0.35);
% \fill[orange!25] (3.5,-0.35) rectangle (7.3,0.35);
% \fill[red!20] (7.3,-0.35) rectangle (10,0.35);

% % boundaries
% \draw[thick] (3.5,-0.45)--(3.5,0.45);
% \draw[thick] (7.3,-0.45)--(7.3,0.45);

% % phase labels
% \node[align=center,font=\small] at (1.75,0)
% {Fully overdamped};

% \node[align=center,font=\small] at (5.4,0)
% {Mixed};

% \node[align=center,font=\small] at (8.65,0)
% {Unstable};

% % threshold labels
% \node[below=6pt] at (3.5,0)
% {$(1-\sqrt{\eta})^2$};

% \node[below=6pt] at (7.3,0)
% {$1-\eta/B^{q}$};

% % endpoints
% \node[below=6pt] at (0,0) {$0$};
% \node[below=6pt] at (10,0) {$1$};

% % % mixed explanation
% % \node[above=10pt,align=center,font=\small]
% % at (5.4,0)
% % {
% % sharp modes: underdamped\\
% % flat modes: overdamped
% % };
% \end{tikzpicture}

% \vspace{-.5em}
% \caption{
% {\bf Damping and stability regimes of momentum}
% The mixed regime consists of sharp underdamped modes and flat overdamped
% modes. Here  \(q=1\) for Polyak  and \(q=\beta\) for Nesterov.
% }
% \label{fig: phase-digram-damping}
% \end{figure}


The overdamped modes evolve with the rescaled learning rate
\[
\eta_\rho:=\frac{\eta}{1-\rho},
\]
which plays the role of the effective learning rate for momentum. The following theorem establishes the  scaling law for Polyak momentum, whose proof is deferred to Appendix~\ref{app:hbm scaling law}.

\begin{highlightedtheorem}
\begin{theorem}[Scaling law for Polyak]
\label{thm:hbm scaling law}
\normalfont
Assume \(\eta\lesssim1\) and \(k\gtrsim\max(\eta^{-1}_\rho,(1-\rho)^{-1})\). Then
\begin{equation}
\label{eq:hb-loss-scaling}
\mathbb E[\cE(\btheta_k)]\eqsim
\begin{cases}
(\eta_\rho k)^{-s}+\dfrac{\sigma^2}{B}\eta_\rho, & 1-\rho\gtrsim\sqrt{\eta},\\[2mm]
P(k)+(\eta_\rho k)^{-s}+\dfrac{\sigma^2}{B}\eta_\rho, & \dfrac{\eta}{B}\lesssim1-\rho\lesssim\sqrt{\eta},
\end{cases}
\end{equation}
where \(P(k)\) denotes the underdamped transient and satisfies \(0\leq P(k)\lesssim\rho^k\). If \(1-\rho\lesssim\eta/B\), the dynamics are not risk stable.
\end{theorem}
\end{highlightedtheorem}


Theorem~\ref{thm:hbm scaling law} identifies three regimes as \(\rho\) increases: fully overdamped SGD-like dynamics with effective learning rate \(\eta_\rho\), mixed damping with an additional transient \(P(k)\), and eventually risk instability.
In the mixed-damping regime, the two risk contributions have a direct spectral interpretation. The power-law term \((\eta_\rho k)^{-s}\) arises from the overdamped flat modes in \(\II_{\mathrm{over}}\), whereas \(P(k)\) comes from the underdamped sharp modes in \(\II_{\mathrm{under}}\) and decays with envelope \(\rho^k\). Increasing \(\rho\) therefore accelerates flat-mode learning through \(\eta_\rho\), while slowing the relaxation of sharp modes. Since \(\rho^k\approx\exp(-(1-\rho)k)\), the latter occurs on the timescale
\[
\tau_\rho:=\frac{1}{1-\rho}\log\left(\frac{\eta}{(1-\rho)^2}\right).
\]


% \paragraph{Long-time equivalence to SGD.}
% \begin{remark}
Once the underdamped transient becomes negligible, the Polyak risk law reduces to the SGD law under the replacement \(\eta\mapsto\eta_\rho\), leaving the signal--noise tradeoff unchanged. The distinction lies in the admissible effective learning rate. By Theorem~\ref{thm:stability}, SGD requires \(\eta\lesssim1\), whereas Polyak allows \(\eta_\rho\lesssim B\), so its admissible range grows linearly with the batch size. As we show in Section~\ref{sec:data-scaling}, this enlarged range increases the critical batch size but does not improve the best data-scaling exponent.
% \end{remark}

For Nesterov, the same damping structure persists, but its look-ahead mechanism modifies both the underdamped transient and noise accumulation in the mixed-damping regime.

\begin{highlightedtheorem}
\begin{theorem}[Scaling law for Nesterov]
\label{thm:nesterov scaling law}
\normalfont
Assume \(\eta<1\) and \(k\gtrsim\max(\eta^{-1}_\rho,(1-\rho)^{-1})\). Then
\begin{equation}
\label{eq:nag-loss-scaling}
\mathbb E[\cE(\btheta_k)]\eqsim
\begin{cases}
(\eta_\rho k)^{-s}+\dfrac{\sigma^2}{B}\eta_\rho, & 1-\rho\gtrsim\sqrt{\eta},\\[2mm]
N(k)+(\eta_\rho k)^{-s}+\dfrac{\sigma^2}{B}\min\{\eta_\rho,\eta_\rho^{1/\beta}\}, & \dfrac{\eta}{B^\beta}\lesssim1-\rho\lesssim\sqrt{\eta},
\end{cases}
\end{equation}
where \(N(k)\) denotes the underdamped transient and satisfies \(0\leq N(k)\lesssim\min\{1,(\eta k)^{-s}\}\rho^k\). If \(1-\rho\lesssim\eta/B^\beta\), the dynamics are not risk stable.
\end{theorem}
\end{highlightedtheorem}

The proof is deferred to Appendix~\ref{app:nesterov scaling law}.
In the fully overdamped regime, Nesterov and Polyak obey the same scaling law, both behaving like SGD with effective learning rate \(\eta_\rho\). Their difference in the mixed-damping regime stems from the curvature-adaptive damping induced by Nesterov's look-ahead mechanism, which damps sharper directions more strongly; see Eq.~\eqref{eqn: curvature-adaptivity}. This suppresses the underdamped transient from
\[
P(k)\lesssim\rho^k
\qquad\textnormal{to}\qquad
N(k)\lesssim\min\{1,(\eta k)^{-s}\}\rho^k,
\]
while leaving the asymptotic power-law signal term unchanged at \((\eta_\rho k)^{-s}\). The same mechanism also suppresses stochastic noise accumulation. When \(\eta_\rho\gg1\), the noise floor scales as
\[
\frac{\sigma^2}{B}\eta_\rho
\qquad\textnormal{for Polyak},\qquad
\frac{\sigma^2}{B}\eta_\rho^{1/\beta}
\qquad\textnormal{for Nesterov}.
\]
For asymptotic data scaling, this reduction in noise accumulation is the key effect.


\section{Optimal data scaling across batch-size regimes}
\label{sec:data-scaling}

We now optimize the final-step risk using the scaling laws derived above, subject to the one-pass constraint $T=D/B$ and the stability constraints. We focus on the large-batch setting, allowing $B$ to grow with the data budget $D$ rather than remain fixed.
\begin{definition}[Batch-size exponent]
\label{def:batch-size-exponent}
For a sequence of training problems with \(D\to\infty\), the batch size has exponent \(b\in[0,1)\) if
\[
B=D^{b+o(1)},
\qquad\text{equivalently}\qquad
b=\lim_{D\to\infty}\frac{\log B}{\log D}.
\]
\end{definition}
% Then  \(b\) quantifies how the batch size scales with the data budget.

\begin{definition}[Batch-size-dependent data-scaling exponent]
\label{def:batch-scaling-exponent}
Fix \(b\in[0,1)\), so that \(B=D^{b+o(1)}\) and \(T=D/B\). For an optimization method with admissible hyperparameter region \(\cA_B\), define its optimal data-scaling exponent \(r(b)\) by
\[
\inf_{(\eta,\rho)\in\cA_B}\EE[\cE(\btheta_T)]
=
D^{-r(b)+o(1)}.
\]
\end{definition}
Here, \(\rho\) is omitted for SGD. For the three methods considered, we write \(r_{\sgd}(b)\), \(r_{\polyak}(b)\), and \(r_{\nesterov}(b)\). A larger \(r(b)\) indicates better data efficiency at batch-size exponent \(b\).

Define the three transition exponents
\begin{equation}
b_1=\frac{s}{s+1},\qquad
b_2=\frac{2s+1/\beta}{2s+1+1/\beta},\qquad
b_3=\frac{2s+1}{2s+2}.
\end{equation}
Since \(\beta>1\) and \(s>0\), \(0<b_1<b_2<b_3<1\). The following theorem gives the optimal data-scaling exponents across batch-size regimes; the proof is deferred to Appendix~\ref{app:sec:data-scaling}.

\begin{highlightedtheorem}
\begin{theorem}[Batch-size-dependent data scaling]
\label{thm:batch-size-scaling}
For \(B=D^{b+o(1)}\) with \(b\in[0,1)\), 
\[
r_{\sgd}(b)=
\begin{cases}
\frac{s}{s+1}, & 0\leq b<b_1,\\[2mm]
s(1-b), & b_1\leq b<1,
\end{cases}
\qquad
r_{\polyak}(b)=
\begin{cases}
\frac{s}{s+1}, & 0\leq b<b_3,\\[2mm]
2s(1-b), & b_3\leq b<1,
\end{cases}
\]
and
\[
r_{\nesterov}(b)=
\begin{cases}
\frac{s}{s+1}, & 0\leq b<b_1,\\[2mm]
\frac{s(1+(\beta-1)b)}{s\beta+1}, & b_1\leq b<b_2,\\[3mm]
2s(1-b), & b_2\leq b<1.
\end{cases}
\]
\end{theorem}
\end{highlightedtheorem}



This theorem shows that the achievable data efficiency depends sharply on batch-size scaling. The resulting phase structure is most clearly seen in Figure~\ref{fig:intro-result} (left), which visualizes the three data-scaling exponents as functions of the batch exponent \(b\) and makes their relative ordering across regimes immediate. The three transition exponents play distinct roles: \(b_1\) is the critical batch exponent of SGD and marks the onset of Nesterov's improvement, \(b_2\) maximizes the Nesterov data-scaling exponent, and \(b_3\) is the critical batch exponent of Polyak.

\begin{itemize}
\item \textbf{Small batch} (\(b\leq b_1\)). All three methods achieve the same optimal data-scaling exponent:
\[
r_{\nesterov}(b)=r_{\polyak}(b)=r_{\sgd}(b).
\]

\item \textbf{Large batch} (\(b_1<b<b_3\)). The three methods separate:
\[
r_{\nesterov}(b)>r_{\polyak}(b)>r_{\sgd}(b),
\]
with Nesterov attaining its best data-scaling exponent at \(b=b_2\).

\item \textbf{Ultra-large batch} (\(b\geq b_3\)). Polyak and Nesterov share the same exponent and outperform SGD:
\[
r_{\nesterov}(b)=r_{\polyak}(b)>r_{\sgd}(b).
\]
\end{itemize}

Thus, Polyak enlarges the batch-size range over which the best small-batch  data scaling is preserved, whereas Nesterov further improves the achievable data-scaling exponent in the large-batch regime.




\section{Numerical validations}
\label{sec:experiment}

In this section, we numerically validate our theoretical predictions.
% \footnote{Code is available at \faGithub~\url{https://github.com/alexhuang13/Momentum_scaling}.} 
Experimental details and additional results are provided in Appendix~\ref{app:experiment}.


\paragraph{Risk stability boundaries.}
We validate the stability predictions of Theorem~\ref{thm:stability} for \((s,\beta)=(3,2)\). Figure~\ref{fig:numerical-summary} (left) fixes \(\rho\) and varies \(B\), confirming \(\eta_\polyak^\crit\eqsim B\) and \(\eta_\nesterov^\crit\eqsim B^\beta\) before saturation, while $\eta_{\sgd}^{\crit}$ remains order one. Figure~\ref{fig:numerical-summary} (middle) fixes \(B=32\) and varies \(1-\rho\), confirming the predicted linear dependence of both momentum stability boundaries on \(1-\rho\). Additional results are provided in Appendix~\ref{app:stability-experiment}.
% \paragraph{Mean-square stability boundaries.}
% We first test the stability conditions derived in
% Section~\ref{sec:stability}. Fixing the momentum factor \(\rho\), we vary the batch size \(B\) and
% scan the learning rate \(\eta\) for SGD, Polyak, and Nesterov. For each
% method and batch size, we estimate the empirical stability boundary as
% the largest learning rate for which the loss remains bounded over the
% prescribed training horizon. The precise stability criterion and
% experimental protocol are provided in
% Appendix~\ref{app:stability-experiment}.  
% Each marker in Figure~\ref{fig:numerical-summary} (left) reports the mean estimate over five independent random seeds. As predicted, the stability boundary of SGD remains of order one, while those of Polyak and Nesterov grow as \(\eta_\polyak^\crit\eqsim B\) and \(\eta_\nesterov^\crit\eqsim B^\beta\), respectively, before saturating at their order-one  ceilings. These results confirm the distinct batch-size scalings  predicted by Theorem~\ref{thm:stability}. Additional results for different choices of $s$ and $\beta$ are reported in Appendix~\ref{app:stability-experiment}.

\paragraph{Risk-dynamics scaling laws.}
We next test the full risk dynamics predicted by Theorems~\ref{thm:hbm scaling law} and~\ref{thm:nesterov scaling law}. For fixed \((\eta,\rho,B)\) and \((s,\beta)=(0.3,4)\), Figure~\ref{fig:numerical-summary} (right) verifies three key predictions: an underdamped transient whose transition timescale has dominant $(1-\rho)^{-1}$ dependence, a common power-law signal decay \((\eta_\rho k)^{-s}\), and distinct asymptotic noise floors. In particular, once the transient becomes negligible, Polyak and Nesterov follow the same power law, while Nesterov saturates at a substantially lower noise floor. These are consistent with our theoretical predictions.
 Additional results across \((s,\beta)\) and hyperparameter settings are provided in Appendix~\ref{app:scaling-law-experiments}.

\begin{figure}[!h]
    \centering
    \begin{subfigure}[t]{0.32\textwidth}
        \centering
        \includegraphics[width=\linewidth]{figures/stability/beta2s3_plot_main.pdf}
        % \label{fig:stability-beta15}
    \end{subfigure}
    \hspace{-.2em}
    \begin{subfigure}[t]{0.32\textwidth}
        \centering
        \includegraphics[width=\linewidth]{figures/stability/new_stability.pdf}
    \end{subfigure}
    \hspace{-.2em}
    \begin{subfigure}[t]{0.32\textwidth}
        \centering
        \includegraphics[width=\linewidth]{figures/scaling_law_new/total_loss_experiment_vs_theory.pdf}
        % \label{fig:scaling-law-main}
    \end{subfigure}
    \vspace{-.5em}
  \caption{\textbf{Numerical validation of stability and risk-dynamics scaling laws.}
\textbf{Left:} Critical learning rate versus batch size at fixed momentum. Markers show numerical estimates and dashed curves the theoretical predictions.
\textbf{Middle:} Critical learning rate versus \(1-\rho\) at fixed batch size, confirming the predicted linear scaling before the order-one ceiling.
\textbf{Right:} Polyak and Nesterov risk dynamics, with empirical results shown by thin translucent lines and theoretical predictions by thick solid lines.  The theory captures the underdamped transient, power-law decay, and noise floor. The inset confirms the dominant \((1-\rho)^{-1}\) dependence of the predicted transition timescale \(\tau_\rho\eqsim(1-\rho)^{-1}\log\!\big(\eta/(1-\rho)^2\big)\), with fitted log-log slopes \(-1.00\) for Polyak and \(-1.01\) for Nesterov.}
    \label{fig:numerical-summary}
    \vspace{-1em}
\end{figure}


\paragraph{Batch-size phase diagram at a fixed data budget.} We first test the three-regime structure predicted by Theorem~\ref{thm:batch-size-scaling} at a finite data budget. We fix \(D=2^{14}\) and, for each batch size \(B\), tune \((\eta,\rho)\) to minimize the final excess risk. Figure~\ref{fig:intro-result} (right) exhibits a progression across the three regimes. At small batch sizes, SGD, Polyak, and Nesterov achieve identical risk. As \(B\) increases, SGD deteriorates first, while Polyak maintains the small-batch performance over a wider range, reflecting its enlarged critical batch size. Nesterov goes further: in the large-batch regime, its risk continues to decrease beyond the best level attained by SGD and Polyak, before eventually increasing at ultra-large batch sizes. Thus, Polyak enlarges the data-efficient batch-size range, while Nesterov further improves data efficiency.

\paragraph{Data-scaling exponents across batch-size regimes.}
Having verified the global phase structure, we next test the predicted data-scaling exponents at representative points in each regime. We continue with \((s,\beta)=(0.3,4)\) and choose \(B\eqsim1\), \(B\eqsim D^{0.4}\), and \(B\eqsim D^{0.8}\), representing small batch, large batch, and ultra-large batch, respectively. For each data budget \(D\), we set \((\eta,\rho)\) according to the theoretically prescribed scalings. Figure~\ref{fig:three-regime-plk} shows that the empirical slopes closely match Theorem~\ref{thm:batch-size-scaling}: all three methods share the same exponent for small batch; for large batch, Nesterov outperforms Polyak, which outperforms SGD; and for ultra-large batch, Polyak and Nesterov again share the same exponent and both outperform SGD. These results quantitatively validate the batch-dependent data-scaling exponents underlying the phase diagram. Experimental details and additional settings are provided in Appendix~\ref{app:data-scaling-experiments}.

\begin{figure}[H]
\centering
\includegraphics[width=0.32\textwidth]{figures/small_batch.pdf}
\includegraphics[width=0.32\textwidth]{figures/large_batch.pdf}
\includegraphics[width=0.32\textwidth]{figures/super_large_batch.pdf}
\vspace{-.5em}
\caption{\textbf{Data-scaling exponents across the three batch-size regimes.} We choose representative batch-size scalings \(B\eqsim1\) \textbf{(left)}, \(B\eqsim D^{0.4}\) \textbf{(middle)}, and \(B\eqsim D^{0.8}\) \textbf{(right)}, and vary the data budget \(D\). Markers show empirical final excess risks, while dashed lines indicate the scaling exponents predicted by Theorem~\ref{thm:batch-size-scaling}. For these representative choices, the predicted ordering is recovered: Nesterov \(=\) Polyak \(=\) SGD at small batch, Nesterov \(>\) Polyak \(>\) SGD at large batch, and Nesterov \(=\) Polyak \(>\) SGD at ultra-large batch.}
% \vspace{-1em}
\label{fig:three-regime-plk}
\end{figure}

% From Figure \ref{fig: three-regime-plk} (middle, right), we can see that,  SGD no longer attains its data-optimal rate, since the batch size grows beyond the critical batch size. In contrast, Polyak momentum can compensate by appropriately scaling the momentum parameter and retains the rate \(D^{-\frac{s}{s+1}}\), while Nesterov momentum exhibits the steeper slope predicted by its improved data scaling theory. These results illustrate that momentum methods can make effective use of larger batch sizes, with Nesterov momentum achieving the fastest data scaling among the three methods.


% \begin{figure}[t]
%     \centering
%     \begin{subfigure}[t]{0.49\textwidth}
%         \centering
%         \includegraphics[width=\linewidth]{figures/mom_large_batch.pdf}
%         \caption{\(s=0.3\) and \(\beta=4\).}
%         \label{fig:polyak data scaling}
%     \end{subfigure}
%     \hfill
%     \begin{subfigure}[t]{0.49\textwidth}
%         \centering
%         \includegraphics[width=\linewidth]{figures/nag_acc.pdf}
%         \caption{\(s=0.8\) and \(\beta=4\).}
%         \label{fig:nesterov data scaling}
%     \end{subfigure}
%     \caption{Empirical comparison of the data scaling behavior of SGD, Polyak momentum, and Nesterov momentum. Markers show the simulation results, and dashed lines indicate the corresponding theoretical predictions. \textbf{Left:} For \(s=0.3\) and \(\beta=4\), SGD and Polyak momentum are compared under the batch size scalings \(B\eqsim D^{s/(s+1)}\) and \(B\eqsim D^{2s/(s+1)}\). Both methods achieve the optimal rate at the smaller batch size, whereas at the larger batch size SGD slows down while Polyak momentum preserves the rate \(D^{-s/(s+1)}\). \textbf{Right:} For \(s=0.8\) and \(\beta=4\), all three algorithms use the Nesterov-optimal batch size scaling \(B\eqsim D^{2s\beta/(s\beta+2\beta-1)}\). Under this scaling, SGD is slower than its data-optimal rate, Polyak momentum retains the rate \(D^{-s/(s+1)}\), and Nesterov momentum achieves the faster rate predicted by our theory.}
%     \label{fig:data scaling}
% \end{figure}




% \section{Proof sketch}
% \paragraph{Signal--noise decomposition.}
% Theorems \ref{thm:hbm scaling law} and \ref{thm:nesterov scaling law} show that the risk dynamics of momentum methods also admit a natural \textbf{signal-learning} and \textbf{noise-forgetting} interpretation:
% \begin{equation}
% \label{eq:momentum-signal-noise-convolution}
% \EE[\cE(\btheta_k)]
% \eqsim
% \underbrace{\cS(k)}_{\text{signal learning}}
% +
% \underbrace{
% \frac{\eta^2}{B}
% \sum_{i=0}^{k-1}
% \mathcal K(k-1-i)
% (\EE[\cE_i]
% +\sigma^2)}_{\text{noise accumulation}},
% \end{equation}
% where
% \[
% \cS(k)
% :=
% \frac12\sum_{j\geq1}\lambda_j
% \bigl(e_{k}^{\mathrm{det}}(j)\bigr)^2,
% \qquad
% \mathcal K(\tau)
% :=
% \sum_{j\geq1}\lambda_j^2G_{\tau}(j)^2.
% \]
% Here $\cS(k)$ denotes the deterministic excess risk, and $\mathcal K(\tau)$ is the forgetting kernel, with $G_{\tau}(j)$ denoting the coordinate-wise Green function. The signal term $\cS(k)$ characterizes how the deterministic dynamics learn the target signal across the eigendirections. The kernel $\mathcal K(\tau)$ measures how much of a noise perturbation injected $\tau$ iterations earlier is retained by the current iterate, so the noise contribution takes the form of a convolution between the historical noise and the forgetting kernel. The regime-dependent scalings of $\cS$ and $\mathcal K$ are established in Appendix~\ref{app:hbm scaling law} and Appendix \ref{app:nesterov scaling law}.

% \paragraph{Fully overdamped regime.}
% When \(\eta\lesssim(1-\rho)^2\), all eigendirections are overdamped. In this regime, both Polyak and Nesterov momentum behave like SGD with effective learning rate \(\eta_{\mathrm{eff}}=\frac{\eta}{1-\rho}\). Their loss therefore has the same signal-learning and noise-accumulation structure as SGD after replacing \(\frac{\eta}{1-\rho}\) by \(\eta_{\mathrm{eff}}\). 
% % Compared with SGD using the same raw learning rate, momentum learns the signal faster but increases label noise by the same factor. 

% \paragraph{Mixed overdamped--underdamped regime.}
% When \(\eta\gtrsim(1-\rho)^2\), the spectrum separates into sharp underdamped directions and flat overdamped directions. Correspondingly, the deterministic signal term separates into an exponentially decaying transient and a polynomially decaying tail. 

% \paragraph{Spectral adaptivity accelerates noise forgetting for Nesterov momentum.}
% Nesterov's look-ahead correction makes the memory of an underdamped direction depend on its curvature: historical noise decays as \(\rho^i\) under Polyak, but as \([\rho(1-\eta\lambda_j)]^i\) under Nesterov. Thus all underdamped Polyak directions forget noise on the same momentum timescale \((1-\rho)^{-1}\), whereas the Nesterov forgetting time is \((1-\rho+\eta\lambda_j)^{-1}\) and becomes shorter as \(\lambda_j\) increases. This difference clips the accumulated additive label noise in each direction:
% \[
% \textnormal{Polyak:}\quad
% \cN_{\polyak,j}\eqsim\frac{\sigma^2}{B}\frac{\eta\lambda_j}{1-\rho},
% \qquad
% \textnormal{Nesterov:}\quad
% \cN_{\nesterov,j}\eqsim\frac{\sigma^2}{B}\frac{\eta\lambda_j}{1-\rho+\eta\lambda_j}.
% \]
% In flat directions with \(\eta\lambda_j\ll1-\rho\), the two contributions are comparable. In sharp directions with \(\eta\lambda_j\gtrsim1-\rho\), however, the Nesterov contribution saturates at order one, while the Polyak contribution continues to grow with curvature. Consequently, summing over the power-law spectrum gives
% \[
% \cN_{\polyak}^{\mathrm{label}}
% \eqsim
% \frac{\sigma^2}{B}\frac{\eta}{1-\rho},
% \qquad
% \cN_{\nesterov}^{\mathrm{label}}
% \eqsim
% \frac{\sigma^2}{B}
% \min\left\{
% \frac{\eta}{1-\rho},
% \left(\frac{\eta}{1-\rho}\right)^{1/\beta}
% \right\}.
% \]


\section{Conclusion and discussion}

We developed a unified analysis of Polyak and Nesterov momentum in large-batch, one-pass training, characterizing their risk stability, risk dynamics, and optimal data scaling. A central conclusion is that the benefit of momentum depends critically on the batch size. At small batch sizes, neither method improves the data-scaling exponent of vanilla SGD. As the batch size grows, however, the two methods provide distinct benefits: Polyak preserves the small-batch optimal data-scaling exponent over a substantially wider range of batch sizes than SGD, whereas Nesterov further improves the data-scaling exponent in the large-batch regime. At ultra-large batch sizes, training becomes limited by the number of sequential updates, and the two momentum methods again exhibit the same signal-limited scaling.

\paragraph{Implications for momentum scheduling.}
Our scaling analysis allows the batch size \(B\), learning rate \(\eta\), and momentum factor \(\rho\) to scale with the data budget \(D\), while keeping them fixed within each training run. A complete theory of time-varying schedules is therefore beyond the present analysis. Nevertheless, our stability results already have direct implications for momentum scheduling. Recent work has increasingly explored time-varying momentum schedules, while also observing that increasing momentum during training can become unstable under stochastic gradients and may require warmup or additional damping~\citep{pagliardini2025ademamix,yarotsky2025sgd,yarotsky2026corner,ferbach2026logarithmic}. Our results provide a complementary quantitative perspective: for the classical Polyak and Nesterov updates, the stability boundary predicts how long an increasing-momentum schedule can remain stable and how this horizon scales with the batch size.

A particularly instructive example is the classical scaling \(1-\rho_k\eqsim k^{-1}\), which underlies Nesterov acceleration in deterministic optimization~\citep{nesterov1983method}. Although Theorem~\ref{thm:stability} is derived for constant momentum, applying its stability boundary locally along this slowly varying schedule predicts
\begin{equation}\label{eqn: stable-schedule}
\eta\lesssim\min\left\{1,\frac{B}{k}\right\}\quad\text{for Polyak},\qquad \eta\lesssim\min\left\{1,\frac{B^\beta}{k}\right\}\quad\text{for Nesterov}.
\end{equation}
Thus, for fixed \(\eta\) and \(B\), the schedule may remain stable for a substantial fraction of training but eventually cross the stochastic stability boundary as \(\rho_k\to1\). The corresponding instability horizons are predicted to scale as
\[
k_{\polyak}^*\eqsim\frac{B}{\eta},\qquad k_{\nesterov}^*\eqsim\frac{B^\beta}{\eta}.
\]
Prior work has observed that increasing the batch size can delay the instability of aggressive momentum schedules~\citep{yarotsky2025sgd}; to our knowledge, the explicit batch-size scaling of this delayed-instability horizon, and its distinction between  Polyak and Nesterov, has not been characterized previously.

\begin{wrapfigure}{r}{0.36\textwidth}
\centering
\vspace{-1.5em}
{\hspace{2.3em}\footnotesize \textbf{Fixed $\eta$, $\rho_k=1-1/(k+1)$}}\\[-.1em]
\includegraphics[width=\linewidth]{figures/blowup_time.pdf}
\vspace{-3.5em}
\end{wrapfigure}
The right panel illustrates this delayed-instability phenomenon for PLK regression with \((s,\beta)=(0.3,4)\) under the schedule \(\rho_k=1-1/(k+1)\); experimental details are provided in Appendix~\ref{app:subsec:blow-up}. Both methods remain stable and benefit from the increasing momentum before eventually becoming unstable. Consistent with the predicted horizons, Nesterov remains stable for substantially more iterations than Polyak because its stability horizon scales as \(B^\beta\) rather than \(B\). This allows Nesterov to exploit the increasing-momentum schedule over a longer training interval and reach a lower risk before instability occurs.

This example highlights a constraint on momentum scheduling that is absent from deterministic acceleration theory: the benefit of increasing momentum depends not only on how strongly it accelerates optimization, but also on how long the resulting trajectory remains inside the stochastic stability region. It also suggests a simple tuning principle. Momentum, learning rate, and batch size should be scheduled jointly so that \(1-\rho_k\) remains above the stochastic stability scale, namely \(1-\rho_k\gtrsim\eta_k/B_k\) for Polyak and \(1-\rho_k\gtrsim\eta_k/B_k^\beta\) for Nesterov. Increasing momentum can therefore be supported either by increasing the batch size or by decreasing the learning rate, suggesting a direct route from the present stability theory to practical momentum-schedule design.

\paragraph{Outlook.}
Our results suggest that learning rate, batch size, and momentum should be designed jointly rather than tuned in isolation. The stability conditions determine how aggressively momentum can be increased at a given learning rate and batch size, while the scaling laws quantify the resulting tradeoff between faster signal learning and stronger noise accumulation. Extending this framework to time-varying \(\eta_k\), \(B_k\), and \(\rho_k\) could therefore provide principled guidance for jointly scaling and scheduling these hyperparameters in large-scale training~\citep{li2026functional,li2026optimal,wang2025sharpness,wang2026fast}. We expect such a theory to provide a practical route toward designing more stable and data-efficient training schedules, rather than relying primarily on empirical tuning.

\section*{Use of AI}
We used AI tools for drafting and language editing throughout the manuscript. In Sections B--F, they were also used to accelerate algebraic calculations and asymptotic estimates. The proof strategies and mathematical arguments were developed by the authors, who verified all AI-assisted calculations and estimates.

\section*{Acknowledgement}
Lei Wu is supported in part by the National Natural Science Foundation of China (NSFC12522120, NSFC92470122, and NSFC12288101). Kairui Li is partially supported by the elite undergraduate training program of School of Mathematical Sciences in Peking University. We thank  Sanxing Cao, Zilin Wang, and Jinbo Wang for helpful discussions. 

\bibliography{ref.bib}
\bibliographystyle{bibstyle}


\newpage


\appendix
\part{Appendix}
\parttoc

\newpage
\section{Technical preliminaries}
\subsection{Equivalence to kernel regression}

We briefly relate the PLK regression in Section~\ref{sec:plkr} to standard kernel regression. Given a positive semidefinite kernel \(K:\cX\times\cX\to\RR\), kernel regression learns over its associated reproducing kernel Hilbert space (RKHS) \(\cH_K\).  Given an input distribution $\cD$, under mild conditions, Mercer's theorem ensures that the kernel admits the following eigendecomposition
\[
K(\bx,\bx')
=
\sum_{j=1}^\infty \lambda_j e_j(\bx) e_j(\bx'),
\]
where \(\{e_j\}_{j\geq1}\) are orthonormal in \(L^2(\cD)\). The standard capacity condition assumes \(\lambda_j\eqsim j^{-\beta}\). Target regularity is commonly described through the interpolation spaces
\[
\cH_K^r
=
\left\{
\sum_{j=1}^\infty a_j e_j:
\sum_{j=1}^\infty \frac{a_j^2}{\lambda_j^r}<\infty
\right\},
\]
with a source condition of the form \(f^*\in\cH_K^r\).

Our formulation is the corresponding feature-map representation. Under Assumption~\ref{ass:diagonal-covariance}, define \(e_j:=\phi_j/\lambda_j^{1/2}\). Then \(\{e_j\}_{j\geq1}\) is orthonormal in \(L^2(\cD)\), and
\[
K_\phi(\bx,\bx')
=
\sum_{j=1}^\infty \phi_j(\bx)\phi_j(\bx')
=
\sum_{j=1}^\infty \lambda_j e_j(\bx)e_j(\bx').
\]
Thus, Assumption~\ref{ass:power-law}, \(\lambda_j\eqsim j^{-\beta}\), is precisely the standard capacity condition.

For the target function,
\[
f^*
=
\sum_{j=1}^\infty \theta_j^*\phi_j
=
\sum_{j=1}^\infty a_j^*e_j,
\qquad
a_j^*:=\lambda_j^{1/2}\theta_j^*.
\]
Our task-difficulty assumption gives \(|a_j^*|^2\eqsim j^{-1}\lambda_j^s\). Hence, for any \(\delta\in(0,s)\),
\[
\sum_{j=1}^\infty
\frac{|a_j^*|^2}{\lambda_j^{s-\delta}}
\eqsim
\sum_{j=1}^\infty j^{-1}\lambda_j^\delta
\eqsim
\sum_{j=1}^\infty j^{-1-\beta\delta}
<\infty.
\]
Therefore, \(f^*\in\cH_{K_\phi}^{s-\delta}\) for every \(\delta\in(0,s)\). Thus, \(\beta\) corresponds to the standard capacity exponent, while \(s\) characterizes the source regularity of the target.

\subsection{Rotational equivariance of SGD and momentum methods}
\label{app:subsec:rotation invariant}

We show that Assumption~\ref{ass:diagonal-covariance} is without loss of generality. Since \(\bH=\EE[\bphi(\bx)\otimes\bphi(\bx)]\) is positive, self-adjoint, and trace class, it admits a spectral decomposition
\[
\bH=\bm U\bm\Lambda\bm U^*,
\qquad
\bm\Lambda=\diag(\lambda_1,\lambda_2,\ldots),
\]
for some unitary operator \(\bm U\). Define the rotated variables
\[
\btheta'=\bm U^*\btheta,\qquad
\btheta^{*\prime}=\bm U^*\btheta^*,\qquad
\bphi'(\bx)=\bm U^*\bphi(\bx).
\]
Unitary invariance of the inner product gives
\[
f'(\bx;\btheta')=\langle\bphi'(\bx),\btheta'\rangle
=\langle\bphi(\bx),\btheta\rangle=f(\bx;\btheta),
\]
and hence, for every mini-batch \(\cB_k\),
\[
\bg'(\btheta';\cB_k)=\bm U^*\bg(\btheta;\cB_k).
\]
It follows immediately that the SGD and Polyak recursions are preserved under \(\btheta_k'=\bm U^*\btheta_k\). For Nesterov, the look-ahead point transforms as
\[
\btheta_k^{\la\prime}
=\bm U^*\btheta_k^{\la}
=\btheta_k'+\rho(\btheta_k'-\btheta_{k-1}'),
\]
so its recursion is preserved as well. Finally,
\[
\bH'=\EE[\bphi'(\bx)\otimes\bphi'(\bx)]
=\bm U^*\bH\bm U
=\bm\Lambda.
\]
Thus, SGD, Polyak, and Nesterov are equivariant under unitary changes of basis, while their predictions and risks are unchanged. We may therefore work in the eigenbasis of \(\bH\), where the covariance operator is diagonal, without loss of generality.

\subsection{Gaussian fourth-moment identities}
\label{app:subsec:gaussian-fourth-moment}

We record several Gaussian fourth-moment identities used repeatedly below.

\begin{lemma}[Gaussian fourth-moment identities]
\label{app:lem:gaussian-fourth-moment}
Let \(\bx\sim\mathcal N(\bm0,\bH)\) in a separable Hilbert space \(\mathcal H\), where \(\bH\) is a positive trace-class covariance operator. For any \(\bu,\bv\in\mathcal H\),
\begin{equation}
\label{app:eq:gaussian-fourth-moment-directional}
\EE\!\left[\langle\bx,\bu\rangle^2\langle\bx,\bv\rangle^2\right]
=
\langle\bu,\bH\bu\rangle\langle\bv,\bH\bv\rangle
+
2\langle\bu,\bH\bv\rangle^2.
\end{equation}
Equivalently, for any self-adjoint trace-class operator \(\bA\),
\begin{equation}
\label{app:eq:gaussian-fourth-moment}
\EE\!\left[(\bx\otimes\bx)\bA(\bx\otimes\bx)\right]
=
2\bH\bA\bH+\bH\Tr(\bH\bA).
\end{equation}
% 
\end{lemma}

\begin{proof}
The directional identity follows directly from Isserlis' formula. For any \(\bu,\bv\in\mathcal H\), the random variables \(\langle\bx,\bu\rangle\) and \(\langle\bx,\bv\rangle\) are jointly Gaussian, and hence
\[
\EE\!\left[\langle\bx,\bu\rangle^2\langle\bx,\bv\rangle^2\right]
=
\langle\bu,\bH\bu\rangle\langle\bv,\bH\bv\rangle
+
2\langle\bu,\bH\bv\rangle^2.
\]

We next prove the operator identity. Let \(\bA=\sum_j a_j\be_j\otimes\be_j\) be the spectral decomposition of the self-adjoint trace-class operator \(\bA\). For arbitrary \(\bu,\bv\in\mathcal H\), Isserlis' formula gives
\begin{align*}
&\left\langle
\bu,
\EE\!\left[(\bx\otimes\bx)\bA(\bx\otimes\bx)\right]
\bv
\right\rangle\\
&\qquad=
\sum_j a_j
\EE\!\left[
\langle\bu,\bx\rangle
\langle\be_j,\bx\rangle^2
\langle\bx,\bv\rangle
\right]\\
&\qquad=
2\sum_j a_j
\langle\bu,\bH\be_j\rangle
\langle\be_j,\bH\bv\rangle
+
\langle\bu,\bH\bv\rangle
\sum_j a_j\langle\be_j,\bH\be_j\rangle\\
&\qquad=
2\langle\bu,\bH\bA\bH\bv\rangle
+
\langle\bu,\bH\bv\rangle\Tr(\bH\bA).
\end{align*}
Since this holds for all \(\bu,\bv\in\mathcal H\),
\[
\EE\!\left[(\bx\otimes\bx)\bA(\bx\otimes\bx)\right]
=
2\bH\bA\bH+\bH\Tr(\bH\bA).
\]
\end{proof}


\begin{corollary}[Empirical covariance fluctuation]
\label{app:cor:gaussian-empirical-covariance}
Let
\[
\widehat{\bH}:=\frac1B\sum_{b=1}^B\bx_b\otimes\bx_b,
\qquad
\bSigma:=\widehat{\bH}-\bH,
\qquad
\bx_b\overset{\mathrm{i.i.d.}}{\sim}\mathcal N(\bm0,\bH).
\]
Then, for any self-adjoint trace-class operator \(\bA\),
\begin{equation}
\label{app:eq:gaussian-empirical-covariance}
\EE[\bSigma\bA\bSigma]
=
\frac1B\left(\bH\bA\bH+\bH\Tr(\bH\bA)\right).
\end{equation}
\end{corollary}

\begin{proof}
Write \(\bSigma=B^{-1}\sum_{b=1}^B\bZ_b\), where \(\bZ_b:=\bx_b\otimes\bx_b-\bH\) are independent and centered. Hence
\[
\EE[\bSigma\bA\bSigma]
=
\frac1B\EE[\bZ\bA\bZ]
=
\frac1B\left(\EE[(\bx\otimes\bx)\bA(\bx\otimes\bx)]-\bH\bA\bH\right),
\]
and the result follows from Lemma~\ref{app:lem:gaussian-fourth-moment}.
\end{proof}



\subsection{Noise--curvature alignment}
\label{app:hessian alignment}

We prove Proposition~\ref{prop:noise-geometry}. For a mini-batch of size \(B\), write
$
\bxi_k(\btheta)=\frac1B\sum_{i=1}^B\bzeta_{k,i}(\btheta),
$
where \(\bzeta_{k,i}(\btheta)\) are conditionally independent, centered single-sample gradient noises. Hence, for any unit vector \(\bnu\in\mathcal H\),
\[
\EE\!\left[\left.|\langle\bxi_k(\btheta),\bnu\rangle|^2\right|\btheta\right]
=
\frac1B\EE\!\left[\left.|\langle\bzeta(\btheta),\bnu\rangle|^2\right|\btheta\right].
\]

Let \(\bdelta:=\btheta-\btheta^*\) and \(\bphi:=\bphi(\bx)\). Since \(y=\langle\bphi,\btheta^*\rangle+\epsilon\), the single-sample gradient is
$
\bg_{\bx,y}(\btheta)
=
\bigl(\langle\bphi,\bdelta\rangle-\epsilon\bigr)\bphi,
$
while \(\nabla\cR(\btheta)=\bH\bdelta\). Therefore,
\[
\bzeta(\btheta)
=
\langle\bphi,\bdelta\rangle\bphi-\bH\bdelta-\epsilon\bphi.
\]
Using the independence of \(\epsilon\) and \(\bphi\), \(\EE[\epsilon]=0\), and $\EE[\bphi\bphi]=\bH$, we obtain
\[
\EE\!\left[\left.|\langle\bzeta(\btheta),\bnu\rangle|^2\right|\btheta\right]
=
\EE\!\left[\langle\bphi,\bdelta\rangle^2\langle\bphi,\bnu\rangle^2\right]
-\langle\bdelta,\bH\bnu\rangle^2
+\sigma^2\langle\bnu,\bH\bnu\rangle.
\]
Applying Lemma~\ref{app:lem:gaussian-fourth-moment}  gives
\[
\EE\!\left[\langle\bphi,\bdelta\rangle^2\langle\bphi,\bnu\rangle^2\right]
=
\langle\bdelta,\bH\bdelta\rangle\langle\bnu,\bH\bnu\rangle
+
2\langle\bdelta,\bH\bnu\rangle^2.
\]
Thus,
\begin{equation}
\label{app:eq:noise-direction-exact}
\EE\!\left[\left.|\langle\bzeta(\btheta),\bnu\rangle|^2\right|\btheta\right]
=
\bigl(\sigma^2+\langle\bdelta,\bH\bdelta\rangle\bigr)\langle\bnu,\bH\bnu\rangle
+
\langle\bdelta,\bH\bnu\rangle^2.
\end{equation}
The last term is nonnegative and, by Cauchy--Schwarz in the \(\bH\)-inner product,
\[
\langle\bdelta,\bH\bnu\rangle^2
\leq
\langle\bdelta,\bH\bdelta\rangle
\langle\bnu,\bH\bnu\rangle.
\]
Since \(\cE(\btheta)=\frac12\langle\bdelta,\bH\bdelta\rangle\), Eq.~\eqref{app:eq:noise-direction-exact} implies
\[
\EE\!\left[\left.|\langle\bzeta(\btheta),\bnu\rangle|^2\right|\btheta\right]
\eqsim
\bigl(\cE(\btheta)+\sigma^2\bigr)\langle\bnu,\bH\bnu\rangle.
\]
Combining this with the \(1/B\) variance reduction from mini-batching and setting \(\btheta=\btheta_k\) yields
\[
\EE\!\left[\left.|\langle\bxi_k(\btheta_k),\bnu\rangle|^2\right|\btheta_k\right]
\eqsim
\frac1B\bigl(\cE(\btheta_k)+\sigma^2\bigr)\langle\bnu,\bH\bnu\rangle,
\]
which proves Proposition~\ref{prop:noise-geometry}.



\section{Risk stability analysis}
\label{app:sec:stability}

\subsection{Renewal formulation of risk stability}
\label{app:subsec:label-noise-stability}

We first recall a general boundedness criterion for renewal equations~\citep{feller1991introduction}. Consider a nonnegative sequence \((x_k)_{k\geq0}\) satisfying
\[
x_k=a_k+\sum_{j=0}^{k-1}K_{k-1-j}x_j,
\]
where \(a_k\geq0\) is the input and \(K=(K_k)_{k\geq0}\) is the renewal kernel. Its total mass \(\|K\|_1:=\sum_{k\geq0}K_k\) measures the overall feedback strength of the recursion.

\begin{lemma}[Boundedness of a renewal equation]
\label{app:lem:renewal-stability}
Suppose \(0<\underline a:=\inf_{k\geq0}a_k\leq\sup_{k\geq0}a_k=:\bar a<\infty\). Then
\[
\sup_{k\geq0}x_k<\infty\quad\Longleftrightarrow\quad \|K\|_1<1.
\]
\end{lemma}

\begin{proof}
Suppose first that \(\|K\|_1<1\), and let \(M_k:=\max_{0\leq j\leq k}x_j\). Then \(x_k\leq\bar a+\|K\|_1M_{k-1}\), which implies
\[
\sup_{k\geq0}x_k\leq\max\left\{x_0,\frac{\bar a}{1-\|K\|_1}\right\}<\infty.
\]
Conversely, suppose that \((x_k)\) is bounded and let \(L:=\liminf_{k\to\infty}x_k\). For every fixed \(N\), nonnegativity gives
\[
x_k\geq\underline a+\sum_{m=0}^{N}K_mx_{k-1-m}.
\]
Taking \(\liminf_{k\to\infty}\) yields \(L\geq\underline a+L\sum_{m=0}^{N}K_m\). Letting \(N\to\infty\), boundedness is impossible if \(\sum_{m\geq0}K_m\geq1\).
\end{proof}

For each method \(\cA\in\{\sgd,\polyak,\nesterov\}\), we will derive a nonnegative risk moment \(\ell_k^{(\cA)}\), equal to \(2\EE[\cE(\btheta_k)]\) for SGD and Polyak and comparable to it for Nesterov, satisfying
\begin{equation}
\label{app:eq:generic-risk-renewal}
\ell_k^{(\cA)}=h_k^{(\cA)}+\sum_{j=0}^{k-1}K_{k-1-j}^{(\cA)}\bigl(\ell_j^{(\cA)}+\sigma^2\bigr).
\end{equation}
Setting \(x_k=\ell_k^{(\cA)}+\sigma^2\) and \(a_k=h_k^{(\cA)}+\sigma^2\) reduces~\eqref{app:eq:generic-risk-renewal} to the renewal equation above. Hence, once the coordinate-wise dynamics are stable and \(h_k^{(\cA)}\) is bounded, Lemma~\ref{app:lem:renewal-stability} gives the global feedback condition
\[
\kappa_{\cA}(\eta,\rho,B):=\|K^{(\cA)}\|_1=\sum_{k\geq0}K_k^{(\cA)}<1.
\]
The kernel \(K^{(\cA)}\) describes how perturbations propagate through the optimizer dynamics and depends on the method and hyperparameters \((\eta,\rho,B)\), but not on the label-noise variance \(\sigma^2\). Additive label noise enters as the constant external forcing \(\sigma^2\), whereas multiplicative noise feeds the current risk back through \(K^{(\cA)}\). The following subsections derive \(K^{(\cA)}\), compute \(\kappa_{\cA}\), and combine \(\kappa_{\cA}<1\) with the corresponding coordinate-wise stability condition to obtain the critical learning rates.
\subsection{Risk stability of SGD}
\label{app:subsec:sgd-stability}

The SGD error recursion, including label noise, is
\begin{equation}
\label{app:eq:sgd-stability-recursion}
\be_{k+1}
=
(\bI-\eta\widehat{\bH}_k)\be_k
+
\eta\widehat{\boldsymbol\zeta}_k.
\end{equation}

\begin{theorem}[Risk stability of SGD]
\label{app:thm:sgd-stability}
Let the Hessian be a positive trace-class operator with spectrum
\[
\bH=\diag(\lambda_1,\lambda_2,\ldots),
\qquad
\lambda_1=\lambda_{\max}\geq\lambda_2\geq\cdots>0,
\]
and suppose
\(\sigma>0.\)
Then SGD is risk stable if and only if
\begin{equation}
\label{app:eq:sgd-local-stability}
\eta<\frac{2B}{(B+1)\lambda_{\max}}
\end{equation}
and
\begin{equation}
\label{app:eq:sgd-global-stability}
\kappa_{\sgd}(\eta)
:=
\sum_{i\geq1}
\frac{\eta\lambda_i}{2B-(B+1)\eta\lambda_i}
<1.
\end{equation}
\end{theorem}

\begin{proof}
Define the covariance operator by
\[
\bC_k:=\EE[\be_k\otimes\be_k].
\]
We work with its spectral coordinates. For every finite-risk initialization, the relevant risk moment is
\[
\ell_k
:=
\Tr(\bH\bC_k)
=
\sum_{i\geq1}\lambda_i
\EE\bigl[|\langle\mathbf e_i,\be_k\rangle|^2\bigr]
=
2\EE[\cE(\btheta_k)].
\]
In particular, \(\ell_0<\infty\). The covariance identities below may be justified first on finite spectral
truncations and then passed to the limit by monotone convergence. Expanding
\eqref{app:eq:sgd-stability-recursion}, using the vanishing cross terms, and applying
Lemma~\ref{app:lem:gaussian-fourth-moment} gives
\begin{equation}
\label{app:eq:sgd-covariance}
\bC_{k+1}
=
(\bI-\eta\bH)\bC_k(\bI-\eta\bH)
+
\frac{\eta^2}{B}
\left(
\bH\bC_k\bH
+
\bH\Tr(\bH\bC_k)
+
\sigma^2\bH
\right).
\end{equation}
Define the coordinate variance by
\[
u_{i,k}:=\langle\mathbf e_i,\bC_k\mathbf e_i\rangle.
\]
Taking the diagonal entries of \eqref{app:eq:sgd-covariance} yields
\begin{equation}
\label{app:eq:sgd-diagonal-recursion}
u_{i,k+1}
=
a_i u_{i,k}
+
\frac{\eta^2}{B}\lambda_i\bigl(\ell_k+\sigma^2\bigr),
\qquad
a_i:=(1-\eta\lambda_i)^2+\frac{\eta^2}{B}\lambda_i^2.
\end{equation}

We first determine the local condition. Since
\[
1-a_i
=
\eta\lambda_i
\left[
2-\left(1+\frac1B\right)\eta\lambda_i
\right],
\]
the coordinatewise contraction condition
\[
a_i<1
\qquad\text{for every }i
\]
is equivalent to
\eqref{app:eq:sgd-local-stability}. If this condition fails, then
\eqref{app:eq:sgd-diagonal-recursion} contains an eigendirection satisfying
\[
u_{i,k+1}
\geq
u_{i,k}+\frac{\eta^2\sigma^2}{B}\lambda_i,
\]
and therefore the risk is unbounded. Hence the local condition is necessary.

Assume henceforth that \eqref{app:eq:sgd-local-stability} holds. Iterating
\eqref{app:eq:sgd-diagonal-recursion}, multiplying the \(i\)-th equation by \(\lambda_i\), and summing
over all eigendirections gives
\begin{equation}
\label{app:eq:sgd-renewal}
\ell_k
=
h_k
+
\sum_{s=0}^{k-1}
K_{k-1-s}\bigl(\ell_s+\sigma^2\bigr),
\end{equation}
where
\[
h_k:=\sum_{i\geq1}\lambda_i a_i^k u_{i,0},
\qquad
K_m:=\frac{\eta^2}{B}\sum_{i\geq1}\lambda_i^2a_i^m.
\]
The local condition gives
\[
0\leq a_i<1,
\qquad
0\leq h_k\leq \ell_0.
\]
Moreover,
\[
\sum_{m\geq0}K_m
=
\frac{\eta^2}{B}
\sum_{i\geq1}\frac{\lambda_i^2}{1-a_i}
=
\sum_{i\geq1}
\frac{\eta\lambda_i}{2B-(B+1)\eta\lambda_i}
=
\kappa_{\sgd}(\eta).
\]
Lemma~\ref{app:lem:renewal-stability} therefore gives
\[
\sup_{k\geq0}\ell_k<\infty
\quad\Longleftrightarrow\quad
\kappa_{\sgd}(\eta)<1,
\]
which completes the proof.

\end{proof}

\begin{corollary}[Critical learning rate of SGD]
\label{app:cor:sgd-critical-rate}
There exists a unique critical learning rate satisfying
\[
0<\eta_{\sgd}^{\crit}<\frac{2B}{(B+1)\lambda_{\max}}
\]
and
\begin{equation}
\label{app:eq:sgd-critical-equation}
\sum_{i\geq1}
\frac{
\eta_{\sgd}^{\crit}\lambda_i
}{
2B-(B+1)\eta_{\sgd}^{\crit}\lambda_i
}
=1.
\end{equation}
The SGD iteration is risk stable if and only if
\[
0<\eta<\eta_{\sgd}^{\crit}.
\]
Moreover,
\begin{equation}
\label{app:eq:sgd-critical-bounds}
\frac{2B}{(B+1)\lambda_{\max}+\Tr(\bH)}
\leq
\eta_{\sgd}^{\crit}
\leq
\min\left\{
\frac{2B}{(B+1)\lambda_{\max}},
\frac{2B}{\Tr(\bH)}
\right\}.
\end{equation}
Consequently, under the fixed-exponent power-law assumptions
\(
\lambda_i\eqsim i^{-\beta},
\beta>1
\), we have
\[
\eta_{\sgd}^{\crit}\eqsim1.
\]
\end{corollary}

\begin{proof}
The renewal mass \(\kappa_{\sgd}(\eta)\) is continuous and strictly increasing on the local stability interval. It vanishes at \(\eta=0\) and diverges as \(\eta\) approaches the local boundary, because the denominator of the leading eigendirection tends to zero. Hence there is a unique solution of \(\kappa_{\sgd}(\eta)=1\), which gives the stated critical learning rate and stability interval.

At the critical point, every denominator in \eqref{app:eq:sgd-critical-equation} is at most \(2B\). Therefore
\(1\geq\frac{\eta_{\sgd}^{\crit}}{2B}\Tr(\bH),\)
which yields
\[
\eta_{\sgd}^{\crit}
\leq
\frac{2B}{\Tr(\bH)}.
\]
The local condition gives the other upper bound in \eqref{app:eq:sgd-critical-bounds}. Conversely, every denominator is at least
\[
2B-(B+1)\eta_{\sgd}^{\crit}\lambda_{\max}.
\]
Using the critical equation then gives
\[
1
\leq
\frac{
\eta_{\sgd}^{\crit}\Tr(\bH)
}{
2B-(B+1)\eta_{\sgd}^{\crit}\lambda_{\max}
},
\]
and rearranging yields the lower bound in \eqref{app:eq:sgd-critical-bounds}.

Finally, under \(\lambda_i\eqsim i^{-\beta}\) with \(\beta>1\), we have \(\lambda_{\max}\eqsim1\) and \(\Tr(\bH)\eqsim1\). The two-sided bounds therefore imply \(\eta_{\sgd}^{\crit}\eqsim1\).
\end{proof}



\subsection{Risk stability of Polyak}
\label{app:subsec:polyak-stability}

The Polyak error recursion, including label noise, is
\begin{equation}
\label{app:eq:polyak-stability-recursion}
\be_{k+1}
=
(\bI-\eta\widehat{\bH}_k)\be_k
+
\rho(\be_k-\be_{k-1})
+
\eta\widehat{\boldsymbol\zeta}_k.
\end{equation}
Writing
\(\bD=(1+\rho)\bI-\eta\bH,\)
Equation \eqref{app:eq:polyak-stability-recursion} becomes
\[
\be_{k+1}
=
\bD\be_k
-
\rho\be_{k-1}
-
\eta({\bH}_k-\bH)\be_k
+
\eta\widehat{\boldsymbol\zeta}_k.
\]

\begin{theorem}[Risk stability of Polyak]
\label{app:thm:polyak-stability}
Let the Hessian be a positive trace-class operator and assume
\[
\bH=\diag(\lambda_1,\lambda_2,\ldots),
\qquad
\lambda_1=\lambda_{\max}\geq\lambda_2\geq\cdots>0,
\qquad
0\leq\rho<1,
\qquad
\sigma>0.
\]
Then Polyak momentum is risk stable if and only if
\begin{equation}
\label{app:eq:polyak-local-stability}
\eta
<
\frac{
2B(1-\rho^2)
}{
\lambda_{\max}\left[B(1-\rho)+(1+\rho)\right]
}
\end{equation}
and
\begin{equation}
\label{app:eq:polyak-global-stability}
\kappa_{\polyak}(\eta,\rho)
:=
\sum_{i\geq1}
\frac{
\eta\lambda_i(1+\rho)
}{
2B(1-\rho^2)
-
\eta\lambda_i\left[B(1-\rho)+(1+\rho)\right]
}
<1.
\end{equation}
\end{theorem}

\begin{proof}
Define the coordinate second moments
\[
\begin{aligned}
\bU_k&:=\EE[\be_k\otimes\be_k],
&
\bP_k&:=\EE[\be_k\otimes\be_{k-1}],
\\
\bQ_k&:=\EE[\be_{k-1}\otimes\be_k],
&
\bR_k&:=\EE[\be_{k-1}\otimes\be_{k-1}].
\end{aligned}
\]
The risk moment relevant to risk stability is
\[
\ell_k
:=
\Tr(\bH\bU_k)
=
2\EE[\cE(\btheta_k)].
\]
Expanding \eqref{app:eq:polyak-stability-recursion}, using the vanishing cross terms, and
applying Lemma~\ref{app:lem:gaussian-fourth-moment} gives
\begin{equation}
\label{app:eq:polyak-covariance-recursion}
\left\{
\begin{aligned}
\bU_{k+1}
={}&
\bD\bU_k\bD
-
\rho\bD\bP_k
-
\rho\bQ_k\bD
+
\rho^2\bR_k
\\
&+
\frac{\eta^2}{B}
\left(
\bH\bU_k\bH
+
\bH\Tr(\bH\bU_k)
+
\sigma^2\bH
\right),
\\
\bP_{k+1}
={}&
\bD\bU_k-\rho\bQ_k,
\\
\bQ_{k+1}
={}&
\bU_k\bD-\rho\bP_k,
\\
\bR_{k+1}
={}&
\bU_k.
\end{aligned}
\right.
\end{equation}
In the eigenbasis of the Hessian, define
\[
u_{i,k}:=\langle\mathbf e_i,\bU_k\mathbf e_i\rangle,
\qquad
v_{i,k}:=\langle\mathbf e_i,\bP_k\mathbf e_i\rangle,
\qquad
r_{i,k}:=\langle\mathbf e_i,\bR_k\mathbf e_i\rangle.
\]
The diagonal entries of the two cross-covariance operators coincide because
\(
\bQ_k=\bP_k^*.
\)
Set
\[
d_i:=1+\rho-\eta\lambda_i,
\qquad
\bX_{i,k}:=(u_{i,k},v_{i,k},r_{i,k})^\top,
\qquad
\bb:=(1,0,0)^\top,
\]
and
\[
\bM_i
:=
\begin{pmatrix}
d_i^2+\dfrac{\eta^2}{B}\lambda_i^2
&
-2\rho d_i
&
\rho^2
\\
d_i
&
-\rho
&
0
\\
1
&
0
&
0
\end{pmatrix}.
\]
Taking the diagonal entries of \eqref{app:eq:polyak-covariance-recursion} yields
\begin{equation}
\label{app:eq:polyak-block-recursion}
\bX_{i,k+1}
=
\bM_i\bX_{i,k}
+
\frac{\eta^2}{B}\lambda_i\bb\bigl(\ell_k+\sigma^2\bigr).
\end{equation}

After removing the global risk-coupling term, the characteristic polynomial of the isolated block
is
\[
p_i(\mu)
=
\mu^3
+
\left(
\rho-d_i^2-\frac{\eta^2}{B}\lambda_i^2
\right)\mu^2
+
\rho
\left(
d_i^2-\frac{\eta^2}{B}\lambda_i^2-\rho
\right)\mu
-
\rho^3.
\]
The Jury stability criterion shows that all roots lie strictly inside the unit disk if and only if
\[
2B(1-\rho^2)
-
\eta\lambda_i\left[B(1-\rho)+(1+\rho)\right]
>0.
\]
Requiring this inequality in every eigendirection gives
\eqref{app:eq:polyak-local-stability}. If the local condition fails, the isolated second-moment
system has a mode on or outside the unit circle.
Because the label noise forcing enters its current-error coordinate directly and has strictly
positive variance, the corresponding coordinate risk is unbounded. Thus, the local condition is
necessary.

Assume henceforth that the local condition holds, and define
\(g_{i,m}:=\bb^\top\bM_i^m\bb.\)
This quantity is nonnegative. Iterating
\eqref{app:eq:polyak-block-recursion}, taking the first component, multiplying by the
corresponding eigenvalue, and summing over all eigendirections gives
\begin{equation}
\label{app:eq:polyak-renewal}
\ell_k
=
h_k
+
\sum_{s=0}^{k-1}
K_{k-1-s}\bigl(\ell_s+\sigma^2\bigr),
\end{equation}
where
\[
h_k
:=
\sum_{i\geq1}
\lambda_i\bb^\top\bM_i^k\bX_{i,0},
\qquad
K_m
:=
\frac{\eta^2}{B}
\sum_{i\geq1}\lambda_i^2g_{i,m}.
\]
The strict local condition gives the uniform power bound
\(\sup_{i,k}\|\bM_i^k\|<\infty.\)
Positivity of the initial coordinate covariance also gives
\(|v_{i,0}|\leq\frac{u_{i,0}+r_{i,0}}{2}.\)
Together, these estimates imply
\[
0\leq h_k
\leq
C\sum_{i\geq1}\lambda_i(u_{i,0}+r_{i,0})
<\infty.
\]
A direct calculation gives
\begin{equation}
\label{app:eq:polyak-resolvent}
\bb^\top(\bI-\bM_i)^{-1}
=
\frac{
B\bigl(1+\rho,-2\rho d_i,\rho^2(1+\rho)\bigr)
}{
\eta\lambda_i
\left\{
2B(1-\rho^2)
-
\eta\lambda_i\left[B(1-\rho)+(1+\rho)\right]
\right\}
}.
\end{equation}
Consequently,
\[
\sum_{m\geq0}K_m
=
\frac{\eta^2}{B}
\sum_{i\geq1}
\lambda_i^2\bb^\top(\bI-\bM_i)^{-1}\bb
=
\kappa_{\polyak}(\eta,\rho).
\]
Lemma~\ref{app:lem:renewal-stability} therefore gives
\[
\sup_{k\geq0}\ell_k<\infty
\quad\Longleftrightarrow\quad
\kappa_{\polyak}(\eta,\rho)<1,
\]
which completes the proof.

\end{proof}

\begin{corollary}[Critical learning rate of Polyak]
\label{app:cor:polyak-critical-rate}
There exists a unique critical learning rate satisfying
\[
0
<
\eta_{\polyak}^{\crit}
<
\frac{
2B(1-\rho^2)
}{
\lambda_{\max}\left[B(1-\rho)+(1+\rho)\right]
}
\]
and the critical equation
\begin{equation}
\label{app:eq:polyak-critical-equation}
\sum_{i\geq1}
\frac{
\eta_{\polyak}^{\crit}\lambda_i(1+\rho)
}{
2B(1-\rho^2)
-
\eta_{\polyak}^{\crit}\lambda_i
\left[B(1-\rho)+(1+\rho)\right]
}
=1.
\end{equation}
The Polyak iteration is risk stable if and only if
\[
0<\eta<\eta_{\polyak}^{\crit}.
\]
Moreover,
\begin{equation}
\label{app:eq:polyak-critical-bounds}
\begin{aligned}
\frac{
2B(1-\rho^2)
}{
\lambda_{\max}\left[B(1-\rho)+(1+\rho)\right]
+
(1+\rho)\Tr(\bH)
}
\leq
\eta_{\polyak}^{\crit},
\\
\eta_{\polyak}^{\crit}
\leq
\min\left\{
\frac{
2B(1-\rho^2)
}{
\lambda_{\max}\left[B(1-\rho)+(1+\rho)\right]
},
\frac{2B(1-\rho)}{\Tr(\bH)}
\right\}.
\end{aligned}
\end{equation}
Under the power-law assumption
\(\lambda_i\eqsim i^{-\beta},\beta>1
\), when \(B\gg1,1-\rho\ll1,\) we have
\[
\eta_{\polyak}^{\crit}
\eqsim
\min\{1,B(1-\rho)\}.
\]
\end{corollary}

\begin{proof}
The renewal mass \(\kappa_{\polyak}(\eta,\rho)\) is continuous and strictly increasing throughout the local stability interval. It
vanishes at the origin and diverges at the local boundary. Theorem~\ref{app:thm:polyak-stability}
therefore gives a unique critical learning rate and the stated stability interval.

At the critical point, every denominator in
\eqref{app:eq:polyak-critical-equation} is at most
\(2B(1-\rho^2)\)
and hence
\[
\eta_{\polyak}^{\crit}
\leq
\frac{2B(1-\rho)}{\Tr(\bH)},
\]
while the local condition gives the other upper bound in
\eqref{app:eq:polyak-critical-bounds}. Conversely, every denominator is at least
\(
2B(1-\rho^2)
-
\eta_{\polyak}^{\crit}\lambda_{\max}
\left[B(1-\rho)+(1+\rho)\right].
\)
Substitution into the critical equation gives
\[
1
\leq
\frac{
\eta_{\polyak}^{\crit}(1+\rho)\Tr(\bH)
}{
2B(1-\rho^2)
-
\eta_{\polyak}^{\crit}\lambda_{\max}
\left[B(1-\rho)+(1+\rho)\right]
}.
\]
Rearranging this inequality gives the lower bound in
\eqref{app:eq:polyak-critical-bounds}.

Finally, the power-law spectrum and the high-momentum regime give
\[
\lambda_{\max}\eqsim1,
\qquad
\Tr(\bH)\eqsim1,
\qquad
1-\rho^2\eqsim1-\rho.
\]
Consequently, both sides of \eqref{app:eq:polyak-critical-bounds} scale as
\[
\frac{B(1-\rho)}{B(1-\rho)+1}
\eqsim
\min\{1,B(1-\rho)\}.
\]
\end{proof}

\subsection{Risk stability of Nesterov}
\label{app:subsec:nesterov-stability}

Define the look-ahead error
\[
\be_k^{\la}
:=
(1+\rho)\be_k-\rho\be_{k-1}.
\]
The Nesterov error recursion, including label noise, is
\begin{equation}
\label{app:eq:nesterov-stability-recursion}
\be_{k+1}
=
(\bI-\eta\widehat{\bH}_k)\be_k^{\la}
+
\eta\widehat{\boldsymbol\zeta}_k.
\end{equation}
With the deterministic transition
\(\bD=\bI-\eta\bH,\)
the recursion becomes
\[
\be_{k+1}
=
\bD\be_k^{\la}
-
\eta({\bH}_k-\bH)\be_k^{\la}
+
\eta\widehat{\boldsymbol\zeta}_k.
\]

\begin{theorem}[Risk stability of Nesterov]
\label{app:thm:nesterov-stability}
Let the Hessian be a positive trace-class operator and assume
\[
\bH=\diag(\lambda_1,\lambda_2,\ldots),
\qquad
\lambda_1=\lambda_{\max}\geq\lambda_2\geq\cdots>0,
\qquad
0\leq\rho<1,
\qquad
\sigma>0.
\]
Define
\[
\Delta_{\nesterov}(x)
:=
2B(1-\rho^2)
+
\left[4B\rho^2+(B-1)\rho-(B+1)\right]x
-
(B+1)\rho(1+2\rho)x^2.
\]
Then Nesterov momentum is risk stable if and only if
\begin{equation}
\label{app:eq:nesterov-local-stability}
\Delta_{\nesterov}(\eta\lambda_{\max})>0
\end{equation}
and
\begin{equation}
\label{app:eq:nesterov-global-stability}
\kappa_{\nesterov}(\eta,\rho)
:=
\sum_{i\geq1}
\frac{
\eta\lambda_i
\left[1+\rho+\rho(1+2\rho)\eta\lambda_i\right]
}{
\Delta_{\nesterov}(\eta\lambda_i)
}
<1.
\end{equation}
\end{theorem}

\begin{proof}
Define
\[
\bU_k:=\EE[\be_k\otimes\be_k],
\qquad
\bV_k:=\EE[\be_k^{\la}\otimes\be_k^{\la}],
\qquad
\bP_k:=\EE[\be_k\otimes\be_{k-1}],
\]
\[
\bQ_k:=\EE[\be_{k-1}\otimes\be_k],
\qquad
\bR_k:=\EE[\be_{k-1}\otimes\be_{k-1}].
\]
Expanding \eqref{app:eq:nesterov-stability-recursion}, using the vanishing cross terms, and
applying Lemma~\ref{app:lem:gaussian-fourth-moment} gives
\begin{equation}
\label{app:eq:nesterov-covariance-recursion}
\left\{
\begin{aligned}
\bU_{k+1}
={}&
\bD\bV_k\bD
+
\frac{\eta^2}{B}
\left(
\bH\bV_k\bH
+
\bH\Tr(\bH\bV_k)
+
\sigma^2\bH
\right),
\\
\bV_k
={}&
(1+\rho)^2\bU_k
-
\rho(1+\rho)(\bP_k+\bQ_k)
+
\rho^2\bR_k,
\\
\bP_{k+1}
={}&
\bD\bigl((1+\rho)\bU_k-\rho\bQ_k\bigr),
\\
\bQ_{k+1}
={}&
\bigl((1+\rho)\bU_k-\rho\bP_k\bigr)\bD,
\\
\bR_{k+1}
={}&
\bU_k.
\end{aligned}
\right.
\end{equation}

The renewal equation below naturally controls the look-ahead loss moment
\[
\ell_k^{\la}
:=
\Tr(\bH\bV_k)
=
\EE\|\be_k^{\la}\|_{\bH}^2.
\]
This is equivalent to controlling the risk at the actual iterate. Define its loss moment by
\[
\ell_k
:=
\Tr(\bH\bU_k)
=
2\EE[\cE(\btheta_k)].
\]
Then
\begin{equation}
\label{app:eq:nesterov-risk-equivalence-forward}
\ell_k^{\la}
\leq
2(1+\rho)^2\ell_k+2\rho^2\ell_{k-1}.
\end{equation}
Conversely, since \(\bD\) commutes with \(\bH\) and \(\bH\) is trace class,
the first equation in \eqref{app:eq:nesterov-covariance-recursion} gives
\begin{equation}
\label{app:eq:nesterov-risk-equivalence-backward}
\ell_{k+1}
\leq
C_{\eta,\rho,\bH,B}
\bigl(\ell_k^{\la}+\sigma^2\bigr).
\end{equation}
Therefore
\[
\sup_k\ell_k^{\la}<\infty
\quad\Longleftrightarrow\quad
\sup_k\EE[\cE(\btheta_k)]<\infty.
\]

Let
\[
u_{i,k}:=\langle\mathbf e_i,\bU_k\mathbf e_i\rangle,
\qquad
p_{i,k}:=\langle\mathbf e_i,\bP_k\mathbf e_i\rangle,
\qquad
r_{i,k}:=\langle\mathbf e_i,\bR_k\mathbf e_i\rangle,
\]
and
\[
v_{i,k}:=\langle\mathbf e_i,\bV_k\mathbf e_i\rangle.
\]
The cross-covariance operators satisfy
\[
\bQ_k=\bP_k^*.
\]
Therefore,
\[
v_{i,k}
=
(1+\rho)^2u_{i,k}
-
2\rho(1+\rho)p_{i,k}
+
\rho^2r_{i,k}.
\]
Set
\[
d_i:=1-\eta\lambda_i,
\qquad
a_i:=d_i^2+\frac{\eta^2}{B}\lambda_i^2,
\qquad
\bX_{i,k}:=(u_{i,k},p_{i,k},r_{i,k})^\top,
\]
\[
\bb:=(1,0,0)^\top,
\qquad
\bc:=\bigl((1+\rho)^2,-2\rho(1+\rho),\rho^2\bigr)^\top,
\]
and
\[
\bM_i
:=
\begin{pmatrix}
a_i(1+\rho)^2
&
-2a_i\rho(1+\rho)
&
a_i\rho^2
\\
d_i(1+\rho)
&
-d_i\rho
&
0
\\
1
&
0
&
0
\end{pmatrix}.
\]
Then
\[
v_{i,k}=\bc^\top\bX_{i,k},
\qquad
\ell_k^{\la}=\sum_{i\geq1}\lambda_i\bc^\top\bX_{i,k},
\]
and the diagonal part of \eqref{app:eq:nesterov-covariance-recursion} becomes
\begin{equation}
\label{app:eq:nesterov-block-recursion}
\bX_{i,k+1}
=
\bM_i\bX_{i,k}
+
\frac{\eta^2}{B}\lambda_i\bb\bigl(\ell_k^{\la}+\sigma^2\bigr).
\end{equation}

The characteristic polynomial of the isolated block is
\[
p_i(\mu)
=
\mu^3
-
\left[a_i(1+\rho)^2-d_i\rho\right]\mu^2
+
a_i\rho\left[d_i(1+\rho)^2-\rho\right]\mu
-
a_id_i\rho^3.
\]
The Jury stability criterion shows that all roots lie strictly inside the unit disk if and only if
\[
\Delta_{\nesterov}(\eta\lambda_i)>0.
\]
The polynomial is a concave quadratic satisfying
\[
\Delta_{\nesterov}(0)=2B(1-\rho^2)>0,
\]
and it has exactly one positive root. Imposing the coordinate condition in every eigendirection is
therefore equivalent to \eqref{app:eq:nesterov-local-stability}. If it fails,
the nondegenerate label-noise forcing excites a nonstable isolated mode, and the corresponding
coordinate risk is unbounded.

Assume henceforth that the local condition holds. Iterating
\eqref{app:eq:nesterov-block-recursion}, multiplying the \(i\)-th coordinate contribution by \(\lambda_i\), and summing over all
eigendirections gives
\begin{equation}
\label{app:eq:nesterov-renewal}
\ell_k^{\la}
=
h_k
+
\sum_{s=0}^{k-1}
K_{k-1-s}\bigl(\ell_s^{\la}+\sigma^2\bigr),
\end{equation}
where
\[
h_k
:=
\sum_{i\geq1}
\lambda_i\bc^\top\bM_i^k\bX_{i,0},
\qquad
K_m
:=
\frac{\eta^2}{B}
\sum_{i\geq1}
\lambda_i^2\bc^\top\bM_i^m\bb.
\]
Both sequences are nonnegative: their summands are look-ahead variances generated by isolated
coordinate systems. The local blocks are uniformly power bounded, and positivity of the initial
coordinate covariance gives
\[
0\leq h_k
\leq
C\sum_{i\geq1}\lambda_i(u_{i,0}+r_{i,0})
<\infty.
\]
A direct calculation gives
\begin{equation}
\label{app:eq:nesterov-resolvent}
\bc^\top(\bI-\bM_i)^{-1}
=
\frac{
B\bigl(
1+\rho+\rho(1+2\rho)\eta\lambda_i,
-2\rho(1+\rho),
\rho^2[1+\rho(1-\eta\lambda_i)]
\bigr)
}{
\eta\lambda_i\Delta_{\nesterov}(\eta\lambda_i)
}.
\end{equation}
Its first entry yields
\[
\sum_{m\geq0}K_m
=
\frac{\eta^2}{B}
\sum_{i\geq1}
\lambda_i^2\bc^\top(\bI-\bM_i)^{-1}\bb
=
\kappa_{\nesterov}(\eta,\rho).
\]
Lemma~\ref{app:lem:renewal-stability}, together with
\eqref{app:eq:nesterov-risk-equivalence-forward}--
\eqref{app:eq:nesterov-risk-equivalence-backward}, gives
\[
\sup_{k\geq0}\EE[\cE(\btheta_k)]<\infty
\quad\Longleftrightarrow\quad
\kappa_{\nesterov}(\eta,\rho)<1,
\]
which completes the proof.

\end{proof}

\begin{corollary}[Critical learning rate of Nesterov]
\label{app:cor:nesterov-critical-rate}
There exists a unique critical learning rate satisfying
\[
\kappa_{\nesterov}
\bigl(\eta_{\nesterov}^{\crit},\rho\bigr)
=1
\]
before the local stability boundary. The Nesterov iteration is risk stable if and only if
\[
0<\eta<\eta_{\nesterov}^{\crit}.
\]
Under the power-law and high-momentum assumptions
\(
\lambda_i\eqsim i^{-\beta},
\beta>1,
B\gg1,
1-\rho\ll1,
\)
\begin{equation}
\label{app:eq:nesterov-critical-scaling}
\eta_{\nesterov}^{\crit}
\eqsim
\min\{1,B^\beta(1-\rho)\}.
\end{equation}
\end{corollary}

\begin{proof}
First consider the local condition. The polynomial is concave and satisfies
\[
\Delta_{\nesterov}(0)>0.
\]
It has a unique positive root, denoted by \(x_+(B,\rho)\), and hence the local condition is
equivalent to
\[
\eta\lambda_{\max}<x_+(B,\rho).
\]
In the high-momentum regime,
\[
B\gg1,
\qquad
\rho\to1^-,
\]
the root remains of constant order. Indeed,
\[
x_+(B,\rho)\eqsim1,
\qquad
x_+(B,1)
=
\frac{4B-2}{3(B+1)}
\longrightarrow
\frac43.
\]
Together with the power-law normalization
\(\lambda_{\max}\eqsim1,\)
this gives the order-one local ceiling
\[
\eta\lesssim1.
\]

We next consider the global renewal mass \(\kappa_{\nesterov}(\eta,\rho)\). Each summand is continuous and strictly increasing up to
the local boundary, and the leading eigendirection diverges at that boundary. Hence
Theorem~\ref{app:thm:nesterov-stability} gives a unique solution of the critical equation and the
stated stability interval.

On every fixed subinterval strictly inside the local stability region,
\[
\kappa_{\nesterov}(\eta,\rho)
\eqsim
\frac1B
\sum_{i\geq1}
\frac{\eta\lambda_i}{1-\rho+\eta\lambda_i}.
\]
For the power-law spectrum, the spectral sum satisfies
\[
\sum_{i\geq1}
\frac{\eta\lambda_i}{1-\rho+\eta\lambda_i}
\eqsim
\begin{cases}
\dfrac{\eta}{1-\rho},
&
\eta\lesssim1-\rho,
\\[3mm]
\left(\dfrac{\eta}{1-\rho}\right)^{1/\beta},
&
\eta\gtrsim1-\rho.
\end{cases}
\]
At a global critical point, the first regime would require
\(
\frac{\eta}{1-\rho}\eqsim B,
\)
contradicting
\(
\frac{\eta}{1-\rho}\lesssim1
\)
when the batch size is large. Therefore the global critical point lies in the second regime, where
\[
\left(\frac{\eta}{1-\rho}\right)^{1/\beta}
\eqsim
B.
\]
Thus the global condition permits learning rates up to
\[
\eta\eqsim B^\beta(1-\rho).
\]
Combining this scale with the order-one local ceiling yields
\eqref{app:eq:nesterov-critical-scaling}.
\end{proof}



% \subsection{Stability of SGD}
% \label{app:subsec:sgd-stability}
% With \(\sigma=0\), the SGD error recursion is
% \begin{equation}
% \label{app:eq:sgd-stability-recursion}
% \be_{k+1}=(\bI-\eta\widehat{\bH}_k)\be_k.
% \end{equation}

% \begin{theorem}[Mean-square stability of SGD]
% \label{app:thm:sgd-stability}
% Let \(\bH=\diag(\lambda_1,\lambda_2,\ldots)\) be a positive trace-class operator, where \(\lambda_1=\lambda_{\max}\geq\lambda_2\geq\cdots>0\). Then the homogeneous SGD recursion \eqref{app:eq:sgd-stability-recursion} is risk stable if and only if
% \begin{equation}
% \label{app:eq:sgd-local-stability}
% \eta<\frac{2B}{(B+1)\lambda_{\max}}
% \end{equation}
% and
% \begin{equation}
% \label{app:eq:sgd-global-stability}
% \mathcal R_{\sgd}(\eta)
% :=
% \sum_{i\geq1}
% \frac{\eta\lambda_i}{2B-(B+1)\eta\lambda_i}
% <1.
% \end{equation}
% \end{theorem}

% \begin{proof}
% Define the covariance operator \(\bC_k:=\EE[\be_k\otimes\be_k]\). Since \(\be_k\) is square-integrable, \(\bC_k\) is positive and trace class, with
% \[
% \Tr(\bC_k)=\EE\lVert\be_k\rVert^2.
% \]
% Because the mini-batch at iteration \(k\) is independent of the past and \(\EE[\bSigma_k]=0\), expanding \eqref{app:eq:sgd-stability-recursion} gives
% \begin{equation}
% \label{app:eq:sgd-covariance-preliminary}
% \bC_{k+1}
% =
% (\bI-\eta\bH)\bC_k(\bI-\eta\bH)
% +
% \eta^2\EE[\bSigma_k\bC_k\bSigma_k].
% \end{equation}
% Applying Lemma~\ref{app:lem:gaussian-fourth-moment} with \(\bA=\bC_k\) gives
% \[
% \EE[\bSigma_k\bC_k\bSigma_k]
% =
% \frac{1}{B}\left(\bH\bC_k\bH+\bH\Tr(\bH\bC_k)\right).
% \]
% Substituting this identity into \eqref{app:eq:sgd-covariance-preliminary}, we obtain
% \begin{equation}
% \label{app:eq:sgd-covariance}
% \bC_{k+1}
% =
% (\bI-\eta\bH)\bC_k(\bI-\eta\bH)
% +
% \frac{\eta^2}{B}\left(\bH\bC_k\bH+\bH\Tr(\bH\bC_k)\right).
% \end{equation}
% Strong risk stability is equivalent to \(\Tr(\bC_k)\to0\). Since \(\bC_k\) is positive, its trace is completely determined by its diagonal entries. Moreover, because \(\bH\) is diagonal, the diagonal entries of \eqref{app:eq:sgd-covariance} form a closed recursion that does not depend on the off-diagonal entries. Therefore, it is sufficient to analyze the diagonal covariance dynamics.
% Let \(\{\mathbf e_i\}_{i\geq1}\) denote the eigenbasis of \(\bH\), and define
% \[
% u_{i,k}:=\langle\mathbf e_i,\bC_k\mathbf e_i\rangle,
% \qquad
% S_k:=\Tr(\bH\bC_k)=\sum_{i\geq1}\lambda_i u_{i,k}.
% \]
% Since \(\bC_k\) is positive and trace class,
% \[
% \bu_k:=(u_{1,k},u_{2,k},\ldots)\in\ell_+^1,
% \qquad
% \lVert\bu_k\rVert_{\ell^1}=\Tr(\bC_k).
% \]
% Taking the diagonal entries of \eqref{app:eq:sgd-covariance} gives
% \begin{equation}
% \label{app:eq:sgd-diagonal-recursion}
% u_{i,k+1}
% =
% a_i u_{i,k}
% +
% \frac{\eta^2}{B}\lambda_iS_k,
% \qquad
% a_i:=(1-\eta\lambda_i)^2+\frac{\eta^2}{B}\lambda_i^2.
% \end{equation}
% Equivalently,
% \[
% \bu_{k+1}=\bM_\eta\bu_k,
% \qquad
% \bM_\eta=\bD_\eta+\frac{\eta^2}{B}\bm\lambda\otimes\bm\lambda,
% \]
% where
% \[
% \bD_\eta=\diag(a_1,a_2,\ldots),
% \qquad
% \bm\lambda=(\lambda_1,\lambda_2,\ldots)^\top.
% \]
% We first determine the local stability condition. A necessary condition is \(a_i<1\) for every \(i\). Since
% \[
% 1-a_i
% =
% \eta\lambda_i
% \left[
% 2-\left(1+\frac{1}{B}\right)\eta\lambda_i
% \right],
% \]
% this is equivalent to
% \[
% \eta<\frac{2B}{(B+1)\lambda_{\max}},
% \]
% which proves \eqref{app:eq:sgd-local-stability}.
% We now assume that \eqref{app:eq:sgd-local-stability} holds. Iterating \eqref{app:eq:sgd-diagonal-recursion} gives
% \begin{equation}
% \label{app:eq:sgd-diagonal-iteration}
% u_{i,k}
% =
% a_i^k u_{i,0}
% +
% \frac{\eta^2}{B}\lambda_i
% \sum_{s=0}^{k-1}
% a_i^{k-1-s}S_s.
% \end{equation}
% Multiplying by \(\lambda_i\) and summing over \(i\geq1\), we obtain the renewal equation
% \begin{equation}
% \label{app:eq:sgd-renewal}
% S_k=h_k+\sum_{s=0}^{k-1}K_{k-1-s}S_s,
% \end{equation}
% where
% \[
% h_k:=\sum_{i\geq1}\lambda_i a_i^k u_{i,0},
% \qquad
% K_m:=\frac{\eta^2}{B}\sum_{i\geq1}\lambda_i^2a_i^m.
% \]
% The forcing sequence \((h_k)_{k\geq0}\) is summable. Indeed,
% \[
% \sum_{k\geq0}h_k
% =
% \sum_{i\geq1}\frac{\lambda_i u_{i,0}}{1-a_i}
% \leq
% \frac{\lVert\bu_0\rVert_{\ell^1}}
% {\eta\left[2-\left(1+\frac{1}{B}\right)\eta\lambda_{\max}\right]}
% <\infty.
% \]
% Moreover,
% \[
% \sum_{m\geq0}K_m
% =
% \frac{\eta^2}{B}\sum_{i\geq1}\frac{\lambda_i^2}{1-a_i}
% =
% \sum_{i\geq1}\frac{\eta\lambda_i}{2B-(B+1)\eta\lambda_i}
% =
% \mathcal R_{\sgd}(\eta).
% \]
% Suppose first that \(\mathcal R_{\sgd}(\eta)<1\). Since all terms in \eqref{app:eq:sgd-renewal} are nonnegative, summing the renewal equation and applying Tonelli's theorem gives
% \[
% \sum_{k\geq0}S_k
% \leq
% \frac{\sum_{k\geq0}h_k}
% {1-\mathcal R_{\sgd}(\eta)}
% <\infty.
% \]
% Summing \eqref{app:eq:sgd-diagonal-iteration} over \(i\) gives
% \begin{equation}
% \label{app:eq:sgd-trace-recursion}
% \lVert\bu_k\rVert_{\ell^1}
% =
% \sum_{i\geq1}a_i^k u_{i,0}
% +
% \frac{\eta^2}{B}
% \sum_{s=0}^{k-1}G_{k-1-s}S_s,
% \end{equation}
% where
% \[
% G_m:=\sum_{i\geq1}\lambda_i a_i^m.
% \]
% For every fixed \(i\), \(a_i^k\to0\). Since \(\bu_0\in\ell^1\), dominated convergence gives
% \[
% \sum_{i\geq1}a_i^k u_{i,0}\to0.
% \]
% Similarly, since \(\lambda\in\ell^1\),
% \[
% G_m\to0,
% \qquad
% \sup_{m\geq0}G_m\leq\Tr(\bH)<\infty.
% \]
% Because \((S_k)_{k\geq0}\in\ell^1\), the convolution term in \eqref{app:eq:sgd-trace-recursion} also converges to zero. Therefore,
% \[
% \EE\lVert\be_k\rVert^2
% =
% \Tr(\bC_k)
% =
% \lVert\bu_k\rVert_{\ell^1}
% \to0.
% \]
% This proves sufficiency.
% We next prove necessity. If \eqref{app:eq:sgd-local-stability} fails, then \(a_i\geq1\) for some \(i\). Choosing an initialization supported on this eigendirection gives
% \[
% u_{i,k+1}\geq a_i u_{i,k},
% \]
% so the second moment does not converge to zero. Hence the local condition is necessary.
% Now assume that the local condition holds but \(\mathcal R_{\sgd}(\eta)>1\). For \(\mu>1\), define
% \[
% F(\mu)
% :=
% \frac{\eta^2}{B}
% \sum_{i\geq1}
% \frac{\lambda_i^2}{\mu-a_i}.
% \]
% The function \(F\) is continuous and strictly decreasing, with
% \[
% \lim_{\mu\downarrow1}F(\mu)
% =
% \mathcal R_{\sgd}(\eta)
% >1,
% \qquad
% \lim_{\mu\to\infty}F(\mu)=0.
% \]
% Hence there exists \(\mu>1\) such that \(F(\mu)=1\). Define
% \[
% v_i:=\frac{\lambda_i}{\mu-a_i}.
% \]
% Since \(\mu>1\) and \(\lambda\in\ell^1\), we have \(\bv=(v_1,v_2,\ldots)\in\ell_+^1\). Moreover,
% \[
% \bM_\eta\bv=\mu\bv.
% \]
% Choosing the deterministic initialization \((\be_0)_i=\sqrt{v_i}\) gives \(\bu_0=\bv\), and hence
% \[
% \bu_k=\mu^k\bv.
% \]
% Consequently,
% \[
% \EE\lVert\be_k\rVert^2
% =
% \mu^k\lVert\bv\rVert_{\ell^1}
% \to\infty.
% \]
% It remains to consider the boundary \(\mathcal R_{\sgd}(\eta)=1\). Define
% \[
% w_i
% :=
% \frac{\lambda_i}{1-a_i}
% =
% \frac{1}
% {\eta\left[2-\left(1+\frac{1}{B}\right)\eta\lambda_i\right]}.
% \]
% The local condition implies \(\bw=(w_1,w_2,\ldots)\in\ell_+^\infty\), and \(\mathcal R_{\sgd}(\eta)=1\) implies
% \[
% \bM_\eta^*\bw=\bw.
% \]
% Thus, for every nonzero \(\bu_0\in\ell_+^1\),
% \[
% \langle\bw,\bu_k\rangle
% =
% \langle\bw,\bu_0\rangle
% >0
% \qquad
% \text{for every }k.
% \]
% If \(\lVert\bu_k\rVert_{\ell^1}\to0\), then
% \[
% \langle\bw,\bu_k\rangle
% \leq
% \lVert\bw\rVert_{\ell^\infty}
% \lVert\bu_k\rVert_{\ell^1}
% \to0,
% \]
% which is a contradiction. Hence the critical boundary is not strongly risk asymptotically stable.
% \end{proof}

% \begin{corollary}[Critical learning rate of SGD]
% \label{app:cor:sgd-critical-rate}
% There exists a unique critical learning rate \(\eta_{\sgd}^{\crit}\) satisfying
% \[
% 0<\eta_{\sgd}^{\crit}<\frac{2B}{(B+1)\lambda_{\max}}
% \]
% and
% \begin{equation}
% \label{app:eq:sgd-critical-equation}
% \sum_{i\geq1}
% \frac{
% \eta_{\sgd}^{\crit}\lambda_i
% }{
% 2B-(B+1)\eta_{\sgd}^{\crit}\lambda_i
% }
% =1.
% \end{equation}
% The homogeneous SGD dynamics are risk stable if and only if
% \[
% 0<\eta<\eta_{\sgd}^{\crit}.
% \]
% Moreover,
% \begin{equation}
% \label{app:eq:sgd-critical-bounds}
% \frac{2B}
% {(B+1)\lambda_{\max}+\Tr(\bH)}
% \leq
% \eta_{\sgd}^{\crit}
% \leq
% \min\left\{
% \frac{2B}{(B+1)\lambda_{\max}},
% \frac{2B}{\Tr(\bH)}
% \right\}.
% \end{equation}
% Consequently, under \(\lambda_i\eqsim i^{-\beta}\), \(\beta>1\), with \(\beta\) regarded as fixed,
% \[
% \eta_{\sgd}^{\crit}\eqsim1.
% \]
% \end{corollary}

% \begin{proof}
% The function \(\mathcal R_{\sgd}(\eta)\) is continuous and strictly increasing on
% \[
% 0<\eta<\frac{2B}{(B+1)\lambda_{\max}}.
% \]
% Moreover, \(\mathcal R_{\sgd}(0)=0\), while \(\mathcal R_{\sgd}(\eta)\to\infty\) as \(\eta\) approaches the local stability boundary. Hence the critical learning rate exists and is unique.
% At the critical learning rate, using
% \[
% 2B-(B+1)\eta_{\sgd}^{\crit}\lambda_i\leq2B
% \]
% in \eqref{app:eq:sgd-critical-equation} gives
% \[
% 1
% \geq
% \frac{
% \eta_{\sgd}^{\crit}\Tr(\bH)
% }{2B},
% \]
% and hence
% \[
% \eta_{\sgd}^{\crit}
% \leq
% \frac{2B}{\Tr(\bH)}.
% \]
% The local condition gives the other upper bound in \eqref{app:eq:sgd-critical-bounds}.
% Conversely,
% \[
% 2B-(B+1)\eta_{\sgd}^{\crit}\lambda_i
% \geq
% 2B-(B+1)\eta_{\sgd}^{\crit}\lambda_{\max}.
% \]
% Therefore,
% \[
% 1
% \leq
% \frac{
% \eta_{\sgd}^{\crit}\Tr(\bH)
% }{
% 2B-(B+1)\eta_{\sgd}^{\crit}\lambda_{\max}
% },
% \]
% which yields the lower bound in \eqref{app:eq:sgd-critical-bounds}.
% Finally, when \(\lambda_i\eqsim i^{-\beta}\) with fixed \(\beta>1\),
% \[
% \lambda_{\max}\eqsim1,
% \qquad
% \Tr(\bH)=\sum_{i\geq1}\lambda_i\eqsim1.
% \]
% The two-sided bound \eqref{app:eq:sgd-critical-bounds} therefore gives
% \[
% \eta_{\sgd}^{\crit}\eqsim1.
% \]
% \end{proof}

% \begin{remark}[Absence of a uniform spectral gap]
% \label{app:rem:no-uniform-spectral-gap}
% Although \(a_i<1\) for every fixed eigendirection in the stable regime, the power-law spectrum satisfies \(\lambda_i\to0\), and hence
% \[
% a_i\to1
% \qquad
% \text{as }i\to\infty.
% \]
% Thus, the infinite-dimensional moment dynamics generally do not possess a uniform spectral gap. In particular, the spectral radius of the global covariance operator need not be strictly smaller than one even though every admissible covariance trajectory converges strongly to zero. Stability here is strong risk asymptotic stability: the convergence holds in trace norm for each admissible initial covariance, but need not be uniform or exponential over all spectral directions.
% \end{remark}

% \subsection{Stability of Polyak}
% \label{app:subsec:polyak-stability}
% With \(\sigma=0\), the Polyak error recursion is
% \begin{equation}
% \label{app:eq:polyak-stability-recursion}
% \be_{k+1}
% =
% (\bI-\eta\widehat{\bH}_k)\be_k
% +
% \rho(\be_k-\be_{k-1}).
% \end{equation}
% Writing
% \[
% \widehat{\bH}_k=\bH+\bSigma_k,
% \qquad
% \bD=(1+\rho)\bI-\eta\bH,
% \]
% the recursion becomes
% \[
% \be_{k+1}
% =
% \bD\be_k-\rho\be_{k-1}-\eta\bSigma_k\be_k.
% \]

% \begin{theorem}[Mean-square stability of Polyak]
% \label{app:thm:polyak-stability}
% Let \(\bH=\diag(\lambda_1,\lambda_2,\ldots)\) be a positive trace-class operator, where \(\lambda_1=\lambda_{\max}\geq\lambda_2\geq\cdots>0\), and let \(0\leq\rho<1\). Then the homogeneous Polyak recursion \eqref{app:eq:polyak-stability-recursion} is risk stable if and only if
% \begin{equation}
% \label{app:eq:polyak-local-stability}
% \eta
% <
% \frac{
% 2B(1-\rho^2)
% }{
% \lambda_{\max}
% \left[B(1-\rho)+(1+\rho)\right]
% }
% \end{equation}
% and
% \begin{equation}
% \label{app:eq:polyak-global-stability}
% \mathcal R_{\polyak}(\eta,\rho)
% :=
% \sum_{i\geq1}
% \frac{
% \eta\lambda_i(1+\rho)
% }{
% 2B(1-\rho^2)
% -
% \eta\lambda_i
% \left[B(1-\rho)+(1+\rho)\right]
% }
% <1.
% \end{equation}
% \end{theorem}

% \begin{proof}
% Define
% \[
% \bU_k:=\EE[\be_k\otimes\be_k],
% \qquad
% \bP_k:=\EE[\be_k\otimes\be_{k-1}],
% \qquad
% \bQ_k:=\EE[\be_{k-1}\otimes\be_k],
% \qquad
% \bR_k:=\EE[\be_{k-1}\otimes\be_{k-1}].
% \]
% These operators are trace class whenever the augmented state has a finite second moment, and
% \[
% \EE\lVert\bz_k\rVert^2
% =
% \Tr(\bU_k)+\Tr(\bR_k).
% \]
% Because the mini-batch at iteration \(k\) is independent of the past and \(\EE[\bSigma_k]=0\), expanding \eqref{app:eq:polyak-stability-recursion} and applying Lemma~\ref{app:lem:gaussian-fourth-moment} gives
% \begin{equation}
% \label{app:eq:polyak-covariance-recursion}
% \left\{
% \begin{aligned}
% \bU_{k+1}
% ={}&
% \bD\bU_k\bD
% -\rho\bD\bP_k
% -\rho\bQ_k\bD
% +\rho^2\bR_k
% +
% \frac{\eta^2}{B}
% \left(
% \bH\bU_k\bH
% +
% \bH\Tr(\bH\bU_k)
% \right),
% \\
% \bP_{k+1}
% ={}&
% \bD\bU_k-\rho\bQ_k,
% \\
% \bQ_{k+1}
% ={}&
% \bU_k\bD-\rho\bP_k,
% \\
% \bR_{k+1}
% ={}&
% \bU_k.
% \end{aligned}
% \right.
% \end{equation}
% Strong risk stability is equivalent to
% \[
% \Tr(\bU_k)+\Tr(\bR_k)\to0.
% \]
% Since \(\bR_{k+1}=\bU_k\), it is sufficient and necessary to prove \(\Tr(\bU_k)\to0\). Moreover, because \(\bH\) and \(\bD\) are diagonal, the diagonal entries of \eqref{app:eq:polyak-covariance-recursion} form a closed recursion.
% Let \(\{\mathbf e_i\}_{i\geq1}\) denote the eigenbasis of \(\bH\), and define
% \[
% u_{i,k}:=\langle\mathbf e_i,\bU_k\mathbf e_i\rangle,
% \qquad
% v_{i,k}:=\langle\mathbf e_i,\bP_k\mathbf e_i\rangle,
% \qquad
% r_{i,k}:=\langle\mathbf e_i,\bR_k\mathbf e_i\rangle,
% \]
% and
% \[
% S_k:=\Tr(\bH\bU_k)=\sum_{i\geq1}\lambda_i u_{i,k}.
% \]
% Since \(\bQ_k=\bP_k^*\), the diagonal entries of \(\bP_k\) and \(\bQ_k\) coincide. Setting
% \[
% d_i:=1+\rho-\eta\lambda_i,
% \]
% we obtain
% \begin{equation}
% \label{app:eq:polyak-diagonal-recursion}
% \left\{
% \begin{aligned}
% u_{i,k+1}
% ={}&
% \left(
% d_i^2+\frac{\eta^2}{B}\lambda_i^2
% \right)u_{i,k}
% -
% 2\rho d_i v_{i,k}
% +
% \rho^2r_{i,k}
% +
% \frac{\eta^2}{B}\lambda_iS_k,
% \\
% v_{i,k+1}
% ={}&
% d_i u_{i,k}-\rho v_{i,k},
% \\
% r_{i,k+1}
% ={}&
% u_{i,k}.
% \end{aligned}
% \right.
% \end{equation}
% Define
% \[
% \bX_{i,k}:=(u_{i,k},v_{i,k},r_{i,k})^\top,
% \qquad
% \bb:=(1,0,0)^\top,
% \]
% and
% \[
% \bM_i
% :=
% \begin{pmatrix}
% d_i^2+\dfrac{\eta^2}{B}\lambda_i^2
% &
% -2\rho d_i
% &
% \rho^2
% \\
% d_i
% &
% -\rho
% &
% 0
% \\
% 1
% &
% 0
% &
% 0
% \end{pmatrix}.
% \]
% Then \eqref{app:eq:polyak-diagonal-recursion} becomes
% \begin{equation}
% \label{app:eq:polyak-block-recursion}
% \bX_{i,k+1}
% =
% \bM_i\bX_{i,k}
% +
% \frac{\eta^2}{B}\lambda_i\bb S_k.
% \end{equation}
% We first determine the local stability condition. Ignoring the global trace-coupling term \(S_k\), the characteristic polynomial of \(\bM_i\) is
% \[
% p_i(\mu)
% =
% \mu^3
% +
% \left(
% \rho-d_i^2-\frac{\eta^2}{B}\lambda_i^2
% \right)\mu^2
% +
% \rho
% \left(
% d_i^2-\frac{\eta^2}{B}\lambda_i^2-\rho
% \right)\mu
% -
% \rho^3.
% \]
% The Jury stability criterion shows that all roots lie strictly inside the unit disk if and only if
% \[
% 2B(1-\rho^2)
% -
% \eta\lambda_i
% \left[B(1-\rho)+(1+\rho)\right]
% >0.
% \]
% Requiring this condition for every \(i\) is equivalent to \eqref{app:eq:polyak-local-stability}.
% We now assume that the local condition holds. Define
% \[
% g_{i,m}:=\bb^\top\bM_i^m\bb.
% \]
% This quantity is nonnegative because it is the variance generated by the isolated \(i\)-th coordinate dynamics from the covariance state \(\bb\). For each fixed \(i\), \(g_{i,m}\to0\). Moreover, the family of local blocks is uniformly power bounded under the strict local condition, so there exists \(C<\infty\), independent of \(i\) and \(m\), such that
% \[
% 0\leq g_{i,m}\leq C.
% \]
% Iterating \eqref{app:eq:polyak-block-recursion}, taking the first component, multiplying by \(\lambda_i\), and summing over \(i\geq1\) gives
% \begin{equation}
% \label{app:eq:polyak-renewal}
% S_k=h_k+\sum_{s=0}^{k-1}K_{k-1-s}S_s,
% \end{equation}
% where
% \[
% h_k
% :=
% \sum_{i\geq1}
% \lambda_i\bb^\top\bM_i^k\bX_{i,0},
% \qquad
% K_m
% :=
% \frac{\eta^2}{B}
% \sum_{i\geq1}
% \lambda_i^2g_{i,m}.
% \]
% The sequence \((h_k)_{k\geq0}\) is nonnegative and summable. Indeed,
% \[
% \sum_{k\geq0}h_k
% =
% \sum_{i\geq1}
% \lambda_i
% \bb^\top(\bI-\bM_i)^{-1}\bX_{i,0}.
% \]
% A direct calculation gives
% \begin{equation}
% \label{app:eq:polyak-resolvent}
% \bb^\top(\bI-\bM_i)^{-1}
% =
% \frac{
% B
% \bigl(
% 1+\rho,\,
% -2\rho d_i,\,
% \rho^2(1+\rho)
% \bigr)
% }{
% \eta\lambda_i
% \left\{
% 2B(1-\rho^2)
% -
% \eta\lambda_i
% \left[B(1-\rho)+(1+\rho)\right]
% \right\}
% }.
% \end{equation}
% The initial coordinate covariance satisfies
% \[
% \begin{pmatrix}
% u_{i,0} & v_{i,0}\\
% v_{i,0} & r_{i,0}
% \end{pmatrix}
% \succeq0.
% \]
% Therefore,
% \[
% v_{i,0}^2\leq u_{i,0}r_{i,0},
% \qquad
% |v_{i,0}|
% \leq
% \frac{u_{i,0}+r_{i,0}}{2}.
% \]
% Define
% \[
% \delta_{\polyak}
% :=
% 2B(1-\rho^2)
% -
% \eta\lambda_{\max}
% \left[B(1-\rho)+(1+\rho)\right]
% >0.
% \]
% Using \eqref{app:eq:polyak-resolvent} and \(|d_i|\leq1+\rho+\eta\lambda_{\max}\), we obtain
% \[
% \lambda_i
% \bb^\top(\bI-\bM_i)^{-1}\bX_{i,0}
% \leq
% C(u_{i,0}+r_{i,0}),
% \]
% where \(C<\infty\) is independent of \(i\). Hence
% \[
% \sum_{k\geq0}h_k
% \leq
% C\sum_{i\geq1}(u_{i,0}+r_{i,0})
% <\infty.
% \]
% By Tonelli's theorem,
% \[
% \sum_{m\geq0}K_m
% =
% \frac{\eta^2}{B}
% \sum_{i\geq1}
% \lambda_i^2
% \bb^\top(\bI-\bM_i)^{-1}\bb.
% \]
% The first entry of \eqref{app:eq:polyak-resolvent} gives
% \[
% \bb^\top(\bI-\bM_i)^{-1}\bb
% =
% \frac{
% B(1+\rho)
% }{
% \eta\lambda_i
% \left\{
% 2B(1-\rho^2)
% -
% \eta\lambda_i
% \left[B(1-\rho)+(1+\rho)\right]
% \right\}
% },
% \]
% and consequently
% \[
% \sum_{m\geq0}K_m
% =
% \mathcal R_{\polyak}(\eta,\rho).
% \]
% Suppose first that \(\mathcal R_{\polyak}(\eta,\rho)<1\). Summing \eqref{app:eq:polyak-renewal} and applying Tonelli's theorem gives
% \[
% \sum_{k\geq0}S_k
% \leq
% \frac{\sum_{k\geq0}h_k}
% {1-\mathcal R_{\polyak}(\eta,\rho)}
% <\infty.
% \]
% Taking the first component of the iterated block recursion and summing over \(i\) yields
% \[
% \sum_{i\geq1}u_{i,k}
% =
% H_k
% +
% \frac{\eta^2}{B}
% \sum_{s=0}^{k-1}
% G_{k-1-s}S_s,
% \]
% where
% \[
% H_k
% :=
% \sum_{i\geq1}\bb^\top\bM_i^k\bX_{i,0},
% \qquad
% G_m
% :=
% \sum_{i\geq1}\lambda_i g_{i,m}.
% \]
% For every fixed \(i\), \(\bb^\top\bM_i^k\bX_{i,0}\to0\). The uniform power bound and
% \[
% \sum_{i\geq1}(u_{i,0}+r_{i,0})<\infty
% \]
% allow us to apply dominated convergence, giving \(H_k\to0\). Similarly,
% \[
% G_m\to0,
% \qquad
% \sup_{m\geq0}G_m
% \leq
% C\Tr(\bH)
% <\infty.
% \]
% Since \((S_k)_{k\geq0}\in\ell^1\), the convolution term also converges to zero. Therefore,
% \[
% \Tr(\bU_k)
% =
% \sum_{i\geq1}u_{i,k}
% \to0.
% \]
% Since \(\bR_{k+1}=\bU_k\),
% \[
% \EE\lVert\bz_k\rVert^2
% =
% \Tr(\bU_k)+\Tr(\bR_k)
% \to0.
% \]
% This proves sufficiency.
% We next prove necessity. If \eqref{app:eq:polyak-local-stability} fails, at least one local block \(\bM_i\) is not asymptotically stable. The isolated moment map preserves the cone of positive semidefinite covariance matrices, so there exists a deterministic augmented initialization supported on this eigendirection whose second moment does not converge to zero. Hence the local condition is necessary.
% Now suppose that the local condition holds but \(\mathcal R_{\polyak}(\eta,\rho)>1\). For \(\mu>1\), define
% \[
% F(\mu)
% :=
% \frac{\eta^2}{B}
% \sum_{i\geq1}
% \lambda_i^2
% \bb^\top(\mu\bI-\bM_i)^{-1}\bb.
% \]
% The Neumann-series representation shows that \(F\) is continuous and strictly decreasing, with
% \[
% \lim_{\mu\downarrow1}F(\mu)
% =
% \mathcal R_{\polyak}(\eta,\rho)
% >1,
% \qquad
% \lim_{\mu\to\infty}F(\mu)=0.
% \]
% Hence there exists \(\mu>1\) such that \(F(\mu)=1\). Define
% \[
% \bW_i
% :=
% \lambda_i
% (\mu\bI-\bM_i^*)^{-1}\bb.
% \]
% Then \(\bW=(\bW_i)_{i\geq1}\) is an eigenvector of the adjoint global moment operator with eigenvalue \(\mu\). Moreover, for every positive covariance state \(\bX_i\),
% \[
% \langle\bW_i,\bX_i\rangle
% =
% \lambda_i
% \sum_{m\geq0}
% \mu^{-m-1}
% \bb^\top\bM_i^m\bX_i
% \geq0.
% \]
% Choose a deterministic initialization such that \(\langle\bW,\bX_0\rangle>0\). Then
% \[
% \langle\bW,\bX_k\rangle
% =
% \mu^k\langle\bW,\bX_0\rangle,
% \]
% which cannot converge to zero. Hence the Polyak dynamics are not strongly risk asymptotically stable.
% It remains to consider \(\mathcal R_{\polyak}(\eta,\rho)=1\). Define
% \[
% \bW_i
% :=
% \lambda_i(\bI-\bM_i^*)^{-1}\bb.
% \]
% The local condition and \eqref{app:eq:polyak-resolvent} imply that \(\bW\) is bounded, while \(\mathcal R_{\polyak}(\eta,\rho)=1\) gives
% \[
% \bM_{\polyak}^*\bW=\bW,
% \]
% where \(\bM_{\polyak}\) denotes the global diagonal moment operator. Choosing an initialization with \(\langle\bW,\bX_0\rangle>0\) gives
% \[
% \langle\bW,\bX_k\rangle
% =
% \langle\bW,\bX_0\rangle
% >0
% \]
% for every \(k\), contradicting convergence of the augmented covariance to zero. Hence the critical boundary is not strongly risk asymptotically stable.
% \end{proof}

% \begin{corollary}[Critical learning rate of Polyak]
% \label{app:cor:polyak-critical-rate}
% There exists a unique critical learning rate \(\eta_{\polyak}^{\crit}\) satisfying
% \[
% 0
% <
% \eta_{\polyak}^{\crit}
% <
% \frac{
% 2B(1-\rho^2)
% }{
% \lambda_{\max}
% \left[B(1-\rho)+(1+\rho)\right]
% }
% \]
% and
% \begin{equation}
% \label{app:eq:polyak-critical-equation}
% \sum_{i\geq1}
% \frac{
% \eta_{\polyak}^{\crit}
% \lambda_i(1+\rho)
% }{
% 2B(1-\rho^2)
% -
% \eta_{\polyak}^{\crit}\lambda_i
% \left[B(1-\rho)+(1+\rho)\right]
% }
% =1.
% \end{equation}
% The homogeneous Polyak dynamics are risk stable if and only if
% \[
% 0<\eta<\eta_{\polyak}^{\crit}.
% \]
% Moreover,
% \begin{equation}
% \label{app:eq:polyak-critical-bounds}
% \frac{
% 2B(1-\rho^2)
% }{
% \lambda_{\max}
% \left[B(1-\rho)+(1+\rho)\right]
% +
% (1+\rho)\Tr(\bH)
% }
% \leq
% \eta_{\polyak}^{\crit}
% \leq
% \min\left\{
% \frac{
% 2B(1-\rho^2)
% }{
% \lambda_{\max}
% \left[B(1-\rho)+(1+\rho)\right]
% },
% \frac{
% 2B(1-\rho)
% }{
% \Tr(\bH)
% }
% \right\}.
% \end{equation}
% Consequently, under \(\lambda_i\eqsim i^{-\beta}\) with fixed \(\beta>1\), in the scaling regime \(B\gg1\) and \(1-\rho\ll1\),
% \[
% \eta_{\polyak}^{\crit}
% \eqsim
% \min\{1,B(1-\rho)\}.
% \]
% \end{corollary}

% \begin{proof}
% The function \(\mathcal R_{\polyak}(\eta,\rho)\) is continuous and strictly increasing on the local stability interval. It vanishes at \(\eta=0\) and diverges as \(\eta\) approaches the local stability boundary. Hence the critical learning rate exists and is unique.
% At the critical learning rate, bounding each denominator in \eqref{app:eq:polyak-critical-equation} from above by \(2B(1-\rho^2)\) gives
% \[
% \eta_{\polyak}^{\crit}
% \leq
% \frac{2B(1-\rho)}{\Tr(\bH)}.
% \]
% The local condition gives the other upper bound in \eqref{app:eq:polyak-critical-bounds}.
% Conversely, bounding each denominator from below by
% \[
% 2B(1-\rho^2)
% -
% \eta_{\polyak}^{\crit}\lambda_{\max}
% \left[B(1-\rho)+(1+\rho)\right]
% \]
% gives the lower bound in \eqref{app:eq:polyak-critical-bounds}.
% Finally,
% \[
% \lambda_{\max}\eqsim1,
% \qquad
% \Tr(\bH)\eqsim1,
% \qquad
% 1-\rho^2\eqsim1-\rho.
% \]
% Both sides of \eqref{app:eq:polyak-critical-bounds} therefore scale as
% \[
% \frac{B(1-\rho)}{B(1-\rho)+1}
% \eqsim
% \min\{1,B(1-\rho)\}.
% \]
% \end{proof}

% \subsection{Stability of Nesterov}
% \label{app:subsec:nesterov-stability}
% Define the look-ahead error
% \[
% \be_k^{\la}
% :=
% (1+\rho)\be_k-\rho\be_{k-1}.
% \]
% With \(\sigma=0\), the Nesterov error recursion is
% \begin{equation}
% \label{app:eq:nesterov-stability-recursion}
% \be_{k+1}
% =
% (\bI-\eta\widehat{\bH}_k)\be_k^{\la}.
% \end{equation}
% Writing
% \[
% \widehat{\bH}_k=\bH+\bSigma_k,
% \qquad
% \bD=\bI-\eta\bH,
% \]
% the recursion becomes
% \[
% \be_{k+1}
% =
% \bD\be_k^{\la}
% -
% \eta\bSigma_k\be_k^{\la}.
% \]

% \begin{theorem}[Mean-square stability of Nesterov]
% \label{app:thm:nesterov-stability}
% Let \(\bH=\diag(\lambda_1,\lambda_2,\ldots)\) be a positive trace-class operator, where \(\lambda_1=\lambda_{\max}\geq\lambda_2\geq\cdots>0\), and let \(0\leq\rho<1\). Define
% \[
% \Delta_{\nesterov}(x)
% :=
% 2B(1-\rho^2)
% +
% \left[
% 4B\rho^2+(B-1)\rho-(B+1)
% \right]x
% -
% (B+1)\rho(1+2\rho)x^2.
% \]
% Then the homogeneous Nesterov recursion \eqref{app:eq:nesterov-stability-recursion} is risk stable if and only if
% \begin{equation}
% \label{app:eq:nesterov-local-stability}
% \Delta_{\nesterov}(\eta\lambda_{\max})>0
% \end{equation}
% and
% \begin{equation}
% \label{app:eq:nesterov-global-stability}
% \mathcal R_{\nesterov}(\eta,\rho)
% :=
% \sum_{i\geq1}
% \frac{
% \eta\lambda_i
% \left[
% 1+\rho+\rho(1+2\rho)\eta\lambda_i
% \right]
% }{
% \Delta_{\nesterov}(\eta\lambda_i)
% }
% <1.
% \end{equation}
% \end{theorem}

% \begin{proof}
% Define
% \[
% \bU_k:=\EE[\be_k\otimes\be_k],
% \qquad
% \bV_k:=\EE[\be_k^{\la}\otimes\be_k^{\la}],
% \qquad
% \bP_k:=\EE[\be_k\otimes\be_{k-1}],
% \]
% \[
% \bQ_k:=\EE[\be_{k-1}\otimes\be_k],
% \qquad
% \bR_k:=\EE[\be_{k-1}\otimes\be_{k-1}].
% \]
% Because the mini-batch at iteration \(k\) is independent of the past and \(\EE[\bSigma_k]=0\), expanding \eqref{app:eq:nesterov-stability-recursion} and applying Lemma~\ref{app:lem:gaussian-fourth-moment} gives
% \begin{equation}
% \label{app:eq:nesterov-covariance-recursion}
% \left\{
% \begin{aligned}
% \bU_{k+1}
% ={}&
% \bD\bV_k\bD
% +
% \frac{\eta^2}{B}
% \left(
% \bH\bV_k\bH
% +
% \bH\Tr(\bH\bV_k)
% \right),
% \\
% \bV_k
% ={}&
% (1+\rho)^2\bU_k
% -
% \rho(1+\rho)(\bP_k+\bQ_k)
% +
% \rho^2\bR_k,
% \\
% \bP_{k+1}
% ={}&
% \bD\bigl((1+\rho)\bU_k-\rho\bQ_k\bigr),
% \\
% \bQ_{k+1}
% ={}&
% \bigl((1+\rho)\bU_k-\rho\bP_k\bigr)\bD,
% \\
% \bR_{k+1}
% ={}&
% \bU_k.
% \end{aligned}
% \right.
% \end{equation}
% Since
% \[
% \EE\lVert\bz_k\rVert^2
% =
% \Tr(\bU_k)+\Tr(\bR_k),
% \]
% strong risk stability is equivalent to
% \[
% \Tr(\bU_k)+\Tr(\bR_k)\to0.
% \]
% Because \(\bR_{k+1}=\bU_k\), it is sufficient and necessary to prove \(\Tr(\bU_k)\to0\). Moreover, since \(\bH\) and \(\bD\) are diagonal, the diagonal entries of \eqref{app:eq:nesterov-covariance-recursion} form a closed recursion.
% Let
% \[
% u_{i,k}:=\langle\mathbf e_i,\bU_k\mathbf e_i\rangle,
% \qquad
% p_{i,k}:=\langle\mathbf e_i,\bP_k\mathbf e_i\rangle,
% \qquad
% r_{i,k}:=\langle\mathbf e_i,\bR_k\mathbf e_i\rangle,
% \]
% and
% \[
% v_{i,k}:=\langle\mathbf e_i,\bV_k\mathbf e_i\rangle.
% \]
% Since \(\bQ_k=\bP_k^*\),
% \[
% v_{i,k}
% =
% (1+\rho)^2u_{i,k}
% -
% 2\rho(1+\rho)p_{i,k}
% +
% \rho^2r_{i,k}.
% \]
% Set
% \[
% S_k:=\Tr(\bH\bV_k)=\sum_{i\geq1}\lambda_i v_{i,k},
% \qquad
% d_i:=1-\eta\lambda_i,
% \qquad
% a_i:=d_i^2+\frac{\eta^2}{B}\lambda_i^2.
% \]
% Taking the diagonal entries of \eqref{app:eq:nesterov-covariance-recursion} gives
% \begin{equation}
% \label{app:eq:nesterov-diagonal-recursion}
% \left\{
% \begin{aligned}
% u_{i,k+1}
% ={}&
% a_i
% \left[
% (1+\rho)^2u_{i,k}
% -
% 2\rho(1+\rho)p_{i,k}
% +
% \rho^2r_{i,k}
% \right]
% +
% \frac{\eta^2}{B}\lambda_iS_k,
% \\
% p_{i,k+1}
% ={}&
% d_i
% \left[
% (1+\rho)u_{i,k}-\rho p_{i,k}
% \right],
% \\
% r_{i,k+1}
% ={}&
% u_{i,k}.
% \end{aligned}
% \right.
% \end{equation}
% Define
% \[
% \bX_{i,k}:=(u_{i,k},p_{i,k},r_{i,k})^\top,
% \qquad
% \bb:=(1,0,0)^\top,
% \qquad
% \bc:=\bigl((1+\rho)^2,-2\rho(1+\rho),\rho^2\bigr)^\top,
% \]
% and
% \[
% \bM_i
% :=
% \begin{pmatrix}
% a_i(1+\rho)^2
% &
% -2a_i\rho(1+\rho)
% &
% a_i\rho^2
% \\
% d_i(1+\rho)
% &
% -d_i\rho
% &
% 0
% \\
% 1
% &
% 0
% &
% 0
% \end{pmatrix}.
% \]
% Then
% \[
% v_{i,k}=\bc^\top\bX_{i,k},
% \qquad
% S_k=\sum_{i\geq1}\lambda_i\bc^\top\bX_{i,k},
% \]
% and \eqref{app:eq:nesterov-diagonal-recursion} becomes
% \begin{equation}
% \label{app:eq:nesterov-block-recursion}
% \bX_{i,k+1}
% =
% \bM_i\bX_{i,k}
% +
% \frac{\eta^2}{B}\lambda_i\bb S_k.
% \end{equation}
% We first determine the local stability condition. The characteristic polynomial of the isolated block \(\bM_i\) is
% \[
% p_i(\mu)
% =
% \mu^3
% -
% \left[
% a_i(1+\rho)^2-d_i\rho
% \right]\mu^2
% +
% a_i\rho
% \left[
% d_i(1+\rho)^2-\rho
% \right]\mu
% -
% a_id_i\rho^3.
% \]
% The Jury stability criterion shows that all roots lie strictly inside the unit disk if and only if
% \[
% \Delta_{\nesterov}(\eta\lambda_i)>0.
% \]
% Since \(\Delta_{\nesterov}\) is a concave quadratic with
% \[
% \Delta_{\nesterov}(0)=2B(1-\rho^2)>0
% \]
% and exactly one positive root, requiring this condition for every \(i\) is equivalent to \eqref{app:eq:nesterov-local-stability}.
% We now assume that the local condition holds. Iterating \eqref{app:eq:nesterov-block-recursion} gives
% \[
% \bX_{i,k}
% =
% \bM_i^k\bX_{i,0}
% +
% \frac{\eta^2}{B}\lambda_i
% \sum_{s=0}^{k-1}
% \bM_i^{k-1-s}\bb S_s.
% \]
% Multiplying by \(\lambda_i\bc^\top\) and summing over \(i\geq1\), we obtain
% \begin{equation}
% \label{app:eq:nesterov-renewal}
% S_k=h_k+\sum_{s=0}^{k-1}K_{k-1-s}S_s,
% \end{equation}
% where
% \[
% h_k
% :=
% \sum_{i\geq1}
% \lambda_i\bc^\top\bM_i^k\bX_{i,0},
% \qquad
% K_m
% :=
% \frac{\eta^2}{B}
% \sum_{i\geq1}
% \lambda_i^2\bc^\top\bM_i^m\bb.
% \]
% Both sequences are nonnegative: \(\bc^\top\bM_i^k\bX_{i,0}\) is the look-ahead variance generated by the isolated \(i\)-th coordinate dynamics, while \(\bc^\top\bM_i^m\bb\) is the look-ahead variance generated from the covariance state \(\bb\).
% We next show that \((h_k)_{k\geq0}\) is summable. Since every local block is stable,
% \[
% \sum_{k\geq0}h_k
% =
% \sum_{i\geq1}
% \lambda_i
% \bc^\top(\bI-\bM_i)^{-1}\bX_{i,0}.
% \]
% A direct calculation gives
% \begin{equation}
% \label{app:eq:nesterov-resolvent}
% \bc^\top(\bI-\bM_i)^{-1}
% =
% \frac{
% B
% \bigl(
% 1+\rho+\rho(1+2\rho)\eta\lambda_i,\,
% -2\rho(1+\rho),\,
% \rho^2[1+\rho(1-\eta\lambda_i)]
% \bigr)
% }{
% \eta\lambda_i
% \Delta_{\nesterov}(\eta\lambda_i)
% }.
% \end{equation}
% The initial coordinate covariance satisfies
% \[
% \begin{pmatrix}
% u_{i,0} & p_{i,0}\\
% p_{i,0} & r_{i,0}
% \end{pmatrix}
% \succeq0,
% \]
% and therefore
% \[
% |p_{i,0}|
% \leq
% \frac{u_{i,0}+r_{i,0}}{2}.
% \]
% Since \(\Delta_{\nesterov}\) is positive on \([0,\eta\lambda_{\max}]\), there exists \(\delta>0\) such that
% \[
% \Delta_{\nesterov}(\eta\lambda_i)\geq\delta
% \qquad
% \text{for every }i.
% \]
% All coefficients in the numerator of \eqref{app:eq:nesterov-resolvent} are uniformly bounded. Hence
% \[
% \lambda_i
% \bc^\top(\bI-\bM_i)^{-1}\bX_{i,0}
% \leq
% C(u_{i,0}+r_{i,0})
% \]
% for some \(C<\infty\) independent of \(i\). Therefore,
% \[
% \sum_{k\geq0}h_k
% \leq
% C\sum_{i\geq1}(u_{i,0}+r_{i,0})
% <\infty.
% \]
% By Tonelli's theorem,
% \[
% \sum_{m\geq0}K_m
% =
% \frac{\eta^2}{B}
% \sum_{i\geq1}
% \lambda_i^2
% \bc^\top(\bI-\bM_i)^{-1}\bb.
% \]
% The first entry of \eqref{app:eq:nesterov-resolvent} gives
% \[
% \bc^\top(\bI-\bM_i)^{-1}\bb
% =
% \frac{
% B
% \left[
% 1+\rho+\rho(1+2\rho)\eta\lambda_i
% \right]
% }{
% \eta\lambda_i
% \Delta_{\nesterov}(\eta\lambda_i)
% },
% \]
% and consequently
% \[
% \sum_{m\geq0}K_m
% =
% \mathcal R_{\nesterov}(\eta,\rho).
% \]
% Suppose first that \(\mathcal R_{\nesterov}(\eta,\rho)<1\). Summing \eqref{app:eq:nesterov-renewal} and applying Tonelli's theorem gives
% \[
% \sum_{k\geq0}S_k
% \leq
% \frac{\sum_{k\geq0}h_k}
% {1-\mathcal R_{\nesterov}(\eta,\rho)}
% <\infty.
% \]
% Taking the first component of the iterated block recursion and summing over \(i\) gives
% \[
% \sum_{i\geq1}u_{i,k}
% =
% H_k
% +
% \frac{\eta^2}{B}
% \sum_{s=0}^{k-1}
% G_{k-1-s}S_s,
% \]
% where
% \[
% H_k
% :=
% \sum_{i\geq1}
% \bb^\top\bM_i^k\bX_{i,0},
% \qquad
% G_m
% :=
% \sum_{i\geq1}
% \lambda_i\bb^\top\bM_i^m\bb.
% \]
% For every fixed \(i\), the local condition gives
% \[
% \bb^\top\bM_i^k\bX_{i,0}\to0,
% \qquad
% \bb^\top\bM_i^m\bb\to0.
% \]
% The family of local blocks is uniformly power bounded on
% \[
% 0\leq\eta\lambda_i\leq\eta\lambda_{\max},
% \]
% so dominated convergence gives
% \[
% H_k\to0,
% \qquad
% G_m\to0,
% \qquad
% \sup_{m\geq0}G_m<\infty.
% \]
% Since \((S_k)_{k\geq0}\in\ell^1\), the convolution term converges to zero. Therefore,
% \[
% \Tr(\bU_k)
% =
% \sum_{i\geq1}u_{i,k}
% \to0.
% \]
% Because \(\bR_{k+1}=\bU_k\),
% \[
% \EE\lVert\bz_k\rVert^2\to0.
% \]
% This proves sufficiency.
% We next prove necessity. If \eqref{app:eq:nesterov-local-stability} fails, then at least one local block \(\bM_i\) is not asymptotically stable. Since the isolated moment map preserves the cone of positive semidefinite covariance matrices, there exists a positive semidefinite initial covariance whose second moment does not converge to zero. Decomposing this covariance into rank-one matrices shows that at least one deterministic augmented initialization supported on this eigendirection is not strongly risk asymptotically stable. Hence the local condition is necessary.
% Now assume that the local condition holds but \(\mathcal R_{\nesterov}(\eta,\rho)>1\). For \(\mu>1\), define
% \[
% F(\mu)
% :=
% \frac{\eta^2}{B}
% \sum_{i\geq1}
% \lambda_i^2
% \bc^\top(\mu\bI-\bM_i)^{-1}\bb.
% \]
% Using
% \[
% (\mu\bI-\bM_i)^{-1}
% =
% \sum_{m\geq0}
% \mu^{-m-1}\bM_i^m,
% \]
% we see that \(F\) is continuous and strictly decreasing, with
% \[
% \lim_{\mu\downarrow1}F(\mu)
% =
% \mathcal R_{\nesterov}(\eta,\rho)
% >1,
% \qquad
% \lim_{\mu\to\infty}F(\mu)=0.
% \]
% Hence there exists \(\mu>1\) such that \(F(\mu)=1\). For each \(i\), define
% \[
% \bW_i
% :=
% \lambda_i
% (\mu\bI-\bM_i^*)^{-1}\bc.
% \]
% Then \(\bW=(\bW_i)_{i\geq1}\) is an eigenvector of the adjoint global moment operator with eigenvalue \(\mu\). Moreover, for every positive semidefinite covariance state \(\bX_i\),
% \[
% \langle\bW_i,\bX_i\rangle
% =
% \lambda_i
% \sum_{m\geq0}
% \mu^{-m-1}
% \bc^\top\bM_i^m\bX_i
% \geq0.
% \]
% Choosing a deterministic initialization such that \(\langle\bW,\bX_0\rangle>0\) gives
% \[
% \langle\bW,\bX_k\rangle
% =
% \mu^k\langle\bW,\bX_0\rangle,
% \]
% which cannot converge to zero. Hence the dynamics are not strongly risk asymptotically stable.
% It remains to consider \(\mathcal R_{\nesterov}(\eta,\rho)=1\). Define
% \[
% \bW_i
% :=
% \lambda_i(\bI-\bM_i^*)^{-1}\bc.
% \]
% The local condition and \eqref{app:eq:nesterov-resolvent} imply that \(\bW\) is bounded, while \(\mathcal R_{\nesterov}(\eta,\rho)=1\) gives
% \[
% \bM_{\nesterov}^*\bW=\bW,
% \]
% where \(\bM_{\nesterov}\) denotes the global diagonal moment operator. Choosing an initialization with \(\langle\bW,\bX_0\rangle>0\) gives
% \[
% \langle\bW,\bX_k\rangle
% =
% \langle\bW,\bX_0\rangle
% >0
% \]
% for every \(k\), contradicting convergence of the augmented covariance to zero. Hence the critical boundary is not strongly risk asymptotically stable.
% \end{proof}

% \begin{corollary}[Critical learning rate of Nesterov]
% \label{app:cor:nesterov-critical-rate}
% There exists a unique critical learning rate \(\eta_{\nesterov}^{\crit}\) satisfying
% \[
% \mathcal R_{\nesterov}
% \bigl(
% \eta_{\nesterov}^{\crit},
% \rho
% \bigr)
% =1
% \]
% before the local stability boundary. The homogeneous Nesterov dynamics are risk stable if and only if
% \[
% 0<\eta<\eta_{\nesterov}^{\crit}.
% \]
% Under \(\lambda_i\eqsim i^{-\beta}\) with fixed \(\beta>1\), in the scaling regime \(B\gg1\) and \(1-\rho\ll1\),
% \begin{equation}
% \label{app:eq:nesterov-critical-scaling}
% \eta_{\nesterov}^{\crit}
% \eqsim
% \min\{1,B^\beta(1-\rho)\}.
% \end{equation}
% \end{corollary}

% \begin{proof}
% Since \(\Delta_{\nesterov}\) is a concave quadratic with \(\Delta_{\nesterov}(0)>0\), the local condition
% \[
% \Delta_{\nesterov}(\eta\lambda_{\max})>0
% \]
% is equivalent to
% \[
% \eta\lambda_{\max}<x_+(B,\rho),
% \]
% where \(x_+(B,\rho)\) is the unique positive root of \(\Delta_{\nesterov}\). In the regime \(B\gg1\) and \(\rho\to1^-\),
% \[
% x_+(B,\rho)\eqsim1.
% \]
% Indeed,
% \[
% x_+(B,1)
% =
% \frac{4B-2}{3(B+1)}
% \to
% \frac{4}{3}.
% \]
% Since \(\lambda_{\max}\eqsim1\), the local condition gives
% \[
% \eta\lesssim1.
% \]
% On every fixed subinterval strictly inside the local stability region,
% \[
% \mathcal R_{\nesterov}(\eta,\rho)
% \eqsim
% \frac{1}{B}
% \sum_{i\geq1}
% \frac{\eta\lambda_i}
% {1-\rho+\eta\lambda_i}.
% \]
% Under \(\lambda_i\eqsim i^{-\beta}\),
% \[
% \sum_{i\geq1}
% \frac{\eta\lambda_i}
% {1-\rho+\eta\lambda_i}
% \eqsim
% \begin{cases}
% \dfrac{\eta}{1-\rho},
% &
% \eta\lesssim1-\rho,
% \\[3mm]
% \left(
% \dfrac{\eta}{1-\rho}
% \right)^{1/\beta},
% &
% \eta\gtrsim1-\rho.
% \end{cases}
% \]
% In the first regime, the critical equation \(\mathcal R_{\nesterov}(\eta,\rho)\eqsim1\) would require
% \[
% \frac{\eta}{1-\rho}\eqsim B,
% \]
% which contradicts
% \[
% \frac{\eta}{1-\rho}\lesssim1
% \]
% when \(B\gg1\). Hence the largest stable learning rate lies in the second regime, where the critical equation becomes
% \[
% \left(
% \frac{\eta}{1-\rho}
% \right)^{1/\beta}
% \eqsim B.
% \]
% Therefore,
% \[
% \eta\eqsim B^\beta(1-\rho).
% \]
% Combining this global condition with the local condition \(\eta\lesssim1\) gives
% \[
% \eta_{\nesterov}^{\crit}
% \eqsim
% \min\{1,B^\beta(1-\rho)\}.
% \]
% \end{proof}

% \newpage
\section{SGD scaling law}
\label{app:sgd_scaling_law}

\begin{theorem}[SGD scaling law \citep{li2026functional}]
\label{thm:sgd_scaling_law}
Under Assumptions~\ref{ass:diagonal-covariance} and~\ref{ass:power-law}, if the stability condition \(\eta\lesssim1\) holds, then for $k \gtrsim 1/\eta$, the expected excess risk satisfies:
\begin{equation}
\label{eq:sgd loss}
\EE[\cE(\btheta_k)]
\eqsim
\left(\eta k\right)^{-s}
+
\frac{\eta\sigma^2}{B}.
\end{equation}
\end{theorem}

\begin{proof}
\citet[Theorem 5.2]{li2026functional} establishes that for a finite-width model of capacity $M$, the expected excess risk under constant learning rate $\eta$ and batch size $B$ scales as:
\[
    \EE[\cE(\btheta_k)] \eqsim M^{-s\beta} + (\eta k)^{-s} + \frac{\eta\sigma^2}{B}.
\]
Our formulation (Section~\ref{sec:plkr}) considers the infinite-dimensional Hilbert space $\mathcal{H}$. Taking the limit $M \to \infty$, the approximation error $M^{-s\beta} \to 0$. Under the matching noise covariance structure (Proposition~\ref{prop:noise-geometry}), the result immediately follows.
\end{proof}







% \newpage
% \section{SGD scaling law}
% \label{app:sgd_scaling_law}

% We provide the complete proof of the SGD scaling law. Following Appendix~\ref{app:hbm scaling law}, we decompose the parameter error into deterministic signal and stochastic noise contributions, then bound the deterministic bias, convolution kernel, label-noise floor, and multiplicative noise over distinct temporal stages.

% \subsection{Signal--noise decomposition} 
% \label{app:sgd_signal_noise_decomp}

% Let 
% \[
% \be_k := \btheta_k - \btheta^*
% \]
% denote the parameter error. 
% Using the gradient decomposition 
% \[
% \bg(\btheta_k; \cB_k) = \bH\be_k + \bxi_k(\btheta_k),
% \]
% we can subtract $\btheta^*$ from both sides of the SGD update rule \eqref{eq:sgd-update} to obtain the error recursion:
% \[
%     \be_{k+1} = (\bI - \eta \bH)\be_k - \eta \bxi_k(\btheta_k).
% \]

% In the eigenbasis of $\bH = \diag(\lambda_1,\lambda_2,\ldots)$, the $j$-th coordinate $e_k(j)$ satisfies:
% \[
%     e_{k+1}(j) = (1-\eta\lambda_j)e_k(j) - \eta\xi_k(j),
% \]
% where $\xi_k(j)$ is the $j$-th coordinate of $\bxi_k(\btheta_k)$.

% The discrete Green function is:
% \[
%     G_k(j) := (1-\eta\lambda_j)^k, \qquad k \ge 0.
% \]

% Unrolling from $e_0(j) = -\theta_j^*$ yields the exact impulse-response representation:
% \begin{equation} \label{eq:sgd_impulse_response}
%     e_k(j)
%     =
%     e_k^{\mathrm{det}}(j)
%     -
%     \eta\sum_{m=0}^{k-1}G_{k-1-m}(j)\xi_m(j),
% \end{equation}
% where $e_k^{\mathrm{det}}(j) = -(1-\eta\lambda_j)^k \theta_j^*$ is the unperturbed deterministic trajectory. We define the deterministic bias and convolution kernel as:
% \[
%     \cS_{\sgd}(k) := \frac{1}{2}\sum_{j\ge1}\lambda_j\,e_k^{\mathrm{det}}(j)^2,
%     \qquad
%     \mathcal K_{\sgd}(k) := \sum_{j\ge1}\lambda_j^2\,G_k(j)^2.
% \]

% \begin{proposition}[Risk recursion for SGD]
% \label{prop:sgd_risk_recursion}
% Under the gradient-noise covariance assumption, the expected excess risk obeys the two-sided recursion:
% \begin{equation} \label{eq:sgd_risk_recursion_form}
%     \EE[\cE(\btheta_k)]
%     \eqsim
%     \cS_{\sgd}(k)
%     +
%     \frac{\eta^2\sigma^2}{B}\sum_{\tau=0}^{k-1}\mathcal{K}_{\sgd}(\tau)
%     +
%     \frac{\eta^2}{B}\sum_{\tau=0}^{k-1}\mathcal{K}_{\sgd}(\tau)\EE[\cE(\btheta_{k-1-\tau})].
% \end{equation}
% \end{proposition}
% \begin{proof}
% Substituting \eqref{eq:sgd_impulse_response} into the excess risk, since the gradient noise is conditionally mean-zero and independent across time, cross-terms vanish:
% \[
%     \EE[e_k(j)^2]
%     =
%     e_k^{\mathrm{det}}(j)^2
%     +
%     \eta^2\sum_{m=0}^{k-1}G_{k-1-m}(j)^2\EE[\xi_m(j)^2].
% \]
% Using Proposition~\ref{prop:noise-geometry} for the variance $\EE[\xi_m(j)^2] \eqsim \frac{1}{B}\lambda_j(\EE[\cE(\btheta_m)] + \sigma^2)$, multiplying by $\frac{1}{2}\lambda_j$, and summing over $j$ yields the stated recursion. 
% \end{proof}

% \subsection{Deterministic risk calculation}
% \label{app:sgd_bias_kernel}

% \begin{proposition}[Deterministic bias of SGD]
% \label{prop:sgd_deterministic_bias}
% Under Assumptions \ref{ass:diagonal-covariance} and \ref{ass:power-law}, assuming the local stability condition $\eta\lambda_1 \le c_0$ for a sufficiently small constant $c_0 \in (0, 1)$, we have:
% \begin{equation*}
%     \cS_{\sgd}(k) \eqsim 
%     \begin{cases}
%     1, & \qquad k \lesssim \frac{1}{\eta}, \\
%     (\eta k)^{-s}, & \qquad k \gtrsim \frac{1}{\eta}.
%     \end{cases}
% \end{equation*}
% \end{proposition}
% \begin{proof}
% By definition, $\cS_{\sgd}(k) = \frac{1}{2}\sum_{j\ge 1} \lambda_j (\theta_j^*)^2 (1-\eta\lambda_j)^{2k}$. Using $\lambda_j (\theta_j^*)^2 \eqsim j^{-1-\beta s}$ and $\lambda_j \eqsim j^{-\beta}$, we evaluate the sum in two stages.

% \textbf{Stage 1: $k \lesssim \frac{1}{\eta}$.} \\
% For $j = \mathcal{O}(1)$, we have $\eta\lambda_j k \lesssim 1$. Thus,
% \[
%     (1-\eta\lambda_j)^{2k} \ge (1-\eta\lambda_1)^{2k} \eqsim \exp(-2\eta\lambda_1 k) \eqsim 1.
% \]
% Since $(1-\eta\lambda_j)^{2k} \le 1$, we have $(1-\eta\lambda_j)^{2k} \eqsim 1$. The bias is dominated by these modes:
% \[
%     \cS_{\sgd}(k) \eqsim \sum_{j\ge 1} j^{-1-\beta s} \eqsim 1,
% \]
% which converges since $\beta s > 0$.

% \textbf{Stage 2: $k \gtrsim \frac{1}{\eta}$.} \\
% Using $(1-\eta\lambda_j)^{2k} \eqsim \exp(-2\eta k j^{-\beta})$:
% \[
%     \cS_{\sgd}(k) \eqsim \sum_{j\ge 1} j^{-1-\beta s} \exp(-2\eta k j^{-\beta}).
% \]
% For the lower bound, restricting to $j \ge (\eta k)^{1/\beta}$ gives $\eta k j^{-\beta} \le 1$, yielding:
% \[
% \sum_{j \ge (\eta k)^{1/\beta}} j^{-1-\beta s} \eqsim (\eta k)^{-s}.
% \]
% For the upper bound, letting $u = 2\eta k x^{-\beta}$:
% \[
%     \int_1^\infty x^{-1-\beta s} \exp(-2\eta k x^{-\beta}) dx = \frac{1}{\beta} (2\eta k)^{-s} \int_0^{2\eta k} u^{s-1} e^{-u} du \lesssim (\eta k)^{-s}.
% \]
% Thus, $\cS_{\sgd}(k) \eqsim (\eta k)^{-s}$.
% \end{proof}







% % \begin{lemma}[Convolution Kernel of SGD]
% % \label{lem:sgd_kernel}
% % Under the same analytical assumptions, the SGD convolution kernel scales as:
% % \begin{equation*}
% %     \mathcal{K}_{\sgd}(k) \eqsim 
% %     \begin{cases}
% %     1, & \qquad k \lesssim \frac{1}{\eta}, \\
% %     (\eta k)^{-\left(2-\frac{1}{\beta}\right)}, & \qquad k \gtrsim \frac{1}{\eta}.
% %     \end{cases}
% % \end{equation*}
% % \end{lemma}
% % \begin{proof}
% % By definition, $\mathcal{K}_{\sgd}(k) \eqsim \sum_{j\ge 1} j^{-2\beta} (1-\eta\lambda_j)^{2k}$.

% % \textbf{Stage 1: $k \lesssim \frac{1}{\eta}$.} \\
% % Since $(1-\eta\lambda_j)^{2k} \eqsim 1$ for leading modes and $\beta > 1$:
% % \[
% %     \mathcal{K}_{\sgd}(k) \eqsim \sum_{j\ge 1} j^{-2\beta} \eqsim 1.
% % \]

% % \textbf{Stage 2: $k \gtrsim \frac{1}{\eta}$.} \\
% % Using an integral approximation with $u = 2\eta k x^{-\beta}$:
% % \[
% %     \sum_{j\ge 1} j^{-2\beta} e^{-2\eta k j^{-\beta}} \eqsim \int_1^\infty x^{-2\beta} e^{-2\eta k x^{-\beta}} dx = \frac{1}{\beta} (2\eta k)^{-\left(2-\frac{1}{\beta}\right)} \int_0^{2\eta k} u^{1-\frac{1}{\beta}} e^{-u} du.
% % \]
% % Since $\beta > 1$, $\int_0^{\infty} u^{1-1/\beta} e^{-u} du = \Gamma(2 - 1/\beta) \eqsim 1$, yielding $\mathcal{K}_{\sgd}(k) \eqsim (\eta k)^{-\left(2-\frac{1}{\beta}\right)}$.
% % \end{proof}

% % \subsection{Multiplicative noise and total loss}
% % \label{app:sgd_noise_bounds}

% % \begin{lemma}[Label-Noise Contribution]
% % \label{lem:sgd_label_noise}
% % For $k \gtrsim \frac{1}{\eta}$, the accumulated label-noise contribution establishes the additive noise floor:
% % \begin{equation*}
% %     \frac{\eta^2 \sigma^2}{B} \sum_{\tau=0}^{k-1} \mathcal{K}_{\sgd}(\tau) \eqsim \frac{\eta \sigma^2}{B}.
% % \end{equation*}
% % \end{lemma}
% % \begin{proof}
% % The saturated sum for each coordinate is:
% % \[
% %     \sum_{\tau=0}^{\infty} G_\tau(j)^2 = \frac{1}{1 - (1-\eta\lambda_j)^2} = \frac{1}{2\eta\lambda_j - \eta^2\lambda_j^2} \eqsim \frac{1}{\eta\lambda_j},
% % \]
% % assuming the stability condition $\eta\lambda_j \ll 1$. Summing over all coordinates gives:
% % \[
% %     \frac{\eta^2 \sigma^2}{B} \sum_{j\ge 1} \lambda_j^2 \sum_{\tau=0}^{\infty} G_\tau(j)^2 \eqsim \frac{\eta^2 \sigma^2}{B} \sum_{j\ge 1} \frac{\lambda_j^2}{\eta\lambda_j} = \frac{\eta \sigma^2}{B} \sum_{j\ge 1} \lambda_j \eqsim \frac{\eta \sigma^2}{B},
% % \]
% % since $\sum \lambda_j \eqsim 1$ for $\beta > 1$. For $k \gtrsim 1/\eta$, the finite sum captures a constant fraction of this infinite series, yielding $\frac{\eta\sigma^2}{B}$.
% % \end{proof}

% % \begin{lemma}[Deterministic Convolution for SGD]
% % \label{lem:sgd_conv}
% % Under the stability condition $\eta \lesssim 1$, for $k \gtrsim \frac{1}{\eta}$, the convolution between the deterministic bias sequence and the SGD kernel perfectly satisfies:
% % \begin{equation*}
% %     \sum_{m=0}^{k-1} \cS_{\sgd}(m) \mathcal{K}_{\sgd}(k-1-m) \eqsim \frac{1}{\eta} (\eta k)^{-\left(s \wedge \left(2-\frac{1}{\beta}\right)\right)}.
% % \end{equation*}
% % \end{lemma}
% % \begin{proof}
% % We evaluate the discrete convolution $\sum_{m=0}^{k-1} \cS_{\sgd}(m) \mathcal{K}_{\sgd}(k-1-m)$ by systematically splitting the summation into four distinct temporal stages based on the piecewise asymptotic behaviors of $\cS_{\sgd}$ and $\mathcal{K}_{\sgd}$.

% % \textbf{Stage 1: $0 \le m \lesssim \frac{1}{\eta}$.} \\
% % In this initial phase, the deterministic bias has not yet decayed, giving $\cS_{\sgd}(m) \eqsim 1$. Since $k \gtrsim 1/\eta$, the kernel index maps to $n = k-1-m \eqsim k \gtrsim 1/\eta$, placing the kernel in its polynomial decay regime: $\mathcal{K}_{\sgd}(n) \eqsim (\eta k)^{-\left(2-\frac{1}{\beta}\right)}$. Summing exactly over this discrete sequence of length $\approx 1/\eta$ outputs:
% % \[
% %     \sum_{m=0}^{1/\eta} 1 \cdot (\eta k)^{-\left(2-\frac{1}{\beta}\right)} \eqsim \frac{1}{\eta} (\eta k)^{-\left(2-\frac{1}{\beta}\right)}.
% % \]

% % \textbf{Stage 2: $\frac{1}{\eta} \lesssim m \le \frac{k}{2}$.} \\
% % Moving deeper into the iteration count, both sequences are in their polynomial regimes: $\cS_{\sgd}(m) \eqsim (\eta m)^{-s}$, and $\mathcal{K}_{\sgd}(k-1-m) \eqsim (\eta k)^{-\left(2-\frac{1}{\beta}\right)}$ since $k-1-m \eqsim k$. The convolution overlap extracts to:
% % \[
% %     \sum_{m=1/\eta}^{k/2} (\eta m)^{-s} (\eta k)^{-\left(2-\frac{1}{\beta}\right)} = \eta^{-s} (\eta k)^{-\left(2-\frac{1}{\beta}\right)} \sum_{m=1/\eta}^{k/2} m^{-s}.
% % \]
% % To evaluate the sub-polynomial sum over $m$, we perform a detailed case discussion based on the source exponent $s$:
% % \begin{itemize}
% %     \item \textbf{When $s < 1$:} The sum diverges and is dominated by its upper limit, computing as $\sum_{m=1/\eta}^{k/2} m^{-s} \eqsim k^{1-s}$. The corresponding sequence term becomes:
% %     \[
% %         \eta^{-s} (\eta k)^{-\left(2-\frac{1}{\beta}\right)} k^{1-s} = \frac{1}{\eta} (\eta k)^{1-s-\left(2-\frac{1}{\beta}\right)}.
% %     \]
% %     Since the capacity condition guarantees $\beta > 1$, we explicitly have $2-\frac{1}{\beta} > 1$. This implies $1-s-\left(2-\frac{1}{\beta}\right) < -s$. Considering $\eta k \gtrsim 1$ in this long-time regime, the exponent is strictly more negative than $-s$, meaning this cross-boundary term is provably bounded by $\frac{1}{\eta}(\eta k)^{-s}$ (which will be structurally generated in Stage 3).
    
% %     \item \textbf{When $s > 1$:} The sum is convergent and is strictly dominated by its lower limit evaluation boundary point: $\sum_{m=1/\eta}^{k/2} m^{-s} \eqsim (1/\eta)^{1-s} = \eta^{s-1}$. The sequence item simplifies down to:
% %     \[
% %         \eta^{-s} (\eta k)^{-\left(2-\frac{1}{\beta}\right)} \eta^{s-1} = \frac{1}{\eta} (\eta k)^{-\left(2-\frac{1}{\beta}\right)}.
% %     \]
    
% %     \item \textbf{When $s = 1$:} The sum yields a logarithmic factor $\ln(\eta k)$. Since $2-\frac{1}{\beta} > 1$, $\frac{1}{\eta} (\eta k)^{-\left(2-\frac{1}{\beta}\right)} \ln(\eta k)$ is asymptotically absorbed and bounded by $\frac{1}{\eta} (\eta k)^{-1}$.
% % \end{itemize}
% % Across all cases, the contribution from Stage 2 is rigorously bounded by $\frac{1}{\eta} (\eta k)^{-\left(s \wedge \left(2-\frac{1}{\beta}\right)\right)}$.

% % \textbf{Stage 3: $\frac{k}{2} < m \le k - \frac{1}{\eta}$.} \\
% % Entering the reversed proximity condition where $m \eqsim k$, we have $\cS_{\sgd}(m) \eqsim (\eta k)^{-s}$. Introducing the index transformation $n = k-1-m$, $n$ ranges inversely from $1/\eta$ up to $k/2$. Here, the kernel resides in its active decay regime: $\mathcal{K}_{\sgd}(n) \eqsim (\eta n)^{-\left(2-\frac{1}{\beta}\right)}$.
% % \begin{align*}
% %     \sum_{n=1/\eta}^{k/2} \cS_{\sgd}(k-1-n) \mathcal{K}_{\sgd}(n) 
% %     &\eqsim (\eta k)^{-s} \eta^{-\left(2-\frac{1}{\beta}\right)} \sum_{n=1/\eta}^{k/2} n^{-\left(2-\frac{1}{\beta}\right)}.
% % \end{align*}
% % Because of the restriction $\beta > 1 \implies 2-\frac{1}{\beta} > 1$, the secondary sum mapped over $n$ is convergent and identically subsumed by its starting lower boundary node evaluation: $\sum n^{-\left(2-\frac{1}{\beta}\right)} \eqsim (1/\eta)^{1-\left(2-\frac{1}{\beta}\right)} = \eta^{1-\frac{1}{\beta}}$. Substituting this back, the fragment equates to:
% % \[
% %     (\eta k)^{-s} \eta^{-\left(2-\frac{1}{\beta}\right)} \eta^{1-\frac{1}{\beta}} = \eta^{-1} (\eta k)^{-s} = \frac{1}{\eta} (\eta k)^{-s}.
% % \]

% % \textbf{Stage 4: $k - \frac{1}{\eta} < m \le k-1$.} \\
% % In this final interval, $m \eqsim k$, yielding $\cS_{\sgd}(m) \eqsim (\eta k)^{-s}$. The transformed kernel index $n = k-1-m$ ranges strictly from $0$ up to $1/\eta$. In this short-time regime for the kernel, the exponential decay has not yet activated, giving $\mathcal{K}_{\sgd}(n) \eqsim 1$.
% % Summing this constant structural factor over the discrete interval of length $\approx 1/\eta$ yields:
% % \[
% %     \sum_{n=0}^{1/\eta} (\eta k)^{-s} \cdot 1 \eqsim \frac{1}{\eta} (\eta k)^{-s}.
% % \]

% % To summarize and resolve the full sequence sum framework, the cumulative total overlap convolution restricts exclusively to tracing the maximum bound defined natively across the four stage-wise evaluations. Stages 1 and 2 contribute the $\frac{1}{\eta} (\eta k)^{-\left(2-\frac{1}{\beta}\right)}$ term, while Stages 3 and 4 rigidly define the $\frac{1}{\eta} (\eta k)^{-s}$ term. Summing them seamlessly establishes:
% % \[
% %     \sum_{m=0}^{k-1} \cS_{\sgd}(m) \mathcal{K}_{\sgd}(k-1-m) \eqsim \frac{1}{\eta} (\eta k)^{-\left(2-\frac{1}{\beta}\right)} + \frac{1}{\eta} (\eta k)^{-s} \eqsim \frac{1}{\eta} (\eta k)^{-\left(s \wedge \left(2-\frac{1}{\beta}\right)\right)}.
% % \]
% % \end{proof}

% % \begin{proposition}[Multiplicative Noise Bound for SGD]
% % \label{prop:sgd_self_noise_bounds}
% % Under the stability condition $\eta \lesssim 1$, the aggregate multiplicative noise feedback $\mathcal{M}_k = \frac{\eta^2}{B} \sum_{m=0}^{k-1} \mathcal{K}_{\sgd}(k-1-m) \EE[\cE(\btheta_m)]$ fundamentally scales as:
% % \begin{equation*}
% %     \mathcal{M}_k \eqsim \frac{\eta}{B} (\eta k)^{-\left(s \wedge \left(2-\frac{1}{\beta}\right)\right)} + \frac{\eta^2 \sigma^2}{B^2}.
% % \end{equation*}
% % \end{proposition}
% % \begin{proof}
% % Using the risk recursion $\EE[\cE(\btheta_m)] \le C \left[ \cS_{\sgd}(m) + \frac{\eta^2\sigma^2}{B}\sum_{r=0}^{m-1} \mathcal{K}_{\sgd}(r) + \mathcal{M}_m \right]$, we decompose $\mathcal{M}_k$:
% % \begin{align*}
% % \mathcal{M}_k \le\quad & C\frac{\eta^2}{B} \sum_{m=1}^{k} \mathcal{K}_{\sgd}(k-m)\cS_{\sgd}(m) \\
% % &+ C\frac{\eta^4\sigma^2}{B^2} \sum_{m=1}^{k} \mathcal{K}_{\sgd}(k-m) \sum_{r=0}^{m-1} \mathcal{K}_{\sgd}(r) \\
% % &+ C\frac{\eta^2}{B} \sum_{m=1}^{k} \mathcal{K}_{\sgd}(k-m)\mathcal{M}_m.
% % \end{align*}
% % \textbf{1. Label-Noise Convolution:} Using $\sum_{n=0}^\infty \mathcal{K}_{\sgd}(n) \eqsim 1/\eta$:
% % \[
% % \frac{\eta^4\sigma^2}{B^2} \sum_{n=0}^{k-1} \mathcal{K}_{\sgd}(n) \sum_{r=0}^{k-1-n} \mathcal{K}_{\sgd}(r) \le \frac{\eta^4\sigma^2}{B^2} \left(\sum_{n=0}^\infty \mathcal{K}_{\sgd}(n)\right)^2 \eqsim \frac{\eta^2\sigma^2}{B^2}. 
% % \]
% % \textbf{2. Deterministic Convolution:} By Lemma~\ref{lem:sgd_conv}:
% % \[
% %     \frac{\eta^2}{B} \left( \frac{1}{\eta} (\eta k)^{-\left(s \wedge \left(2-\frac{1}{\beta}\right)\right)} \right) = \frac{\eta}{B} (\eta k)^{-\left(s \wedge \left(2-\frac{1}{\beta}\right)\right)}.
% % \]
% % \textbf{3. Self-Feedback Convolution:} Using $\sum_{r=0}^n \mathcal{K}_{\sgd}(n-r)\mathcal{K}_{\sgd}(r) \le C \frac{1}{\eta} \mathcal{K}_{\sgd}(n)$:
% % \[
% %     \frac{\eta^2}{B} \sum_{m=1}^{k} \mathcal{K}_{\sgd}(k-m)\mathcal{M}_m \le C\frac{\eta^2}{B} \frac{1}{\eta} \mathcal{M}_k = C\frac{\eta}{B} \mathcal{M}_k.
% % \]
% % For $\eta \lesssim 1$ and $B \ge 1$, choosing the constants appropriately ensures $C\eta/B \le 1/2$. Thus, this feedback term can be absorbed into the left-hand side.
% % \end{proof}

% % \begin{theorem}[SGD Scaling Law]
% % \label{thm:sgd_scaling_law}
% % Assume that the primary structural stability condition \(\eta\lesssim1\) holds for normal SGD execution. For temporal parameter mapping extending across long-timescale horizons where $k \gtrsim 1/\eta$, we rigorously establish:
% % \begin{equation}
% % \label{eq:sgd loss}
% % \EE[\cE(\btheta_k)]
% % \eqsim
% % \underbrace{
% % \left(\eta k\right)^{-s}
% % }_{\text{signal}}
% % +
% % \underbrace{
% % \frac{\eta\sigma^2}{B}
% % }_{\text{additive noise}}.
% % \end{equation}
% % \end{theorem}
% % \begin{proof}
% % By Proposition~\ref{prop:sgd_risk_recursion}, the expected risk decomposes as:
% % \[
% %     \EE[\cE(\btheta_k)] \eqsim \cS_{\sgd}(k) + \frac{\eta^2\sigma^2}{B}\sum_{\tau=0}^{k-1}\mathcal{K}_{\sgd}(\tau) + \mathcal{M}_k.
% % \]
% % Substituting the asymptotic bounds for $k \gtrsim \frac{1}{\eta}$:
% % \begin{enumerate}
% %     \item \textbf{Deterministic Bias}: $\cS_{\sgd}(k) \eqsim (\eta k)^{-s}.$
% %     \item \textbf{Label Noise}: $\text{LN}(k) := \frac{\eta^2\sigma^2}{B}\sum_{\tau=0}^{k-1}\mathcal{K}_{\sgd}(\tau) \eqsim \frac{\eta\sigma^2}{B}.$
% %     \item \textbf{Multiplicative Noise}: $\mathcal{M}_k \eqsim \frac{\eta}{B} (\eta k)^{-\left(s \wedge (2-\frac{1}{\beta})\right)} + \frac{\eta^2\sigma^2}{B^2}.$
% % \end{enumerate}

% % We now absorb the multiplicative noise fragments. First, the noise floor is absorbed by the label noise since $\eta/B \lesssim 1$:
% % \[
% %     \frac{\eta^2\sigma^2}{B^2} = \left(\frac{\eta}{B}\right) \frac{\eta\sigma^2}{B} \lesssim \frac{\eta\sigma^2}{B}.
% % \]
% % Next, we bound $M_2(k) := \frac{\eta}{B} (\eta k)^{-\left(s \wedge \left(2-\frac{1}{\beta}\right)\right)}$.
% % \begin{itemize}
% %     \item If $s \le 2-\frac{1}{\beta}$, then 
% %     \[
% %         M_2(k) = \frac{\eta}{B} (\eta k)^{-s} \lesssim (\eta k)^{-s} \eqsim \cS_{\sgd}(k).
% %     \]
% %     \item If $s > 2-\frac{1}{\beta}$, then $M_2(k) = \frac{\eta}{B} (\eta k)^{-\left(2-\frac{1}{\beta}\right)}$.
% %     When the bias dominates the label noise $\left((\eta k)^{-s} \ge \frac{\eta}{B}\right)$, we have:
% %     \[
% %         M_2(k) = \frac{\eta}{B} (\eta k)^{s-(2-\frac{1}{\beta})} (\eta k)^{-s} \le \left(\frac{\eta}{B}\right)^{\frac{2-\frac{1}{\beta}}{s}} (\eta k)^{-s} \lesssim \cS_{\sgd}(k).
% %     \]
% %     When the label noise dominates $\left((\eta k)^{-s} \le \frac{\eta}{B}\right)$, since $\eta k \gtrsim 1$, we have $(\eta k)^{-\left(2-\frac{1}{\beta}\right)} \lesssim 1$, yielding:
% %     \[
% %         M_2(k) \lesssim \frac{\eta}{B} \eqsim \text{LN}(k).
% %     \]
% % \end{itemize}
% % In all cases, the multiplicative term is bounded by $\max(\cS_{\sgd}(k), \text{LN}(k))$. Thus, it is absorbed, leaving the final scaling law.
% % \end{proof}


% % \begin{theorem}[SGD Scaling Law, Noiseless Case]
% % \label{thm:sgd_scaling_law_noiseless}
% % Assume that the primary structural stability condition \(\eta\lesssim1\) holds for normal SGD execution. For temporal parameter mapping extending across long-timescale horizons where \(k \gtrsim 1/\eta\), and assuming \(\sigma=0\), the expected excess risk satisfies:
% % \begin{equation}
% % \EE[\cE(\btheta_k)]
% % \eqsim
% % \begin{cases}
% % \underbrace{ \left(\eta k\right)^{-s} }_{\text{signal}}, & 0 < s \le 2-\frac{1}{\beta}, \\[12pt]
% % \underbrace{ \left(\eta k\right)^{-s} }_{\text{signal}} + \underbrace{ \frac{\eta}{B}\left(\eta k\right)^{-(2-\frac{1}{\beta})} }_{\text{multiplicative noise}}, & s > 2-\frac{1}{\beta}.
% % \end{cases}
% % \end{equation}
% % \end{theorem}

% % \begin{proof}
% % By Proposition~\ref{prop:sgd_risk_recursion}, setting \(\sigma=0\) eliminates the additive label-noise accumulation entirely. The expected risk reduces precisely to the summation of the deterministic bias and the multiplicative noise feedback:
% % \[
% %     \EE[\cE(\btheta_k)] \eqsim \cS_{\sgd}(k) + \mathcal{M}_k.
% % \]
% % For the long-time regime \(k \gtrsim 1/\eta\), Proposition~\ref{prop:sgd_deterministic_bias} provides the deterministic bias \(\cS_{\sgd}(k) \eqsim (\eta k)^{-s}\).
% % Furthermore, Proposition~\ref{prop:sgd_self_noise_bounds} evaluated at \(\sigma=0\) yields the multiplicative noise \(\mathcal{M}_k \eqsim \frac{\eta}{B} (\eta k)^{-\left(s \wedge (2-\frac{1}{\beta})\right)}\).

% % We systematically analyze the structural absorption of the multiplicative noise based on the source exponent \(s\):
% % \begin{itemize}
% %     \item When \(0 < s \le 2-\frac{1}{\beta}\), the multiplicative noise evaluates to \(\frac{\eta}{B} (\eta k)^{-s}\). \\
% %     Under the stability condition \(\eta \lesssim 1\), and \(B \ge 1\), the coefficient strictly satisfies \(\frac{\eta}{B} \lesssim 1\). Therefore, the multiplicative noise is universally bounded by \(\cS_{\sgd}(k)\).
% %     \item When \(s > 2-\frac{1}{\beta}\), the multiplicative noise evaluates to \(\frac{\eta}{B} (\eta k)^{-(2-\frac{1}{\beta})}\). 
% % \end{itemize}
% % Synthesizing these two regimes exactly produces the piecewise scaling law.
% % \end{proof}



% \subsection{Uniform boundedness and total risk}
% \label{app:sgd_noise_bounds}

% As long as the additive label noise is strictly positive ($\sigma > 0$), we can prove that the expected risk is uniformly bounded, which allows the multiplicative gradient noise to be completely absorbed by the additive label noise floor.


% \begin{lemma}[Label-noise contribution]
% \label{lem:sgd_label_noise}
% For $k \gtrsim \frac{1}{\eta}$, the accumulated label-noise contribution establishes the additive noise floor:
% \begin{equation*}
%     \frac{\eta^2 \sigma^2}{B} \sum_{\tau=0}^{k-1} \mathcal{K}_{\sgd}(\tau) \eqsim \frac{\eta \sigma^2}{B}.
% \end{equation*}
% \end{lemma}
% \begin{proof}
% The saturated sum for each coordinate is:
% \[
%     \sum_{\tau=0}^{\infty} G_\tau(j)^2 = \frac{1}{1 - (1-\eta\lambda_j)^2} = \frac{1}{2\eta\lambda_j - \eta^2\lambda_j^2} \eqsim \frac{1}{\eta\lambda_j},
% \]
% assuming the stability condition $\eta\lambda_j \ll 1$. Summing over all coordinates gives:
% \[
%     \frac{\eta^2 \sigma^2}{B} \sum_{j\ge 1} \lambda_j^2 \sum_{\tau=0}^{\infty} G_\tau(j)^2 \eqsim \frac{\eta^2 \sigma^2}{B} \sum_{j\ge 1} \frac{\lambda_j^2}{\eta\lambda_j} = \frac{\eta \sigma^2}{B} \sum_{j\ge 1} \lambda_j \eqsim \frac{\eta \sigma^2}{B},
% \]
% since $\sum \lambda_j \eqsim 1$ for $\beta > 1$. For $k \gtrsim 1/\eta$, the finite sum captures a constant fraction of this infinite series, yielding $\frac{\eta\sigma^2}{B}$.
% \end{proof}

% \begin{lemma}[Uniform boundedness of expected risk]
% \label{lem:sgd_loss_bounded}
% Suppose $\sigma > 0$. Under the stability condition $\eta \le c$ for a sufficiently small universal constant $c > 0$, the expected excess risk of SGD is uniformly bounded for all $k \ge 0$:
% \[
%     \EE[\cE(\btheta_k)] \le M,
% \]
% where $M > 0$ is a constant independent of $k$, $\eta$, and $B$.
% \end{lemma}
% \begin{proof}
% We proceed by induction on $k$. By Proposition~\ref{prop:sgd_risk_recursion}, there exist universal constants $C_1, C_2, C_3 > 0$ such that for any $k \ge 1$:
% \begin{equation} \label{eq:sgd_risk_recursion_bound}
%     \EE[\cE(\btheta_k)] \le C_1 \cS_{\sgd}(k) + C_2 \frac{\eta^2\sigma^2}{B}\sum_{\tau=0}^{k-1}\mathcal{K}_{\sgd}(\tau) + C_3 \frac{\eta^2}{B}\sum_{\tau=0}^{k-1}\mathcal{K}_{\sgd}(\tau)\EE[\cE(\btheta_{k-1-\tau})].
% \end{equation}

% First, by Proposition~\ref{prop:sgd_deterministic_bias}, the deterministic bias monotonically decays and is uniformly bounded by its initialization: $\cS_{\sgd}(k) \le \cE(\btheta_0) =: M_1$.

% Second, as shown in the derivation of Lemma~\ref{lem:sgd_label_noise}, the infinite sum of the convolution kernel is bounded by $\sum_{\tau=0}^\infty \mathcal{K}_{\sgd}(\tau) \le \frac{C_K}{\eta}$. Therefore, the label noise contribution is universally bounded:
% \[
%     C_2 \frac{\eta^2\sigma^2}{B}\sum_{\tau=0}^{k-1}\mathcal{K}_{\sgd}(\tau) \le C_2 \sigma^2 \frac{C_K \eta}{B} \le C_2 \sigma^2 C_K c =: M_2,
% \]
% where we used the stability condition $\eta \le c$ and the fact that $B \ge 1 \implies \eta/B \le c$.

% Assuming inductively $\EE[\cE(\btheta_m)] \le M$ for $m < k$, the multiplicative noise term yields:
% \[
%     C_3 \frac{\eta^2}{B}\sum_{\tau=0}^{k-1}\mathcal{K}_{\sgd}(\tau)\EE[\cE(\btheta_{k-1-\tau})] \le C_3 \frac{\eta^2}{B} \left( \frac{C_K}{\eta} \right) M = C_3 C_K \frac{\eta}{B} M \le C_3 C_K c M.
% \]
% Choosing $c \le \frac{1}{2 C_3 C_K}$ limits this term to $M/2$. Substituting into \eqref{eq:sgd_risk_recursion_bound} yields:
% \[
%     \EE[\cE(\btheta_k)] \le M_1 + M_2 + \frac{M}{2}.
% \]
% Setting $M = 2(M_1 + M_2)$ guarantees $\EE[\cE(\btheta_k)] \le M$. The base case $k=0$ holds trivially since $\EE[\cE(\btheta_0)] = \cS_{\sgd}(0) \le M_1 \le M$. This completes the induction.
% \end{proof}

% \begin{corollary}[Absorption of multiplicative noise]
% \label{cor:sgd_noise_absorption}
% For $\sigma > 0$, under the stability condition $\eta \lesssim 1$, the multiplicative gradient noise is strictly bounded by the additive label noise up to a constant factor:
% \[
%     \frac{\eta^2}{B}\sum_{\tau=0}^{k-1}\mathcal{K}_{\sgd}(\tau)\EE[\cE(\btheta_{k-1-\tau})] \le \frac{M}{\sigma^2} \left( \frac{\eta^2\sigma^2}{B}\sum_{\tau=0}^{k-1}\mathcal{K}_{\sgd}(\tau) \right).
% \]
% Consequently, the expected excess risk is asymptotically equivalent to the sum of the deterministic bias and the label noise:
% \[
%     \EE[\cE(\btheta_k)] \eqsim \cS_{\sgd}(k) + \frac{\eta^2\sigma^2}{B}\sum_{\tau=0}^{k-1}\mathcal{K}_{\sgd}(\tau).
% \]
% \end{corollary}


% \begin{theorem}[SGD scaling law]
% \label{thm:sgd_scaling_law}
% Assume that the primary structural stability condition \(\eta\lesssim1\) holds for normal SGD execution. For temporal parameter mapping extending across long-timescale horizons where $k \gtrsim 1/\eta$, we rigorously establish:
% \begin{equation}
% \label{eq:sgd loss}
% \EE[\cE(\btheta_k)]
% \eqsim
% \underbrace{
% \left(\eta k\right)^{-s}
% }_{\text{signal}}
% +
% \underbrace{
% \frac{\eta\sigma^2}{B}
% }_{\text{additive noise}}.
% \end{equation}
% \end{theorem}
% \begin{proof}
% By Corollary~\ref{cor:sgd_noise_absorption}, the expected risk is completely determined by the deterministic bias and the label noise:
% \[
%     \EE[\cE(\btheta_k)] \eqsim \cS_{\sgd}(k) + \frac{\eta^2\sigma^2}{B}\sum_{\tau=0}^{k-1}\mathcal{K}_{\sgd}(\tau).
% \]
% For $k \gtrsim \frac{1}{\eta}$, we substitute the deterministic bias from Proposition~\ref{prop:sgd_deterministic_bias}:
% \[
%     \cS_{\sgd}(k) \eqsim (\eta k)^{-s}.
% \]
% We substitute the label noise floor from Lemma~\ref{lem:sgd_label_noise}:
% \[
%     \frac{\eta^2\sigma^2}{B}\sum_{\tau=0}^{k-1}\mathcal{K}_{\sgd}(\tau) \eqsim \frac{\eta\sigma^2}{B}.
% \]
% Summing these two terms directly yields the exact scaling law.
% \end{proof}





\section{Polyak scaling law}
\label{app:hbm scaling law}

% In this section, we prove Theorem~\ref{thm:hbm scaling law} following the framework of \citet{li2026functional}, by decomposing the risk dynamics into signal learning and noise accumulation.


In this section, we establish Theorem~\ref{thm:hbm scaling law} using the Functional Scaling Law (FSL) framework of \citet{li2026functional}. Our starting point is the following expected excess risk:
\begin{align}
    \EE[\cE(\btheta_k)] 
    &\eqsim \underbrace{\cS_{\polyak}(k)}_{\text{signal learning}} +\underbrace{\sum_{\tau=0}^{k-1}\underbrace{\mathcal{K}_{\polyak}(\tau)}_{\text{forgetting kernel}} \left(\sigma^2 + \EE[\cE(\btheta_{k-1-\tau})]\right) \frac{\eta^2}{B}}_{\text{noise accumulation}} \label{eq:pol_fsl_intro_org}\\
    &\eqsim \cS_{\polyak}(k) + \underbrace{\frac{\eta^2\sigma^2}{B}\sum_{\tau=0}^{k-1}\mathcal{K}_{\polyak}(\tau)}_{\text{additive label noise}} + \underbrace{\frac{\eta^2}{B}\sum_{\tau=0}^{k-1}\mathcal{K}_{\polyak}(\tau)\EE[\cE(\btheta_{k-1-\tau})]}_{\text{multiplicative noise}}. \label{eq:pol_fsl_intro}
\end{align}

Our analysis proceeds in three steps:
\begin{itemize}
\item First, Section~\ref{app:pol_signal_noise_decomp} derives the \textbf{general discrete risk recursion} \eqref{eq:pol_fsl_intro} from an exact coordinate-wise impulse-response representation. Section~\ref{app:pol_decomp_modes} then analyzes the characteristic roots of the second-order dynamics and separates the spectrum into \textbf{overdamped and underdamped modes}.

\item Second, Sections~\ref{app:pol_preparation} and~\ref{app:pol_uniform_boundedness} deal with the \textbf{stochastic terms} in the recursion. Section~\ref{app:pol_preparation} computes the additive label noise, while Section~\ref{app:pol_uniform_boundedness} shows that, under the risk stability condition, the multiplicative noise contribution can be absorbed into the label noise term up to constant factors.

\item Finally, Sections~\ref{app:pol_fully_overdamped_loss} and ~\ref{app:pol_mixed_loss} characterize the \textbf{deterministic signal learning term} \(S_{\mathrm{polyak}}(k)\) in the \textbf{fully overdamped and mixed-damping regimes}, and therefore derive the scaling law expression.
\end{itemize}

% We systematically evaluate each functional component in the following sections:

% \textbf{Section~\ref{app:pol_signal_noise_decomp}:} Establishes the discrete risk recursion \eqref{eq:pol_fsl_intro} via exact coordinate-wise impulse-response representation.
% \textbf{Section~\ref{app:pol_decomp_modes}:} Partitions the spectrum at the crossover index $J_{\max} \eqsim \big(\frac{\eta}{(1-\rho)^2}\big)^{1/\beta}$ derived discriminant of the second-order iteration $\Delta_j=0$.
% \textbf{Section~\ref{app:pol_preparation}:} Evaluates the saturated Green energy to extract the additive label noise floor: $\frac{\eta^2\sigma^2}{B}\sum_{\tau=0}^{k-1}\mathcal{K}_{\polyak}(\tau) \eqsim \frac{\sigma^2\eta}{B(1-\rho)}$.
% \textbf{Section~\ref{app:pol_uniform_boundedness}:} Proves $\EE[\cE(\btheta_k)] \le M$ to absorb multiplicative noise, reducing \eqref{eq:pol_fsl_intro} to $\EE[\cE(\btheta_k)] \eqsim \cS_{\polyak}(k) + \frac{\sigma^2\eta}{B(1-\rho)}$.
% \textbf{Section~\ref{app:pol_fully_overdamped_loss}:} Computes the fully overdamped deterministic signal bias: $\cS_{\polyak}(k) \eqsim \big(\frac{1-\rho}{\eta k}\big)^s$.
% \textbf{Section~\ref{app:pol_mixed_loss}:} Bounds the underdamped transient to establish the mixed-regime bias: $\cS_{\polyak}(k) \eqsim \cS_{\mathrm{under}}(k) + \big(\frac{1-\rho}{\eta k}\big)^s$.


\subsection{Signal--noise decomposition and risk recursion} 
\label{app:pol_signal_noise_decomp}


Recall the parameter error 
\[
\be_k := \btheta_k - \btheta^*.
\]
Substituting the gradient decomposition into the Polyak update \eqref{eq:hb-update}, we obtain the error recursion:
\[
    \be_{k+1} = \bigl((1+\rho)\bI-\eta\bH\bigr)\be_k -\rho\be_{k-1} -\eta\bxi_k(\btheta_k).
\]

In the diagonal eigenbasis of the population Hessian $\bH = \diag(\lambda_1,\lambda_2,\ldots)$, the $j$-th coordinate of the parameter error, denoted as $e_k(j)$, obeys the following second-order scalar difference equation:
\[
    e_{k+1}(j) = (1+\rho-\eta\lambda_j)e_k(j) -\rho e_{k-1}(j) -\eta\xi_k(j),
\]
where $\xi_k(j)$ is the $j$-th coordinate of the stochastic gradient noise vector $\bxi_k(\btheta_k)$.

The homogeneous characteristic polynomial associated with the deterministic part of this coordinate-wise recursion is
\[
    y^2-(1+\rho-\eta\lambda_j)y+\rho=0.
\]

For each coordinate $j$, let $y_{j,+}$ and $y_{j,-}$ denote the two roots of this polynomial. 

We define the discrete Green function (impulse response) for the $j$-th coordinate as
\[
    G_k(j):=\frac{y_{j,+}^{k+1}-y_{j,-}^{k+1}}{y_{j,+}-y_{j,-}},
    \qquad k\ge 0,
\]
with the natural continuous limiting definition when $y_{j,+}=y_{j,-}$ (i.e., critical damping). 

By viewing the stochastic noise sequence $-\eta\xi_k(j)$ as an external forcing term, the stochastic parameter error trajectory can be decomposed exactly into a deterministic signal trajectory and a discrete convolution of the gradient noise. This exact impulse-response representation is formalized in the following proposition.

\begin{proposition}[Impulse-response representation]
\label{prop:pol_impulse_response}
For any step $k \ge 0$, the coordinate-wise error trajectory of Polyak satisfies:
\begin{equation} \label{eq:pol_impulse_response}
    e_k(j)
    =
    e_k^{\mathrm{det}}(j)
    -
    \eta\sum_{m=0}^{k-1}G_{k-1-m}(j)\xi_m(j),
\end{equation}
where $e_k^{\mathrm{det}}(j)$ is the unperturbed deterministic trajectory solving:
\[
    e_{k+1}^{\mathrm{det}}(j)
    =
    (1+\rho-\eta\lambda_j)e_k^{\mathrm{det}}(j)
    -
    \rho e_{k-1}^{\mathrm{det}}(j),
\]
with initializations $e_0^{\mathrm{det}}(j) = e_0(j) = -\theta_j^*$ and $e_{-1}^{\mathrm{det}}(j) = e_{-1}(j) = -\theta_j^*$ (assuming $\btheta_0=\btheta_{-1}=\bm0$).
\end{proposition}

\begin{proof}

Define the stochastic residual error $\tilde{e}_k(j) := e_k(j) - e_k^{\mathrm{det}}(j)$. Subtracting the deterministic trajectory from the stochastic parameter error yields the inhomogeneous system:
\begin{equation} \label{eq:inhomogeneous_system}
    \tilde{e}_{k+1}(j) - (1+\rho-\eta\lambda_j)\tilde{e}_k(j) + \rho\tilde{e}_{k-1}(j) = -\eta\xi_k(j),
\end{equation}
subject to the initial conditions:
\begin{equation*} \label{eq:zero_initial_conditions}
    \tilde{e}_0(j) = \tilde{e}_{-1}(j) = 0.
\end{equation*}

We introduce the discrete linear difference operator $\mathcal{L}_j$:
\begin{equation*}
    (\mathcal{L}_j u)_k := u_{k+1} - (1+\rho-\eta\lambda_j)u_k + \rho u_{k-1}.
\end{equation*}

By the definition of the discrete Green function $G_k(j)$:
\begin{equation*}
    (\mathcal{L}_j G(j))_k = 0 \quad (\forall k \ge 0), \qquad \text{with} \quad G_{-1}(j) = 0, \quad G_0(j) = 1.
\end{equation*}

Let the impulse response to a point source at step $m \ge 0$ be defined as $H_k(j; m) := G_{k-1-m}(j)$, which satisfies:
\begin{equation*}
    (\mathcal{L}_j H(j; m))_k = \delta_{k, m}, \qquad \text{where} \quad \delta_{k,m} = \begin{cases} 1, & k = m, \\ 0, & k \neq m. \end{cases}
\end{equation*}

Applying the superposition principle to the linear operator $\mathcal{L}_j$ under the forcing term $-\eta\xi_k(j)$ in \eqref{eq:inhomogeneous_system} yields:
\begin{align*}
    \tilde{e}_k(j) &= \sum_{m=0}^{k-1} \big(-\eta\xi_m(j)\big) H_k(j; m) \nonumber \\
    &= -\eta \sum_{m=0}^{k-1} G_{k-1-m}(j)\xi_m(j).
\end{align*}

Substituting the definition of $\tilde{e}_k(j)$ back into the equation completes the proof:
\begin{equation*}
    e_k(j) = e_k^{\mathrm{det}}(j) - \eta \sum_{m=0}^{k-1} G_{k-1-m}(j)\xi_m(j).
\end{equation*}
\end{proof}

To quantify the excess risk, we define the deterministic bias (the \textit{signal} component) and the convolution kernel as follows:
\[
    \cS_{\polyak}(k) := \frac{1}{2}\sum_{j\ge1}\lambda_j\,e_k^{\mathrm{det}}(j)^2 = \cE(\btheta_k^{\mathrm{det}}),
    \qquad
    \mathcal K_{\polyak}(k) := \sum_{j\ge1}\lambda_j^2\,G_k(j)^2.
\]

The following proposition establishes that the expected excess risk decomposes, up to constant factors, into the deterministic signal bias, the accumulated additive label noise, and the accumulated multiplicative gradient noise.

\begin{proposition}[Risk recursion for Polyak]
\label{prop:pol_risk_recursion}
Under the gradient-noise covariance assumption, the expected excess risk obeys the two-sided recursion:
\begin{equation} \label{eq:pol_risk_recursion_form}
    \EE[\cE(\btheta_k)]
    \eqsim
    \cS_{\polyak}(k)
    +
    \frac{\eta^2\sigma^2}{B}\sum_{\tau=0}^{k-1}\mathcal{K}_{\polyak}(\tau)
    +
    \frac{\eta^2}{B}\sum_{\tau=0}^{k-1}\mathcal{K}_{\polyak}(\tau)\EE[\cE(\btheta_{k-1-\tau})].
\end{equation}
Equivalently, the right-hand side provides both rigorous upper and lower bounds for the expected excess risk up to universal constant factors.
\end{proposition}

\begin{proof}

By substituting the Green function representation \eqref{eq:pol_impulse_response} into the definition of the excess risk, we analyze the second moment for each error coordinate:
\[
    e_k(j)^2
    =
    \left( e_k^{\mathrm{det}}(j)
    -\eta\sum_{m=0}^{k-1}G_{k-1-m}(j)\xi_m(j) \right)^2.
\]

Because the stochastic gradient noise is conditionally mean-zero ($\EE[\xi_m(j) \mid \btheta_m] = 0$) and independent across distinct time steps, all cross-terms between the deterministic trajectory and the noise, as well as cross-time noise products, vanish upon taking the expectation. Hence,
\[
    \EE[e_k(j)^2]
    =
    e_k^{\mathrm{det}}(j)^2
    +
    \eta^2\sum_{m=0}^{k-1}G_{k-1-m}(j)^2\EE[\xi_m(j)^2].
\]

Multiplying by $\frac{1}{2}\lambda_j$ and summing over all coordinates $j \ge 1$ gives the expected excess risk:
\begin{align*}
    \EE[\cE(\btheta_k)]
    &\eqsim
    \frac{1}{2} \sum_{j\ge1}\lambda_j e_k^{\mathrm{det}}(j)^2
    +
    \frac{\eta^2}{2}\sum_{m=0}^{k-1}\sum_{j\ge1}
    \lambda_j G_{k-1-m}(j)^2\EE[\xi_m(j)^2].
\end{align*}

According to the property of curvature-aligned noise defined in Eq.~\eqref{eqn:noise-geometry}, for the coordinate basis $\bm{u}_j$, the variance satisfies the following:
\[
    \EE[\xi_m(j)^2]
    \eqsim
    \frac{1}{B}\lambda_j\big(\EE[\cE(\btheta_m)] + \sigma^2\big),
\]
where the factor $1/B$ accounts for the mini-batching variance reduction. 

Substituting this into the previous equation yields:
\begin{align*}
    \EE[\cE(\btheta_k)]
    &\eqsim
    \frac{1}{2} \sum_{j\ge1}\lambda_j e_k^{\mathrm{det}}(j)^2
    +
    \frac{\eta^2}{B}\sum_{m=0}^{k-1}
    \big(\sigma^2+\EE[\cE(\btheta_m)]\big)
    \sum_{j\ge1}\lambda_j^2 G_{k-1-m}(j)^2  \\
    &=
    \cS_{\polyak}(k)
    +
    \frac{\eta^2\sigma^2}{B}\sum_{m=0}^{k-1}\mathcal K_{\polyak}(k-1-m)
    +
    \frac{\eta^2}{B}\sum_{m=0}^{k-1}\mathcal K_{\polyak}(k-1-m)\EE[\cE(\btheta_m)].
\end{align*}

Finally, performing a change of variables by setting $\tau = k-1-m$ yields:
\[
    \EE[\cE(\btheta_k)]
    \eqsim
    \cS_{\polyak}(k)
    +
    \frac{\eta^2\sigma^2}{B}\sum_{\tau=0}^{k-1}\mathcal{K}_{\polyak}(\tau)
    +
    \frac{\eta^2}{B}\sum_{\tau=0}^{k-1}\mathcal{K}_{\polyak}(\tau)\EE[\cE(\btheta_{k-1-\tau})].
\]

The two-sided equivalence follows directly from the two-sided comparability in the noise covariance assumption and the non-negativity of $\cS_{\polyak}(k)$ and $\mathcal K_{\polyak}(k)$.
\end{proof}

\subsection{Overdamped and underdamped modes decomposition}
\label{app:pol_decomp_modes}

Since both the deterministic error signal $e_k^{\mathrm{det}}(j)$ and the discrete Green function $G_k(j)$ obey the  same homogeneous difference equation, their dynamic behaviors are simultaneously dictated by the same characteristic roots. Depending on the sign of the corresponding discriminant, these roots transition from real to complex conjugate pairs, driving the system from monotonic decay to oscillatory transients. The following proposition rigorously establishes the damping threshold, defines the crossover index $J_{\max}$, and partitions the spectrum into two distinct global regimes.

\begin{proposition}[Exact spectral decomposition and crossover boundary]
\label{prop:pol_mode_decomposition}
For each coordinate $j$, the characteristic roots of the Polyak recursion are given by
\[
    y_{j, \pm} = \frac{(1+\rho-\eta\lambda_j) \pm \sqrt{\Delta_j}}{2},
\]
where the discriminant is $\Delta_j := (1+\rho-\eta\lambda_j)^2 - 4\rho$. 

The critical transition $\Delta_j = 0$ corresponds exactly to the spectral threshold $\eta\lambda_j = (1-\sqrt{\rho})^2$, which is asymptotically equivalent to $\eta\lambda_j \eqsim (1-\rho)^2$. This threshold classifies the individual eigen-directions into:
\begin{enumerate}
    \item \textbf{Overdamped Modes} ($\eta\lambda_j < (1-\sqrt{\rho})^2$, i.e., $\Delta_j > 0$): The roots are strictly real.
    \item \textbf{Underdamped Modes} ($\eta\lambda_j > (1-\sqrt{\rho})^2$, i.e., $\Delta_j < 0$): The roots are complex conjugates.
\end{enumerate}
In the mixed-damping regime, the boundary index $J_{\max}$ separating the underdamped modes from the overdamped modes is defined as the crossover index satisfying $\eta\lambda_{J_{\max}} \eqsim (1-\rho)^2$:
\begin{equation}
\label{eq:J_max}
    J_{\max} \eqsim \left(\frac{\eta}{(1-\rho)^2}\right)^{1/\beta}.
\end{equation}
Consequently, based on the maximum eigenvalue $\lambda_{\max}$ of the population Hessian $\bH$, the algorithm globally operates in one of two regimes:
\begin{itemize}
    \item \textbf{Fully Overdamped Regime} ($\eta\lambda_{\max} \lesssim (1-\rho)^2$): All spectral modes are overdamped.
    \item \textbf{Mixed-Damping Regime} ($\eta\lambda_{\max} \gtrsim (1-\rho)^2$): The spectrum bifurcates into high-curvature underdamped modes ($j \le J_{\max}$) and flat overdamped modes ($j > J_{\max}$).
\end{itemize}
\end{proposition}

\begin{proof}
The phase transition between real and complex roots occurs exactly at the boundary where the discriminant vanishes ($\Delta_j = 0$). Solving the quadratic equation $(1+\rho-\eta\lambda_j)^2 = 4\rho$ for the relevant stable branch yields the exact threshold:
\[
    \eta\lambda_j = 1 + \rho - 2\sqrt{\rho} = (1-\sqrt{\rho})^2.
\]
By multiplying the numerator and denominator by $(1+\sqrt{\rho})^2$, we can rewrite this exact threshold as: $\eta\lambda_j = \frac{(1-\rho)^2}{(1+\sqrt{\rho})^2}.$
Since $1+\sqrt{\rho} \eqsim 2$ for all $\rho \in (0, 1)$, we obtain the fundamental scaling equivalence for the damping boundary:
\[
    \eta\lambda_j \eqsim (1-\rho)^2.
\]
To locate the exact crossover boundary in the mixed-damping regime, we apply the power-law spectrum $\lambda_j \eqsim j^{-\beta}$. Setting the critical damping boundary condition $\eta\lambda_{J_{\max}} \eqsim (1-\rho)^2$, we obtain: 
$\eta J_{\max}^{-\beta} \eqsim (1-\rho)^2  \Longleftrightarrow  J_{\max}^{-\beta} \eqsim \frac{(1-\rho)^2}{\eta}.$
Taking the $-1/\beta$-th power of both sides yields:
\[
    J_{\max} \eqsim \left(\frac{\eta}{(1-\rho)^2}\right)^{1/\beta},
\]
which establishes that the underdamped modes are exactly indexed by $j \le J_{\max}$ and the overdamped modes by $j > J_{\max}$.

For an overdamped mode where $\eta\lambda_j < (1-\sqrt{\rho})^2$, the discriminant is strictly positive ($\Delta_j > 0$), and the roots $y_{j,+}$ and $y_{j,-}$ are real and distinct.

For an underdamped mode where $\eta\lambda_j > (1-\sqrt{\rho})^2$, the discriminant is strictly negative ($\Delta_j < 0$). The roots become a complex conjugate pair $y_{j,\pm} = \sqrt{\rho} e^{\pm i \omega_j}$, where the exact phase angle $\omega_j \in (0, \pi)$ satisfies:
$\cos\omega_j = \frac{1+\rho-\eta\lambda_j}{2\sqrt{\rho}}.$

Finally, to determine the global behavior of the system, we examine the maximum eigenvalue $\lambda_{\max}$. If $\eta\lambda_{\max} \lesssim (1-\rho)^2$, then every coordinate satisfies $\Delta_j \ge 0$ in the scaling sense, placing the system entirely in the Fully Overdamped Regime. If $\eta\lambda_{\max} \gtrsim (1-\rho)^2$, the spectral modes with $\lambda_j \gtrsim (1-\rho)^2/\eta$ form a non-empty underdamped subset, while the remaining tail $\lambda_j \lesssim (1-\rho)^2/\eta$ remains overdamped, placing the system globally in the mixed-damping regime.
\end{proof}

\subsection{Label noise calculation}
\label{app:pol_preparation}

% To evaluate the summation of exponentially decaying terms and polynomially decaying terms during the error accumulation, we first introduce the following exponential-polynomial convolution estimate.

% \begin{lemma}[Exponential-Polynomial Convolution]
% \label{lem:pol_convolution_estimates1}
% Let $a>0$, $k\ge 2$, and $0<\rho<1$. Assume
% \begin{equation}
% \label{eq:exp_poly_conv_condition}
%     (1-\rho)k\ge 2,
%     \qquad
%     (1-\rho)k\ge 2(a+1)\log k.
% \end{equation}
% Then the exponential--polynomial convolution admits the explicit two-sided bound
% \begin{equation}
% \label{eq:exp_poly_conv_explicit}
%     (1-e^{-1})
%     \frac{1}{1-\rho}\frac{1}{k^a}
%     \le
%     \sum_{m=0}^{k-1}\rho^m\left(\frac{1}{k-m}\right)^a
%     \le
%     \left(2^a+C_a\right)\frac{1}{1-\rho}\frac{1}{k^a},
% \end{equation}
% where
% \begin{equation}
% \label{eq:Ca_exp_poly_conv}
%     C_a:=
%     \begin{cases}
%     \dfrac{2}{e(1-a)}, & 0<a<1,\\[6pt]
%     1, & a=1,\\[6pt]
%     1+\dfrac{1}{a-1}, & a>1.
%     \end{cases}
% \end{equation}
% In particular, under~\eqref{eq:exp_poly_conv_condition},
% \[
%     \sum_{m=0}^{k-1}\rho^m\left(\frac{1}{k-m}\right)^a
%     \sim_{a}
%     \frac{1}{1-\rho}\left(\frac{1}{k}\right)^a .
% \]
% \end{lemma}

% \begin{proof}
% Let
% \[
%     S_k:=\sum_{m=0}^{k-1}\rho^m(k-m)^{-a},
%     \qquad
%     \delta:=1-\rho,
%     \qquad
%     z:=\delta k.
% \]
% We split the sum into
% \[
%     S_k=S_{k,1}+S_{k,2},
% \]
% where
% \[
%     S_{k,1}:=\sum_{0\le m\le k/2}\rho^m(k-m)^{-a},
%     \qquad
%     S_{k,2}:=\sum_{k/2<m\le k-1}\rho^m(k-m)^{-a}.
% \]
% Here and below, $0\le m\le k/2$ means $0\le m\le \lfloor k/2\rfloor$.

% We first bound $S_{k,1}$. On $0\le m\le k/2$, we have $\frac{k}{2}\le k-m\le k,$
% and therefore
% \[
%     k^{-a}\le (k-m)^{-a}\le 2^a k^{-a}.
% \]
% Consequently,
% \begin{align*}
%     S_{k,1}
%     &\ge
%     k^{-a}\sum_{m=0}^{\lfloor k/2\rfloor}\rho^m
%     =
%     k^{-a}\frac{1-\rho^{\lfloor k/2\rfloor+1}}{1-\rho},\\
%     S_{k,1}
%     &\le
%     2^a k^{-a}\sum_{m=0}^{\lfloor k/2\rfloor}\rho^m
%     \le
%     2^a k^{-a}\frac{1}{1-\rho}.
% \end{align*}
% Since $\rho=1-\delta\le e^{-\delta}$ and $\lfloor k/2\rfloor+1\ge k/2$, condition~\eqref{eq:exp_poly_conv_condition} gives
% \[
%     \rho^{\lfloor k/2\rfloor+1}
%     \le
%     e^{-\delta k/2}
%     =e^{-z/2}
%     \le e^{-1}.
% \]
% Thus
% \begin{equation}
% \label{eq:St1_bounds_explicit}
%     (1-e^{-1})\frac{1}{1-\rho}\frac{1}{k^a}
%     \le
%     S_{k,1}
%     \le
%     2^a\frac{1}{1-\rho}\frac{1}{k^a}.
% \end{equation}
% This already proves the desired lower bound for $S_k$.

% It remains to upper-bound $S_{k,2}$. Put $n=k-m$. Then $1\le n<k/2$ and
% \[
%     S_{k,2}
%     =
%     \sum_{1\le n<k/2}\rho^{k-n}n^{-a}.
% \]
% Since $k-n>k/2$ on this range,
% \[
%     \rho^{k-n}\le \rho^{k/2}\le e^{-\delta k/2}=e^{-z/2},
% \]
% so
% \begin{equation*}
%     S_{k,2}
%     \le
%     e^{-z/2}\sum_{1\le n<k/2}n^{-a}.
% \end{equation*}
% We now consider the three possible ranges of $a$.

% If $0<a<1$, then
% \[
%     \sum_{1\le n<k/2}n^{-a}\le \sum_{1\le n\le k}n^{-a}\le \frac{k^{1-a}}{1-a}.
% \]
% Therefore
% \[
%     S_{k,2}
%     \le
%     \frac{e^{-z/2}k^{1-a}}{1-a}
%     =
%     \frac{z e^{-z/2}}{1-a}\frac{1}{1-\rho}\frac{1}{k^a}.
% \]
% The function $z e^{-z/2}$ is bounded above by $2/e$ for $z\ge0$. Hence
% \begin{equation}
% \label{eq:St2_bound_a_less_1}
%     S_{k,2}
%     \le
%     \frac{2}{e(1-a)}\frac{1}{1-\rho}\frac{1}{k^a}.
% \end{equation}

% If $a=1$, then
% \[
%     \sum_{1\le n<k/2}n^{-1}\le 1+\log k.
% \]
% Using $z\ge 2(a+1)\log k=4\log k$ and $z=(1-\rho)k\le k$, we obtain
% \begin{align*}
%     S_{k,2}
%     &\le e^{-z/2}(1+\log k) \
%     &= z e^{-z/2}(1+\log k)\frac{1}{1-\rho}\frac{1}{k} \
%     &\le k\cdot k^{-2}(1+\log k)\frac{1}{1-\rho}\frac{1}{k}
%     \le
%     \frac{1}{1-\rho}\frac{1}{k},
% \end{align*}
% where the last step uses $1+\log k\le k$ for $k\ge2$. Thus
% \begin{equation}
% \label{eq:St2_bound_a_eq_1}
%     S_{k,2}
%     \le
%     \frac{1}{1-\rho}\frac{1}{k}.
% \end{equation}

% If $a>1$, then
% \[
%     \sum_{1\le n<k/2}n^{-a}\le 1+\int_1^\infty x^{-a}\,dx
%     =1+\frac{1}{a-1}.
% \]
% Again using $z\ge 2(a+1)\log k$ and $z\le k$, we have
% \begin{align*}
%     S_{k,2}
%     &\le
%     \left(1+\frac{1}{a-1}\right)e^{-z/2} \\
%     &=
%     \left(1+\frac{1}{a-1}\right)
%     z k^{a-1}e^{-z/2}
%     \frac{1}{1-\rho}\frac{1}{k^a} \\
%     &\le
%     \left(1+\frac{1}{a-1}\right)
%     k^a k^{-(a+1)}
%     \frac{1}{1-\rho}\frac{1}{k^a} \\
%     &\le
%     \left(1+\frac{1}{a-1}\right)
%     \frac{1}{1-\rho}\frac{1}{k^a}.
% \end{align*}
% Therefore
% \begin{equation}
% \label{eq:St2_bound_a_greater_1}
%     S_{k,2}
%     \le
%     \left(1+\frac{1}{a-1}\right)\frac{1}{1-\rho}\frac{1}{k^a}.
% \end{equation}
% Combining~\eqref{eq:St2_bound_a_less_1},~\eqref{eq:St2_bound_a_eq_1}, and~\eqref{eq:St2_bound_a_greater_1}, we have in all cases
% \[
%     S_{k,2}
%     \le
%     C_a\frac{1}{1-\rho}\frac{1}{k^a},
% \]
% with $C_a$ defined in~\eqref{eq:Ca_exp_poly_conv}. Combining this upper bound with~\eqref{eq:St1_bounds_explicit} gives
% \[
%     S_k=S_{k,1}+S_{k,2}
%     \le
%     (2^a+C_a)\frac{1}{1-\rho}\frac{1}{k^a}.
% \]
% Together with the lower bound, this proves~\eqref{eq:exp_poly_conv_explicit}.
% \end{proof}

Before bounding the full multiplicative risk, we must isolate the purely additive label-noise contribution, which establishes the noise floor for both regimes.

\begin{lemma}[Label-noise contribution]
\label{lem:pol_label_noise}
There exist constants $C,c_-,c_+>0$ such that if
\begin{equation}
\label{eq:pol_label_noise_time_condition}
    k\ge
    C
    \max\left\{
        \frac{1}{1-\rho},
        \frac{1-\rho}{\eta}
    \right\},
\end{equation}
then the finite-time label-noise contribution satisfies
\begin{equation}
\label{eq:pol_label_noise_explicit_bounds}
    c_-\frac{\sigma^2\eta}{B(1-\rho)}
    \le
    \frac{\eta^2\sigma^2}{B}\sum_{\tau=0}^{k-1}\mathcal{K}_{\polyak}(\tau)
    \le
    c_+\frac{\sigma^2\eta}{B(1-\rho)}.
\end{equation}
In particular,
\[
    \frac{\eta^2\sigma^2}{B}\sum_{\tau=0}^{k-1}\mathcal{K}_{\polyak}(\tau)
    \eqsim
    \frac{\sigma^2}{B}\frac{\eta}{1-\rho}.
\]
\end{lemma}

\begin{proof}
We first compute the saturated coordinate Green energy exactly. Fix a coordinate $j$ and write
\[
    x_j:=\eta\lambda_j,
    \qquad
    A_j:=1+\rho-x_j.
\]
The Green function satisfies
\[
    G_{k+1}(j)=A_jG_k(j)-\rho G_{k-1}(j),
    \qquad
    G_0(j)=1,
    \qquad
    G_{-1}(j)=0.
\]
Assume the coordinate is stable. Define
\[
    S_j:=\sum_{\tau=0}^\infty G_\tau(j)^2,
    \qquad
    C_j:=\sum_{\tau=0}^\infty G_\tau(j)G_{\tau-1}(j),
\]
where $G_{-1}(j)=0$. Multiplying the recursion by $G_k(j)$ and summing over $k\ge0$ gives
   $C_j=A_jS_j-\rho C_j,$ hence $C_j=\frac{A_j}{1+\rho}S_j.$

Next square the recursion and sum over $k\ge0$:
\[
    S_j-1
    =
    \left(A_j^2+\rho^2\right)S_j-2A_j\rho C_j.
\]
Substituting $C_j=A_jS_j/(1+\rho)$ gives
\begin{align*}
    1
    &=
    S_j\left[1-A_j^2-\rho^2+\frac{2A_j^2\rho}{1+\rho}\right] \\
    &=
    S_j\frac{1-\rho}{1+\rho}
    \left[(1+\rho)^2-(1+\rho-x_j)^2\right] \\
    &=
    S_j\frac{1-\rho}{1+\rho}
    x_j\left(2(1+\rho)-x_j\right).
\end{align*}
Therefore
\begin{equation}
\label{eq:pol_green_energy_exact}
    \sum_{\tau=0}^\infty G_\tau(j)^2
    =
    \frac{1+\rho}{(1-\rho)\eta\lambda_j\left(2(1+\rho)-\eta\lambda_j\right)}.
\end{equation}

We now pass from the saturated sum to the finite-time sum. The root formula for $G_k(j)$ implies the following tail bound: there exist universal constants $C_G,c_G>0$ such that
\begin{equation}
\label{eq:pol_green_tail_bound}
    \sum_{\tau=k}^{\infty}G_\tau(j)^2
    \le
    C_G\exp\left\{-c_G k\min\left(1-\rho,\frac{\eta\lambda_j}{1-\rho}\right)\right\}
    \sum_{\tau=0}^{\infty}G_\tau(j)^2.
\end{equation}
This is the standard Polyak relaxation estimate.

Because $\lambda_j\eqsim j^{-\beta}$ and $\beta>1$, the spectral mass is finite:
$\sum_{j\ge1}\lambda_j<\infty.$\\
Choose a fixed integer $J_*$, depending only on the spectrum constants and $\beta$, such that
$\sum_{j=1}^{J_*}\lambda_j\ge\frac12\sum_{j\ge1}\lambda_j.$\\
Since $J_*$ is fixed, $\lambda_{J_*}$ is a positive constant independent of $\eta,\rho,k,B$. Thus, after increasing the universal constant $C$ in~\eqref{eq:pol_label_noise_time_condition} if necessary, condition~\eqref{eq:pol_label_noise_time_condition} implies that for every $1\le j\le J_*$,
\[
    k\min\left(1-\rho,\frac{\eta\lambda_j}{1-\rho}\right)
    \ge
    \frac{\log(2C_G)}{c_G}.
\]
By~\eqref{eq:pol_green_tail_bound}, for all $j\le J_*$,
\[
    \sum_{\tau=k}^{\infty}G_\tau(j)^2
    \le
    \frac12\sum_{\tau=0}^{\infty}G_\tau(j)^2.
\]
Therefore
\begin{equation}
\label{eq:finite_green_lower_half}
    \sum_{\tau=0}^{k-1}G_\tau(j)^2
    \ge
    \frac12\sum_{\tau=0}^{\infty}G_\tau(j)^2,
    \qquad 1\le j\le J_*.
\end{equation}
For all $j$, the trivial upper bound
\begin{equation}
\label{eq:finite_green_upper_saturated}
    \sum_{\tau=0}^{k-1}G_\tau(j)^2
    \le
    \sum_{\tau=0}^{\infty}G_\tau(j)^2
\end{equation}
holds.

We now insert these bounds into the finite-time label-noise term. Since all terms are nonnegative, we can exchange the $\tau$- and $j$-sums:
\[
    \frac{\eta^2\sigma^2}{B}\sum_{\tau=0}^{k-1}\mathcal K_{\polyak}(\tau)
    =
    \frac{\eta^2\sigma^2}{B}
    \sum_{j\ge1}\lambda_j^2\sum_{\tau=0}^{k-1}G_\tau(j)^2.
\]
For the lower bound, use~\eqref{eq:finite_green_lower_half} and~\eqref{eq:pol_green_energy_exact} on the first $J_*$ modes:
\begin{align*}
    \frac{\eta^2\sigma^2}{B}\sum_{\tau=0}^{k-1}\mathcal K_{\polyak}(\tau)
    &\ge
    \frac12\frac{\eta^2\sigma^2}{B}
    \sum_{j=1}^{J_*}\lambda_j^2
    \frac{1+\rho}{(1-\rho)\eta\lambda_j(2(1+\rho)-\eta\lambda_j)} \\
    &=
    \frac12\frac{\eta\sigma^2(1+\rho)}{B(1-\rho)}
    \sum_{j=1}^{J_*}\frac{\lambda_j}{2(1+\rho)-\eta\lambda_j}.
\end{align*}
Since $2(1+\rho)-\eta\lambda_j\le2(1+\rho)$,
\[
    \frac{1+\rho}{2(1+\rho)-\eta\lambda_j}\ge \frac12.
\]
Thus
\begin{align*}
    \frac{\eta^2\sigma^2}{B}\sum_{\tau=0}^{k-1}\mathcal K_{\polyak}(\tau)
    &\ge
    \frac14\frac{\eta\sigma^2}{B(1-\rho)}
    \sum_{j=1}^{J_*}\lambda_j.
\end{align*}
The finite sum $\sum_{j=1}^{J_*}\lambda_j$ is a positive constant, so this gives the lower bound in~\eqref{eq:pol_label_noise_explicit_bounds} for a suitable constant $c_->0$.

For the upper bound, use~\eqref{eq:finite_green_upper_saturated} and~\eqref{eq:pol_green_energy_exact} over all modes:
\begin{align*}
    \frac{\eta^2\sigma^2}{B}\sum_{\tau=0}^{k-1}\mathcal K_{\polyak}(\tau)
    &\le
    \frac{\eta\sigma^2(1+\rho)}{B(1-\rho)}
    \sum_{j\ge1}\frac{\lambda_j}{2(1+\rho)-\eta\lambda_j} \\
    &\le
    \frac{\eta\sigma^2(1+\rho)}{B(1-\rho)}
    \frac{1}{2(1+\rho)-\eta\lambda_1}
    \sum_{j\ge1}\lambda_j.
\end{align*}
The last factor is finite, and the denominator is bounded below by the stability margin. Hence the upper bound in~\eqref{eq:pol_label_noise_explicit_bounds} holds for a suitable constant $c_+>0$.
\end{proof}

% \subsection{Loss in the fully overdamped regime}
% \label{app:pol_fully_overdamped_loss}

% In the fully overdamped regime, the deterministic error trajectory decays monotonically, yielding the following polynomial bias.

% \begin{proposition}[Deterministic loss, fully overdamped regime]
% \label{prop:pol_deterministic_bias0}
% Under Assumptions~\ref{ass:diagonal-covariance} and~\ref{ass:power-law}, in the fully overdamped regime
% $\eta\lesssim (1-\rho)^2$, we have:
% \begin{equation*}
%     \cS_{\polyak}(k)
%     \eqsim
%     \begin{cases}
%     1,&\qquad k\lesssim \frac{1-\rho}{\eta},\\
%     \left(\frac{1-\rho}{\eta k}\right)^s, &\qquad k\gtrsim\frac{1-\rho}{\eta}.\\
%     \end{cases}
% \end{equation*}
% \end{proposition}

% \begin{proof}
% In the fully overdamped regime, we can assume: $\eta\lambda_1\le c_0 (1-\rho)^2$, with a sufficiently small universal constant $c_0>0$.

% Let the deterministic coordinate multiplier $h_k(j):=e_k^{\mathrm{det}}(j)/e_0^{\mathrm{det}}(j) = -e_k^{\mathrm{det}}(j)/\theta_j^*$ satisfy
% \[
%     h_{k+1}(j)=(1+\rho-\eta\lambda_j)h_k(j)-\rho h_{k-1}(j),
%     \qquad
%     h_0(j)=1,
%     \qquad
%     h_1(j)=1-\eta\lambda_j.
% \]
% The characteristic roots are
% \[
%     r_{j,+},r_{j,-}
%     =
%     \frac{1+\rho-\eta\lambda_j\pm\sqrt{(1+\rho-\eta\lambda_j)^2-4\rho}}{2}.
% \]
% We first give explicit root bounds. Since $\eta\le \frac{c_0}{\lambda_1}(1-\rho)^2$ and $\lambda_j\le\lambda_1$, we have $0\le \eta\lambda_j\le c_0(1-\rho)^2$ for all $j$. Put $r=1-u$. The characteristic polynomial evaluated at $1-u$ is
% \[
%     (1-u)^2-(1+\rho-\eta\lambda_j)(1-u)+\rho
%     =
%     \eta\lambda_j-(1-\rho+\eta\lambda_j)u+u^2.
% \]
% Choosing $c_0>0$ sufficiently small, the two evaluations
% \[
%     \eta\lambda_j-(1-\rho+\eta\lambda_j)\frac{\eta\lambda_j}{2(1-\rho)}+\left(\frac{\eta\lambda_j}{2(1-\rho)}\right)^2>0,
% \]
% and
% \[
%     \eta\lambda_j-(1-\rho+\eta\lambda_j)\frac{2\eta\lambda_j}{1-\rho}+\left(\frac{2\eta\lambda_j}{1-\rho}\right)^2<0
% \]
% hold uniformly for $0<\eta\lambda_j\le c_0(1-\rho)^2$. Therefore the slow root satisfies
% \begin{equation}
% \label{eq:pol_slow_root_bounds_sgd}
%     1-\frac{2\eta\lambda_j}{1-\rho}
%     \le
%     r_{j,+}
%     \le
%     1-\frac{\eta\lambda_j}{2(1-\rho)}.
% \end{equation}
% The fast root satisfies $r_{j,+}r_{j,-}=\rho=1-(1-\rho)$. Since $r_{j,+}\ge1-2c_0(1-\rho)$, after decreasing $c_0$ if needed,
% \begin{equation}
% \label{eq:pol_fast_root_bound_sgd}
%     0\le r_{j,-}\le 1-\frac{1-\rho}{2}.
% \end{equation}

% Now write the solution as
% \[
%     h_k(j)=A_jr_{j,+}^k+B_jr_{j,-}^k.
% \]
% From $h_0=1$ and $h_1=1-\eta\lambda_j$, we obtain the coefficients
% \[
%     A_j=\frac{1-\eta\lambda_j-r_{j,-}}{r_{j,+}-r_{j,-}},
%     \qquad
%     B_j=1-A_j = \frac{r_{j,+}-(1-\eta\lambda_j)}{r_{j,+}-r_{j,-}}.
% \]
% Using~\eqref{eq:pol_slow_root_bounds_sgd}--\eqref{eq:pol_fast_root_bound_sgd}, and taking $c_0$ sufficiently small, these coefficients obey uniform bounds:
% \begin{equation}
% \label{eq:pol_sgd_coeff_bounds}
%     \frac12\le A_j\le2,
%     \qquad
%     |B_j|\le2.
% \end{equation}
% Equivalently, applying the initial conditions, the exact solution can be expressed in the fractional form:
% \begin{equation}
% \label{eq:hk_exact_form}
%     h_k(j) = \frac{r_{j,+} r_{j,-}\left(-r_{j,+}^{k-1} + r_{j,-}^{k-1}\right) + (1-\eta\lambda_j)\left(r_{j,+}^k - r_{j,-}^k\right)}{r_{j,+} - r_{j,-}}.
% \end{equation}

% To evaluate the deterministic bias $\cS_{\polyak}(k) = \sum_{j\ge1}\lambda_j(\theta_j^*)^2 h_k(j)^2$, we divide the time $k$ into three successive stages.

% \textbf{Stage 1: $k \lesssim \frac{1}{1-\rho}$.} \\
% In this extremely short-time regime, taking the limits $r_{j,+}\approx r_{j,-} \approx r \approx 1$ and using $r^2 \approx \rho$ in~\eqref{eq:hk_exact_form} yields the heuristic trajectory:
% \[
%     h_k(j) \simeq -\rho(k-1)r^{k-2} + (1-\eta\lambda_j)k r^{k-1} = k\left((1-\eta\lambda_j)r - \rho \right)r^{k-2} + \rho r^{k-2}.
% \]
% Since $(1-\eta\lambda_j)r - \rho = \mathcal{O}(1-\rho)$, the first term is bounded by $\mathcal{O}(k(1-\rho)) \lesssim 1$, yielding $h_k(j) \simeq 1$. 
% Rigorously, for the leading modes ($j=\mathcal{O}(1)$), since $\eta\lambda_j \le c_0(1-\rho)^2$, the slow root satisfies $r_{j,+} \ge 1 - 2c_0(1-\rho)$. For $k \lesssim \frac{1}{1-\rho}$, we have
% \[
%     r_{j,+}^k \ge \left(1 - 2c_0(1-\rho)\right)^{\frac{C}{1-\rho}} \eqsim 1.
% \]
% Since the dynamics are unconditionally stable ($h_k(j) \le 1$), it follows that $h_k(j) \eqsim 1$ for the leading modes. Consequently, the bias is dominated by these non-decayed modes:
% \[
%     \cS_{\polyak}(k) \eqsim \sum_{j\ge1}\lambda_j(\theta_j^*)^2 \eqsim 1.
% \]

% \textbf{Stage 2: $\frac{1}{1-\rho} \lesssim k \lesssim \frac{1-\rho}{\eta}$.} \\
% For $k \gtrsim \frac{1}{1-\rho}$, the fast root decays exponentially since $r_{j,-}^k \le \left(1-\frac{1-\rho}{2}\right)^k \ll 1$. The multiplier is  dominated by the slow root:
% \[
%     h_k(j) \simeq \frac{1-\eta\lambda_j - r_{j,-}}{r_{j,+} - r_{j,-}} r_{j,+}^k \simeq r_{j,+}^k.
% \]
% For the leading modes with $\lambda_j = \Theta(1)$, within the time window $k \lesssim \frac{1-\rho}{\eta}$, the exponent is bounded by:
% \[
%     r_{j,+}^k \eqsim \left(1 - c\frac{\eta\lambda_j}{1-\rho}\right)^k \eqsim \exp\left(-c\frac{\eta\lambda_j}{1-\rho}k\right) \eqsim 1,
% \]
% where we used the condition $\frac{\eta\lambda_j}{1-\rho}k \lesssim \lambda_j \le \lambda_1 = \mathcal{O}(1)$. Thus, the leading modes have not yet experienced significant exponential decay, preserving $h_k(j) \eqsim 1$. As a result, the deterministic bias remains unchanged:
% \[
%     \cS_{\polyak}(k) \eqsim 1.
% \]

% \textbf{Stage 3: $k \gtrsim \frac{1-\rho}{\eta}$.} \\
% In this long-time regime, the exponential decay of the slow root becomes significant across the spectrum. Combining~\eqref{eq:pol_slow_root_bounds_sgd},~\eqref{eq:pol_fast_root_bound_sgd}, and~\eqref{eq:pol_sgd_coeff_bounds}, there are universal constants $c_1,c_2,c_3,C_1>0$ such that, for all $k\ge C_1(1-\rho)^{-1}$,
% \begin{equation}
% \label{eq:pol_sgd_coordinate_two_sided_decay}
%     c_1\exp\left(-c_2\frac{\eta\lambda_j}{1-\rho}k\right)
%     \le
%     |h_k(j)|
%     \le
%     C_1\exp\left(-c_3\frac{\eta\lambda_j}{1-\rho}k\right).
% \end{equation}
% Substituting~\eqref{eq:pol_sgd_coordinate_two_sided_decay} into the deterministic bias gives
% \begin{equation}
% \label{eq:pol_sgd_bias_spectral_comparison}
%     c\sum_{j\ge1}\lambda_j(\theta_j^*)^2
%     \exp\left(-C\frac{\eta\lambda_j}{1-\rho}k\right)
%     \le
%     \mathcal D_{\polyak}(k)
%     \le
%     C\sum_{j\ge1}\lambda_j(\theta_j^*)^2
%     \exp\left(-c\frac{\eta\lambda_j}{1-\rho}k\right).
% \end{equation}
% By Assumption~\ref{ass:power-law},
% \[
%     \lambda_j(\theta_j^*)^2\eqsim j^{-1-\beta s}, \qquad \lambda_j\eqsim j^{-\beta}.
% \]
% Thus it remains to bound sums of the form
% \[
%     \sum_{j\ge1}j^{-1-\beta s}
%     \exp\left(-c\frac{\eta k}{1-\rho}j^{-\beta}\right).
% \]
% The condition on $k$ implies $\frac{\eta k}{1-\rho}\ge C$.

% For the lower bound, set $J:=\left\lceil \left(\frac{\eta k}{1-\rho}\right)^{1/\beta}\right\rceil$. For $j\ge J$, we have $\frac{\eta k}{1-\rho}j^{-\beta}\le1$, so $\exp(-C\frac{\eta k}{1-\rho}j^{-\beta})\ge e^{-C}$. Therefore,
% \[
%     \sum_{j\ge1}j^{-1-\beta s}\exp\left(-C\frac{\eta k}{1-\rho}j^{-\beta}\right)
%     \ge
%     e^{-C}\sum_{j\ge J}j^{-1-\beta s}
%     \ge
%     c\left(\frac{\eta k}{1-\rho}\right)^{-s}.
% \]
% For the upper bound, the integral comparison with $u=c\frac{\eta k}{1-\rho}x^{-\beta}$ gives
% \begin{align*}
%     \sum_{j\ge1}j^{-1-\beta s}\exp\left(-c\frac{\eta k}{1-\rho}j^{-\beta}\right)
%     &\le
%     C\int_1^\infty x^{-1-\beta s}\exp\left(-c\frac{\eta k}{1-\rho}x^{-\beta}\right)\,dx \\
%     &=
%     C \left(\frac{\eta k}{1-\rho}\right)^{-s}\int_0^{c\frac{\eta k}{1-\rho}}u^{s-1}e^{-u}\,du
%     \le
%     C \left(\frac{\eta k}{1-\rho}\right)^{-s}.
% \end{align*}
% Combining the upper and lower bounds in~\eqref{eq:pol_sgd_bias_spectral_comparison} yields
% \[
%     \mathcal D_{\polyak}(k)
%     \eqsim
%     \left(\frac{1-\rho}{\eta k}\right)^s.
% \]
% Synthesizing the three stages completes the proof.
% \end{proof}

% Without underdamped modes, the Green function kernel exhibits purely overdamped polynomial scaling.

% \begin{lemma}[Convolution Kernel, Fully Overdamped Regime]
% \label{lem:pol_kernel0}
% Under Assumption~\ref{ass:power-law}, we have:
% \begin{equation*}
%     \mathcal{K}_{\polyak}(k)
%     \eqsim \begin{cases}
%     k^2, &\qquad k\lesssim \frac{1}{1-\rho},\\
%     \frac{1}{(1-\rho)^2}, &\qquad \frac{1}{1-\rho}\lesssim k \lesssim \frac{1-\rho}{\eta},\\
%     \frac{1}{(1-\rho)^2}
%     \left(\frac{1-\rho}{\eta k}\right)^{2-\frac1\beta}, &\qquad k \gtrsim  \frac{1-\rho}{\eta}.
% \end{cases}
% \end{equation*}
% \end{lemma}

% \begin{proof}
% We work in the Fully Overdamped Regime, where $\eta\lambda_1\le c_0(1-\rho)^2$ with $c_0>0$ sufficiently small. The Green function $G_k(j)$ satisfies the homogeneous recursion:
% \[
%     G_{k+1}(j)=(1+\rho-\eta\lambda_j)G_k(j)-\rho G_{k-1}(j),
%     \qquad
%     G_0(j)=1,
%     \qquad
%     G_{-1}(j)=0.
% \]
% The exact formula for the Green function is:
% \[
%     G_k(j)=\frac{r_{j,+}^{k+1}-r_{j,-}^{k+1}}{r_{j,+}-r_{j,-}} = \sum_{m=0}^{k} r_{j,+}^{k-m} r_{j,-}^m.
% \]
% From Proposition \ref{prop:pol_deterministic_bias0}, the roots $r_{j,+}$ and $r_{j,-}$ are real and satisfy:
% \[
%     1-\frac{2\eta\lambda_j}{1-\rho}
%     \le r_{j,+}\le
%     1-\frac{\eta\lambda_j}{2(1-\rho)},
%     \qquad
%     0 \le r_{j,-} \le 1-\frac{1-\rho}{2},
% \]
% and their difference is bounded by $r_{j,+}-r_{j,-}\eqsim 1-\rho$. Also, $r_{j,+}r_{j,-} = \rho$, meaning $r_{j,-} \ge \rho = 1-(1-\rho)$.

% By definition, the convolution kernel is:
% \[
%     \mathcal K_{\polyak}(k) = \sum_{j\ge1}\lambda_j^2 G_k(j)^2.
% \]
% We divide the time $k$ into three successive stages to evaluate $\mathcal K_{\polyak}(k)$.

% \textbf{Stage 1: $k \lesssim \frac{1}{1-\rho}$.} \\
% In this extremely short-time regime, the roots are very close to 1. Using the uniform bounds $r_{j,+} \ge 1 - 2c_0(1-\rho)$ and $r_{j,-} \ge 1-(1-\rho)$, we have for all $m \le k \lesssim \frac{1}{1-\rho}$:
% \[
%     r_{j,+}^{k-m} \ge \left(1 - 2c_0(1-\rho)\right)^{\frac{C}{1-\rho}} \eqsim 1, \qquad r_{j,-}^m \ge \left(1 - (1-\rho)\right)^{\frac{C}{1-\rho}} \eqsim 1.
% \]
% Since all terms in the algebraic expansion of $G_k(j)$ are strictly positive and bounded away from zero by a constant, the sum is proportional to the number of terms:
% \[
%     G_k(j) = \sum_{m=0}^{k} r_{j,+}^{k-m} r_{j,-}^m \eqsim \sum_{m=0}^{k} 1 = k+1 \eqsim k.
% \]
% Since this holds for the leading modes which dominate the convergent spectral sum ($\sum_{j\ge1}\lambda_j^2 < \infty$ due to $\beta > 1$), we obtain:
% \[
%     \mathcal K_{\polyak}(k) \eqsim \sum_{j\ge1}\lambda_j^2 k^2 = k^2 \sum_{j\ge1}\lambda_j^2 \eqsim k^2.
% \]

% \textbf{Stage 2: $\frac{1}{1-\rho} \lesssim k \lesssim \frac{1-\rho}{\eta}$.} \\
% In this intermediate regime, $k$ is sufficiently large such that the fast root decays exponentially:
% \[
%     r_{j,-}^{k+1} \le \left(1-\frac{1-\rho}{2}\right)^{k+1} \ll r_{j,+}^{k+1}.
% \]
% Thus, the Green function is  dominated by the slow root:
% \[
%     G_k(j) \simeq \frac{r_{j,+}^{k+1}}{r_{j,+}-r_{j,-}} \eqsim \frac{r_{j,+}^{k+1}}{1-\rho}.
% \]
% For the leading modes ($j = \mathcal{O}(1)$), since $k \lesssim \frac{1-\rho}{\eta}$, the exponent is bounded by:
% \[
%     r_{j,+}^{k+1} \ge \left(1 - \frac{2\eta\lambda_j}{1-\rho}\right)^{k+1} \eqsim \exp\left(-C \frac{\eta\lambda_j}{1-\rho}k\right) \ge \exp(-C \lambda_j) \eqsim 1.
% \]
% This implies the slow root has not yet experienced significant exponential decay for the leading modes, giving $G_k(j) \eqsim \frac{1}{1-\rho}$. For the tail modes, we trivially have $G_k(j) \lesssim \frac{1}{1-\rho}$. Therefore, the kernel is dominated by these non-decayed leading modes:
% \[
%     \mathcal K_{\polyak}(k) = \sum_{j\ge1}\lambda_j^2 G_k(j)^2 \eqsim \sum_{j=1}^{\mathcal{O}(1)}\lambda_j^2 \frac{1}{(1-\rho)^2} \eqsim \frac{1}{(1-\rho)^2}.
% \]

% \textbf{Stage 3: $k \gtrsim \frac{1-\rho}{\eta}$.} \\
% In this long-time regime, the exponential decay of the slow root becomes dominant across the spectrum. Substituting the two-sided root bounds into $G_k(j) \simeq \frac{r_{j,+}^{k+1}}{r_{j,+}-r_{j,-}}$, there exist constants $c, C > 0$ such that:
% \begin{equation*}
%     \frac{c}{1-\rho}
%     \exp\left(-C\frac{\eta\lambda_j}{1-\rho}k\right)
%     \le
%     |G_k(j)|
%     \le
%     \frac{C}{1-\rho}
%     \exp\left(-c\frac{\eta\lambda_j}{1-\rho}k\right).
% \end{equation*}
% Squaring this and summing over the spectrum gives:
% \begin{equation*}
%     \frac{c}{(1-\rho)^2}
%     \sum_{j\ge1}\lambda_j^2
%     \exp\left(-C\frac{\eta\lambda_j}{1-\rho}k\right)
%     \le
%     \mathcal K_{\polyak}(k)
%     \le
%     \frac{C}{(1-\rho)^2}
%     \sum_{j\ge1}\lambda_j^2
%     \exp\left(-c\frac{\eta\lambda_j}{1-\rho}k\right).
% \end{equation*}
% Using $\lambda_j \eqsim j^{-\beta}$, we estimate the corresponding spectral sum:
% \[
%     \sum_{j\ge1}j^{-2\beta}
%     \exp\left(-c\frac{\eta k}{1-\rho}j^{-\beta}\right).
% \]
% For the lower bound, let $J:=\left\lceil \left(\frac{\eta k}{1-\rho}\right)^{1/\beta}\right\rceil$. For $j\ge J$, we have $\frac{\eta k}{1-\rho}j^{-\beta}\le1$, yielding $\exp\left(-C\frac{\eta k}{1-\rho}j^{-\beta}\right)\ge e^{-C}$. Thus:
% \[
%     \sum_{j\ge1}j^{-2\beta}
%     \exp\left(-C\frac{\eta k}{1-\rho}j^{-\beta}\right)
%     \ge
%     e^{-C}\sum_{j\ge J}j^{-2\beta}
%     \ge
%     c J^{1-2\beta}
%     \ge
%     c\left(\frac{\eta k}{1-\rho}\right)^{-\left(2-\frac1\beta\right)}.
% \]
% For the upper bound, an integral comparison with $u=c\frac{\eta k}{1-\rho}x^{-\beta}$ gives:
% \begin{align*}
%     \sum_{j\ge1}j^{-2\beta}
%     \exp\left(-c\frac{\eta k}{1-\rho}j^{-\beta}\right)
%     &\le
%     C\int_1^\infty x^{-2\beta}
%     \exp\left(-c\frac{\eta k}{1-\rho}x^{-\beta}\right)\,dx \\
%     &=
%     C\left(\frac{\eta k}{1-\rho}\right)^{-\left(2-\frac1\beta\right)}
%     \int_0^{c\frac{\eta k}{1-\rho}}u^{1-\frac1\beta}e^{-u}\,du \\
%     &\le
%     C\left(\frac{\eta k}{1-\rho}\right)^{-\left(2-\frac1\beta\right)}.
% \end{align*}
% Combining these bounds yields:
% \[
%     \mathcal K_{\polyak}(k)
%     \eqsim
%     \frac{1}{(1-\rho)^2}
%     \left(\frac{1-\rho}{\eta k}\right)^{2-\frac1\beta}.
% \]
% Synthesizing the three stages completes the proof.
% \end{proof}

% By convolving the deterministic bias and the kernel, we bound the multiplicative noise accumulation.

% \begin{proposition}[Multiplicative Noise Bound, Fully Overdamped Regime]
% \label{prop:pol_self_noise_bounds0}
% In the Fully Overdamped Regime, under the stability condition $\eta\le cB(1-\rho)$ with $c>0$ sufficiently small, and for
% $k\ge C\max\{(1-\rho)^{-1},(1-\rho)/\eta\}$, we have:
% \[
%     \frac{\eta^2}{B}\sum_{m=0}^{k-1}\mathcal{K}_{\polyak}(k-1-m)\EE[\cE(\btheta_m)]
%     \eqsim
%     \frac{\eta}{B(1-\rho)}
%         \left(\frac{1-\rho}{\eta k}\right)^{(2-\frac1\beta)\wedge s}+\frac{\eta^2\sigma^2}{B^2(1-\rho)^2}.
% \]
% \end{proposition}

% \begin{proof}
% Let
% \[
%     \mathcal M_k:=
%     \frac{\eta^2}{B}\sum_{m=0}^{k-1}\mathcal K_{\polyak}(k-1-m)\EE[\cE(\btheta_m)].
% \]
% By Proposition~\ref{prop:pol_risk_recursion}, the expected excess risk is bounded by:
% \[
%     \EE[\cE(\btheta_m)]
%     \le
%     C\left[
%         \mathcal D_{\polyak}(m)
%         +\frac{\eta^2\sigma^2}{B}\sum_{r=0}^{m-1}\mathcal K_{\polyak}(r)
%         +\mathcal M_m
%     \right].
% \]
% Substituting this bound into the definition of $\mathcal M_k$ yields three convolution terms:
% \begin{align}
% \label{eq:prop17_bootstrap_start}
%     \mathcal M_k
%     \le{}&
%     C\frac{\eta^2}{B}\sum_{m=0}^{k-1}
%         \mathcal K_{\polyak}(k-1-m)\mathcal D_{\polyak}(m) \nonumber \\
%     & + C\frac{\eta^4\sigma^2}{B^2}\sum_{m=0}^{k-1}\mathcal K_{\polyak}(k-1-m)\sum_{r=0}^{m-1}\mathcal K_{\polyak}(r) \nonumber \\
%     & + C\frac{\eta^2}{B}\sum_{m=0}^{k-1}
%         \mathcal K_{\polyak}(k-1-m)\mathcal M_m .
% \end{align}

% We first bound the label-noise convolution term. Changing the summation indices and exploiting the non-negativity of the kernel, we have:
% \begin{align*}
%     \frac{\eta^4\sigma^2}{B^2}\sum_{m=0}^{k-1}\mathcal K_{\polyak}(k-1-m)\sum_{r=0}^{m-1}\mathcal K_{\polyak}(r)
%     &= \frac{\eta^4\sigma^2}{B^2} \sum_{n=0}^{k-1} \mathcal K_{\polyak}(n) \sum_{r=0}^{k-2-n} \mathcal K_{\polyak}(r) \\
%     &\le \frac{\eta^4\sigma^2}{B^2} \left( \sum_{n=0}^{\infty} \mathcal K_{\polyak}(n) \right)^2.
% \end{align*}
% From the saturated Green-energy estimate in Lemma~\ref{lem:pol_label_noise}, we have $\sum_{n=0}^{\infty} \mathcal K_{\polyak}(n) \eqsim \frac{1}{\eta(1-\rho)}$. Therefore, this term is bounded by $\mathcal{O}\left( \frac{\eta^2\sigma^2}{B^2(1-\rho)^2} \right)$. \\
% For the lower bound, since $k \ge C\frac{1-\rho}{\eta}$, the sum over the restricted range $n, r \le k/2$ captures a constant fraction of the infinite sum, yielding the matching lower bound $\eqsim \frac{\eta^2\sigma^2}{B^2(1-\rho)^2}$. This non-decaying constant term represents the fundamental noise floor amplified by momentum.

% Next, we rigorously evaluate the deterministic convolution term. For asymptotic simplicity, we write $I_{\mathrm{det}} := \sum_{m=0}^{k-1} \mathcal K_{\polyak}(m)\mathcal D_{\polyak}(k-m)$. Based on the piecewise asymptotic forms of $\mathcal K_{\polyak}$ and $\mathcal D_{\polyak}$ established in Proposition~\ref{prop:pol_deterministic_bias0} and Lemma~\ref{lem:pol_kernel0}, we divide the summation into five distinct stages:

% \textbf{Stage 1:} $0 \le m \lesssim \frac{1}{1-\rho}$. \\
% Here, $\mathcal K_{\polyak}(m) \eqsim m^2$ and $k-m \eqsim k \gtrsim \frac{1-\rho}{\eta}$, giving $\mathcal D_{\polyak}(k-m) \eqsim \left(\frac{1-\rho}{\eta k}\right)^s$.
% \[
%     \sum_{m=0}^{\frac{1}{1-\rho}} m^2 \left(\frac{1-\rho}{\eta k}\right)^s \eqsim \frac{1}{(1-\rho)^3} \left(\frac{1-\rho}{\eta k}\right)^s.
% \]

% \textbf{Stage 2:} $\frac{1}{1-\rho} \lesssim m \lesssim \frac{1-\rho}{\eta}$. \\
% Here, $\mathcal K_{\polyak}(m) \eqsim \frac{1}{(1-\rho)^2}$ and $\mathcal D_{\polyak}(k-m) \eqsim \left(\frac{1-\rho}{\eta k}\right)^s$. \\
% The length of this interval is $\mathcal{O}(\frac{1-\rho}{\eta})$.
% \[
%     \sum_{m=\frac{1}{1-\rho}}^{\frac{1-\rho}{\eta}} \frac{1}{(1-\rho)^2} \left(\frac{1-\rho}{\eta k}\right)^s \eqsim \frac{1-\rho}{\eta} \cdot \frac{1}{(1-\rho)^2} \left(\frac{1-\rho}{\eta k}\right)^s = \frac{1}{\eta(1-\rho)} \left(\frac{1-\rho}{\eta k}\right)^s.
% \]

% \textbf{Stage 3:} $\frac{1-\rho}{\eta} \lesssim m \lesssim k - \frac{1-\rho}{\eta}$. \\
% Here, $\mathcal K_{\polyak}(m) \eqsim \frac{1}{(1-\rho)^2} \left(\frac{1-\rho}{\eta m}\right)^{2-\frac1\beta}$ and $\mathcal D_{\polyak}(k-m) \eqsim \left(\frac{1-\rho}{\eta(k-m)}\right)^s$.
% \begin{align*}
%     &\sum_{m=\frac{1-\rho}{\eta}}^{k - \frac{1-\rho}{\eta}} \frac{1}{(1-\rho)^2} \left(\frac{1-\rho}{\eta m}\right)^{2-\frac1\beta} \left(\frac{1-\rho}{\eta(k-m)}\right)^s \\
%     &\eqsim \frac{1}{(1-\rho)^2} \left(\frac{1-\rho}{\eta}\right)^{2-\frac1\beta+s} \left[ \sum_{m=\frac{1-\rho}{\eta}}^{k/2} \frac{1}{m^{2-\frac1\beta}} \frac{1}{k^s} + \sum_{m=k/2}^{k-\frac{1-\rho}{\eta}} \frac{1}{k^{2-\frac1\beta}} \frac{1}{(k-m)^s} \right] \\
%     &\eqsim \frac{1}{(1-\rho)^2} \left(\frac{1-\rho}{\eta}\right)^{2-\frac1\beta+s} \left[ k^{-s} \left(\frac{1-\rho}{\eta}\right)^{-\left(1-\frac1\beta\right)} + k^{-\left(2-\frac1\beta\right)} \begin{cases} k^{1-s}, & s<1 \\ \left(\frac{1-\rho}{\eta}\right)^{1-s}, & s>1 \end{cases} \right] \\
%     &\eqsim \frac{1}{\eta(1-\rho)} \left(\frac{1-\rho}{\eta k}\right)^s + \frac{1}{\eta(1-\rho)} \left(\frac{1-\rho}{\eta k}\right)^{2-\frac1\beta} \\
%     &\eqsim \frac{1}{\eta(1-\rho)} \left(\frac{1-\rho}{\eta k}\right)^{\left(2-\frac1\beta\right)\wedge s}.
% \end{align*}

% \textbf{Stage 4:} $k - \frac{1-\rho}{\eta} \lesssim m \lesssim k - \frac{1}{1-\rho}$. \\
% Here, $\mathcal K_{\polyak}(m) \eqsim \frac{1}{(1-\rho)^2} \left(\frac{1-\rho}{\eta k}\right)^{2-\frac1\beta}$ and $k-m \in [\frac{1}{1-\rho}, \frac{1-\rho}{\eta}]$, giving $\mathcal D_{\polyak}(k-m) \eqsim 1$.\\
% The length is $\mathcal{O}(\frac{1-\rho}{\eta})$.
% \[
%     \sum_{m=k-\frac{1-\rho}{\eta}}^{k-\frac{1}{1-\rho}} \frac{1}{(1-\rho)^2} \left(\frac{1-\rho}{\eta k}\right)^{2-\frac1\beta} \eqsim \frac{1-\rho}{\eta} \cdot \frac{1}{(1-\rho)^2} \left(\frac{1-\rho}{\eta k}\right)^{2-\frac1\beta} = \frac{1}{\eta(1-\rho)} \left(\frac{1-\rho}{\eta k}\right)^{2-\frac1\beta}.
% \]

% \textbf{Stage 5:} $k - \frac{1}{1-\rho} \lesssim m \le k-1$. \\
% Here, $\mathcal K_{\polyak}(m) \eqsim \frac{1}{(1-\rho)^2} \left(\frac{1-\rho}{\eta k}\right)^{2-\frac1\beta}$ and $k-m \in [1, \frac{1}{1-\rho}]$, giving $\mathcal D_{\polyak}(k-m) \eqsim 1$. \\
% The length is $\mathcal{O}(\frac{1}{1-\rho})$.
% \[
%     \sum_{m=k-\frac{1}{1-\rho}}^{k-1} \frac{1}{(1-\rho)^2} \left(\frac{1-\rho}{\eta k}\right)^{2-\frac1\beta} \eqsim \frac{1}{1-\rho} \cdot \frac{1}{(1-\rho)^2} \left(\frac{1-\rho}{\eta k}\right)^{2-\frac1\beta} = \frac{1}{(1-\rho)^3} \left(\frac{1-\rho}{\eta k}\right)^{2-\frac1\beta}.
% \]

% Since $\eta \lesssim (1-\rho)^2$ in the Fully Overdamped Regime, we have $\frac{1}{\eta(1-\rho)} \gtrsim \frac{1}{(1-\rho)^3}$. Thus, the sums from Stages 1 and 5 are strictly lower order. The dominant contribution to $I_{\mathrm{det}}$ comes from Stages 2, 3, and 4, yielding:
% \[
%     I_{\mathrm{det}} \eqsim \frac{1}{\eta(1-\rho)} \left(\frac{1-\rho}{\eta k}\right)^{\left(2-\frac1\beta\right)\wedge s}.
% \]
% Multiplying by $\frac{\eta^2}{B}$, the deterministic convolution term is bounded by
% \[
% \mathcal{O}\left( \frac{\eta}{B(1-\rho)} \left(\frac{1-\rho}{\eta k}\right)^{\left(2-\frac1\beta\right)\wedge s} \right). 
% \]
% The matching lower bound follows identically from the two-sided nature of the constituent bounds.

% Finally, we absorb the feedback term in~\eqref{eq:prop17_bootstrap_start}. Exchanging the order of summation gives a two-kernel convolution:
% \begin{align*}
%     \frac{\eta^2}{B}\sum_{m=0}^{k-1}\mathcal K_{\polyak}(k-1-m)\mathcal M_m
%     &= \frac{\eta^2}{B}\sum_{\ell=0}^{k-2} \left[ \frac{\eta^2}{B}\sum_{m=\ell+1}^{k-1} \mathcal K_{\polyak}(k-1-m)\mathcal K_{\polyak}(m-1-\ell) \right] \EE[\cE(\btheta_\ell)].
% \end{align*}
% Using the kernel sub-convolution bound 
% \[
% \sum_{r} \mathcal K_{\polyak}(n-r)\mathcal K_{\polyak}(r) \le C \left(\sum_{r=0}^{\infty}\mathcal K_{\polyak}(r)\right) \mathcal K_{\polyak}(n) \le \frac{C}{\eta(1-\rho)}\mathcal K_{\polyak}(n),
% \]
% we have:
% \begin{align*}
%     \frac{\eta^2}{B}\sum_{m=0}^{k-1}\mathcal K_{\polyak}(k-1-m)\mathcal M_m
%     &\le C\frac{\eta}{B(1-\rho)} \frac{\eta^2}{B}\sum_{\ell=0}^{k-2}\mathcal K_{\polyak}(k-2-\ell)\EE[\cE(\btheta_\ell)] \nonumber \\
%     &\le C\frac{\eta}{B(1-\rho)}\mathcal M_{k-1} \le C'\frac{\eta}{B(1-\rho)}\mathcal M_k .
% \end{align*}
% Choosing the stability constant $c>0$ in $\eta\le cB(1-\rho)$ sufficiently small ensures $C'\frac{\eta}{B(1-\rho)} \le 1/2$, allowing this term to be  absorbed into the left-hand side of~\eqref{eq:prop17_bootstrap_start}.

% Combining the bounded deterministic and label-noise convolutions completes the proof:
% \[
%     \mathcal M_k \eqsim \frac{\eta}{B(1-\rho)} \left(\frac{1-\rho}{\eta k}\right)^{\left(2-\frac1\beta\right)\wedge s} + \frac{\eta^2\sigma^2}{B^2(1-\rho)^2}.
% \]
% \end{proof}

% Combining the deterministic bias, label noise, and multiplicative noise yields the complete scaling law for the fully overdamped regime.

% \begin{theorem}[Scaling Law, Fully Overdamped Regime]
% \label{thm:pol_scaling_sgd}
% Under Assumptions~\ref{ass:diagonal-covariance},~\ref{ass:power-law}, the stability condition $\eta \lesssim B(1-\rho)$, and $\eta\lesssim(1-\rho)^2$, $k\gtrsim \frac{1-\rho}{\eta}$, the expected excess risk satisfies:
% \[
%     \EE[\cE(\btheta_k)]
%     \eqsim
%     \left(\frac{1-\rho}{\eta k}\right)^s
%     +
%     \frac{\sigma^2\eta}{B(1-\rho)}.
% \]
% \end{theorem}

% \begin{proof}
% By the expected risk recursion established in Proposition~\ref{prop:pol_risk_recursion}, the total excess risk decomposes into three additive components: the deterministic bias, the accumulated label noise, and the multiplicative noise feedback $\mathcal{M}_k$:
% \[
%     \EE[\cE(\btheta_k)]
%     \eqsim
%     \cS_{\polyak}(k)
%     +
%     \frac{\eta^2\sigma^2}{B}\sum_{\tau=0}^{k-1}\mathcal{K}_{\polyak}(\tau)
%     +
%     \mathcal M_k.
% \]

% For the Fully Overdamped Regime ($k \gtrsim \frac{1-\rho}{\eta}$), we substitute the corresponding asymptotic bounds derived in the preceding propositions:
% \begin{enumerate}
%     \item \textbf{Deterministic Bias} (Proposition~\ref{prop:pol_deterministic_bias0}):
%     \[
%         \cS_{\polyak}(k) \eqsim \left(\frac{1-\rho}{\eta k}\right)^s.
%     \]
%     \item \textbf{Label Noise} (Lemma~\ref{lem:pol_label_noise}):
%     \[
%         \frac{\eta^2\sigma^2}{B}\sum_{\tau=0}^{k-1}\mathcal{K}_{\polyak}(\tau) \eqsim \frac{\sigma^2\eta}{B(1-\rho)}.
%     \]
%     \item \textbf{Multiplicative Noise} (Proposition~\ref{prop:pol_self_noise_bounds0}):
%     \[
%         \mathcal M_k \eqsim \frac{\eta}{B(1-\rho)} \left(\frac{1-\rho}{\eta k}\right)^{\left(2-\frac1\beta\right)\wedge s} + \frac{\eta^2\sigma^2}{B^2(1-\rho)^2}.
%     \]
% \end{enumerate}

% Summing these three terms, we analyze the structural absorption of the multiplicative noise fragments. First, under the stability condition $\eta \lesssim B(1-\rho)$, the constant noise floor originating from the multiplicative noise is strictly bounded by the primary label noise and is thus  absorbed:
% \[
%     \frac{\eta^2\sigma^2}{B^2(1-\rho)^2} = \left(\frac{\eta}{B(1-\rho)}\right) \frac{\sigma^2\eta}{B(1-\rho)} \lesssim \frac{\sigma^2\eta}{B(1-\rho)}.
% \]

% Next, we bound the polynomial tail of the multiplicative noise, denoted $M_3(k) = \frac{\eta}{B(1-\rho)} \left(\frac{1-\rho}{\eta k}\right)^{\left(2-\frac1\beta\right)\wedge s}$:
% \begin{itemize}
%     \item If $s \le 2-\frac{1}{\beta}$, then $M_3(k) = \frac{\eta}{B(1-\rho)} \left(\frac{1-\rho}{\eta k}\right)^s$. Since $\frac{\eta}{B(1-\rho)} \lesssim 1$, we explicitly have $M_3(k) \lesssim \left(\frac{1-\rho}{\eta k}\right)^s \eqsim \cS_{\polyak}(k)$. Thus, it is cleanly absorbed into the deterministic bias.
%     \item If $s > 2-\frac{1}{\beta}$, then $M_3(k) = \frac{\eta}{B(1-\rho)} \left(\frac{1-\rho}{\eta k}\right)^{2-\frac{1}{\beta}}$.
%     When the deterministic bias dominates the label noise $\left( \left(\frac{1-\rho}{\eta k}\right)^{-s} \le \frac{B(1-\rho)}{\eta} \right)$, reversing the inequality provides an upper bound that yields explicit containment:
%     \[
%         M_3(k) = \frac{\eta}{B(1-\rho)} \left(\frac{1-\rho}{\eta k}\right)^{s-\left(s-\left(2-\frac{1}{\beta}\right)\right)} \le \left(\frac{\eta}{B(1-\rho)}\right)^{\frac{2-1/\beta}{s}} \left(\frac{1-\rho}{\eta k}\right)^s \lesssim \cS_{\polyak}(k).
%     \]
%     Conversely, in the long-time limit where label noise dominates $\left( \left(\frac{1-\rho}{\eta k}\right)^{-s} \ge \frac{B(1-\rho)}{\eta} \right)$, since the operational regime dictates $k \gtrsim \frac{1-\rho}{\eta}$, we have $\frac{1-\rho}{\eta k} \lesssim 1$. This structural mapping guarantees $\left(\frac{1-\rho}{\eta k}\right)^{2-\frac{1}{\beta}} \lesssim 1$, enforcing containment:
%     \[
%         M_3(k) \lesssim \frac{\eta}{B(1-\rho)} \cdot 1 \eqsim \text{LN}(k).
%     \]
% \end{itemize}
% In all distinct bounding configurations, the un-filtered multiplicative term is universally bounded by $\max(\cS_{\polyak}(k), \text{LN}(k))$. Thus, it is thoroughly absorbed, yielding the final scaling law:
% \[
%     \EE[\cE(\btheta_k)]
%     \eqsim
%     \left(\frac{1-\rho}{\eta k}\right)^s
%     +
%     \frac{\sigma^2\eta}{B(1-\rho)}.
% \]
% This completes the proof.
% \end{proof}

% \begin{theorem}[Scaling Law, No Label Noise, Fully Overdamped Regime]
% \label{thm:pol_scaling_sgd_noiseless}
% Under Assumptions~\ref{ass:diagonal-covariance} and~\ref{ass:power-law}, with \(\sigma=0\), the stability condition \(\eta \lesssim B(1-\rho)\), and in the fully overdamped regime \(\eta\lesssim(1-\rho)^2\), for \(k \gtrsim \frac{1-\rho}{\eta}\), the expected excess risk satisfies:
% \begin{equation}
% \label{eq:pol_scaling_sgd_noiseless}
%     \EE[\cE(\btheta_k)]
%     \eqsim
%     \left(\frac{1-\rho}{\eta k}\right)^s
%     +
%     \frac{\eta}{B(1-\rho)}
%     \left(\frac{1-\rho}{\eta k}\right)^{s \wedge \left(2-\frac1\beta\right)}.
% \end{equation}
% \end{theorem}

% \begin{proof}
% By the expected risk recursion established in Proposition~\ref{prop:pol_risk_recursion}, the total excess risk decomposes as:
% \[
%     \EE[\cE(\btheta_k)]
%     \eqsim
%     \cS_{\polyak}(k)
%     +
%     \frac{\eta^2\sigma^2}{B}\sum_{\tau=0}^{k-1}\mathcal{K}_{\polyak}(\tau)
%     +
%     \mathcal M_{\polyak, k}.
% \]
% Setting \(\sigma=0\) eliminates the additive label-noise term, reducing the decomposition to:
% \[
%     \EE[\cE(\btheta_k)]
%     \eqsim
%     \cS_{\polyak}(k)
%     +
%     \mathcal M_{\polyak, k}.
% \]

% For the fully overdamped regime \(k \gtrsim \frac{1-\rho}{\eta}\), Proposition~\ref{prop:pol_deterministic_bias0} gives the deterministic bias:
% \[
%     \cS_{\polyak}(k) \eqsim \left(\frac{1-\rho}{\eta k}\right)^s,
% \]
% while Proposition~\ref{prop:pol_self_noise_bounds0} yields the multiplicative noise:
% \[
%     \mathcal M_{\polyak, k}
%     \eqsim
%     \frac{\eta}{B(1-\rho)}
%     \left(\frac{1-\rho}{\eta k}\right)^{s \wedge \left(2-\frac1\beta\right)}.
% \]
% Substituting these two asymptotic bounds directly gives:
% \[
%     \EE[\cE(\btheta_k)]
%     \eqsim
%     \left(\frac{1-\rho}{\eta k}\right)^s
%     +
%     \frac{\eta}{B(1-\rho)}
%     \left(\frac{1-\rho}{\eta k}\right)^{s \wedge \left(2-\frac1\beta\right)}.
% \]
% This proves the stated theorem.
% \end{proof}




% \subsection{Loss in the mixed overdamped--underdamped regime}
% \label{app:pol_mixed_loss}

% When underdamped modes are present, the deterministic bias acquires an exponentially decaying oscillatory transient alongside the polynomial tail. The boundary between them is  $J_{\max}$~\ref{eq:J_max}.

% The following proposition characterizes the deterministic bias of Polyak momentum in the mixed overdamped--underdamped regime.

% \begin{proposition}[Deterministic loss, mixed regime]
% \label{prop:pol_deterministic_bias}
% Under Assumptions~\ref{ass:diagonal-covariance} and~\ref{ass:power-law}, assume the stability condition
% $\eta\lesssim \min\{1,B(1-\rho)\}$. For the mixed regime \(\eta\gtrsim(1-\rho)^2\), we have:
% \begin{equation*}
%     \cS_{\polyak}(k)
%     \eqsim
%     \begin{cases}
%     \rho^k+\left(\frac{(1-\rho)^2}{\eta}\right)^s, &\qquad k\lesssim\frac{1}{1-\rho},\\
%     \rho^k+\left(\frac{1-\rho}{\eta k}\right)^s, &\qquad k\gtrsim\frac{1}{1-\rho}.
%     \end{cases}
% \end{equation*}
% \end{proposition}

% \begin{proof}
% Let the coordinate multiplier $h_k(j):= -e_k^{\mathrm{det}}(j)/\theta_j^*$. Then it satisfies
% \[
%     h_{k+1}(j)=(1+\rho-\eta\lambda_j)h_k(j)-\rho h_{k-1}(j),
%     \qquad
%     h_0(j)=1,
%     \qquad
%     h_1(j)=1-\eta\lambda_j.
% \]
% The characteristic discriminant is negative precisely when
% \[
%     (1-\sqrt\rho)^2<\eta\lambda_j<(1+\sqrt\rho)^2.
% \]
% Thus the comparison between $\eta$ and $(1-\rho)^2$ determines whether a coordinate is overdamped or underdamped. In the mixed regime $\eta \gtrsim (1-\rho)^2$, the crossover index is defined by $\eta\lambda_{J_{\max}} \eqsim (1-\rho)^2$. We decompose the deterministic bias into the overdamped tail ($j > J_{\max}$) and the underdamped leading block ($j \le J_{\max}$).

% For overdamped coordinates in the stable small-step range, the root comparison used in Proposition~\ref{prop:pol_deterministic_bias0} gives constants $c_1,C_1>0$ such that
% \begin{equation}
% \label{eq:pol_momentum_overdamped_multiplier_upper}
%     |h_k(j)|
%     \le
%     C_1\exp\left(-c_1\frac{\eta\lambda_j}{1-\rho}k\right).
% \end{equation}
% Moreover, if $\eta\lambda_j\le c_1(1-\rho)^2$ and $k\ge C_1(1-\rho)^{-1}$, the slow root dominates the fast root and
% \begin{equation*}
%     |h_k(j)|
%     \ge
%     c_1\exp\left(-C_1\frac{\eta\lambda_j}{1-\rho}k\right).
% \end{equation*}
% To evaluate the overdamped contribution to the bias, we discuss two cases for $k$:

% \textbf{Case 1: $k \lesssim \frac{1}{1-\rho}$.} \\
% In this regime, for all overdamped modes $j > J_{\max}$, we have $\eta\lambda_j \lesssim (1-\rho)^2$. \\
% Consequently, the exponent satisfies $\frac{\eta\lambda_j}{1-\rho}k \lesssim \frac{(1-\rho)^2}{1-\rho}\frac{1}{1-\rho} = 1$. The exponential decay is bounded away from zero, so $|h_k(j)| \eqsim 1$. Therefore, the sum is strictly dominated by its lower limit:
% \[
%     \sum_{j > J_{\max}} \lambda_j(\theta_j^*)^2 h_k(j)^2 \eqsim \sum_{j > J_{\max}} j^{-1-\beta s} \eqsim (J_{\max})^{-\beta s} \eqsim \left(\frac{(1-\rho)^2}{\eta}\right)^s.
% \]

% \textbf{Case 2: $k \gtrsim \frac{1}{1-\rho}$.} \\
% Since $\lambda_j\eqsim j^{-\beta}$, the time condition implies that every
% $j\ge \left(\frac{C\eta k}{1-\rho}\right)^{1/\beta}$
% satisfies $\eta\lambda_j\le c_1(1-\rho)^2$, after increasing $C$ if necessary. \\
% Using
% $\lambda_j(e_0(j))^2\eqsim j^{-1-\beta s}$, we get
% \begin{equation}
% \label{eq:pol_momentum_poly_lower}
%     \sum_{j > J_{\max}} \lambda_j(\theta_j^*)^2 h_k(j)^2
%     \ge
%     c\sum_{j\ge (C\eta k/(1-\rho))^{1/\beta}}j^{-1-\beta s}
%     \ge
%     c\left(\frac{1-\rho}{\eta k}\right)^s .
% \end{equation}
% Conversely, by~\eqref{eq:pol_momentum_overdamped_multiplier_upper} and the same integral comparison,
% \begin{equation*}
%     \sum_{j > J_{\max}}
%     \lambda_j(\theta_j^*)^2h_k(j)^2
%     \le
%     C\sum_{j\ge1}j^{-1-\beta s}
%     \exp\left(-c_1\frac{\eta k}{1-\rho}j^{-\beta}\right)
%     \le
%     C\left(\frac{1-\rho}{\eta k}\right)^s .
% \end{equation*}
% Therefore the overdamped tail contributes exactly the polynomial term, up to constants.

% It remains to estimate the underdamped contribution. For $j \le J_{\max}$, write the two roots as
% \[
%     \sqrt{\rho}e^{i\omega_j},
%     \quad
%     \sqrt{\rho}e^{-i\omega_j},
%     \qquad \text{where} \quad
%     \cos \omega_j=\frac{1+\rho-\eta\lambda_j}{2\sqrt \rho}.
% \]
% Solving the two initial conditions gives the exact identity
% \begin{equation}
% \label{eq:pol_underdamped_h_formula}
%     h_k(j)
%     =\rho^{k/2}
%     \left[
%         \cos(k\omega_j)
%         +\frac{1-\rho-\eta\lambda_j}{2\sqrt{\rho}\sin \omega_j}
%         \sin(k\omega_j)
%     \right].
% \end{equation}
% The sine term is retained in the oscillatory estimates below. Set
% \[
%     A_j:=\frac{1-\rho-\eta\lambda_j}{2\sqrt{\rho}\sin \omega_j}.
% \]
% On the separated underdamped band
% \[
%     C_{\rm mom}(1-\rho)^2\le \eta\lambda_j\le c_{\rm st}/2,
% \]
% one has
% \[
%     4\rho\sin^2\omega_j
%     =\bigl(\eta\lambda_j-(1-\sqrt\rho)^2\bigr)
%      \bigl((1+\sqrt\rho)^2-\eta\lambda_j\bigr)
%     \ge c\eta\lambda_j,
% \]
% and
% \[
%     (1-\rho-\eta\lambda_j)^2
%     \le(1-\rho)^2,
% \]
% hence $|A_j|\le C$. 
% Thus~\eqref{eq:pol_underdamped_h_formula} immediately gives the upper bound
% \begin{equation*}
%     \sum_{j \le J_{\max}}
%     \lambda_j(\theta_j^*)^2h_k(j)^2
%     \le
%     C\rho^k,
% \end{equation*}
% because $\sum_{j\ge1}\lambda_j(\theta_j^*)^2<\infty$.

% For the lower bound of the underdamped contribution, we evaluate the sum over the underdamped modes $j \le J_{\max}$, which is given by:
% \[
%     \cS_{\mathrm{under}}(k) = \rho^k \sum_{j \le J_{\max}} \lambda_j (\theta_j^*)^2 (1+A_j^2) \cos^2(k\omega_j - \psi_j).
% \]
% To isolate the oscillatory behavior, we define the positive spectral weights $w_j$ and decompose the cosine-squared terms as follows:

% \[
%     \cS_{\mathrm{under}}(k) =\sum_{j \le J_{\max}} w_j \cos^2(k\omega_j - \psi_j) = \frac{1}{2} \sum_{j \le J_{\max}} w_j + \frac{1}{2} \sum_{j \le J_{\max}} w_j \cos(2k\omega_j - 2\psi_j),
% \]
% where $w_j = \lambda_j (\theta_j^*)^2 (1+A_j^2).$

% The first sum represents the stable static mass, which satisfies $\sum_{j \le J_{\max}} w_j \ge M_w > 0$. To analyze the oscillatory second sum, we view the indices as a continuous variable $x \in [1, J_{\max}]$. For the continuous oscillatory integral $I(k)$, we define:
% \[
%     I(k) = \int_{1}^{J_{\max}} w(x) \cos(2k\omega(x) - 2\psi(x)) \, dx.
% \]
% We introduce the substitution $u = \omega(x)$ to map the continuous domain $[1, J_{\max}]$ to $[u_{\min}, u_{\max}]$, denoting the inverse function as $x(u)$.
% Under this substitution, the oscillatory integral is represented by:
% \[
%     I(k) = \int_{u_{\min}}^{u_{\max}} g(u) \cos(2ku - \phi(u)) \, du, \qquad \text{where}\ g(u) = w(x(u)) \left| x'(u) \right|, \phi(u) = 2\psi(x(u)).
% \]
% To bound the integral, we perform integration by parts using the exact differential form:
% \[
%     d\left(\sin(2ku - \phi(u))\right) = (2k - \phi'(u)) \cos(2ku - \phi(u)) \, du
% \]
% Applying this integration yields:
% \[
%     I(k) = \int_{u_{\min}}^{u_{\max}} \frac{g(u)}{2k - \phi'(u)} \, d\left(\sin(2ku - \phi(u))\right).
% \]
% Let $v_k(u) = \frac{g(u)}{2k - \phi'(u)},$
% and then evaluating this expression completely gives the boundary terms and the residual integral:
% \[
%     I(k) = \left[ v_k(u) \sin(2ku - \phi(u)) \right]_{u_{\min}}^{u_{\max}} - \int_{u_{\min}}^{u_{\max}} v_k'(u) \sin(2ku - \phi(u)) \, du.
% \]
% Assuming that $g(u)$, $\phi(u)$, and their derivatives are continuous and bounded on $[u_{\min}, u_{\max}]$, we define the maximum values:
% \[
%     M_0 = \max |g(u)|, \quad M_1 = \max |g'(u)|, \quad L_1 = \max |\phi'(u)|, \quad L_2 = \max |\phi''(u)|
% \]
% For a sufficiently large step $k$ satisfying $k \ge L_1$, we have $2k - \phi'(u) \ge k$, which allows us to write:
% \[
%     k \ge L_1, \quad 2k - \phi'(u) \ge k
% \]
% \[
%     |v_k(u)| \le \frac{M_0}{k}, \qquad |v_k'(u)| = \left| \frac{g'(u)}{2k - \phi'(u)} + \frac{g(u)\phi''(u)}{(2k - \phi'(u))^2} \right| \le \frac{M_1 + M_0 L_2}{k}.
% \]
% Combining these bounds on the boundary terms and the integral, we bound the absolute value of $I(k)$ as:
% \[
%     |I(k)| \le \frac{2M_0 + (u_{\max} - u_{\min})(M_1 + M_0 L_2)}{k} =: \frac{C_1}{k}.
% \]
% This establishes the temporal decay of the continuous oscillatory component. Thus, \(I(k)=O(1/k)\) as \(k\to\infty\).

% We find that for any step $k$ satisfying:
% \[
%     k \ge C\max\left\{\frac{1}{1-\rho}, \frac{1-\rho}{\eta}\right\}
% \]
% The oscillatory term is strictly dominated by a small fraction of the static mass:
% \[
%     |I(k)| \le \frac{1}{4} \int_{1}^{J_{\max}} w(x) \, dx
% \]
% This guarantees that the oscillatory interference cannot cancel the static component, yielding a strictly positive lower bound for the sum:
% \begin{equation*}
%     \sum_{j \le J_{\max}} \lambda_j (\theta_j^*)^2 (1+A_j^2) \cos^2(k\omega_j - \psi_j) \ge c > 0.
% \end{equation*}
% Thus, the discrete sum is strictly bounded away from zero, yielding $\cS_{\mathrm{under}}(k) \ge c \rho^k$.

% Synthesizing the overdamped and underdamped contributions, we obtain the full deterministic bias $\cS_{\polyak}(k) = \cS_{\mathrm{over}}(k) + \cS_{\mathrm{under}}(k)$:
% \begin{itemize}
%     \item When $k \lesssim \frac{1-\rho}{\eta}$, combining the short-time overdamped term with the underdamped bound gives:
%     \[
%         \cS_{\polyak}(k) \eqsim \rho^k + \left(\frac{(1-\rho)^2}{\eta}\right)^s.
%     \]
%     \item When $k \gtrsim \frac{1-\rho}{\eta}$, combining the long-time overdamped polynomial tail with the underdamped bound gives:
%     \[
%         \cS_{\polyak}(k) \eqsim \rho^k + \left(\frac{1-\rho}{\eta k}\right)^s.
%     \]
% \end{itemize}
% This concludes the proof.
% \end{proof}


% \begin{lemma}[Convolution Kernel, Mixed Regime]
% \label{lem:pol_kernel}
% Under Assumptions~\ref{ass:diagonal-covariance} and~\ref{ass:power-law}, assume the stability condition
% $\eta\lesssim \min\{1,B(1-\rho)\}$. For the mixed regime \(\eta\gtrsim(1-\rho)^2\), write \(\mathcal{K}_{\polyak}(k)=\cK_{\mathrm{under}}(k)+\cK_{\mathrm{over}}(k)\), we have:
% \begin{equation*}
%     \mathcal{K}_{\mathrm{over}}(k)
%     \eqsim
%     \begin{cases}
%     k^2 \left(\frac{(1-\rho)^2}{\eta}\right)^{2-\frac1\beta},&\qquad k\lesssim \frac{1}{1-\rho},\\
%     \frac{1}{(1-\rho)^2}\left(\frac{1-\rho}{\eta k}\right)^{2-\frac1\beta},&\qquad k\gtrsim\frac{1}{1-\rho}.
%     \end{cases}
% \end{equation*}
% \begin{equation*}
%     \mathcal{K}_{\mathrm{under}}(k)
%     \eqsim
% \begin{cases}
%     k^2\rho^k,&\qquad k\lesssim\frac{1}{\sqrt{\eta}},\\
%     \frac{\rho^k}{\eta},&\qquad k\gtrsim\frac{1}{\sqrt{\eta}}.
%     \end{cases}
% \end{equation*}
% \end{lemma}

% \begin{proof}
% Recall that the convolution kernel is defined as
% \[
%     \mathcal K_{\polyak}(k)=\sum_{j\ge1}\lambda_j^2G_k(j)^2,
%     \qquad
%     G_k(j)=\frac{y_{j,+}^{k+1}-y_{j,-}^{k+1}}{y_{j,+}-y_{j,-}}.
% \]
% We split the sum into the overdamped tail $\mathcal{K}_{\mathrm{over}}(k)$ where $\eta\lambda_j\lesssim(1-\rho)^2$, and the underdamped leading block $\mathcal{K}_{\mathrm{under}}(k)$ where $\eta\lambda_j\gtrsim(1-\rho)^2$. The crossover index is $J_{\max} \eqsim \left(\frac{\eta}{(1-\rho)^2}\right)^{1/\beta}$.

% \textbf{1. Overdamped Contribution $\mathcal{K}_{\mathrm{over}}(k)$:} \\
% We divide the evaluation into two temporal stages, mirroring the proof method of Lemma~\ref{lem:pol_kernel0}.

% \textit{Stage 1: $k \lesssim \frac{1}{1-\rho}$.} \\
% For the overdamped coordinates with $\eta\lambda_j\lesssim (1-\rho)^2$, the roots $y_{j,+}, y_{j,-}$ are real and close to 1. In this short-time regime, the exponential decay is negligible, and all terms in the algebraic expansion of $G_k(j)$ are strictly positive, yielding $G_k(j) \eqsim k$. Thus, the overdamped part is evaluated by summing over the tail $j \gtrsim J_{\max}$:
% \[
%     \mathcal{K}_{\mathrm{over}}(k) \eqsim \sum_{\eta\lambda_j\lesssim (1-\rho)^2} \lambda_j^2 k^2 = k^2 \sum_{j \gtrsim J_{\max}} j^{-2\beta} \eqsim k^2 (J_{\max})^{1-2\beta}.
% \]
% Substituting $J_{\max} \eqsim \left(\frac{\eta}{(1-\rho)^2}\right)^{1/\beta}$, we obtain $(J_{\max})^{1-2\beta} \eqsim \left(\frac{\eta}{(1-\rho)^2}\right)^{\frac{1-2\beta}{\beta}} = \left(\frac{(1-\rho)^2}{\eta}\right)^{2-\frac1\beta}$. Therefore:
% \[
%     \mathcal{K}_{\mathrm{over}}(k) \eqsim k^2 \left(\frac{(1-\rho)^2}{\eta}\right)^{2-\frac1\beta}.
% \]

% \textit{Stage 2: $k \gtrsim \frac{1}{1-\rho}$.} \\
% For $k \gtrsim \frac{1}{1-\rho}$, the root comparison gives
% \[
%     |G_k(j)|
%     \eqsim
%     \frac{1}{1-\rho}
%     \exp\left(-c\frac{\eta\lambda_j}{1-\rho}k\right),
% \]
% with constants independent of $j,\eta,\rho,k$. Hence the overdamped part is comparable to
% \[
%     \mathcal{K}_{\mathrm{over}}(k)
%     \eqsim
%     \frac{1}{(1-\rho)^2}
%     \sum_{\eta\lambda_j\lesssim (1-\rho)^2}
%     \lambda_j^2
%     \exp\left(-c\frac{\eta\lambda_j}{1-\rho}k\right).
% \]
% Using $\lambda_j\eqsim j^{-\beta}$, the weighted spectral counting rule is
% \[
%     \sum_j \lambda_j^2 F(\lambda_j)
%     \eqsim
%     \int u^{1-\frac1\beta}F(u)\,du .
% \]
% Therefore the overdamped sum is bounded above and below by constant multiples of
% \[
%     \frac{1}{(1-\rho)^2}
%     \int_0^{(1-\rho)^2/\eta}
%     u^{1-\frac1\beta}
%     \exp\left(-c\frac{\eta k}{1-\rho}u\right)\,du .
% \]
% Since $k \gtrsim \frac{1}{1-\rho}$, the exponential decay activates well within the integration domain. The integral evaluates to $\mathcal{O}\left( \left(\frac{1-\rho}{\eta k}\right)^{2-\frac1\beta} \right)$, giving exactly
% \begin{equation*}
%     \mathcal K_{\mathrm{over}}(k)
%     \eqsim
%     \frac{1}{(1-\rho)^2}
%     \left(\frac{1-\rho}{\eta k}\right)^{2-\frac1\beta}.
% \end{equation*}

% \textbf{2. Underdamped Contribution $\mathcal{K}_{\mathrm{under}}(k)$:} \\
% It remains to compute the underdamped part. For $j \le J_{\max}$, write
% \[
%     y_{j,+}=\sqrt \rho e^{i\omega_j},
%     \quad
%     y_{j,-}=\sqrt \rho e^{-i\omega_j},
%     \qquad \text{where} \quad
%     \cos \omega_j=\frac{1+\rho-\eta\lambda_j}{2\sqrt \rho}.
% \]
% Then
% \begin{equation*}
%     G_k(j)=\rho^{k/2}\frac{\sin((k+1)\omega_j)}{\sin \omega_j}.
% \end{equation*}
% On the separated underdamped region $\eta\lambda_j\gtrsim(1-\rho)^2$, stability gives $\sin^2\omega_j \eqsim \eta\lambda_j$. We evaluate this in two stages:

% \textit{Stage 1: $k \lesssim \frac{1}{\sqrt{\eta}}$.} \\
% Since $\lambda_j = \mathcal{O}(1)$, we have $\omega_j \approx \sin\omega_j \eqsim \sqrt{\eta\lambda_j} \lesssim \sqrt{\eta}$. In this short-time regime $k \lesssim \frac{1}{\sqrt{\eta}}$, the angle $(k+1)\omega_j$ is bounded by a small constant of $\mathcal{O}(1)$, so the small-angle approximation applies: $\sin((k+1)\omega_j) \eqsim (k+1)\omega_j \eqsim k\sin\omega_j$.
% Substituting this into the exact formula gives:
% \[
%     G_k(j) = \rho^{k/2}\frac{\sin((k+1)\omega_j)}{\sin \omega_j} \eqsim \rho^{k/2} \frac{k \sin\omega_j}{\sin\omega_j} = k \rho^{k/2}.
% \]
% Consequently, the underdamped contribution is:
% \[
%     \mathcal K_{\mathrm{under}}(k) = \sum_{j \le J_{\max}} \lambda_j^2 G_k(j)^2 \eqsim \sum_{j \le J_{\max}} \lambda_j^2 k^2 \rho^k = k^2 \rho^k \sum_{j \le J_{\max}} \lambda_j^2.
% \]
% Since $\beta > 1$, the spectral mass $\sum_{j \ge 1} \lambda_j^2$ is a convergent sum and bounded by a positive constant $\mathcal{O}(1)$. As $J_{\max}$ is large, $\sum_{j \le J_{\max}} \lambda_j^2 \eqsim 1$. Thus, we obtain exactly:
% \[
%     \mathcal K_{\mathrm{under}}(k) \eqsim k^2 \rho^k.
% \]

% \textit{Stage 2: $k \gtrsim \frac{1}{\sqrt{\eta}}$.} \\
% On the separated underdamped region $\eta\lambda_j\gtrsim(1-\rho)^2$, stability gives
% \[
%     4\rho\sin^2\omega_j
%     =\bigl(\eta\lambda_j-(1-\sqrt\rho)^2\bigr)
%      \bigl((1+\sqrt\rho)^2-\eta\lambda_j\bigr),
% \]
% and hence
% \[
%     \frac{c}{\eta\lambda_j}
%     \le
%     \frac{1}{\sin^2\omega_j}
%     \le
%     \frac{C}{\eta\lambda_j}.
% \]
% Consequently, the underdamped contribution is controlled by
% \[
%     \frac{\rho^k}{\eta}
%     \sum_{\eta\lambda_j\gtrsim(1-\rho)^2}
%     \lambda_j\sin^2((k+1)\omega_j).
% \]
% The upper bound follows immediately from $\sin^2\le1$:
% \[
%     \mathcal K_{\mathrm{under}}(k)
%     \le
%     C\frac{\rho^k}{\eta}
%     \sum_{\eta\lambda_j\gtrsim(1-\rho)^2}\lambda_j.
% \]
% For the lower bound of the underdamped contribution:
% \[
%     \mathcal{K}_{\mathrm{under}}(k) = \sum_{j \le J_{\max}} \lambda_j^2 G_k(j)^2 = \frac{\rho^k}{\eta} \sum_{j \le J_{\max}} \left(\frac{\eta\lambda_j}{\sin^2\omega_j}\right) \lambda_j \sin^2((k+1)\omega_j).
% \]
% Using the approximation:
% $\frac{\eta\lambda_j}{\sin^2\omega_j} \eqsim 1,$
% we decompose the weighted sum:
% \[
%     S(k) = \sum_{j \le J_{\max}} \lambda_j \sin^2((k+1)\omega_j) = \frac{1}{2} \sum_{j \le J_{\max}} \lambda_j - \frac{1}{2} \sum_{j \le J_{\max}} \lambda_j \cos(2(k+1)\omega_j).
% \]
% With the static bound:
% $\sum_{j \le J_{\max}} \lambda_j \ge M_\lambda > 0,$
% we define the continuous oscillatory integral:
% \[
%     J(k) = \int_{1}^{J_{\max}} \lambda(x) \cos(2(k+1)\omega(x)) \, dx.
% \]
% Using the substitution:
% \[
%     u = \omega(x), \qquad [1,J_{\max}] \to [u_{\min}, u_{\max}]
% \]
% We have
% \[
%     J(k) = \int_{u_{\min}}^{u_{\max}} h(u) \cos(2(k+1)u) \, du, \qquad h(u) = \lambda(x(u)) \left| x'(u) \right|.
% \]
% By integration by parts:
% \begin{multline*}
%     J(k) = \int_{u_{\min}}^{u_{\max}} \frac{h(u)}{2(k+1)} \, d\big(\sin(2(k+1)u)\big) = \left[ \frac{h(u)}{2(k+1)} \sin(2(k+1)u) \right]_{u_{\min}}^{u_{\max}}\\ - \int_{u_{\min}}^{u_{\max}} \frac{h'(u)}{2(k+1)} \sin(2(k+1)u) \, du.
% \end{multline*}
% Defining the limits and bounds:
% \[
%     N_0 = \max |h(u)|, \qquad N_1 = \max |h'(u)|.
% \]
% We estimate the absolute value:
% \[
%     |J(k)| \le \frac{|h(u_{\max})| + |h(u_{\min})|}{2(k+1)} + \int_{u_{\min}}^{u_{\max}} \frac{|h'(u)|}{2(k+1)} \, du \le \frac{2N_0 + (u_{\max} - u_{\min})N_1}{2(k+1)} =: \frac{C_2}{k+1}.
% \]
% The oscillatory integral has temporal decay:
% $\mathcal{O}(1/k).$
% The oscillatory integral satisfies:
% \[
%     |J(k)| \le \frac{1}{4} \int_{1}^{J_{\max}} \lambda(x) \, dx.
% \]
% So we have
% \[
%     \sum_{j \le J_{\max}} \lambda_j \sin^2((k+1)\omega_j) \ge c_0 > 0.
% \]
% Thus
% \begin{equation*}
%     \mathcal K_{\mathrm{under}}(k)
%     \eqsim
%     \frac{\rho^k}{\eta}
%     \sum_{\eta\lambda_j\gtrsim(1-\rho)^2}\lambda_j .
% \end{equation*}
% Finally, under the separated momentum condition, the underdamped spectral mass is bounded above and below by positive constants. Indeed, another use of $\lambda_j\eqsim j^{-\beta}$ gives
% \[
%     \sum_{\eta\lambda_j\gtrsim(1-\rho)^2}\lambda_j
%     \eqsim
%     \int_{(1-\rho)^2/\eta}^{1}u^{-1/\beta}\,du
%     \eqsim
%     1-
%     \left(\frac{(1-\rho)^2}{\eta}\right)^{1-\frac1\beta}
%     \eqsim 1.
% \]
% Therefore
% \begin{equation*}
%     \mathcal K_{\mathrm{under}}(k)
%     \eqsim
%     \frac{\rho^k}{\eta}.
% \end{equation*}
% Synthesizing both the overdamped and underdamped parts concludes the proof.
% \end{proof}


% \begin{lemma}[Convolution with Underdamped Kernel]
% \label{lem:conv_under}
% Under the assumptions of Proposition~\ref{prop:pol_deterministic_bias} and Lemma~\ref{lem:pol_kernel}, and under the stability condition
% $\eta\lesssim B(1-\rho)$, for $k\gtrsim \frac{1}{1-\rho}$, the convolution of the deterministic bias with the underdamped kernel satisfies:
% \begin{equation*}
%     \sum_{m=0}^{k-1} \cS_{\polyak}(m) \mathcal{K}_{\mathrm{under}}(k-1-m) 
%     \eqsim 
%     \frac{k\rho^k}{\eta} + \frac{1}{\eta(1-\rho)} \left(\frac{1-\rho}{\eta k}\right)^s.
% \end{equation*}
% \end{lemma}
% \begin{proof}
% We evaluate the convolution $\cS_{\polyak} * \mathcal{K}_{\mathrm{under}}$ by dividing the summation into three distinct stages based on the temporal bounds.

% \textbf{Stage 1:} For $0 \le m \lesssim \frac{1}{1-\rho}$, we have:
% \begin{align*}
%     \sum_{m=1}^{\frac{1}{1-\rho}} \left(\rho^m + \left(\frac{(1-\rho)^2}{\eta}\right)^s\right) \frac{\rho^{k-1-m}}{\eta}
%     &\simeq \sum_{m=1}^{\frac{1}{1-\rho}} \rho^m \frac{\rho^{k-1-m}}{\eta} + \sum_{m=1}^{\frac{1}{1-\rho}} \left(\frac{(1-\rho)^2}{\eta}\right)^s \frac{\rho^{k-1-m}}{\eta} \\
%     &\simeq \frac{\rho^k}{\eta(1-\rho)} + \frac{\rho^k}{\eta(1-\rho)} \left(\frac{(1-\rho)^2}{\eta}\right)^s \\
%     &\simeq \frac{\rho^k}{\eta(1-\rho)}.
% \end{align*}
% The last step rigorously follows since $\eta \gtrsim (1-\rho)^2$ in the mixed regime, implying $\left(\frac{(1-\rho)^2}{\eta}\right)^s \lesssim 1$.

% \textbf{Stage 2:} For $\frac{1}{1-\rho} \lesssim m \lesssim k - \frac{1}{\sqrt{\eta}}$, we explicitly expand the convolution and split the summation of the polynomial-exponential cross term at $m = \frac{k}{2}$, guided by the exponential-polynomial convolution logic (Lemma~\ref{lem:pol_convolution_estimates1}):
% \begin{align*}
%     &\sum_{m=\frac{1}{1-\rho}}^{k-\frac{1}{\sqrt{\eta}}} \left(\rho^m + \left(\frac{1-\rho}{\eta m}\right)^s\right) \frac{\rho^{k-1-m}}{\eta}\\
%     &\simeq \sum_{m=\frac{1}{1-\rho}}^{k-\frac{1}{\sqrt{\eta}}} \rho^m \frac{\rho^{k-m}}{\eta} + \sum_{m=\frac{1}{1-\rho}}^{k-\frac{1}{\sqrt{\eta}}} \left(\frac{1-\rho}{\eta m}\right)^s \frac{\rho^{k-m}}{\eta} \\
%     &\simeq \sum_{m=\frac{1}{1-\rho}}^{k-\frac{1}{\sqrt{\eta}}} \frac{\rho^k}{\eta} + \frac{1}{\eta} \left(\frac{1-\rho}{\eta}\right)^s \left( \sum_{m=\frac{1}{1-\rho}}^{\frac{k}{2}} \frac{\rho^{k-m}}{m^s} + \sum_{m=\frac{k}{2}}^{k-\frac{1}{\sqrt{\eta}}} \frac{\rho^{k-m}}{m^s} \right).
% \end{align*}
% For the first term, summing the constant summand directly yields $\simeq \frac{k\rho^k}{\eta}$. For the split summation, in the first half ($m \le \frac{k}{2}$), the exponential factor is heavily damped as $\rho^{k-m} \le \rho^{k/2}$, making this part exponentially small and strictly subsumed. In the second half ($m \ge \frac{k}{2}$), the polynomial term is smoothly bounded by $m^{-s} \simeq k^{-s}$. Thus, we can rigorously extract the polynomial component and evaluate the remaining geometric series:
% \begin{align*}
%     &\simeq \frac{k\rho^k}{\eta} + \frac{1}{\eta} \left(\frac{1-\rho}{\eta}\right)^s \left( \rho^{k/2} \sum_{m=\frac{1}{1-\rho}}^{\frac{k}{2}} m^{-s} + k^{-s} \sum_{m=\frac{k}{2}}^{k-\frac{1}{\sqrt{\eta}}} \rho^{k-m} \right) \\
%     &\simeq \frac{k\rho^k}{\eta} + \frac{1}{\eta} \left(\frac{1-\rho}{\eta k}\right)^s \sum_{r=\frac{1}{\sqrt{\eta}}}^{\frac{k}{2}} \rho^r .
% \end{align*}
% Since we operate in the mixed regime where $\eta \gtrsim (1-\rho)^2$, it guarantees that $\frac{1}{\sqrt{\eta}} \lesssim \frac{1}{1-\rho}$. This implies the starting term of the geometric series $\rho^{1/\sqrt{\eta}} \simeq 1$, allowing the sum to evaluate to $\simeq \frac{1}{1-\rho}$. We finally obtain:
% \begin{align*}
%     &\simeq \frac{k\rho^k}{\eta} + \frac{1}{\eta} \left(\frac{1-\rho}{\eta k}\right)^s \cdot \frac{1}{1-\rho} \\
%     &\simeq \frac{k\rho^k}{\eta} + \frac{1}{\eta(1-\rho)} \left(\frac{1-\rho}{\eta k}\right)^s.
% \end{align*}

% \textbf{Stage 3:} For $k - \frac{1}{\sqrt{\eta}} \lesssim m \le k-1$, utilizing the short-time asymptotic form of the underdamped kernel:
% \begin{align*}
%     \sum_{m=k-\frac{1}{\sqrt{\eta}}}^{k-1} \left(\rho^m + \left(\frac{1-\rho}{\eta m}\right)^s\right) (k-1-m)^2 \rho^{k-1-m}
%     &\simeq \sum_{m=k-\frac{1}{\sqrt{\eta}}}^{k-1} \left(\rho^k + \left(\frac{1-\rho}{\eta k}\right)^s\right) (k-m)^2 \cdot 1 \\
%     &\simeq \left(\rho^k + \left(\frac{1-\rho}{\eta k}\right)^s\right) \frac{1}{\eta^{3/2}} \\
%     &\simeq \frac{\rho^k}{\eta^{3/2}} + \frac{1}{\eta^{3/2}} \left(\frac{1-\rho}{\eta k}\right)^s.
% \end{align*}
% To bound this, since $\eta \gtrsim (1-\rho)^2$, it holds that $\sqrt{\eta} \gtrsim 1-\rho$, which guarantees $\frac{1}{\eta^{3/2}} \lesssim \frac{1}{\eta(1-\rho)}$. Thus, this stage is strictly subsumed by the preceding stages.

% Combining all three stages, we obtain the total bound:
% \begin{equation*}
%     \cS_{\polyak} * \mathcal{K}_{\mathrm{under}} \simeq \frac{k\rho^k}{\eta} + \frac{1}{\eta(1-\rho)} \left(\frac{1-\rho}{\eta k}\right)^s.
% \end{equation*}
% \end{proof}


% \begin{lemma}[Convolution with Overdamped Kernel]
% \label{lem:conv_over}
% Under the assumptions of Proposition~\ref{prop:pol_deterministic_bias} and Lemma~\ref{lem:pol_kernel}, and under the stability condition
% $\eta\lesssim B(1-\rho)$, for $k\gtrsim \frac{1}{1-\rho}$, the convolution of the deterministic bias with the overdamped kernel satisfies:
% \begin{align*}
%     &\sum_{m=0}^{k-1} \cS_{\polyak}(m) \mathcal{K}_{\mathrm{over}}(k-1-m) \simeq\quad\\ 
%     &\frac{1}{(1-\rho)^3} \left(\frac{1-\rho}{\eta k}\right)^{2-\frac{1}{\beta}} 
%     + \frac{1}{(1-\rho)^3} \left(\frac{(1-\rho)^2}{\eta}\right)^{s \vee (2-\frac{1}{\beta})} \left(\frac{1-\rho}{\eta k}\right)^{s \wedge (2-\frac{1}{\beta})} \\
%     &+ \frac{1}{(1-\rho)^3} \left(\frac{(1-\rho)^2}{\eta}\right)^{2-\frac{1}{\beta}} \rho^k 
%     + \frac{1}{(1-\rho)^3} \left(\frac{(1-\rho)^2}{\eta}\right)^{2-\frac{1}{\beta}} \left(\frac{1-\rho}{\eta k}\right)^s .
% \end{align*}


% \end{lemma}


% \begin{proof}
% We evaluate the convolution $\cS_{\polyak} * \mathcal{K}_{\mathrm{over}}$ by dividing the sum into three distinct temporal stages.

% \textbf{Stage 1:} For $0 \le m \lesssim \frac{1}{1-\rho}$,  since $\eta \gtrsim (1-\rho)^2$, we apply the short-time bias bounds:
% \begin{align*}
%     &\sum_{m=0}^{\frac{1}{1-\rho}} \left(\rho^m + \left(\frac{(1-\rho)^2}{\eta}\right)^s\right) \frac{1}{(1-\rho)^2} \left(\frac{1-\rho}{\eta (k-1-m)}\right)^{2-\frac{1}{\beta}}\\
%     &\simeq \frac{1}{1-\rho} \cdot 1 \cdot \frac{1}{(1-\rho)^2} \left(\frac{1-\rho}{\eta k}\right)^{2-\frac{1}{\beta}} + \left(\frac{(1-\rho)^2}{\eta}\right)^s \frac{1}{1-\rho} \frac{1}{(1-\rho)^2} \left(\frac{1-\rho}{\eta k}\right)^{2-\frac{1}{\beta}} \\
%     &\simeq \frac{1}{(1-\rho)^3} \left(\frac{1-\rho}{\eta k}\right)^{2-\frac{1}{\beta}}.
% \end{align*}

% \textbf{Stage 2:} For $\frac{1}{1-\rho} \lesssim m \le k - \frac{1}{1-\rho}$, the sum is explicitly expanded into four components:
% \begin{align*}
%     &\sum_{m=\frac{1}{1-\rho}}^{k-\frac{1}{1-\rho}} \left(\rho^m + \left(\frac{1-\rho}{\eta m}\right)^s\right) \frac{1}{(1-\rho)^2} \left(\frac{1-\rho}{\eta (k-1-m)}\right)^{2-\frac{1}{\beta}} \\
%     &\simeq \frac{1}{(1-\rho)^2} \sum_{m=\frac{1}{1-\rho}}^{\frac{k}{2}} \rho^m \left(\frac{1-\rho}{\eta k}\right)^{2-\frac{1}{\beta}} + \frac{1}{(1-\rho)^2} \sum_{m=\frac{k}{2}}^{k-\frac{1}{1-\rho}} \rho^m \left(\frac{1-\rho}{\eta (k-1-m)}\right)^{2-\frac{1}{\beta}} \\
%     &\quad + \frac{1}{(1-\rho)^2} \sum_{m=\frac{1}{1-\rho}}^{\frac{k}{2}} \left(\frac{1-\rho}{\eta m}\right)^s \left(\frac{1-\rho}{\eta k}\right)^{2-\frac{1}{\beta}} + \frac{1}{(1-\rho)^2} \sum_{m=\frac{k}{2}}^{k-\frac{1}{1-\rho}} \left(\frac{1-\rho}{\eta k}\right)^s \left(\frac{1-\rho}{\eta (k-1-m)}\right)^{2-\frac{1}{\beta}}.
% \end{align*}
% For the first two terms, since $\sum_{m=\frac{1}{1-\rho}}^{\frac{k}{2}} \rho^m \simeq \frac{1}{1-\rho}$, the first summary term simplifies to $\simeq \frac{1}{1-\rho} \left(\frac{1-\rho}{\eta k}\right)^{2-\frac{1}{\beta}}$. For the second summary term, by the exponential-polynomial convolution Lemma~\ref{lem:pol_convolution_estimates1}, it is bounded by $\lesssim \frac{1}{1-\rho} \left(\frac{1-\rho}{\eta k}\right)^{2-\frac{1}{\beta}}$. 

% Continuing the evaluation of the polynomial-polynomial cross terms:
% \begin{align*}
%     &\simeq \frac{1}{(1-\rho)^3} \left(\frac{1-\rho}{\eta k}\right)^{2-\frac{1}{\beta}} + \frac{1}{(1-\rho)^2} \left(\frac{1-\rho}{\eta}\right)^{2-\frac{1}{\beta}+s} \cdot \begin{cases} k^{1-s} \cdot k^{-(2-\frac{1}{\beta})} & \text{if } s < 1 \\ \left(\frac{1}{1-\rho}\right)^{1-s} \cdot k^{-(2-\frac{1}{\beta})} & \text{if } s > 1 \end{cases} \\
%     &\quad + \frac{1}{(1-\rho)^2} \left(\frac{1-\rho}{\eta}\right)^{2-\frac{1}{\beta}+s} \cdot \left(\frac{1}{1-\rho}\right)^{2-\frac{1}{\beta}} \cdot \frac{1}{k^s} \\
%     &\simeq \frac{1}{(1-\rho)^3} \left(\frac{1-\rho}{\eta k}\right)^{2-\frac{1}{\beta}} + \frac{1}{(1-\rho)^3} \left(\frac{1-\rho}{\eta}\right)^{2-\frac{1}{\beta}+s} \left(\frac{1}{(1-\rho) k }\right)^{s \wedge (2-\frac{1}{\beta})} \cdot \left(1-\rho\right)^{2-\frac{1}{\beta}+s} \\
%     &\simeq \frac{1}{(1-\rho)^3} \left(\frac{1-\rho}{\eta k}\right)^{2-\frac{1}{\beta}} + \frac{1}{(1-\rho)^3} \left(\frac{(1-\rho)^2}{\eta}\right)^{s \vee (2-\frac{1}{\beta})} \cdot \left(\frac{1-\rho}{\eta k}\right)^{s \wedge (2-\frac{1}{\beta})}.
% \end{align*}

% \textbf{Stage 3:} For $k - \frac{1}{1-\rho} \lesssim m \le k-1$, exploiting the short-time scaling of the overdamped kernel:
% \begin{align*}
%     &\sum_{m=k-\frac{1}{1-\rho}}^{k-1} \left(\rho^m + \left(\frac{1-\rho}{\eta m}\right)^s\right) \frac{1}{(1-\rho)^2} \left(\frac{(1-\rho)^2}{\eta}\right)^{2-\frac{1}{\beta}} \\
%     &\simeq \sum_{m=k-\frac{1}{1-\rho}}^{k-1} \rho^k \frac{1}{(1-\rho)^2} \left(\frac{(1-\rho)^2}{\eta}\right)^{2-\frac{1}{\beta}} \cdot 1 + \frac{1}{(1-\rho)^2} \left(\frac{1-\rho}{\eta k}\right)^s \left(\frac{(1-\rho)^2}{\eta}\right)^{2-\frac{1}{\beta}} \sum_{m=k-\frac{1}{1-\rho}}^{k-1} 1 \\
%     &\simeq \frac{1}{(1-\rho)^3} \left(\frac{(1-\rho)^2}{\eta}\right)^{2-\frac{1}{\beta}} \rho^k + \frac{1}{(1-\rho)^3} \left(\frac{1-\rho}{\eta k}\right)^s \left(\frac{(1-\rho)^2}{\eta}\right)^{2-\frac{1}{\beta}}.
% \end{align*}

% Summing all corresponding parts across the three stages, the total bound is established as:
% \begin{align*}
%     \cS_{\polyak} * \mathcal{K}_{\mathrm{over}} \simeq \quad &\frac{1}{(1-\rho)^3} \left(\frac{1-\rho}{\eta k}\right)^{2-\frac{1}{\beta}} + \frac{1}{(1-\rho)^3} \left(\frac{(1-\rho)^2}{\eta}\right)^{s \vee (2-\frac{1}{\beta})} \cdot \left(\frac{1-\rho}{\eta k}\right)^{s \wedge (2-\frac{1}{\beta})} \\
%     &+ \frac{1}{(1-\rho)^3} \left(\frac{(1-\rho)^2}{\eta}\right)^{2-\frac{1}{\beta}} \rho^k + \frac{1}{(1-\rho)^3} \left(\frac{1-\rho}{\eta k}\right)^s \left(\frac{(1-\rho)^2}{\eta}\right)^{2-\frac{1}{\beta}}.
% \end{align*}
% \end{proof}


% \begin{lemma}[Total Deterministic Convolution]
% \label{lem:conv_total}
% Under the assumptions of Proposition~\ref{prop:pol_deterministic_bias} and Lemma~\ref{lem:pol_kernel}, and under the stability condition
% $\eta\lesssim B(1-\rho)$, for $k\gtrsim \frac{1}{1-\rho}$, the total convolution between the deterministic bias and the Polyak momentum kernel satisfies:
% \begin{equation*}
%     \sum_{m=0}^{k-1} \cS_{\polyak}(m) \mathcal{K}_{\polyak}(k-1-m) \simeq \frac{k\rho^k}{\eta} + \frac{1}{(1-\rho)^3} \left(\frac{1-\rho}{\eta k}\right)^{2-\frac{1}{\beta}} + \frac{1}{\eta(1-\rho)} \left(\frac{1-\rho}{\eta k}\right)^s.
% \end{equation*}
% \end{lemma}
% \begin{proof}
% Synthesizing the results from $\cS_{\polyak} * \mathcal{K}_{\mathrm{over}}$, we first rigorously verify the following crucial bounding relation:
% \begin{align*}
%     &\frac{1}{(1-\rho)^3} \left(\frac{(1-\rho)^2}{\eta}\right)^{2-\frac{1}{\beta}} \lesssim \frac{1}{\eta(1-\rho)} \\
%     & \Leftrightarrow \frac{1}{(1-\rho)^3} \left(\frac{(1-\rho)^2}{\eta}\right)^{2-\frac{1}{\beta}} \cdot \eta(1-\rho) \lesssim 1 \\
%     & \Leftrightarrow \left(\frac{(1-\rho)^2}{\eta}\right)^{1-\frac{1}{\beta}} \lesssim 1.
% \end{align*}
% This inequality holds since $\eta \gtrsim (1-\rho)^2$ and the polynomial slack is naturally absorbed within the relevant spectral limits. Armed with this relation, we can strictly simplify the cross terms:
% \[
%     \frac{1}{(1-\rho)^3} \left(\frac{(1-\rho)^2}{\eta}\right)^{2-\frac{1}{\beta}} \left(\frac{1-\rho}{\eta k}\right)^s \lesssim \frac{1}{\eta(1-\rho)} \left(\frac{1-\rho}{\eta k}\right)^s,
% \]
% and for the exponential part:
% \[
%     \frac{1}{(1-\rho)^3} \left(\frac{(1-\rho)^2}{\eta}\right)^{2-\frac{1}{\beta}} \rho^k \lesssim \frac{1}{\eta(1-\rho)} \rho^k \lesssim \frac{k}{\eta} \rho^k.
% \]
% Furthermore, for the main polynomial cross term, it decomposes gracefully as:
% \begin{align*}
%     \frac{1}{(1-\rho)^3} \left(\frac{(1-\rho)^2}{\eta}\right)^{s \vee (2-\frac{1}{\beta})} \left(\frac{1-\rho}{\eta k}\right)^{s \wedge (2-\frac{1}{\beta})} 
%     &\lesssim \frac{1}{(1-\rho)^3} \left(\frac{1-\rho}{\eta k}\right)^{2-\frac{1}{\beta}} + \frac{1}{\eta(1-\rho)} \left(\frac{1-\rho}{\eta k}\right)^s.
% \end{align*}
% Adding the aggregated contribution from the underdamped convolution $\cS_{\polyak} * \mathcal{K}_{\mathrm{under}}$, the total convolution sum is strictly bounded by:
% \begin{equation*}
%     \cS_{\polyak} * \mathcal{K}_{\polyak} \simeq \frac{k\rho^k}{\eta} + \frac{1}{(1-\rho)^3} \left(\frac{1-\rho}{\eta k}\right)^{2-\frac{1}{\beta}} + \frac{1}{\eta(1-\rho)} \left(\frac{1-\rho}{\eta k}\right)^s.
% \end{equation*}
% This completes the proof.
% \end{proof}






% The mixed-regime multiplicative noise intertwines exponential transients and polynomial tails, yielding a composite bound.

% \begin{proposition}[Multiplicative Noise Bounds, Mixed Regime]
% \label{prop:pol_self_noise_bounds}
% Under the assumptions of Proposition~\ref{prop:pol_deterministic_bias} and Lemma~\ref{lem:pol_kernel}, and under the stability condition $\eta\lesssim B(1-\rho)$, we have:
% \begin{align*}
%     \frac{\eta^2}{B}\sum_{m=0}^{k-1}\mathcal{K}_{\polyak}(k-1-m)\EE[\cE(\btheta_m)]
%     \eqsim
%     \frac{\eta^2}{B} \left[ \frac{k\rho^k}{\eta} + \frac{1}{(1-\rho)^3} \left(\frac{1-\rho}{\eta k}\right)^{2-\frac{1}{\beta}} + \frac{1}{\eta(1-\rho)} \left(\frac{1-\rho}{\eta k}\right)^s \right]\\
%     + \frac{\eta^2\sigma^2}{B^2(1-\rho)^2}.
% \end{align*}
% \end{proposition}

% \begin{proof}
% Let
% \[
%     \mathcal M_k:=
%     \frac{\eta^2}{B}\sum_{m=0}^{k-1}\mathcal K_{\polyak}(k-1-m)\EE[\cE(\btheta_m)].
% \]
% By Proposition~\ref{prop:pol_risk_recursion}, the expected excess risk is bounded by:
% \[
%     \EE[\cE(\btheta_m)]
%     \le
%     C\left[
%         \mathcal D_{\polyak}(m)
%         +\frac{\eta^2\sigma^2}{B}\sum_{r=0}^{m-1}\mathcal K_{\polyak}(r)
%         +\mathcal M_m
%     \right].
% \]
% Substituting this bound into the definition of $\mathcal M_k$ yields three convolution terms:
% \begin{align}
% \label{eq:prop_mixed_bootstrap_start}
%     \mathcal M_k
%     \le{}&
%     C\frac{\eta^2}{B}\sum_{m=0}^{k-1}
%         \mathcal K_{\polyak}(k-1-m)\mathcal D_{\polyak}(m) \nonumber \\
%     & + C\frac{\eta^4\sigma^2}{B^2}\sum_{m=0}^{k-1}\mathcal K_{\polyak}(k-1-m)\sum_{r=0}^{m-1}\mathcal K_{\polyak}(r) \nonumber \\
%     & + C\frac{\eta^2}{B}\sum_{m=0}^{k-1}
%         \mathcal K_{\polyak}(k-1-m)\mathcal M_m .
% \end{align}

% We first evaluate the label-noise convolution term. By changing the summation indices and exploiting the non-negativity of the kernel, we have:
% \begin{align*}
%     \frac{\eta^4\sigma^2}{B^2}\sum_{m=0}^{k-1}\mathcal K_{\polyak}(k-1-m)\sum_{r=0}^{m-1}\mathcal K_{\polyak}(r)
%     &= \frac{\eta^4\sigma^2}{B^2} \sum_{n=0}^{k-1} \mathcal K_{\polyak}(n) \sum_{r=0}^{k-2-n} \mathcal K_{\polyak}(r) \\
%     &\le \frac{\eta^4\sigma^2}{B^2} \left( \sum_{n=0}^{\infty} \mathcal K_{\polyak}(n) \right)^2.
% \end{align*}
% From the exact saturated Green-energy identity in Lemma~\ref{lem:pol_label_noise}, the infinite sum over the mixed regime kernel converges to $\sum_{n=0}^{\infty} \mathcal K_{\polyak}(n) \eqsim \frac{1}{\eta(1-\rho)}$. Therefore, this term is tightly bounded by $\mathcal{O}\left( \frac{\eta^2\sigma^2}{B^2(1-\rho)^2} \right)$. Due to the positive definiteness of the sums over restricted ranges capturing a constant fraction for large $k$, the matching lower bound strictly holds.

% Second, we substitute the deterministic convolution term directly using Lemma~\ref{lem:conv_total}. The total convolution between the deterministic bias and the Polyak momentum kernel satisfies:
% \begin{align*}
%     \frac{\eta^2}{B}\sum_{m=0}^{k-1}\mathcal K_{\polyak}(k-1-m)\mathcal D_{\polyak}(m)
%     &\simeq \frac{\eta^2}{B} \left[ \frac{k\rho^k}{\eta} + \frac{1}{(1-\rho)^3} \left(\frac{1-\rho}{\eta k}\right)^{2-\frac{1}{\beta}} + \frac{1}{\eta(1-\rho)} \left(\frac{1-\rho}{\eta k}\right)^s \right].
% \end{align*}

% Finally, we absorb the feedback term in~\eqref{eq:prop_mixed_bootstrap_start}. Exchanging the order of summation gives a two-kernel convolution:
% \begin{align*}
%     \frac{\eta^2}{B}\sum_{m=0}^{k-1}\mathcal K_{\polyak}(k-1-m)\mathcal M_m
%     &= \frac{\eta^2}{B}\sum_{\ell=0}^{k-2} \left[ \frac{\eta^2}{B}\sum_{m=\ell+1}^{k-1} \mathcal K_{\polyak}(k-1-m)\mathcal K_{\polyak}(m-1-\ell) \right] \EE[\cE(\btheta_\ell)].
% \end{align*}
% Applying the sub-convolution bound $\sum_{r} \mathcal K_{\polyak}(n-r)\mathcal K_{\polyak}(r) \le C \left(\sum_{r=0}^{\infty}\mathcal K_{\polyak}(r)\right) \mathcal K_{\polyak}(n) \le \frac{C}{\eta(1-\rho)}\mathcal K_{\polyak}(n)$, we have:
% \begin{align*}
%     \frac{\eta^2}{B}\sum_{m=0}^{k-1}\mathcal K_{\polyak}(k-1-m)\mathcal M_m
%     &\le C\frac{\eta}{B(1-\rho)} \frac{\eta^2}{B}\sum_{\ell=0}^{k-2}\mathcal K_{\polyak}(k-2-\ell)\EE[\cE(\btheta_\ell)] \\
%     &\le C\frac{\eta}{B(1-\rho)}\mathcal M_{k-1} \le C'\frac{\eta}{B(1-\rho)}\mathcal M_k .
% \end{align*}
% By enforcing the stability condition $\eta\le cB(1-\rho)$ with a sufficiently small $c>0$, we ensure $C'\frac{\eta}{B(1-\rho)} \le 1/2$. This permits the feedback term to be  subtracted and absorbed into the left-hand side of~\eqref{eq:prop_mixed_bootstrap_start}.

% Combining the rigorously bounded deterministic convolution and the label-noise convolution, and noting that the lower bound follows identically from the two-sided nature of the constituent bounds, completes the proof:
% \[
%     \mathcal M_k \eqsim \frac{\eta^2}{B} \left[ \frac{k\rho^k}{\eta} + \frac{1}{(1-\rho)^3} \left(\frac{1-\rho}{\eta k}\right)^{2-\frac{1}{\beta}} + \frac{1}{\eta(1-\rho)} \left(\frac{1-\rho}{\eta k}\right)^s \right] + \frac{\eta^2\sigma^2}{B^2(1-\rho)^2}.
% \]
% \end{proof}

% Synthesizing the components provides the complete two-sided scaling law for Polyak momentum in the mixed regime.

% \begin{theorem}[Two-Sided Scaling Law, Mixed Regime]
% \label{thm:pol_two_sided_scaling}
% Under Assumptions~\ref{ass:diagonal-covariance} and~\ref{ass:power-law}, the stability condition $\eta\lesssim \min\{1,B(1-\rho)\}$, and for the mixed regime $\eta\gtrsim(1-\rho)^2$ with $k \gtrsim \frac{1}{1-\rho}$, the expected excess risk satisfies:
% \[
%     \EE[\cE(\btheta_k)]
%     \eqsim
%     \left(1+\frac{\eta k}{B}\right) \rho^k
%     +
%     \left(\frac{1-\rho}{\eta k}\right)^s
%     +
%     \frac{\sigma^2\eta}{B(1-\rho)}.
% \]
% \end{theorem}

% \begin{proof}
% By the expected risk recursion established in Proposition~\ref{prop:pol_risk_recursion}, the total expected excess risk decomposes into the deterministic bias, the accumulated label noise, and the multiplicative noise feedback $\mathcal{M}_k$:
% \[
%     \EE[\cE(\btheta_k)]
%     \eqsim
%     \cS_{\polyak}(k)
%     +
%     \frac{\eta^2\sigma^2}{B}\sum_{\tau=0}^{k-1}\mathcal{K}_{\polyak}(\tau)
%     +
%     \mathcal M_k.
% \]

% Substituting the asymptotic bounds derived in the preceding propositions for the mixed regime ($k \gtrsim \frac{1-\rho}{\eta}$), we assemble the three main components:
% \begin{enumerate}
%     \item \textbf{Deterministic Bias} (Proposition~\ref{prop:pol_deterministic_bias}):
%     \[
%         \cS_{\polyak}(k) \eqsim \rho^k + \left(\frac{1-\rho}{\eta k}\right)^s.
%     \]
%     \item \textbf{Label Noise} (Lemma~\ref{lem:pol_label_noise}):
%     \[
%         \frac{\eta^2\sigma^2}{B}\sum_{\tau=0}^{k-1}\mathcal{K}_{\polyak}(\tau) \eqsim \frac{\sigma^2\eta}{B(1-\rho)}.
%     \]
%     \item \textbf{Multiplicative Noise} (Proposition~\ref{prop:pol_self_noise_bounds}):
%     \[
%         \mathcal M_k \eqsim \frac{\eta k\rho^k}{B} + \frac{\eta^2}{B(1-\rho)^3} \left(\frac{1-\rho}{\eta k}\right)^{2-\frac{1}{\beta}} + \frac{\eta}{B(1-\rho)} \left(\frac{1-\rho}{\eta k}\right)^s + \frac{\eta^2\sigma^2}{B^2(1-\rho)^2}.
%     \]
% \end{enumerate}

% Summing these components, we systematically absorb the lower-order terms under the stability condition $\eta \le cB(1-\rho)$. 

% First, the constant noise floor arising from the multiplicative noise is strictly bounded by the primary label noise term and is thus absorbed:
% \[
%     \frac{\eta^2\sigma^2}{B^2(1-\rho)^2} = \left(\frac{\eta}{B(1-\rho)}\right) \frac{\sigma^2\eta}{B(1-\rho)} \lesssim \frac{\sigma^2\eta}{B(1-\rho)}.
% \]

% Second, the polynomial tail from the multiplicative noise that shares the source exponent $s$ is absorbed into the deterministic polynomial bias:
% \[
%     \frac{\eta}{B(1-\rho)} \left(\frac{1-\rho}{\eta k}\right)^s \lesssim \left(\frac{1-\rho}{\eta k}\right)^s.
% \]

% Third, we examine the polynomial tail from the multiplicative noise with exponent $2-\frac{1}{\beta}$, denoted $M_3(k) = \frac{\eta^2}{B(1-\rho)^3}\left(\frac{1-\rho}{\eta k}\right)^{2-\frac1\beta}$.
% Since the mixed regime specifies $k \gtrsim \frac{1}{1-\rho}$, we have $\eta k \gtrsim \frac{\eta}{1-\rho}$, which implies $\frac{1-\rho}{\eta k} \lesssim \frac{(1-\rho)^2}{\eta}$. 
% Substituting this upper bound into $M_3(k)$, we obtain:
% \begin{align*}
%     M_3(k) &\lesssim \frac{\eta^2}{B(1-\rho)^3} \left(\frac{(1-\rho)^2}{\eta}\right)^{2-\frac{1}{\beta}} \\
%     &= \frac{\eta}{B(1-\rho)} \frac{\eta}{(1-\rho)^2} \left(\frac{(1-\rho)^2}{\eta}\right)^{2-\frac{1}{\beta}} \\
%     &= \frac{\eta}{B(1-\rho)} \left(\frac{(1-\rho)^2}{\eta}\right)^{1-\frac{1}{\beta}}.
% \end{align*}
% In the mixed regime, $\eta \gtrsim (1-\rho)^2$, meaning $\frac{(1-\rho)^2}{\eta} \lesssim 1$. Because $\beta > 1$, the exponent $1-\frac{1}{\beta} > 0$. Thus, $\left(\frac{(1-\rho)^2}{\eta}\right)^{1-\frac{1}{\beta}} \lesssim 1$.
% Consequently, this yields:
% \begin{equation*}
%     M_3(k) \lesssim \frac{\eta}{B(1-\rho)}.
% \end{equation*}
% Treating the label noise variance $\sigma^2$ as a constant $\Theta(1)$, this term is unconditionally bounded by the primary label noise $\text{LN}(k) \eqsim \frac{\sigma^2\eta}{B(1-\rho)}$, and is therefore entirely absorbed.

% Finally, grouping the remaining unabsorbable terms—the underdamped exponential transients and the distinct polynomial tails—yields the exact scaling law:
% \begin{align*}
%     \EE[\cE(\btheta_k)]
%     &\eqsim \left( \rho^k + \frac{\eta k\rho^k}{B} \right) + \left(\frac{1-\rho}{\eta k}\right)^s + \frac{\sigma^2\eta}{B(1-\rho)} \\
%     &= \rho^k\left(1+\frac{\eta k}{B}\right) + \left(\frac{1-\rho}{\eta k}\right)^s + \frac{\sigma^2\eta}{B(1-\rho)}.
% \end{align*}
% This completes the proof.
% \end{proof}


% \begin{theorem}[Two-Sided Scaling Law, No Label Noise, Mixed Regime]
% \label{thm:pol_two_sided_scaling_noiseless}
% Under Assumptions~\ref{ass:diagonal-covariance} and~\ref{ass:power-law}, with \(\sigma=0\), the stability condition \(\eta\lesssim \min\{1,B(1-\rho)\}\), and in the mixed regime \(\eta\gtrsim(1-\rho)^2\), for \(k \gtrsim \frac{1}{1-\rho}\), the expected excess risk satisfies:
% \begin{equation}
% \label{eq:pol_two_sided_scaling_noiseless}
%     \EE[\cE(\btheta_k)]
%     \eqsim
%     \left(
%         1+\frac{\eta k}{B}
%     \right)\rho^k
%     +
%     \left(\frac{1-\rho}{\eta k}\right)^s
%     +
%     \frac{\eta^2}{B(1-\rho)^3}
%     \left(\frac{1-\rho}{\eta k}\right)^{2-\frac1\beta}.
% \end{equation}
% \end{theorem}

% \begin{proof}
% Following Proposition~\ref{prop:pol_risk_recursion} and setting \(\sigma=0\), the expected risk simplifies to:
% \[
%     \EE[\cE(\btheta_k)]
%     \eqsim
%     \cS_{\polyak}(k)
%     +
%     \mathcal M_{\polyak, k}.
% \]
% In the mixed regime \(k \gtrsim \frac{1}{1-\rho}\), Proposition~\ref{prop:pol_deterministic_bias} provides the deterministic bias:
% \[
%     \cS_{\polyak}(k)
%     \eqsim
%     \rho^k
%     +
%     \left(\frac{1-\rho}{\eta k}\right)^s,
% \]
% and Proposition~\ref{prop:pol_self_noise_bounds} gives the multiplicative noise:
% \[
%     \mathcal M_{\polyak, k}
%     \eqsim
%     \frac{\eta k}{B}\rho^k
%     +
%     \frac{\eta^2}{B(1-\rho)^3}
%     \left(\frac{1-\rho}{\eta k}\right)^{2-\frac1\beta}
%     +
%     \frac{\eta}{B(1-\rho)}
%     \left(\frac{1-\rho}{\eta k}\right)^s.
% \]
% Combining the deterministic bias with the multiplicative noise and grouping terms with identical asymptotic forms yields:
% \[
%     \EE[\cE(\btheta_k)]
%     \eqsim
%     \left(1+\frac{\eta k}{B}\right)\rho^k
%     +
%     \left(
%         1+\frac{\eta}{B(1-\rho)}
%     \right)
%     \left(\frac{1-\rho}{\eta k}\right)^s
%     +
%     \frac{\eta^2}{B(1-\rho)^3}
%     \left(\frac{1-\rho}{\eta k}\right)^{2-\frac1\beta}.
% \]
% By the stability condition \(\eta \lesssim B(1-\rho)\), we have \(\frac{\eta}{B(1-\rho)} \lesssim 1\). Hence the factor \(1+\frac{\eta}{B(1-\rho)}\) is bounded by a universal constant and can be absorbed into the asymptotic equivalence, giving:
% \[
%     1+\frac{\eta}{B(1-\rho)} \eqsim 1.
% \]
% Substituting this absorption into the previous display completes the proof.
% \end{proof}



\subsection{Uniform boundedness of expected risk}
\label{app:pol_uniform_boundedness}

As long as the additive label noise is strictly positive ($\sigma > 0$), the scaling law derivation can be  simplified. We can prove that the expected risk is uniformly bounded, which allows the multiplicative gradient noise to be completely absorbed by the additive label noise floor.

% \begin{lemma}[Uniform Boundedness of Expected Risk]
% \label{lem:pol_loss_bounded}
% Suppose $\sigma > 0$. Under the stability condition $\eta \le c B(1-\rho)$ for a sufficiently small universal constant $c > 0$, the expected excess risk of Polyak momentum is uniformly bounded for all $k \ge 0$:
% \[
%     \EE[\cE(\btheta_k)] \le M,
% \]
% where $M > 0$ is a constant independent of $k$, $\eta$, $B$, and $\rho$.
% \end{lemma}
% \begin{proof}
% We proceed by induction on $k$. By Proposition~\ref{prop:pol_risk_recursion}, there exist universal constants $C_1, C_2, C_3 > 0$ such that for any $k \ge 1$:
% \begin{equation} \label{eq:pol_risk_recursion_bound}
%     \EE[\cE(\btheta_k)] \le C_1 \cS_{\polyak}(k) + C_2 \frac{\eta^2\sigma^2}{B}\sum_{\tau=0}^{k-1}\mathcal{K}_{\polyak}(\tau) + C_3 \frac{\eta^2}{B}\sum_{\tau=0}^{k-1}\mathcal{K}_{\polyak}(\tau)\EE[\cE(\btheta_{k-1-\tau})].
% \end{equation}

% First, by Proposition~\ref{prop:pol_mode_decomposition}, all characteristic roots $y_{j,\pm}$ have moduli strictly bounded by 1 for all stable modes. Thus, the unperturbed deterministic trajectory is unconditionally stable. The deterministic multiplier $h_k(j) = -e_k^{\det}(j)/\theta_j^*$ is uniformly bounded by a universal constant $C_h$ for all $k$ and $j$. Consequently, the deterministic bias is bounded by the initial risk:
% \[
%     \cS_{\polyak}(k) \le C_h^2 \cE(\btheta_0) =: M_1.
% \]

% Second, by the saturated Green-energy estimate in Lemma~\ref{lem:pol_label_noise} (from Appendix~\ref{app:pol_preparation}), the infinite sum of the convolution kernel is bounded by $\sum_{\tau=0}^\infty \mathcal{K}_{\polyak}(\tau) \le \frac{C_K}{\eta(1-\rho)}$. Therefore, the label noise contribution is universally bounded:
% \[
%     C_2 \frac{\eta^2\sigma^2}{B}\sum_{\tau=0}^{k-1}\mathcal{K}_{\polyak}(\tau) \le C_2 \sigma^2 \frac{C_K \eta}{B(1-\rho)}.
% \]
% Under the stability condition $\eta \le c B(1-\rho)$, this term is bounded by $C_2 \sigma^2 C_K c =: M_2$.

% Now, assume inductively that $\EE[\cE(\btheta_m)] \le M$ for all $m < k$. For the multiplicative noise term, we have:
% \[
%     C_3 \frac{\eta^2}{B}\sum_{\tau=0}^{k-1}\mathcal{K}_{\polyak}(\tau)\EE[\cE(\btheta_{k-1-\tau})] \le C_3 \frac{\eta^2}{B} \left( \frac{C_K}{\eta(1-\rho)} \right) M = C_3 C_K \frac{\eta}{B(1-\rho)} M \le C_3 C_K c M.
% \]
% By choosing the stability tuning constant $c$ sufficiently small such that $c \le \frac{1}{2 C_3 C_K}$, the multiplicative noise term is strictly bounded by $M/2$. Substituting these into \eqref{eq:pol_risk_recursion_bound} yields:
% \[
%     \EE[\cE(\btheta_k)] \le M_1 + M_2 + \frac{M}{2}.
% \]
% Setting $M = 2(M_1 + M_2)$ ensures that $\EE[\cE(\btheta_k)] \le M$. The base case $k=0$ holds trivially since $\EE[\cE(\btheta_0)] = \cS_{\polyak}(0) \le M_1 \le M$. This completes the induction.
% \end{proof}




\begin{lemma}[Uniform boundedness of expected risk]
\label{lem:pol_loss_bounded}
Suppose $\sigma > 0$. Under the stability condition $\eta \lesssim 1$ and $\eta \le c B(1-\rho)$ for a sufficiently small universal constant $c > 0$, the expected excess risk of Polyak is uniformly bounded for all $k \ge 0$:
\[
    \EE[\cE(\btheta_k)] \le M,
\]
where $M > 0$ is a constant independent of $k$, $\eta$, $B$, and $\rho$.
\end{lemma}

\begin{proof}
We proceed by induction on $k$. By Proposition~\ref{prop:pol_risk_recursion}, there exist universal constants $C_1, C_2, C_3 > 0$ such that for any $k \ge 1$:
\begin{equation} \label{eq:pol_risk_recursion_bound}
    \EE[\cE(\btheta_k)] \le C_1 \cS_{\polyak}(k) + C_2 \frac{\eta^2\sigma^2}{B}\sum_{\tau=0}^{k-1}\mathcal{K}_{\polyak}(\tau) + C_3 \frac{\eta^2}{B}\sum_{\tau=0}^{k-1}\mathcal{K}_{\polyak}(\tau)\EE[\cE(\btheta_{k-1-\tau})].
\end{equation}

First, Propositions~\ref{prop:pol_deterministic_bias0} and \ref{prop:pol_deterministic_bias} in the following subsections establish that the deterministic bias monotonically decays or is uniformly bounded by its initialization across all regimes. Thus:
\[
    \cS_{\polyak}(k) \le C_h^2 \cE(\btheta_0) =: M_1.
\]

Second, applying $\sum_{\tau=0}^\infty \mathcal{K}_{\polyak}(\tau) \le \frac{C_K}{\eta(1-\rho)}$ from Lemma~\ref{lem:pol_label_noise} bounds the label noise:
\[
    C_2 \frac{\eta^2\sigma^2}{B}\sum_{\tau=0}^{k-1}\mathcal{K}_{\polyak}(\tau) \le C_2 \sigma^2 \frac{C_K \eta}{B(1-\rho)}.
\]
Under $\eta \le c B(1-\rho)$, this is uniformly bounded by $C_2 \sigma^2 C_K c =: M_2$.

Assuming inductively $\EE[\cE(\btheta_m)] \le M$ for $m < k$, the multiplicative noise term yields:
\[
    C_3 \frac{\eta^2}{B}\sum_{\tau=0}^{k-1}\mathcal{K}_{\polyak}(\tau)\EE[\cE(\btheta_{k-1-\tau})] \le C_3 \frac{\eta^2}{B} \left( \frac{C_K}{\eta(1-\rho)} \right) M = C_3 C_K \frac{\eta}{B(1-\rho)} M \le C_3 C_K c M.
\]
Choosing $c \le \frac{1}{2 C_3 C_K}$ limits this term to $M/2$. Substituting into \eqref{eq:pol_risk_recursion_bound} yields:
\[
    \EE[\cE(\btheta_k)] \le M_1 + M_2 + \frac{M}{2}.
\]
Setting $M = 2(M_1 + M_2)$ guarantees $\EE[\cE(\btheta_k)] \le M$. The base case $k=0$ holds trivially since $\EE[\cE(\btheta_0)] = \cS_{\polyak}(0) \le M_1 \le M$. This completes the induction.
\end{proof}




\begin{corollary}[Absorption of multiplicative noise]
\label{cor:pol_noise_absorption}
For $\sigma > 0$, under the stability condition $\eta \lesssim \min\{1,B(1-\rho)\}$, the multiplicative gradient noise is strictly bounded by the additive label noise up to a constant factor:
\[
    \frac{\eta^2}{B}\sum_{\tau=0}^{k-1}\mathcal{K}_{\polyak}(\tau)\EE[\cE(\btheta_{k-1-\tau})] \le \frac{M}{\sigma^2} \left( \frac{\eta^2\sigma^2}{B}\sum_{\tau=0}^{k-1}\mathcal{K}_{\polyak}(\tau) \right).
\]
Consequently, the multiplicative noise can be absorbed, and the expected excess risk is asymptotically equivalent to the sum of the deterministic bias and the label noise:
\[
    \EE[\cE(\btheta_k)] \eqsim \cS_{\polyak}(k) + \frac{\eta^2\sigma^2}{B}\sum_{\tau=0}^{k-1}\mathcal{K}_{\polyak}(\tau).
\]
\end{corollary}

%===========================================================
\subsection{Risk dynamics in the fully overdamped regime}
\label{app:pol_fully_overdamped_loss}

In the fully overdamped regime, the deterministic error trajectory decays monotonically, yielding a polynomial bias.

\begin{proposition}[Deterministic risk, fully overdamped regime]
\label{prop:pol_deterministic_bias0}
Under Assumptions~\ref{ass:diagonal-covariance} and~\ref{ass:power-law}, in the fully overdamped regime $\eta\lesssim (1-\rho)^2$, we have:
\begin{equation*}
    \cS_{\polyak}(k)
    \eqsim
    \begin{cases}
    1,&\qquad k\lesssim \frac{1-\rho}{\eta},\\
    \left(\frac{1-\rho}{\eta k}\right)^s, &\qquad k\gtrsim\frac{1-\rho}{\eta}.\\
    \end{cases}
\end{equation*}
\end{proposition}
\begin{proof}
In the fully overdamped regime, we can assume: $\eta\lambda_1\le c_0 (1-\rho)^2$, with a sufficiently small universal constant $c_0>0$.

Let the deterministic coordinate multiplier $h_k(j):=e_k^{\mathrm{det}}(j)/e_0^{\mathrm{det}}(j) = -e_k^{\mathrm{det}}(j)/\theta_j^*$ satisfy
\[
    h_{k+1}(j)=(1+\rho-\eta\lambda_j)h_k(j)-\rho h_{k-1}(j),
    \qquad
    h_0(j)=1,
    \qquad
    h_1(j)=1-\eta\lambda_j.
\]
The characteristic roots are
\[
    r_{j,+},r_{j,-}
    =
    \frac{1+\rho-\eta\lambda_j\pm\sqrt{(1+\rho-\eta\lambda_j)^2-4\rho}}{2}.
\]
We first give explicit root bounds. Since $\eta\le \frac{c_0}{\lambda_1}(1-\rho)^2$ and $\lambda_j\le\lambda_1$, we have $0\le \eta\lambda_j\le c_0(1-\rho)^2$ for all $j$. Put $r=1-u$. The characteristic polynomial evaluated at $1-u$ is
\[
    (1-u)^2-(1+\rho-\eta\lambda_j)(1-u)+\rho
    =
    \eta\lambda_j-(1-\rho+\eta\lambda_j)u+u^2.
\]
Choosing $c_0>0$ sufficiently small, the two evaluations
\[
    \eta\lambda_j-(1-\rho+\eta\lambda_j)\frac{\eta\lambda_j}{2(1-\rho)}+\left(\frac{\eta\lambda_j}{2(1-\rho)}\right)^2>0,
\]
and
\[
    \eta\lambda_j-(1-\rho+\eta\lambda_j)\frac{2\eta\lambda_j}{1-\rho}+\left(\frac{2\eta\lambda_j}{1-\rho}\right)^2<0
\]
hold uniformly for $0<\eta\lambda_j\le c_0(1-\rho)^2$. Therefore the slow root satisfies
\begin{equation}
\label{eq:pol_slow_root_bounds_sgd}
    1-\frac{2\eta\lambda_j}{1-\rho}
    \le
    r_{j,+}
    \le
    1-\frac{\eta\lambda_j}{2(1-\rho)}.
\end{equation}
The fast root satisfies $r_{j,+}r_{j,-}=\rho=1-(1-\rho)$. Since $r_{j,+}\ge1-2c_0(1-\rho)$, after decreasing $c_0$ if needed,
\begin{equation}
\label{eq:pol_fast_root_bound_sgd}
    0\le r_{j,-}\le 1-\frac{1-\rho}{2}.
\end{equation}

Now write the solution as
\[
    h_k(j)=A_jr_{j,+}^k+B_jr_{j,-}^k.
\]
From $h_0=1$ and $h_1=1-\eta\lambda_j$, we obtain the coefficients
\[
    A_j=\frac{1-\eta\lambda_j-r_{j,-}}{r_{j,+}-r_{j,-}},
    \qquad
    B_j=1-A_j = \frac{r_{j,+}-(1-\eta\lambda_j)}{r_{j,+}-r_{j,-}}.
\]
Using~\eqref{eq:pol_slow_root_bounds_sgd}--\eqref{eq:pol_fast_root_bound_sgd}, and taking $c_0$ sufficiently small, these coefficients obey uniform bounds:
\begin{equation}
\label{eq:pol_sgd_coeff_bounds}
    \frac12\le A_j\le2,
    \qquad
    |B_j|\le2.
\end{equation}
Equivalently, applying the initial conditions, the exact solution can be expressed in the fractional form:
\begin{equation}
\label{eq:hk_exact_form}
    h_k(j) = \frac{r_{j,+} r_{j,-}\left(-r_{j,+}^{k-1} + r_{j,-}^{k-1}\right) + (1-\eta\lambda_j)\left(r_{j,+}^k - r_{j,-}^k\right)}{r_{j,+} - r_{j,-}}.
\end{equation}

To evaluate the deterministic bias $\cS_{\polyak}(k) = \sum_{j\ge1}\lambda_j(\theta_j^*)^2 h_k(j)^2$, we divide the time $k$ into three successive stages.

\textbf{Stage 1: $k \lesssim \frac{1}{1-\rho}$.} \\
In this extremely short-time regime, taking the limits $r_{j,+}\approx r_{j,-} \approx r \approx 1$ and using $r^2 \approx \rho$ in~\eqref{eq:hk_exact_form} yields the heuristic trajectory:
\[
    h_k(j) \simeq -\rho(k-1)r^{k-2} + (1-\eta\lambda_j)k r^{k-1} = k\left((1-\eta\lambda_j)r - \rho \right)r^{k-2} + \rho r^{k-2}.
\]
Since $(1-\eta\lambda_j)r - \rho = \mathcal{O}(1-\rho)$, the first term is bounded by $\mathcal{O}(k(1-\rho)) \lesssim 1$, yielding $h_k(j) \simeq 1$. 
Rigorously, for the leading modes ($j=\mathcal{O}(1)$), since $\eta\lambda_j \le c_0(1-\rho)^2$, the slow root satisfies $r_{j,+} \ge 1 - 2c_0(1-\rho)$. For $k \lesssim \frac{1}{1-\rho}$, we have
\[
    r_{j,+}^k \ge \left(1 - 2c_0(1-\rho)\right)^{\frac{C}{1-\rho}} \eqsim 1.
\]
Since the dynamics are unconditionally stable ($h_k(j) \le 1$), it follows that $h_k(j) \eqsim 1$ for the leading modes. Consequently, the bias is dominated by these non-decayed modes:
\[
    \cS_{\polyak}(k) \eqsim \sum_{j\ge1}\lambda_j(\theta_j^*)^2 \eqsim 1.
\]

\textbf{Stage 2: $\frac{1}{1-\rho} \lesssim k \lesssim \frac{1-\rho}{\eta}$.} \\
For $k \gtrsim \frac{1}{1-\rho}$, the fast root decays exponentially since $r_{j,-}^k \le \left(1-\frac{1-\rho}{2}\right)^k \ll 1$. The multiplier is  dominated by the slow root:
\[
    h_k(j) \simeq \frac{1-\eta\lambda_j - r_{j,-}}{r_{j,+} - r_{j,-}} r_{j,+}^k \simeq r_{j,+}^k.
\]
For the leading modes with $\lambda_j = \Theta(1)$, within the time window $k \lesssim \frac{1-\rho}{\eta}$, the exponent is bounded by:
\[
    r_{j,+}^k \eqsim \left(1 - c\frac{\eta\lambda_j}{1-\rho}\right)^k \eqsim \exp\left(-c\frac{\eta\lambda_j}{1-\rho}k\right) \eqsim 1,
\]
where we used the condition $\frac{\eta\lambda_j}{1-\rho}k \lesssim \lambda_j \le \lambda_1 = \mathcal{O}(1)$. Thus, the leading modes have not yet experienced significant exponential decay, preserving $h_k(j) \eqsim 1$. As a result, the deterministic bias remains unchanged:
\[
    \cS_{\polyak}(k) \eqsim 1.
\]

\textbf{Stage 3: $k \gtrsim \frac{1-\rho}{\eta}$.} \\
In this long-time regime, the exponential decay of the slow root becomes significant across the spectrum. Combining~\eqref{eq:pol_slow_root_bounds_sgd},~\eqref{eq:pol_fast_root_bound_sgd}, and~\eqref{eq:pol_sgd_coeff_bounds}, there are universal constants $c_1,c_2,c_3,C_1>0$ such that, for all $k\ge C_1(1-\rho)^{-1}$,
\begin{equation}
\label{eq:pol_sgd_coordinate_two_sided_decay}
    c_1\exp\left(-c_2\frac{\eta\lambda_j}{1-\rho}k\right)
    \le
    |h_k(j)|
    \le
    C_1\exp\left(-c_3\frac{\eta\lambda_j}{1-\rho}k\right).
\end{equation}
Substituting~\eqref{eq:pol_sgd_coordinate_two_sided_decay} into the deterministic bias gives
\begin{equation}
\label{eq:pol_sgd_bias_spectral_comparison}
    c\sum_{j\ge1}\lambda_j(\theta_j^*)^2
    \exp\left(-C\frac{\eta\lambda_j}{1-\rho}k\right)
    \le
    \cS_{\polyak}(k)
    \le
    C\sum_{j\ge1}\lambda_j(\theta_j^*)^2
    \exp\left(-c\frac{\eta\lambda_j}{1-\rho}k\right).
\end{equation}
By Assumption~\ref{ass:power-law},
\[
    \lambda_j(\theta_j^*)^2\eqsim j^{-1-\beta s}, \qquad \lambda_j\eqsim j^{-\beta}.
\]
Thus it remains to bound sums of the form
\[
    \sum_{j\ge1}j^{-1-\beta s}
    \exp\left(-c\frac{\eta k}{1-\rho}j^{-\beta}\right).
\]
The condition on $k$ implies $\frac{\eta k}{1-\rho}\ge C$.

For the lower bound, set $J:=\left\lceil \left(\frac{\eta k}{1-\rho}\right)^{1/\beta}\right\rceil$. For $j\ge J$, we have $\frac{\eta k}{1-\rho}j^{-\beta}\le1$, so $\exp(-C\frac{\eta k}{1-\rho}j^{-\beta})\ge e^{-C}$. Therefore,
\[
    \sum_{j\ge1}j^{-1-\beta s}\exp\left(-C\frac{\eta k}{1-\rho}j^{-\beta}\right)
    \ge
    e^{-C}\sum_{j\ge J}j^{-1-\beta s}
    \ge
    c\left(\frac{\eta k}{1-\rho}\right)^{-s}.
\]
For the upper bound, the integral comparison with $u=c\frac{\eta k}{1-\rho}x^{-\beta}$ gives
\begin{align*}
    \sum_{j\ge1}j^{-1-\beta s}\exp\left(-c\frac{\eta k}{1-\rho}j^{-\beta}\right)
    &\le
    C\int_1^\infty x^{-1-\beta s}\exp\left(-c\frac{\eta k}{1-\rho}x^{-\beta}\right)\,dx \\
    &=
    C \left(\frac{\eta k}{1-\rho}\right)^{-s}\int_0^{c\frac{\eta k}{1-\rho}}u^{s-1}e^{-u}\,du
    \le
    C \left(\frac{\eta k}{1-\rho}\right)^{-s}.
\end{align*}
Combining the upper and lower bounds in~\eqref{eq:pol_sgd_bias_spectral_comparison} yields
\[
    \cS_{\polyak}(k)
    \eqsim
    \left(\frac{1-\rho}{\eta k}\right)^s.
\]
Synthesizing the three stages completes the proof.
\end{proof}

\begin{theorem}[Scaling law, fully overdamped regime]
\label{thm:pol_scaling_sgd}
Under Assumptions~\ref{ass:diagonal-covariance} and~\ref{ass:power-law}, and the stability condition $\eta \lesssim \min\{1,B(1-\rho)\}$, for the fully overdamped regime $\eta\lesssim(1-\rho)^2$, when $k\gtrsim \frac{1-\rho}{\eta}$, the expected excess risk satisfies:
\[
    \EE[\cE(\btheta_k)]
    \eqsim
    \left(\frac{1-\rho}{\eta k}\right)^s
    +
    \frac{\sigma^2\eta}{B(1-\rho)}.
\]
\end{theorem}
\begin{proof}
By Corollary~\ref{cor:pol_noise_absorption}, the expected risk is completely determined by the deterministic bias and the label noise:
\[
    \EE[\cE(\btheta_k)] \eqsim \cS_{\polyak}(k) + \frac{\eta^2\sigma^2}{B}\sum_{\tau=0}^{k-1}\mathcal{K}_{\polyak}(\tau).
\]
For $k \gtrsim \frac{1-\rho}{\eta}$, we substitute the deterministic bias from Proposition~\ref{prop:pol_deterministic_bias0}:
\[
    \cS_{\polyak}(k) \eqsim \left(\frac{1-\rho}{\eta k}\right)^s.
\]
We substitute the label noise floor from Lemma~\ref{lem:pol_label_noise}:
\[
    \frac{\eta^2\sigma^2}{B}\sum_{\tau=0}^{k-1}\mathcal{K}_{\polyak}(\tau) \eqsim \frac{\sigma^2\eta}{B(1-\rho)}.
\]
Summing these two terms directly yields the scaling law.
\end{proof}











\subsection{Risk dynamics in the mixed-damping regime}
\label{app:pol_mixed_loss}

In the mixed-damping regime, the spectrum is partitioned into underdamped and overdamped modes by the crossover index $J_{\max} \eqsim \left(\frac{\eta}{(1-\rho)^2}\right)^{1/\beta}$. Consequently, the deterministic bias decomposes into an overdamped polynomial tail and an underdamped oscillatory transient, denoted by $\cS_{\mathrm{over}}(k)$ and $\cS_{\mathrm{under}}(k)$. Due to destructive interference among the underdamped modes, $\cS_{\mathrm{under}}(k)$ lacks a strictly positive pointwise lower bound. 

\begin{proposition}[Deterministic bias, mixed-damping regime]
\label{prop:pol_deterministic_bias}
Under Assumptions~\ref{ass:diagonal-covariance} and~\ref{ass:power-law}, for the mixed-damping regime \(\eta\gtrsim(1-\rho)^2\), the deterministic bias decomposes into $\cS_{\polyak}(k) = \cS_{\mathrm{over}}(k) + \cS_{\mathrm{under}}(k)$, where the overdamped tail satisfies the two-sided bounds:
\begin{equation*}
    \cS_{\mathrm{over}}(k)
    \eqsim
    \begin{cases}
    \left(\frac{(1-\rho)^2}{\eta}\right)^s, &\qquad k\lesssim\frac{1}{1-\rho},\\
    \left(\frac{1-\rho}{\eta k}\right)^s, &\qquad k\gtrsim\frac{1}{1-\rho},
    \end{cases}
\end{equation*}
and $\cS_{\mathrm{under}}(k)$ represents the underdamped transient whose properties are formally established in Lemma~\ref{lem:pol_underdamped_transient_properties}.
\end{proposition}

\begin{proof}
Let the coordinate multiplier $h_k(j):= -e_k^{\mathrm{det}}(j)/\theta_j^*$. Then it satisfies
\[
    h_{k+1}(j)=(1+\rho-\eta\lambda_j)h_k(j)-\rho h_{k-1}(j),
    \qquad
    h_0(j)=1,
    \qquad
    h_1(j)=1-\eta\lambda_j.
\]
The characteristic discriminant is negative precisely when
\[
    (1-\sqrt\rho)^2<\eta\lambda_j<(1+\sqrt\rho)^2.
\]
Thus the comparison between $\eta$ and $(1-\rho)^2$ determines whether a coordinate is overdamped or underdamped. In the mixed-damping regime $\eta \gtrsim (1-\rho)^2$, the crossover index is defined by $\eta\lambda_{J_{\max}} \eqsim (1-\rho)^2$. We decompose the deterministic bias into the overdamped tail ($j > J_{\max}$) and the underdamped leading block ($j \le J_{\max}$).

For overdamped coordinates ($j > J_{\max}$) in the stable small-step range, the root comparison used in Proposition~\ref{prop:pol_deterministic_bias0} gives constants $c_1,C_1>0$ such that
\begin{equation}
\label{eq:pol_momentum_overdamped_multiplier_upper}
    |h_k(j)|
    \le
    C_1\exp\left(-c_1\frac{\eta\lambda_j}{1-\rho}k\right).
\end{equation}
Moreover, if $\eta\lambda_j\le c_1(1-\rho)^2$ and $k\ge C_1(1-\rho)^{-1}$, the slow root dominates the fast root and
\begin{equation*}
    |h_k(j)|
    \ge
    c_1\exp\left(-C_1\frac{\eta\lambda_j}{1-\rho}k\right).
\end{equation*}
To evaluate the overdamped contribution $\cS_{\mathrm{over}}(k) = \frac{1}{2} \sum_{j > J_{\max}} \lambda_j(\theta_j^*)^2 h_k(j)^2$, we discuss two cases for $k$:

\textbf{Case 1: $k \lesssim \frac{1}{1-\rho}$.} \\
In this regime, for all overdamped modes $j > J_{\max}$, we have $\eta\lambda_j \lesssim (1-\rho)^2$. 
Consequently, the exponent satisfies $\frac{\eta\lambda_j}{1-\rho}k \lesssim \frac{(1-\rho)^2}{1-\rho}\frac{1}{1-\rho} = 1$. The exponential decay is bounded away from zero, so $|h_k(j)| \eqsim 1$. Therefore, the sum is dominated by its lower limit:
\[
    \cS_{\mathrm{over}}(k) \eqsim \sum_{j > J_{\max}} j^{-1-\beta s} \eqsim (J_{\max})^{-\beta s} \eqsim \left(\frac{(1-\rho)^2}{\eta}\right)^s.
\]

\textbf{Case 2: $k \gtrsim \frac{1}{1-\rho}$.} \\
Since $\lambda_j\eqsim j^{-\beta}$, the time condition implies that every
$j\ge \left(\frac{C\eta k}{1-\rho}\right)^{1/\beta}$
satisfies $\eta\lambda_j\le c_1(1-\rho)^2$, after increasing $C$ if necessary.
Using $\lambda_j(\theta_j^*)^2\eqsim j^{-1-\beta s}$, we get
\begin{equation}
\label{eq:pol_momentum_poly_lower}
    \cS_{\mathrm{over}}(k)
    \ge
    c\sum_{j\ge (C\eta k/(1-\rho))^{1/\beta}}j^{-1-\beta s}
    \ge
    c\left(\frac{1-\rho}{\eta k}\right)^s .
\end{equation}
Conversely, by~\eqref{eq:pol_momentum_overdamped_multiplier_upper} and the integral comparison,
\begin{equation*}
    \cS_{\mathrm{over}}(k)
    \le
    C\sum_{j\ge1}j^{-1-\beta s}
    \exp\left(-c_1\frac{\eta k}{1-\rho}j^{-\beta}\right)
    \le
    C\left(\frac{1-\rho}{\eta k}\right)^s .
\end{equation*}
Therefore the overdamped tail contributes exactly the polynomial term, up to constants. 

The remaining deterministic bias comes from the underdamped modes, defined as $\cS_{\mathrm{under}}(k) := \frac{1}{2} \sum_{j \le J_{\max}} \lambda_j (\theta_j^*)^2 h_k(j)^2$, which will be analyzed in Lemma~\ref{lem:pol_underdamped_transient_properties}.
\end{proof}


\begin{lemma}[Properties of the underdamped transient]
\label{lem:pol_underdamped_transient_properties}
Let $\cS_{\mathrm{under}}(k)$ denote the deterministic bias contributed by the underdamped modes ($j \le J_{\max}$). Under the conditions of Proposition~\ref{prop:pol_deterministic_bias}, for every step $k \ge 0$, we have:
\begin{equation*}
    0 \le \cS_{\mathrm{under}}(k) \le C \rho^k.
\end{equation*}
\end{lemma}

\begin{proof}
For $j \le J_{\max}$, write the two roots as
\[
    \sqrt{\rho}e^{i\omega_j},
    \quad
    \sqrt{\rho}e^{-i\omega_j},
    \qquad \text{where} \quad
    \cos \omega_j=\frac{1+\rho-\eta\lambda_j}{2\sqrt \rho}.
\]
Solving the two initial conditions gives the exact identity
\begin{equation}
\label{eq:pol_underdamped_h_formula}
    h_k(j)
    =\rho^{k/2}
    \left[
        \cos(k\omega_j)
        +\frac{1-\rho-\eta\lambda_j}{2\sqrt{\rho}\sin \omega_j}
        \sin(k\omega_j)
    \right].
\end{equation}
Set
\[
    A_j:=\frac{1-\rho-\eta\lambda_j}{2\sqrt{\rho}\sin \omega_j}.
\]
On the separated underdamped band
\[
    C_{\rm mom}(1-\rho)^2\le \eta\lambda_j\le c_{\rm st}/2,
\]
one has
\[
    4\rho\sin^2\omega_j
    =\bigl(\eta\lambda_j-(1-\sqrt\rho)^2\bigr)
     \bigl((1+\sqrt\rho)^2-\eta\lambda_j\bigr)
    \ge c\eta\lambda_j,
\]
and
\[
    (1-\rho-\eta\lambda_j)^2
    \le(1-\rho)^2,
\]
hence $|A_j|\le C$. 
Thus~\eqref{eq:pol_underdamped_h_formula} immediately gives the upper bound
\begin{equation*}
    \cS_{\mathrm{under}}(k) = \frac{1}{2} \sum_{j \le J_{\max}}
    \lambda_j(\theta_j^*)^2h_k(j)^2
    \le
    C\rho^k,
\end{equation*}
because $\sum_{j\ge1}\lambda_j(\theta_j^*)^2<\infty$. The lower bound is trivially zero due to the potential for destructive interference of the oscillatory terms at specific iterations.
\end{proof}

\begin{theorem}[Scaling law, mixed-damping regime]
\label{thm:pol_two_sided_scaling}
Under Assumptions~\ref{ass:diagonal-covariance} and~\ref{ass:power-law}, and the stability condition $\eta \lesssim \min\{1,B(1-\rho)\}$, for the mixed-damping regime $\eta\gtrsim(1-\rho)^2$ with $k \gtrsim \max\left\{\frac{1}{1-\rho}, \frac{1-\rho}{\eta}\right\}$, the expected excess risk satisfies:
\begin{equation}
\label{eq:pol_mixed_scaling_law}
    \EE[\cE(\btheta_k)] \eqsim \cS_{\mathrm{under}}(k) + \left(\frac{1-\rho}{\eta k}\right)^s + \frac{\sigma^2\eta}{B(1-\rho)},
\end{equation}
where the properties of the underdamped transient $\cS_{\mathrm{under}}(k)$ are established in Lemma~\ref{lem:pol_underdamped_transient_properties}.
\end{theorem}

\begin{proof}
By Corollary~\ref{cor:pol_noise_absorption}, the multiplicative noise is absorbed by the additive label noise floor, yielding:
\begin{equation}
\label{eq:pol_total_risk_mixed}
    \EE[\cE(\btheta_k)] 
    \eqsim 
    \cS_{\mathrm{under}}(k) + \cS_{\mathrm{over}}(k) + \frac{\eta^2\sigma^2}{B}\sum_{\tau=0}^{k-1}\mathcal{K}_{\polyak}(\tau).
\end{equation}

By Lemma~\ref{lem:pol_label_noise}, for $k \gtrsim \max\left\{\frac{1}{1-\rho}, \frac{1-\rho}{\eta}\right\}$, the label noise summation captures the saturated floor:
\begin{equation*}
    \frac{\eta^2\sigma^2}{B}\sum_{\tau=0}^{k-1}\mathcal{K}_{\polyak}(\tau) 
    \eqsim 
    \frac{\sigma^2\eta}{B(1-\rho)}.
\end{equation*}

By Proposition~\ref{prop:pol_deterministic_bias}, since $k \gtrsim \frac{1}{1-\rho}$, the overdamped tail reaches its terminal polynomial decay phase:
\begin{equation*}
    \cS_{\mathrm{over}}(k) \eqsim \left(\frac{1-\rho}{\eta k}\right)^s.
\end{equation*}

Substituting these components directly into \eqref{eq:pol_total_risk_mixed} establishes the equivalence \eqref{eq:pol_mixed_scaling_law}.
\end{proof}



% %===========================================================
% \subsection{Risk dynamics in the mixed regime}
% \label{app:pol_mixed_loss}

% When underdamped modes are present, the deterministic bias acquires an exponentially decaying oscillatory transient alongside the polynomial tail.


% \begin{proposition}[Deterministic loss, mixed regime]
% \label{prop:pol_deterministic_bias}
% Under Assumptions~\ref{ass:diagonal-covariance} and~\ref{ass:power-law}, assume the stability condition
% $\eta\lesssim \min\{1,B(1-\rho)\}$. For the mixed regime \(\eta\gtrsim(1-\rho)^2\), we have:
% \begin{equation*}
%     \cS_{\polyak}(k)
%     \eqsim
%     \begin{cases}
%     \rho^k+\left(\frac{(1-\rho)^2}{\eta}\right)^s, &\qquad k\lesssim\frac{1}{1-\rho},\\
%     \rho^k+\left(\frac{1-\rho}{\eta k}\right)^s, &\qquad k\gtrsim\frac{1}{1-\rho}.
%     \end{cases}
% \end{equation*}
% \end{proposition}
% \begin{proof}
% Let the coordinate multiplier $h_k(j):= -e_k^{\mathrm{det}}(j)/\theta_j^*$. Then it satisfies
% \[
%     h_{k+1}(j)=(1+\rho-\eta\lambda_j)h_k(j)-\rho h_{k-1}(j),
%     \qquad
%     h_0(j)=1,
%     \qquad
%     h_1(j)=1-\eta\lambda_j.
% \]
% The characteristic discriminant is negative precisely when
% \[
%     (1-\sqrt\rho)^2<\eta\lambda_j<(1+\sqrt\rho)^2.
% \]
% Thus the comparison between $\eta$ and $(1-\rho)^2$ determines whether a coordinate is overdamped or underdamped. In the mixed regime $\eta \gtrsim (1-\rho)^2$, the crossover index is defined by $\eta\lambda_{J_{\max}} \eqsim (1-\rho)^2$. We decompose the deterministic bias into the overdamped tail ($j > J_{\max}$) and the underdamped leading block ($j \le J_{\max}$).

% For overdamped coordinates in the stable small-step range, the root comparison used in Proposition~\ref{prop:pol_deterministic_bias0} gives constants $c_1,C_1>0$ such that
% \begin{equation}
% \label{eq:pol_momentum_overdamped_multiplier_upper}
%     |h_k(j)|
%     \le
%     C_1\exp\left(-c_1\frac{\eta\lambda_j}{1-\rho}k\right).
% \end{equation}
% Moreover, if $\eta\lambda_j\le c_1(1-\rho)^2$ and $k\ge C_1(1-\rho)^{-1}$, the slow root dominates the fast root and
% \begin{equation*}
%     |h_k(j)|
%     \ge
%     c_1\exp\left(-C_1\frac{\eta\lambda_j}{1-\rho}k\right).
% \end{equation*}
% To evaluate the overdamped contribution to the bias, we discuss two cases for $k$:

% \textbf{Case 1: $k \lesssim \frac{1}{1-\rho}$.} \\
% In this regime, for all overdamped modes $j > J_{\max}$, we have $\eta\lambda_j \lesssim (1-\rho)^2$. \\
% Consequently, the exponent satisfies $\frac{\eta\lambda_j}{1-\rho}k \lesssim \frac{(1-\rho)^2}{1-\rho}\frac{1}{1-\rho} = 1$. The exponential decay is bounded away from zero, so $|h_k(j)| \eqsim 1$. Therefore, the sum is strictly dominated by its lower limit:
% \[
%     \sum_{j > J_{\max}} \lambda_j(\theta_j^*)^2 h_k(j)^2 \eqsim \sum_{j > J_{\max}} j^{-1-\beta s} \eqsim (J_{\max})^{-\beta s} \eqsim \left(\frac{(1-\rho)^2}{\eta}\right)^s.
% \]

% \textbf{Case 2: $k \gtrsim \frac{1}{1-\rho}$.} \\
% Since $\lambda_j\eqsim j^{-\beta}$, the time condition implies that every
% $j\ge \left(\frac{C\eta k}{1-\rho}\right)^{1/\beta}$
% satisfies $\eta\lambda_j\le c_1(1-\rho)^2$, after increasing $C$ if necessary. \\
% Using
% $\lambda_j(e_0(j))^2\eqsim j^{-1-\beta s}$, we get
% \begin{equation}
% \label{eq:pol_momentum_poly_lower}
%     \sum_{j > J_{\max}} \lambda_j(\theta_j^*)^2 h_k(j)^2
%     \ge
%     c\sum_{j\ge (C\eta k/(1-\rho))^{1/\beta}}j^{-1-\beta s}
%     \ge
%     c\left(\frac{1-\rho}{\eta k}\right)^s .
% \end{equation}
% Conversely, by~\eqref{eq:pol_momentum_overdamped_multiplier_upper} and the same integral comparison,
% \begin{equation*}
%     \sum_{j > J_{\max}}
%     \lambda_j(\theta_j^*)^2h_k(j)^2
%     \le
%     C\sum_{j\ge1}j^{-1-\beta s}
%     \exp\left(-c_1\frac{\eta k}{1-\rho}j^{-\beta}\right)
%     \le
%     C\left(\frac{1-\rho}{\eta k}\right)^s .
% \end{equation*}
% Therefore the overdamped tail contributes exactly the polynomial term, up to constants.

% It remains to estimate the underdamped contribution. For $j \le J_{\max}$, write the two roots as
% \[
%     \sqrt{\rho}e^{i\omega_j},
%     \quad
%     \sqrt{\rho}e^{-i\omega_j},
%     \qquad \text{where} \quad
%     \cos \omega_j=\frac{1+\rho-\eta\lambda_j}{2\sqrt \rho}.
% \]
% Solving the two initial conditions gives the exact identity
% \begin{equation}
% \label{eq:pol_underdamped_h_formula}
%     h_k(j)
%     =\rho^{k/2}
%     \left[
%         \cos(k\omega_j)
%         +\frac{1-\rho-\eta\lambda_j}{2\sqrt{\rho}\sin \omega_j}
%         \sin(k\omega_j)
%     \right].
% \end{equation}
% The sine term is retained in the oscillatory estimates below. Set
% \[
%     A_j:=\frac{1-\rho-\eta\lambda_j}{2\sqrt{\rho}\sin \omega_j}.
% \]
% On the separated underdamped band
% \[
%     C_{\rm mom}(1-\rho)^2\le \eta\lambda_j\le c_{\rm st}/2,
% \]
% one has
% \[
%     4\rho\sin^2\omega_j
%     =\bigl(\eta\lambda_j-(1-\sqrt\rho)^2\bigr)
%      \bigl((1+\sqrt\rho)^2-\eta\lambda_j\bigr)
%     \ge c\eta\lambda_j,
% \]
% and
% \[
%     (1-\rho-\eta\lambda_j)^2
%     \le(1-\rho)^2,
% \]
% hence $|A_j|\le C$. 
% Thus~\eqref{eq:pol_underdamped_h_formula} immediately gives the upper bound
% \begin{equation*}
%     \sum_{j \le J_{\max}}
%     \lambda_j(\theta_j^*)^2h_k(j)^2
%     \le
%     C\rho^k,
% \end{equation*}
% because $\sum_{j\ge1}\lambda_j(\theta_j^*)^2<\infty$.

% For the lower bound of the underdamped contribution, we evaluate the sum over the underdamped modes $j \le J_{\max}$, which is given by:
% \[
%     \cS_{\mathrm{under}}(k) = \rho^k \sum_{j \le J_{\max}} \lambda_j (\theta_j^*)^2 (1+A_j^2) \cos^2(k\omega_j - \psi_j).
% \]
% To isolate the oscillatory behavior, we define the positive spectral weights $w_j$ and decompose the cosine-squared terms as follows:

% \[
%     \cS_{\mathrm{under}}(k) =\sum_{j \le J_{\max}} w_j \cos^2(k\omega_j - \psi_j) = \frac{1}{2} \sum_{j \le J_{\max}} w_j + \frac{1}{2} \sum_{j \le J_{\max}} w_j \cos(2k\omega_j - 2\psi_j),
% \]
% where $w_j = \lambda_j (\theta_j^*)^2 (1+A_j^2).$

% The first sum represents the stable static mass, which satisfies $\sum_{j \le J_{\max}} w_j \ge M_w > 0$. To analyze the oscillatory second sum, we view the indices as a continuous variable $x \in [1, J_{\max}]$. For the continuous oscillatory integral $I(k)$, we define:
% \[
%     I(k) = \int_{1}^{J_{\max}} w(x) \cos(2k\omega(x) - 2\psi(x)) \, dx.
% \]
% We introduce the substitution $u = \omega(x)$ to map the continuous domain $[1, J_{\max}]$ to $[u_{\min}, u_{\max}]$, denoting the inverse function as $x(u)$.
% Under this substitution, the oscillatory integral is represented by:
% \[
%     I(k) = \int_{u_{\min}}^{u_{\max}} g(u) \cos(2ku - \phi(u)) \, du, \qquad \text{where}\ g(u) = w(x(u)) \left| x'(u) \right|, \phi(u) = 2\psi(x(u)).
% \]
% To bound the integral, we perform integration by parts using the exact differential form:
% \[
%     d\left(\sin(2ku - \phi(u))\right) = (2k - \phi'(u)) \cos(2ku - \phi(u)) \, du
% \]
% Applying this integration yields:
% \[
%     I(k) = \int_{u_{\min}}^{u_{\max}} \frac{g(u)}{2k - \phi'(u)} \, d\left(\sin(2ku - \phi(u))\right).
% \]
% Let $v_k(u) = \frac{g(u)}{2k - \phi'(u)},$
% and then evaluating this expression completely gives the boundary terms and the residual integral:
% \[
%     I(k) = \left[ v_k(u) \sin(2ku - \phi(u)) \right]_{u_{\min}}^{u_{\max}} - \int_{u_{\min}}^{u_{\max}} v_k'(u) \sin(2ku - \phi(u)) \, du.
% \]
% Assuming that $g(u)$, $\phi(u)$, and their derivatives are continuous and bounded on $[u_{\min}, u_{\max}]$, we define the maximum values:
% \[
%     M_0 = \max |g(u)|, \quad M_1 = \max |g'(u)|, \quad L_1 = \max |\phi'(u)|, \quad L_2 = \max |\phi''(u)|
% \]
% For a sufficiently large step $k$ satisfying $k \ge L_1$, we have $2k - \phi'(u) \ge k$, which allows us to write:
% \[
%     k \ge L_1, \quad 2k - \phi'(u) \ge k
% \]
% \[
%     |v_k(u)| \le \frac{M_0}{k}, \qquad |v_k'(u)| = \left| \frac{g'(u)}{2k - \phi'(u)} + \frac{g(u)\phi''(u)}{(2k - \phi'(u))^2} \right| \le \frac{M_1 + M_0 L_2}{k}.
% \]
% Combining these bounds on the boundary terms and the integral, we bound the absolute value of $I(k)$ as:
% \[
%     |I(k)| \le \frac{2M_0 + (u_{\max} - u_{\min})(M_1 + M_0 L_2)}{k} =: \frac{C_1}{k}.
% \]
% This establishes the temporal decay of the continuous oscillatory component. Thus, \(I(k)=O(1/k)\) as \(k\to\infty\).

% We find that for any step $k$ satisfying:
% \[
%     k \ge C\max\left\{\frac{1}{1-\rho}, \frac{1-\rho}{\eta}\right\}
% \]
% The oscillatory term is strictly dominated by a small fraction of the static mass:
% \[
%     |I(k)| \le \frac{1}{4} \int_{1}^{J_{\max}} w(x) \, dx
% \]
% This guarantees that the oscillatory interference cannot cancel the static component, yielding a strictly positive lower bound for the sum:
% \begin{equation*}
%     \sum_{j \le J_{\max}} \lambda_j (\theta_j^*)^2 (1+A_j^2) \cos^2(k\omega_j - \psi_j) \ge c > 0.
% \end{equation*}
% Thus, the discrete sum is strictly bounded away from zero, yielding $\cS_{\mathrm{under}}(k) \ge c \rho^k$.

% Synthesizing the overdamped and underdamped contributions, we obtain the full deterministic bias $\cS_{\polyak}(k) = \cS_{\mathrm{over}}(k) + \cS_{\mathrm{under}}(k)$:
% \begin{itemize}
%     \item When $k \lesssim \frac{1-\rho}{\eta}$, combining the short-time overdamped term with the underdamped bound gives:
%     \[
%         \cS_{\polyak}(k) \eqsim \rho^k + \left(\frac{(1-\rho)^2}{\eta}\right)^s.
%     \]
%     \item When $k \gtrsim \frac{1-\rho}{\eta}$, combining the long-time overdamped polynomial tail with the underdamped bound gives:
%     \[
%         \cS_{\polyak}(k) \eqsim \rho^k + \left(\frac{1-\rho}{\eta k}\right)^s.
%     \]
% \end{itemize}
% This concludes the proof.
% \end{proof}

% \begin{theorem}[Two-sided scaling law, mixed regime]
% \label{thm:pol_two_sided_scaling}
% Under Assumptions~\ref{ass:diagonal-covariance} and~\ref{ass:power-law}, the stability condition $\eta\lesssim \min\{1,B(1-\rho)\}$, and for the mixed regime $\eta\gtrsim(1-\rho)^2$ with $k \gtrsim \frac{1-\rho}{\eta}$, the expected excess risk satisfies:
% \[
%     \EE[\cE(\btheta_k)]
%     \eqsim
%     \rho^k
%     +
%     \left(\frac{1-\rho}{\eta k}\right)^s
%     +
%     \frac{\sigma^2\eta}{B(1-\rho)}.
% \]
% \end{theorem}

% \begin{proof}
% By Corollary~\ref{cor:pol_noise_absorption}, the multiplicative gradient noise is strictly absorbed by the additive label noise floor, yielding the asymptotic equivalence:
% \begin{equation}
% \label{eq:pol_total_risk_mixed}
%     \EE[\cE(\btheta_k)] 
%     \eqsim 
%     \cS_{\polyak}(k) + \frac{\eta^2\sigma^2}{B}\sum_{\tau=0}^{k-1}\mathcal{K}_{\polyak}(\tau).
% \end{equation}
% By Lemma~\ref{lem:pol_label_noise}, the condition $k \gtrsim \frac{1-\rho}{\eta}$ ensures the label noise summation captures the saturated floor:
% \begin{equation}
% \label{eq:pol_label_noise_saturated}
%     \frac{\eta^2\sigma^2}{B}\sum_{\tau=0}^{k-1}\mathcal{K}_{\polyak}(\tau) 
%     \eqsim 
%     \frac{\sigma^2\eta}{B(1-\rho)}.
% \end{equation}

% To evaluate the deterministic bias $\cS_{\polyak}(k)$ over the entire domain $k \gtrsim \frac{1-\rho}{\eta}$, we analyze the bounds derived in Proposition~\ref{prop:pol_deterministic_bias} by splitting the domain into two consecutive sub-regimes at the momentum relaxation timescale $\frac{1}{1-\rho}$.

% For $\frac{1-\rho}{\eta} \lesssim k \lesssim \frac{1}{1-\rho}$, Proposition~\ref{prop:pol_deterministic_bias} gives:
% \[
%     \cS_{\polyak}(k) \eqsim \rho^k + \left(\frac{(1-\rho)^2}{\eta}\right)^s.
% \]
% Since $k(1-\rho) \lesssim 1$, the exponential transient satisfies:
% \[
%     \rho^k = (1-(1-\rho))^k \eqsim e^{-k(1-\rho)} \eqsim 1.
% \]
% In the mixed regime $\eta \gtrsim (1-\rho)^2$, the second term satisfies $\left(\frac{(1-\rho)^2}{\eta}\right)^s \lesssim 1$. Consequently, the deterministic bias is dominated by the non-decayed transient:
% \begin{equation}
% \label{eq:pol_bias_short}
%     \cS_{\polyak}(k) \eqsim 1.
% \end{equation}
% Concurrently, we evaluate the target analytical formula $\rho^k + \left(\frac{1-\rho}{\eta k}\right)^s$ in this same sub-regime. Utilizing the condition $k \gtrsim \frac{1-\rho}{\eta}$, we obtain:
% \[
%     \left(\frac{1-\rho}{\eta k}\right)^s \lesssim 1.
% \]
% Combining this with $\rho^k \eqsim 1$, the target formula satisfies:
% \begin{equation}
% \label{eq:pol_target_short}
%     \rho^k + \left(\frac{1-\rho}{\eta k}\right)^s \eqsim 1 + \mathcal{O}(1) \eqsim 1.
% \end{equation}
% Equating \eqref{eq:pol_bias_short} and \eqref{eq:pol_target_short} establishes the equivalence $\cS_{\polyak}(k) \eqsim \rho^k + \left(\frac{1-\rho}{\eta k}\right)^s$ for $\frac{1-\rho}{\eta} \lesssim k \lesssim \frac{1}{1-\rho}$, where the polynomial decay is rigorously absorbed by the $\mathcal{O}(1)$ exponential term.

% For $k \gtrsim \frac{1}{1-\rho}$, Proposition~\ref{prop:pol_deterministic_bias} explicitly gives:
% \[
%     \cS_{\polyak}(k) \eqsim \rho^k + \left(\frac{1-\rho}{\eta k}\right)^s.
% \]

% Synthesizing both sub-regimes, the following unified equivalence holds uniformly for all $k \gtrsim \frac{1-\rho}{\eta}$:
% \[
%     \cS_{\polyak}(k) \eqsim \rho^k + \left(\frac{1-\rho}{\eta k}\right)^s.
% \]
% Substituting this and \eqref{eq:pol_label_noise_saturated} into \eqref{eq:pol_total_risk_mixed} yields:
% \[
%     \EE[\cE(\btheta_k)] \eqsim \rho^k + \left(\frac{1-\rho}{\eta k}\right)^s + \frac{\sigma^2\eta}{B(1-\rho)}.
% \]
% This completes the proof.
% \end{proof}




% \subsection{Loss in the Fully Overdamped Regime} 

% \subsection{Loss in the Mixed Overdamped--Underdamped Regime}
% \newpage
\section{Nesterov scaling law}
\label{app:nesterov scaling law}

% In this section, we prove the Nesterov scaling law. Following the structure of Appendix~\ref{app:hbm scaling law}, we first decompose the error into deterministic signal and stochastic noise contributions, then classify the spectrum into overdamped and underdamped modes. We finally evaluate the loss in the fully overdamped regime and in the mixed overdamped--underdamped regime. 

% The main distinction from Polyak is that the squared modulus of an underdamped Nesterov root is $\rho(1-\eta\lambda_j)$ rather than $\rho$, which changes both the oscillatory transient and the label-noise floor.


As for Polyak momentum, we establish Theorem~\ref{thm:nesterov scaling law} using the Functional Scaling Law (FSL) framework of \citet{li2026functional}. Our starting point is the following expected excess risk:
\begin{align}
    \EE[\cE(\btheta_k)] 
    &\eqsim \underbrace{\cS_{\nesterov}(k)}_{\text{signal learning}} +\underbrace{\sum_{\tau=0}^{k-1}\underbrace{\mathcal{K}_{\nesterov}(\tau)}_{\text{forgetting kernel}} \left(\sigma^2 + \EE[\cE(\btheta_{k-\tau})]\right) \frac{\eta^2}{B}}_{\text{noise accumulation}} \label{eq:nes_fsl_intro_org}\\
    &\eqsim \cS_{\nesterov}(k) + \underbrace{\frac{\eta^2\sigma^2}{B}\sum_{\tau=0}^{k-1}\mathcal{K}_{\nesterov}(\tau)}_{\text{additive label noise}} + \underbrace{\frac{\eta^2}{B}\sum_{\tau=0}^{k-1}\mathcal{K}_{\nesterov}(\tau)\EE[\cE(\btheta_{k-\tau})]}_{\text{multiplicative noise}}. \label{eq:nes_fsl_intro}
\end{align}

Our analysis proceeds in three steps.

First, Section~\ref{app:nes_signal_noise_decomp} derives the \textbf{general discrete risk recursion} \eqref{eq:nes_fsl_intro} from an exact coordinate-wise impulse-response representation. Section~\ref{app:nes_decomp_modes} then analyzes the characteristic roots of the second-order dynamics and separates the spectrum into \textbf{overdamped and underdamped modes}.

Second, Sections~\ref{app:nes_preparation} and~\ref{app:nes_uniform_boundedness} deal with the \textbf{stochastic terms} in the recursion. Section~\ref{app:nes_preparation} computes the additive label-noise contribution, which yields a lower noise floor than Polyak; Section~\ref{app:nes_uniform_boundedness} shows that, under the risk stability condition, the multiplicative noise contribution can be absorbed into the label noise term up to constant factors.

Finally, Sections~\ref{app:nes_fully_overdamped_loss} and~\ref{app:nes_mixed_loss} characterize the \textbf{deterministic signal learning term} \(\cS_{\nesterov}(k)\) in the \textbf{fully overdamped and mixed-damping regimes}, and therefore derive the scaling law expression.


\subsection{Signal--noise decomposition and risk recursion}
\label{app:nes_signal_noise_decomp}

Recall the parameter error 
\[
\be_k := \btheta_k - \btheta^*.
\]
For Nesterov, the gradient is evaluated at the look-ahead point $\btheta_k^{\la} = \btheta_k + \rho(\btheta_k - \btheta_{k-1})$. 
Let 
\[
\be_k^{\la} := \btheta_k^{\la} - \btheta^*
\]
denote the look-ahead error. Substituting the gradient decomposition into the Nesterov update \eqref{eq:nag-update}, the error recursion in vector form is:
\[
    \be_{k+1} = (\bI - \eta \bH)\be_k^{\la} - \eta \bxi_k(\btheta_k^{\la}).
\]
In the diagonal eigenbasis of $\bH$, since $\be_k^{\la} = (1+\rho)\be_k - \rho \be_{k-1}$, the $j$-th coordinate of the Nesterov error recursion satisfies:
\[
    e_{k+1}(j) = (1+\rho)(1-\eta\lambda_j)e_k(j) -\rho(1-\eta\lambda_j)e_{k-1}(j) -\eta\xi_k^{\la}(j),
\]
where $\xi_k^{\la}(j)$ is the $j$-th coordinate of the stochastic gradient noise $\bxi_k(\btheta_k^{\la})$ evaluated at the look-ahead point.

Its homogeneous characteristic polynomial is
\[
    y^2-(1+\rho)(1-\eta\lambda_j)y+\rho(1-\eta\lambda_j)=0.
\]

Let \(y_{j,+}\) and \(y_{j,-}\) be the two roots of
\[
    y^2-(1+\rho)(1-\eta\lambda_j)y+\rho(1-\eta\lambda_j)=0.
\]
Define
\[
    G_k^{\nesterov}(j)
    :=
    \frac{y_{j,+}^{k+1}-y_{j,-}^{k+1}}{y_{j,+}-y_{j,-}},
    \qquad k\ge0,
\]
with the usual limiting definition at a repeated root.  Then
\begin{equation}
\label{eq:nes_impulse_response}
    e_k(j)
    =
    e_k^{\mathrm{det}}(j)
    -\eta\sum_{m=0}^{k-1}G_{k-1-m}^{\nesterov}(j)\xi_m^{\la}(j),
\end{equation}
where $e_k^{\mathrm{det}}(j)$ is the deterministic Nesterov error trajectory.

\begin{proposition}[Impulse-response representation]
\label{prop:nes_impulse_response}
The Green function is the unique solution of
\[
    G_{k+1}^{\nesterov}(j)
    =
    (1+\rho)(1-\eta\lambda_j)G_k^{\nesterov}(j)
    -\rho(1-\eta\lambda_j)G_{k-1}^{\nesterov}(j),
    \qquad
    G_0^{\nesterov}(j)=1,
    \qquad
    G_{-1}^{\nesterov}(j)=0.
\]
Moreover, if $e_0^{\mathrm{det}}(j)=-\theta_j^*$ and
$e_1^{\mathrm{det}}(j)=-(1-\eta\lambda_j)\theta_j^*$, then
\[
    e_k^{\mathrm{det}}(j)
    =
    -\frac{(1-\eta\lambda_j-y_{j,-})y_{j,+}^{k}
    -(1-\eta\lambda_j-y_{j,+})y_{j,-}^{k}}
    {y_{j,+}-y_{j,-}}\,\theta_j^* .
\]
\end{proposition}
\begin{proof}
Fix a coordinate $j$. The inhomogeneous scalar recursion is
\[
    e_{k+1}(j)
    =
    (1+\rho)(1-\eta\lambda_j)e_k(j)
    -\rho(1-\eta\lambda_j)e_{k-1}(j)
    -\eta \xi_k^{\la}(j).
\]
The impulse response for this second-order recursion is the sequence $G_k^{\nesterov}(j)$ satisfying
\[
    G_{k+1}^{\nesterov}(j)
    =
    (1+\rho)(1-\eta\lambda_j)G_k^{\nesterov}(j)
    -\rho(1-\eta\lambda_j)G_{k-1}^{\nesterov}(j),
    \qquad
    G_0^{\nesterov}(j)=1,
    \qquad
    G_{-1}^{\nesterov}(j)=0.
\]
Indeed, if an impulse $-\eta \xi_m^{\la}(j)$ is inserted at time $m$, then its contribution to
$e_{m+1}(j)$ is $-\eta\xi_m^{\la}(j)$, its contribution to $e_{m+2}(j)$ is
$-\eta(1+\rho)(1-\eta\lambda_j)\xi_m^{\la}(j)$, and subsequent contributions obey the same
homogeneous recursion. Hence the contribution at time $k$ is
$-\eta G_{k-1-m}^{\nesterov}(j)\xi_m^{\la}(j)$ for $0\le m\le k-1$. Summing over all previous
impulses gives
\[
    e_k(j)
    =
    e_k^{\mathrm{det}}(j)
    -\eta\sum_{m=0}^{k-1}G_{k-1-m}^{\nesterov}(j)\xi_m^{\la}(j),
\]
where $e_k^{\mathrm{det}}(j)$ solves the homogeneous recursion with the same initialization.

It remains to compute the closed form of $G_k^{\nesterov}(j)$. If $y_{j,+}\neq y_{j,-}$, then the general
solution of the homogeneous recurrence is a linear combination of $y_{j,+}^k$ and $y_{j,-}^k$. The initial
conditions $G_{-1}^{\nesterov}(j)=0$ and $G_0^{\nesterov}(j)=1$ give
\[
    G_k^{\nesterov}(j)
    =
    \frac{y_{j,+}^{k+1}-y_{j,-}^{k+1}}{y_{j,+}-y_{j,-}},
    \qquad k\ge0.
\]
If $y_{j,+}=y_{j,-}$, this formula is interpreted by continuity, giving
$G_k^{\nesterov}(j)=(k+1)y_{j,+}^k$. Thus the displayed formula is exactly the unique Green function of
the Nesterov coordinate recursion.

Now consider the deterministic trajectory. It satisfies
\[
    e_{k+1}^{\mathrm{det}}(j)
    =
    (1+\rho)(1-\eta\lambda_j)e_k^{\mathrm{det}}(j)
    -\rho(1-\eta\lambda_j)e_{k-1}^{\mathrm{det}}(j),
    \qquad
    e_0^{\mathrm{det}}(j)=-\theta_j^*,
    \qquad
    e_1^{\mathrm{det}}(j)=-(1-\eta\lambda_j)\theta_j^*.
\]
For $y_{j,+}\neq y_{j,-}$, write the deterministic solution as a linear combination of $y_{j,+}^k$ and
$y_{j,-}^k$. Solving the two-by-two linear system determined by the two displayed initial conditions gives
\[
    e_k^{\mathrm{det}}(j)
    =
    -\frac{(1-\eta\lambda_j-y_{j,-})y_{j,+}^{k}
    -(1-\eta\lambda_j-y_{j,+})y_{j,-}^{k}}
    {y_{j,+}-y_{j,-}}\,\theta_j^* .
\]
The repeated-root case follows by taking the same continuous limit. This proves the claim.
\end{proof}

To quantify the excess risk, we define the deterministic bias (the \textit{signal} component) and the convolution kernel as follows:
\[
    \cS_{\nesterov}(k)
    :=
    \frac12\sum_{j\ge1}\lambda_j\bigl(e_k^{\mathrm{det}}(j)\bigr)^2,
    \qquad
    \mathcal K_{\nesterov}(k)
    :=
    \sum_{j\ge1}\lambda_j^2\bigl(G_k^{\nesterov}(j)\bigr)^2.
\]
Because the stochastic gradient is evaluated at the look-ahead point, define
\[
    \be_k^{\la}:=\be_k+\rho(\be_k-\be_{k-1}),
    \qquad
    \mathcal L_k:=\frac12\|\be_k^{\la}\|_{\bH}^2.
\]

\begin{lemma}[Nesterov look-ahead risk control]
\label{lem:nes_lookahead_loss_control}
Assume the stable small-step condition \(\eta\lambda_1\le c_{\rm st}\), where \(c_{\rm st}>0\) is sufficiently small.  Then the look-ahead risk admits the coordinate expansion
\[
    \mathcal L_m
    =
    \frac12\sum_{j\ge1}\lambda_j
    \left((1+\rho)e_m(j)-\rho e_{m-1}(j)\right)^2.
\]
Moreover, the deterministic look-ahead risk is comparable to the next deterministic Nesterov risk:
\[
    \frac12\sum_{j\ge1}\lambda_j
    \left((1+\rho)e_m^{\mathrm{det}}(j)-\rho e_{m-1}^{\mathrm{det}}(j)\right)^2
    \eqsim
    \cS_{\nesterov}(m+1).
\]
For the stochastic iterates, the corresponding expectation-level control is
\[
    c\,\mathbb{E}\mathcal E_{m+1}-C\frac{\eta^2\sigma^2}{B}
    \le
    \mathbb{E}\mathcal L_m
    \le
    C\,\mathbb{E}\mathcal E_{m+1},
\]
where the constants are independent of \(m,\eta,\rho,B\).
\end{lemma}

\begin{proof}

The coordinate expansion follows directly from \(\be_m^{\la}=(1+\rho)\be_m-\rho\be_{m-1}\) and the definition of the \(\operatorname{\bm H}\)-norm.

For the deterministic comparison, the deterministic Nesterov update gives, for each coordinate \(j\),
\[
    e_{m+1}^{\mathrm{det}}(j)
    =
    (1-\eta\lambda_j)
    \left((1+\rho)e_m^{\mathrm{det}}(j)-\rho e_{m-1}^{\mathrm{det}}(j)\right).
\]
Since \(0<1-\eta\lambda_j\le1\) and \(1-\eta\lambda_j\ge1-c_{\rm st}\), we have
\[
    \left((1+\rho)e_m^{\mathrm{det}}(j)-\rho e_{m-1}^{\mathrm{det}}(j)\right)^2
    \eqsim
    \bigl(e_{m+1}^{\mathrm{det}}(j)\bigr)^2.
\]
Multiplying by \(\lambda_j\) and summing over \(j\ge1\) yields the deterministic claim.

For the stochastic comparison, the actual Nesterov update in coordinate \(j\) is
\[
    e_{m+1}(j)
    =
    (1-\eta\lambda_j)
    \left((1+\rho)e_m(j)-\rho e_{m-1}(j)\right)
    -\eta\xi_m^{\la}(j).
\]
Conditioning on the past, the noise is mean zero at the look-ahead point.  Therefore the mixed term vanishes and
\[
\begin{aligned}
    \mathbb{E}\mathcal E_{m+1}
    \eqsim{}&
    \sum_{j\ge1}\lambda_j(1-\eta\lambda_j)^2
    \mathbb{E}\left((1+\rho)e_m(j)-\rho e_{m-1}(j)\right)^2 \\
    &+
    \eta^2\sum_{j\ge1}\lambda_j\mathbb{E}\bigl(\xi_m^{\la}(j)\bigr)^2 .
\end{aligned}
\]
The first term is comparable to \(\mathbb{E}\mathcal L_m\), because \(1-\eta\lambda_j\) is bounded above and below by positive constants.  For the second term, Proposition~\ref{prop:noise-geometry} gives
\[
    \sum_{j\ge1}\lambda_j\mathbb{E}\bigl(\xi_m^{\la}(j)\bigr)^2
    \lesssim
    \frac{1}{B}\bigl(\sigma^2+\mathbb{E}\mathcal L_m\bigr)
    \sum_{j\ge1}\lambda_j^2.
\]




% The coordinate expansion follows directly from \(\tilde\be_m=(1+\rho)\be_m-\rho\be_{m-1}\) and the definition of the \(\operatorname{\bm H}\)-norm.

% For the deterministic comparison, the deterministic Nesterov update gives, for each coordinate \(j\),
% \[
%     e_{m+1}^{\mathrm{det}}(j)
%     =
%     (1-\eta\lambda_j)
%     \left((1+\rho)e_m^{\mathrm{det}}(j)-\rho e_{m-1}^{\mathrm{det}}(j)\right).
% \]
% Since \(0<1-\eta\lambda_j\le1\) and \(1-\eta\lambda_j\ge1-c_{\rm st}\), we have
% \[
%     \left((1+\rho)e_m^{\mathrm{det}}(j)-\rho e_{m-1}^{\mathrm{det}}(j)\right)^2
%     \eqsim
%     \bigl(e_{m+1}^{\mathrm{det}}(j)\bigr)^2.
% \]
% Multiplying by \(\lambda_j\) and summing over \(j\ge1\) yields the deterministic claim.

% For the stochastic comparison, the actual Nesterov update in coordinate \(j\) is
% \[
%     e_{m+1}(j)
%     =
%     (1-\eta\lambda_j)
%     \left((1+\rho)e_m(j)-\rho e_{m-1}(j)\right)
%     -\eta\xi_m^{\nesterov}(j).
% \]
% Conditioning on the past, the noise is mean zero at the look-ahead point.  Therefore the mixed term vanishes and
% \[
% \begin{aligned}
%     \mathbb{E}\mathcal E_{m+1}
%     \eqsim{}&
%     \sum_{j\ge1}\lambda_j(1-\eta\lambda_j)^2
%     \mathbb{E}\left((1+\rho)e_m(j)-\rho e_{m-1}(j)\right)^2 \\
%     &+
%     \eta^2\sum_{j\ge1}\lambda_j\mathbb{E}\bigl(\xi_m^{\nesterov}(j)\bigr)^2 .
% \end{aligned}
% \]
% The first term is comparable to \(\mathbb{E}\mathcal L_m\), because \(1-\eta\lambda_j\) is bounded above and below by positive constants.  For the second term, Proposition~\ref{prop:noise-geometry} gives
% \[
%     \sum_{j\ge1}\lambda_j\mathbb{E}\bigl(\xi_m^{\nesterov}(j)\bigr)^2
%     \lesssim
%     \frac{1}{B}\bigl(\sigma^2+\mathbb{E}\mathcal L_m\bigr)
%     \sum_{j\ge1}\lambda_j^2.
% \]
Since \(\sum_{j\ge1}\lambda_j^2<\infty\) and \(c_{\rm st}\) is chosen small, the term proportional to \(\eta^2\mathbb{E}\mathcal L_m/B\) is absorbed into the comparable first term.  Thus
\[
    \mathbb{E}\mathcal E_{m+1}
    \eqsim
    \mathbb{E}\mathcal L_m
    +O\left(\frac{\eta^2\sigma^2}{B}\right),
\]
which gives
\[
    c\,\mathbb{E}\mathcal E_{m+1}-C\frac{\eta^2\sigma^2}{B}
    \le
    \mathbb{E}\mathcal L_m
    \le
    C\,\mathbb{E}\mathcal E_{m+1}.
\]
This proves the lemma.
\end{proof}

\begin{proposition}[Nesterov risk recursion]
\label{prop:nes_risk_recursion}
Under the covariance relation in Proposition~\ref{prop:noise-geometry}, the expected excess risk satisfies the two-sided
recursion
\begin{equation}
\label{eq:nes_risk_recursion_form}
    \mathbb{E}\mathcal E_k
    \eqsim
    \cS_{\nesterov}(k)
    +
    \frac{\eta^2\sigma^2}{B}\sum_{\tau=0}^{k-1}\mathcal K_{\nesterov}(\tau)
    +
    \frac{\eta^2}{B}\sum_{\tau=0}^{k-1}\mathcal K_{\nesterov}(\tau)
    \mathbb{E}\mathcal L_{k-1-\tau}.
\end{equation}
Moreover, Lemma~\ref{lem:nes_lookahead_loss_control} gives the required control of the look-ahead risk by the
next-step risk.  In particular, under the stable small-step condition,
\[
    c\,\mathbb{E}\mathcal E_{m+1}-C\frac{\eta^2\sigma^2}{B}
    \le
    \mathbb{E}\mathcal L_m
    \le
    C\,\mathbb{E}\mathcal E_{m+1}.
\]
\end{proposition}

\begin{proof}
Substitute the Nesterov impulse-response representation~\eqref{eq:nes_impulse_response} into the coordinate-wise excess risk.  For every coordinate
\(j\),
\[
e_k(j)
    =
    e_k^{\mathrm{det}}(j)
    -\eta\sum_{m=0}^{k-1}G_{k-1-m}^{\nesterov}(j)\xi_m^{\la}(j).
\]
The noise \(\xi_m^{\la}(j)\) is conditionally mean zero given the past and the look-ahead point
\(\be_m^{\la}\), and the fresh mini-batches are independent across time.  Therefore all mixed
deterministic--stochastic and cross-time stochastic terms vanish after taking expectation.  Hence
\[
    \mathbb{E}[e_k(j)^2]
    =
    \bigl(e_k^{\mathrm{det}}(j)\bigr)^2
    +
    \eta^2\sum_{m=0}^{k-1}
    \bigl(G_{k-1-m}^{\nesterov}(j)\bigr)^2
    \mathbb{E}\bigl(\xi_m^{\la}(j)\bigr)^2,
\]
up to the same universal constants in the covariance comparison of Proposition~\ref{prop:noise-geometry}.  Multiplying by
\(\lambda_j\) and summing over \(j\ge1\) gives
\[
\begin{aligned}
    \mathbb{E}\mathcal E_k
    \eqsim{}&
    \sum_{j\ge1}\lambda_j\bigl(e_k^{\mathrm{det}}(j)\bigr)^2
    +\eta^2\sum_{m=0}^{k-1}\sum_{j\ge1}
    \lambda_j
    \bigl(G_{k-1-m}^{\nesterov}(j)\bigr)^2
    \mathbb{E}\bigl(\xi_m^{\la}(j)\bigr)^2 .
\end{aligned}
\]
Since Nesterov evaluates the stochastic gradient at the look-ahead point \(\be_m^{\la}\),
Proposition~\ref{prop:noise-geometry} gives the coordinate-wise variance relation
\[
    \mathbb{E}\bigl(\xi_m^{\la}(j)\bigr)^2
    \eqsim
    \frac{1}{B}\lambda_j\bigl(\sigma^2+\mathbb{E}\mathcal L_m\bigr).
\]
Substituting this into the previous display yields
\[
\begin{aligned}
    \mathbb{E}\mathcal E_k
    \eqsim{}&
    \sum_{j\ge1}\lambda_j\bigl(e_k^{\mathrm{det}}(j)\bigr)^2
    +\frac{\eta^2}{B}\sum_{m=0}^{k-1}
    \bigl(\sigma^2+\mathbb{E}\mathcal L_m\bigr)
    \sum_{j\ge1}\lambda_j^2
    \bigl(G_{k-1-m}^{\nesterov}(j)\bigr)^2 .
\end{aligned}
\]
Using the definitions of \(\cS_{\nesterov}\) and \(\mathcal K_{\nesterov}\), and then changing
variables \(\tau=k-1-m\), we obtain
\[
    \mathbb{E}\mathcal E_k
    \eqsim
    \cS_{\nesterov}(k)
    +
    \frac{\eta^2\sigma^2}{B}\sum_{\tau=0}^{k-1}\mathcal K_{\nesterov}(\tau)
    +
    \frac{\eta^2}{B}\sum_{\tau=0}^{k-1}\mathcal K_{\nesterov}(\tau)
    \mathbb{E}\mathcal L_{k-1-\tau}.
\]
This proves the claimed two-sided recursion.

It remains to control the look-ahead risk.  This is exactly Lemma~\ref{lem:nes_lookahead_loss_control},
which expands \(\mathcal L_m\) in coordinates and then uses the Nesterov update from step \(m\) to step
\(m+1\) to compare the look-ahead risk with the next-step risk.
\end{proof}


\subsection{Overdamped and underdamped modes decomposition}
\label{app:nes_decomp_modes}

Since both the deterministic error signal $e_k^{\mathrm{det}}(j)$ and the discrete Green function $G_k^{\nesterov}(j)$ obey the same homogeneous difference equation, their behavior is dictated by the same characteristic roots. Depending on the sign of the discriminant, these roots transition from real roots to a complex-conjugate pair. The following proposition establishes the exact Nesterov damping threshold, defines the crossover index $J_{\max}$, and partitions the spectrum into the two global regimes.

\begin{proposition}[Exact spectral decomposition and crossover boundary]
\label{prop:nes_mode_decomposition}
For each coordinate $j$, the characteristic roots of the Nesterov recursion are
\[
    y_{j,\pm}
    =
    \frac{(1+\rho)(1-\eta\lambda_j)\pm\sqrt{\Delta_j^{\nesterov}}}{2},
\]
where
\[
    \Delta_j^{\nesterov}
    :=
    (1+\rho)^2(1-\eta\lambda_j)^2
    -4\rho(1-\eta\lambda_j).
\]
In the stable small-step range $0<\eta\lambda_j<1$, the critical transition $\Delta_j^{\nesterov}=0$ corresponds exactly to
\[
    \eta\lambda_j
    =
    \frac{(1-\rho)^2}{(1+\rho)^2},
\]
which is asymptotically equivalent to $\eta\lambda_j\eqsim(1-\rho)^2$. Thus the individual eigen-directions are classified as follows:
\begin{enumerate}
    \item \textbf{Overdamped Modes} $\bigl(\eta\lambda_j<(1-\rho)^2/(1+\rho)^2\bigr)$: the roots are real.
    \item \textbf{Underdamped Modes} $\bigl(\eta\lambda_j>(1-\rho)^2/(1+\rho)^2\bigr)$: the roots form a complex-conjugate pair.
\end{enumerate}
In the mixed-damping regime, the boundary index $J_{\max}$ separating the underdamped modes from the overdamped modes is defined by the crossover relation $\eta\lambda_{J_{\max}}\eqsim(1-\rho)^2$:
\begin{equation}
\label{eq:J_max_nes}
    J_{\max}
    \eqsim
    \left(\frac{\eta}{(1-\rho)^2}\right)^{1/\beta}.
\end{equation}
Consequently, based on the maximum eigenvalue $\lambda_{\max}$, Nesterov operates in one of two global regimes:
\begin{itemize}
    \item \textbf{Fully Overdamped Regime} $\bigl(\eta\lambda_{\max}\lesssim(1-\rho)^2\bigr)$: all spectral modes are overdamped.
    \item \textbf{Mixed-Damping Regime} $\bigl(\eta\lambda_{\max}\gtrsim(1-\rho)^2\bigr)$: the high-curvature modes $j\le J_{\max}$ are underdamped, whereas the flat modes $j>J_{\max}$ are overdamped.
\end{itemize}
\end{proposition}

\begin{proof}
In the stable small-step range, $1-\eta\lambda_j\in(0,1)$. The condition $\Delta_j^{\nesterov}<0$ is
\[
    (1+\rho)^2(1-\eta\lambda_j)^2
    -4\rho(1-\eta\lambda_j)<0.
\]
Since $1-\eta\lambda_j>0$, this is equivalent to
$1-\eta\lambda_j<\frac{4\rho}{(1+\rho)^2},$
and hence
$\eta\lambda_j>1-\frac{4\rho}{(1+\rho)^2}=\frac{(1-\rho)^2}{(1+\rho)^2}.$
Therefore the equality case gives the exact damping boundary. Since $(1+\rho)^2\eqsim1$ for $0\le\rho<1$, this boundary has the scaling form $\eta\lambda_j\eqsim(1-\rho)^2$.

To locate the crossover index in the mixed-damping regime, use $\lambda_j\eqsim j^{-\beta}$ and impose
$\eta\lambda_{J_{\max}}\eqsim(1-\rho)^2.$
Then
$\eta J_{\max}^{-\beta}\eqsim(1-\rho)^2,$
which gives
\[
    J_{\max}
    \eqsim
    \left(\frac{\eta}{(1-\rho)^2}\right)^{1/\beta}.
\]
Because the eigenvalues are non-increasing, the modes $j\le J_{\max}$ lie above the damping threshold and are underdamped, while the modes $j>J_{\max}$ lie below the threshold and are overdamped. Examining whether this crossover occurs below the leading eigenvalue yields the two global regimes stated above.
\end{proof}


\subsection{Label noise calculation}
\label{app:nes_preparation}

The convolution estimates in Appendix~\ref{app:pol_preparation} will be used throughout this section. It remains to identify the Nesterov label-noise floor and to replace the look-ahead risk in the multiplicative-noise recursion.

\begin{lemma}[Nesterov label-noise contribution]
\label{lem:nes_label_noise}
Assume the stable small-step condition \(\eta\lambda_1\le c_{\rm st}\), where \(c_{\rm st}>0\) is a
sufficiently small universal constant.  There exist constants \(C,c_-,c_+>0\) such that if
\[
    k\ge
    C\max\left\{\frac1{1-\rho},\frac{1-\rho}{\eta}\right\},
\]
then
\[
    c_-\frac{\sigma^2}{B}
    \min\left\{
        \left(\frac{\eta}{1-\rho}\right)^{\frac1\beta},
        \frac{\eta}{1-\rho}
    \right\}
    \le
    \frac{\eta^2\sigma^2}{B}\sum_{\tau=0}^{k-1}\mathcal K_{\nesterov}(\tau)
    \le
    c_+\frac{\sigma^2}{B}
    \min\left\{
        \left(\frac{\eta}{1-\rho}\right)^{\frac1\beta},
        \frac{\eta}{1-\rho}
    \right\}.
\]
Hence the Nesterov label-noise floor is
\[
    \frac{\eta^2\sigma^2}{B}\sum_{\tau=0}^{k-1}\mathcal K_{\nesterov}(\tau)
    \eqsim
    \frac{\sigma^2}{B}
    \min\left\{
        \left(\frac{\eta}{1-\rho}\right)^{\frac1\beta},
        \frac{\eta}{1-\rho}
    \right\}.
\]
\end{lemma}

\begin{proof}
We first compute the saturated Green energy in one eigendirection.  The Nesterov Green function satisfies
\[
    G_{k+1}^{\nesterov}(j)
    =
    (1+\rho)(1-\eta\lambda_j)G_k^{\nesterov}(j)
    -\rho(1-\eta\lambda_j)G_{k-1}^{\nesterov}(j),
    \qquad
    G_0^{\nesterov}(j)=1,
    \qquad
    G_{-1}^{\nesterov}(j)=0.
\]
Multiplying this recursion by \(G_k^{\nesterov}(j)\), summing over \(k\ge0\), and then squaring the
recursion and summing over \(k\ge0\) gives the standard second-order energy identity
\[
    \sum_{k=0}^{\infty}\bigl(G_k^{\nesterov}(j)\bigr)^2
    =
    \frac{1+\rho(1-\eta\lambda_j)}
    {\bigl(1-\rho(1-\eta\lambda_j)\bigr)
    \left(\bigl(1+\rho(1-\eta\lambda_j)\bigr)^2
    -(1+\rho)^2(1-\eta\lambda_j)^2\right)}.
\]
The second factor in the denominator simplifies as
\[
    \bigl(1+\rho(1-\eta\lambda_j)\bigr)^2
    -(1+\rho)^2(1-\eta\lambda_j)^2
    =
    \eta\lambda_j\left(1+(1+2\rho)(1-\eta\lambda_j)\right).
\]
Under \(\eta\lambda_1\le c_{\rm st}\), the remaining factors are bounded above and below by constants
except for \(1-\rho+\eta\lambda_j\) and \(\eta\lambda_j\).  Therefore
\[
    \sum_{k=0}^{\infty}\bigl(G_k^{\nesterov}(j)\bigr)^2
    \eqsim
    \frac{1}{\eta\lambda_j(1-\rho+\eta\lambda_j)}.
\]
Consequently the saturated label-noise contribution is
\[
\begin{aligned}
    \frac{\eta^2\sigma^2}{B}
    \sum_{j\ge1}\lambda_j^2
    \sum_{k=0}^{\infty}\bigl(G_k^{\nesterov}(j)\bigr)^2
    &\eqsim
    \frac{\eta\sigma^2}{B}
    \sum_{j\ge1}\frac{\lambda_j}{1-\rho+\eta\lambda_j}  \\
    &=
    \frac{\sigma^2}{B}
    \sum_{j\ge1}
    \frac{\lambda_j}{\lambda_j+\frac{1-\rho}{\eta}} .
\end{aligned}
\]
Here Nesterov differs from Polyak: the effective
root product is \(\rho(1-\eta\lambda_j)\), so the saturated spectral weight contains
\(\lambda_j/(\lambda_j+(1-\rho)/\eta)\), rather than simply \(\lambda_j\).

We next estimate the spectral sum.  If \(\eta/(1-\rho)\le1\), then \((1-\rho)/\eta\gtrsim1\), and hence
\[
    \sum_{j\ge1}
    \frac{\lambda_j}{\lambda_j+\frac{1-\rho}{\eta}}
    \eqsim
    \frac{\eta}{1-\rho}\sum_{j\ge1}\lambda_j
    \eqsim
    \frac{\eta}{1-\rho}.
\]
If \(\eta/(1-\rho)\ge1\), then the spectral cutoff is
$\lambda_j\eqsim\frac{1-\rho}{\eta}.$
Since \(\lambda_j\eqsim j^{-\beta}\), the number of indices above this cutoff is
$\left(\frac{\eta}{1-\rho}\right)^{\frac1\beta}.$
The modes above the cutoff each contribute a constant amount to the summand, while the tail below the
cutoff contributes the same order:
\[
    \sum_{\lambda_j<\frac{1-\rho}{\eta}}
    \frac{\lambda_j}{\frac{1-\rho}{\eta}}
    \eqsim
    \left(\frac{\eta}{1-\rho}\right)^{\frac1\beta}.
\]
Thus, in all cases,
\[
    \sum_{j\ge1}
    \frac{\lambda_j}{\lambda_j+\frac{1-\rho}{\eta}}
    \eqsim
    \min\left\{
        \left(\frac{\eta}{1-\rho}\right)^{\frac1\beta},
        \frac{\eta}{1-\rho}
    \right\}.
\]

It remains to pass from the saturated energy to the finite-time sum.  The root formula for
\(G_k^{\nesterov}(j)\) gives a uniform tail estimate: after increasing \(C\) if necessary, the time
condition
\[
    k\ge
    C\max\left\{\frac1{1-\rho},\frac{1-\rho}{\eta}\right\}
\]
ensures that the finite partial sum captures a fixed positive fraction of the saturated Green energy on the
spectral region responsible for the lower bounds above.  When \(\eta/(1-\rho)\le1\), this region may be
taken to be a fixed finite block of leading eigendirections, whose relaxation time is of order
\((1-\rho)/\eta\).  When \(\eta/(1-\rho)\ge1\), it may be taken to be the block
\(\lambda_j\gtrsim(1-\rho)/\eta\), whose relaxation time is at most of order \((1-\rho)^{-1}\).
Therefore
\[
    \sum_{\tau=0}^{k-1}\mathcal K_{\nesterov}(\tau)
    \eqsim
    \sum_{j\ge1}\lambda_j^2
    \sum_{\tau=0}^{\infty}\bigl(G_\tau^{\nesterov}(j)\bigr)^2
\]
up to constants for the purpose of the two-sided label-noise estimate.  Combining this finite-time
comparison with the saturated computation above proves the lemma.
\end{proof}

% \begin{lemma}[Nesterov look-ahead replacement after label-noise absorption]
% \label{lem:nes_replace_lookahead_by_next_loss}
% Assume the stable small-step condition of Lemma~\ref{lem:nes_lookahead_loss_control}.  After the label-noise
% term in Lemma~\ref{lem:nes_label_noise} is included, the Nesterov risk recursion may be written as
% \[
%     \mathbb{E}\mathcal E_k
%     \eqsim
%     \cS_{\nesterov}(k)
%     +
%     \frac{\eta^2\sigma^2}{B}\sum_{\tau=0}^{k-1}\mathcal K_{\nesterov}(\tau)
%     +
%     \frac{\eta^2}{B}\sum_{\tau=0}^{k-1}\mathcal K_{\nesterov}(\tau)
%     \mathbb{E}\mathcal E_{k-\tau}.
% \]
% Equivalently, in all subsequent Nesterov self-noise estimates one may replace
% \(\mathbb{E}\mathcal L_{k-1-\tau}\) by \(\mathbb{E}\mathcal E_{k-\tau}\).
% \end{lemma}

% \begin{proof}
% By Lemma~\ref{lem:nes_lookahead_loss_control}, for \(m=k-1-\tau\),
% \[
%     c\,\mathbb{E}\mathcal E_{k-\tau}-C\frac{\eta^2\sigma^2}{B}
%     \le
%     \mathbb{E}\mathcal L_{k-1-\tau}
%     \le
%     C\,\mathbb{E}\mathcal E_{k-\tau}.
% \]
% The upper bound immediately gives
% \[
%     \frac{\eta^2}{B}
%     \sum_{\tau=0}^{k-1}\mathcal K_{\nesterov}(\tau)
%     \mathbb{E}\mathcal L_{k-1-\tau}
%     \lesssim
%     \frac{\eta^2}{B}
%     \sum_{\tau=0}^{k-1}\mathcal K_{\nesterov}(\tau)
%     \mathbb{E}\mathcal E_{k-\tau}.
% \]
% For the lower bound, the same comparison gives
% \[
% \begin{aligned}
%     \frac{\eta^2}{B}
%     \sum_{\tau=0}^{k-1}\mathcal K_{\nesterov}(\tau)
%     \mathbb{E}\mathcal L_{k-1-\tau}
%     \ge{}&
%     c\frac{\eta^2}{B}
%     \sum_{\tau=0}^{k-1}\mathcal K_{\nesterov}(\tau)
%     \mathbb{E}\mathcal E_{k-\tau}  \\
%     &-
%     C\frac{\eta^2}{B}
%     \frac{\eta^2\sigma^2}{B}
%     \sum_{\tau=0}^{k-1}\mathcal K_{\nesterov}(\tau).
% \end{aligned}
% \]
% The final term is a factor \(\eta^2/B\) times the label-noise term.  Under the stable small-step condition
% \(\eta\lambda_1\le c_{\rm st}\) and \(B\ge1\), this factor is bounded by a constant, and by choosing
% \(c_{\rm st}\) sufficiently small it is absorbed into the label-noise contribution already present in the
% risk recursion.  Hence the recursion with \(\mathbb{E}\mathcal L_{k-1-\tau}\) is equivalent, up to constants, to
% the recursion with \(\mathbb{E}\mathcal E_{k-\tau}\).
% \end{proof}







% \subsection{Loss in the fully overdamped regime}
% \label{app:nes_fully_overdamped_loss}

% In the fully overdamped regime, every characteristic root is real and the deterministic bias has no oscillatory component.


% \begin{proposition}[Deterministic Bias, Fully Overdamped Regime]
% \label{prop:nes_deterministic_bias0}
% Under Assumptions~\ref{ass:diagonal-covariance} and~\ref{ass:power-law}, in the fully overdamped regime $\eta\lambda_1\le c_0(1-\rho)^2$ with a sufficiently small universal constant $c_0>0$, we have:
% \begin{equation*}
%     \cS_{\nesterov}(k)
%     \eqsim
%     \begin{cases}
%     1,&\qquad k\lesssim \frac{1-\rho}{\eta},\\
%     \left(\frac{1-\rho}{\eta k}\right)^s, &\qquad k\gtrsim\frac{1-\rho}{\eta}.\\
%     \end{cases}
% \end{equation*}
% \end{proposition}

% \begin{proof}
% For each coordinate, write the deterministic multiplier as
% \[
%     e_k^{\mathrm{det}}(j)=-h_k^{\nesterov}(j)\theta_j^*.
% \]
% The deterministic Nesterov coordinate recursion gives
% \[
%     h_{k+1}^{\nesterov}(j)
%     =
%     (1+\rho)(1-\eta\lambda_j)h_k^{\nesterov}(j)
%     -\rho(1-\eta\lambda_j)h_{k-1}^{\nesterov}(j),
%     \qquad
%     h_0^{\nesterov}(j)=1,
%     \qquad
%     h_1^{\nesterov}(j)=1-\eta\lambda_j.
% \]
% Its characteristic roots are
% \[
%     y_{j,+},y_{j,-}
%     =
%     \frac{(1+\rho)(1-\eta\lambda_j)
%     \pm
%     \sqrt{(1+\rho)^2(1-\eta\lambda_j)^2-4\rho(1-\eta\lambda_j)}}{2}.
% \]
% In the fully overdamped regime, \(\eta\lambda_j\le c_0(1-\rho)^2\) for every \(j\).  To locate the two roots, evaluate the characteristic polynomial at \(1-u\).  A direct expansion gives
% \[
% \begin{aligned}
%     &(1-u)^2
%     -(1+\rho)(1-\eta\lambda_j)(1-u)
%     +\rho(1-\eta\lambda_j) \\
%     &\qquad=
%     \eta\lambda_j
%     -\bigl(1-\rho+(1+\rho)\eta\lambda_j\bigr)u
%     +u^2.
% \end{aligned}
% \]
% Taking \(c_0>0\) sufficiently small, the two evaluations
% \[
%     \eta\lambda_j
%     -\bigl(1-\rho+(1+\rho)\eta\lambda_j\bigr)
%     \frac{\eta\lambda_j}{2(1-\rho)}
%     +
%     \left(\frac{\eta\lambda_j}{2(1-\rho)}\right)^2
%     >0
% \]
% and
% \[
%     \eta\lambda_j
%     -\bigl(1-\rho+(1+\rho)\eta\lambda_j\bigr)
%     \frac{2\eta\lambda_j}{1-\rho}
%     +
%     \left(\frac{2\eta\lambda_j}{1-\rho}\right)^2
%     <0
% \]
% hold uniformly for all \(0<\eta\lambda_j\le c_0(1-\rho)^2\).  Hence the slow root satisfies
% \[
%     1-\frac{2\eta\lambda_j}{1-\rho}
%     \le
%     y_{j,+}
%     \le
%     1-\frac{\eta\lambda_j}{2(1-\rho)}.
% \]
% Moreover,
% \[
%     y_{j,+}y_{j,-}=\rho(1-\eta\lambda_j).
% \]
% Using the lower bound on \(y_{j,+}\), and decreasing \(c_0\) if necessary, the fast root obeys
% \[
%     0\le y_{j,-}
%     \le
%     1-\frac{1-\rho}{2}.
% \]

% The solution of the two-point initial value problem is
% \[
%     h_k^{\nesterov}(j)
%     =
%     \frac{1-\eta\lambda_j-y_{j,-}}{y_{j,+}-y_{j,-}}y_{j,+}^{k}
%     +
%     \frac{y_{j,+}-(1-\eta\lambda_j)}{y_{j,+}-y_{j,-}}y_{j,-}^{k}.
% \]
% The root bounds above imply
% \[
%     \frac12
%     \le
%     \frac{1-\eta\lambda_j-y_{j,-}}{y_{j,+}-y_{j,-}}
%     \le
%     2,
%     \qquad
%     \left|
%     \frac{y_{j,+}-(1-\eta\lambda_j)}{y_{j,+}-y_{j,-}}
%     \right|
%     \le
%     Cc_0.
% \]
% Indeed, \(y_{j,+}-y_{j,-}\eqsim1-\rho\), while \(1-\eta\lambda_j-y_{j,-}\eqsim1-\rho\), and
% \[
%     \left|y_{j,+}-(1-\eta\lambda_j)\right|
%     \lesssim
%     \frac{\eta\lambda_j}{1-\rho}
%     \le
%     c_0(1-\rho).
% \]
% Thus the slow component has a uniformly positive coefficient and the fast component has a coefficient of order \(c_0\).

% To evaluate the deterministic bias $\cS_{\nesterov}(k) = \frac12\sum_{j\ge1}\lambda_j(\theta_j^*)^2 \left(h_k^{\nesterov}(j)\right)^2$, we divide the time $k$ into three successive stages.

% \textbf{Stage 1: $k \lesssim \frac{1}{1-\rho}$.} \\
% In this extremely short-time regime, for the leading modes ($j=\mathcal{O}(1)$), since $\eta\lambda_j \le c_0(1-\rho)^2$, the slow root satisfies $y_{j,+} \ge 1 - 2c_0(1-\rho)$. For $k \lesssim \frac{1}{1-\rho}$, we have
% \[
%     y_{j,+}^k \ge \left(1 - 2c_0(1-\rho)\right)^{\frac{C}{1-\rho}} \eqsim 1.
% \]
% Since the dynamics are unconditionally stable ($h_k^{\nesterov}(j) \le 1$), it follows that $h_k^{\nesterov}(j) \eqsim 1$ for the leading modes. Consequently, the bias is dominated by these non-decayed modes:
% \[
%     \cS_{\nesterov}(k) \eqsim \sum_{j\ge1}\lambda_j(\theta_j^*)^2 \eqsim 1.
% \]

% \textbf{Stage 2: $\frac{1}{1-\rho} \lesssim k \lesssim \frac{1-\rho}{\eta}$.} \\
% For $k \gtrsim \frac{1}{1-\rho}$, the fast root decays exponentially since $y_{j,-}^k \le \left(1-\frac{1-\rho}{2}\right)^k \ll 1$. The multiplier is  dominated by the slow root:
% \[
%     h_k^{\nesterov}(j) \simeq \frac{1-\eta\lambda_j - y_{j,-}}{y_{j,+} - y_{j,-}} y_{j,+}^k \simeq y_{j,+}^k.
% \]
% For the leading modes with $\lambda_j = \Theta(1)$, within the time window $k \lesssim \frac{1-\rho}{\eta}$, the exponent is bounded by:
% \[
%     y_{j,+}^k \eqsim \left(1 - c\frac{\eta\lambda_j}{1-\rho}\right)^k \eqsim \exp\left(-c\frac{\eta\lambda_j}{1-\rho}k\right) \eqsim 1,
% \]
% where we used the condition $\frac{\eta\lambda_j}{1-\rho}k \lesssim \lambda_j \le \lambda_1 = \mathcal{O}(1)$. Thus, the leading modes have not yet experienced significant exponential decay, preserving $h_k^{\nesterov}(j) \eqsim 1$. As a result, the deterministic bias remains unchanged:
% \[
%     \cS_{\nesterov}(k) \eqsim 1.
% \]

% \textbf{Stage 3: $k \gtrsim \frac{1-\rho}{\eta}$.} \\
% In this long-time regime, the exponential decay of the slow root becomes significant across the spectrum. For the slow root, the preceding bounds give
% \[
%     \exp\left(-C\frac{\eta\lambda_j}{1-\rho}k\right)
%     \le
%     y_{j,+}^{k}
%     \le
%     \exp\left(-c\frac{\eta\lambda_j}{1-\rho}k\right).
% \]
% For the fast root, since \(\eta\lambda_j/(1-\rho)\le c_0(1-\rho)\), the ratio of the fast component to the slow component is bounded by
% \[
%     Cc_0
%     \exp\left(-c(1-\rho)k+C\frac{\eta\lambda_j}{1-\rho}k\right)
%     \le
%     Cc_0\exp\left(-c(1-\rho)k\right),
% \]
% after decreasing \(c_0\).  Therefore, for all \(k\ge C(1-\rho)^{-1}\), the fast component is at most a fixed fraction of the slow component.  Consequently there exist constants \(c,C>0\) such that
% \[
%     c\exp\left(-C\frac{\eta\lambda_j}{1-\rho}k\right)
%     \le
%     \left|h_k^{\nesterov}(j)\right|
%     \le
%     C\exp\left(-c\frac{\eta\lambda_j}{1-\rho}k\right).
% \]

% Substituting this coordinate estimate into the deterministic bias gives
% \[
%     c\sum_{j\ge1}\lambda_j(\theta_j^*)^2
%     \exp\left(-C\frac{\eta\lambda_j}{1-\rho}k\right)
%     \le
%     \mathcal D_{\nesterov}(k)
%     \le
%     C\sum_{j\ge1}\lambda_j(\theta_j^*)^2
%     \exp\left(-c\frac{\eta\lambda_j}{1-\rho}k\right).
% \]
% By Assumption~\ref{ass:power-law},
% \[
%     \lambda_j(\theta_j^*)^2\eqsim j^{-1-\beta s},
%     \qquad
%     \lambda_j\eqsim j^{-\beta}.
% \]
% It remains to estimate
% \[
%     \sum_{j\ge1}j^{-1-\beta s}
%     \exp\left(-c\frac{\eta k}{1-\rho}j^{-\beta}\right).
% \]
% The condition $k \gtrsim \frac{1-\rho}{\eta}$ gives \(\frac{\eta k}{1-\rho}\ge C\). \\
% For the lower bound, take the tail
% $j\ge
%     \left\lceil
%     \left(\frac{\eta k}{1-\rho}\right)^{1/\beta}
%     \right\rceil.$
% On this tail the exponential factor is bounded below by a positive constant, hence
% \[
%     \sum_{j\ge1}j^{-1-\beta s}
%     \exp\left(-C\frac{\eta k}{1-\rho}j^{-\beta}\right)
%     \ge
%     c\left(\frac{\eta k}{1-\rho}\right)^{-s}.
% \]
% For the upper bound, an integral comparison gives
% \[
% \begin{aligned}
%     \sum_{j\ge1}j^{-1-\beta s}
%     \exp\left(-c\frac{\eta k}{1-\rho}j^{-\beta}\right)
%     &\le
%     C\int_1^\infty x^{-1-\beta s}
%     \exp\left(-c\frac{\eta k}{1-\rho}x^{-\beta}\right)\,dx \\
%     &\le
%     C\left(\frac{\eta k}{1-\rho}\right)^{-s}.
% \end{aligned}
% \]
% Combining the spectral upper and lower bounds yields
% \[
%     \mathcal D_{\nesterov}(k)
%     \eqsim
%     \left(\frac{1-\rho}{\eta k}\right)^s.
% \]
% Synthesizing the three stages completes the proof.
% \end{proof}

% \begin{lemma}[Convolution Kernel, Fully Overdamped Regime]
% \label{lem:nes_kernel0}
% Under Assumption~\ref{ass:power-law} and $\eta\lambda_1\le c_0(1-\rho)^2$ with sufficiently small \(c_0>0\), we have:
% \begin{equation*}
%     \mathcal{K}_{\nesterov}(k)
%     \eqsim \begin{cases}
%     k^2, &\qquad k\lesssim \frac{1}{1-\rho},\\
%     \frac{1}{(1-\rho)^2}, &\qquad \frac{1}{1-\rho}\lesssim k \lesssim \frac{1-\rho}{\eta},\\
%     \frac{1}{(1-\rho)^2}
%     \left(\frac{1-\rho}{\eta k}\right)^{2-\frac1\beta}, &\qquad k \gtrsim \frac{1-\rho}{\eta}.
% \end{cases}
% \end{equation*}
% \end{lemma}

% \begin{proof}
% The Nesterov Green function satisfies
% \[
%     G_{k+1}^{\nesterov}(j)
%     =
%     (1+\rho)(1-\eta\lambda_j)G_k^{\nesterov}(j)
%     -\rho(1-\eta\lambda_j)G_{k-1}^{\nesterov}(j),
%     \qquad
%     G_0^{\nesterov}(j)=1,
%     \qquad
%     G_{-1}^{\nesterov}(j)=0.
% \]
% The characteristic roots are the same roots used in Proposition~\ref{prop:nes_deterministic_bias0}:
% \[
%     y_{j,+},y_{j,-}
%     =
%     \frac{(1+\rho)(1-\eta\lambda_j)
%     \pm
%     \sqrt{(1+\rho)^2(1-\eta\lambda_j)^2-4\rho(1-\eta\lambda_j)}}{2}.
% \]
% In the fully overdamped regime, \(0\le\eta\lambda_j\le c_0(1-\rho)^2\). The root calculation from Proposition~\ref{prop:nes_deterministic_bias0} gives
% \[
%     1-\frac{2\eta\lambda_j}{1-\rho}
%     \le
%     y_{j,+}
%     \le
%     1-\frac{\eta\lambda_j}{2(1-\rho)},
%     \qquad
%     0\le y_{j,-}\le1-\frac{1-\rho}{2},
% \]
% and
% \[
%     y_{j,+}-y_{j,-}
%     \eqsim
%     1-\rho.
% \]
% The last comparison follows because \(y_{j,+}\) remains within \(O(\eta\lambda_j/(1-\rho))\) of one, while \(y_{j,-}\) remains within \(O(\eta\lambda_j/(1-\rho))\) of \(\rho\).

% Solving the impulse initial conditions gives the exact formula
% \[
%     G_k^{\nesterov}(j)
%     =
%     \frac{y_{j,+}^{k+1}-y_{j,-}^{k+1}}{y_{j,+}-y_{j,-}} = \sum_{m=0}^{k} y_{j,+}^{k-m} y_{j,-}^m.
% \]

% By definition, the convolution kernel is:
% \[
%     \mathcal K_{\nesterov}(k) = \sum_{j\ge1}\lambda_j^2 G_k^{\nesterov}(j)^2.
% \]
% We divide the time $k$ into three successive stages to evaluate $\mathcal K_{\nesterov}(k)$.

% \textbf{Stage 1: $k \lesssim \frac{1}{1-\rho}$.} \\
% In this extremely short-time regime, the roots are very close to 1 and $\rho$ respectively. Using the uniform bounds $y_{j,+} \ge 1 - 2c_0(1-\rho)$ and $y_{j,-} \ge y_{j,+}y_{j,-} = \rho(1-\eta\lambda_j) \ge \rho(1-c_0(1-\rho)^2)$, we have for all $m \le k \lesssim \frac{1}{1-\rho}$:
% \[
%     y_{j,+}^{k-m} \ge \left(1 - 2c_0(1-\rho)\right)^{\frac{C}{1-\rho}} \eqsim 1, \qquad y_{j,-}^m \ge \rho^m \left(1-c_0(1-\rho)^2\right)^m \ge \rho^{\frac{C}{1-\rho}} \cdot 1 \eqsim e^{-C} \eqsim 1.
% \]
% Since all terms in the algebraic expansion of $G_k^{\nesterov}(j)$ are strictly positive and bounded away from zero by a constant, the sum is proportional to the number of terms:
% \[
%     G_k^{\nesterov}(j) = \sum_{m=0}^{k} y_{j,+}^{k-m} y_{j,-}^m \eqsim \sum_{m=0}^{k} 1 = k+1 \eqsim k.
% \]
% Since this holds for the leading modes which dominate the convergent spectral sum ($\sum_{j\ge1}\lambda_j^2 < \infty$ due to $\beta > 1$), we obtain:
% \[
%     \mathcal K_{\nesterov}(k) \eqsim \sum_{j\ge1}\lambda_j^2 k^2 = k^2 \sum_{j\ge1}\lambda_j^2 \eqsim k^2.
% \]

% \textbf{Stage 2: $\frac{1}{1-\rho} \lesssim k \lesssim \frac{1-\rho}{\eta}$.} \\
% In this intermediate regime, $k$ is sufficiently large such that the fast root decays exponentially:
% \[
%     y_{j,-}^{k+1} \le \left(1-\frac{1-\rho}{2}\right)^{k+1} \ll y_{j,+}^{k+1}.
% \]
% Thus, the Green function is  dominated by the slow root:
% \[
%     G_k^{\nesterov}(j) \simeq \frac{y_{j,+}^{k+1}}{y_{j,+}-y_{j,-}} \eqsim \frac{y_{j,+}^{k+1}}{1-\rho}.
% \]
% For the leading modes ($j = \mathcal{O}(1)$), since $k \lesssim \frac{1-\rho}{\eta}$, the exponent is bounded by:
% \[
%     y_{j,+}^{k+1} \ge \left(1 - \frac{2\eta\lambda_j}{1-\rho}\right)^{k+1} \eqsim \exp\left(-C \frac{\eta\lambda_j}{1-\rho}k\right) \ge \exp(-C \lambda_j) \eqsim 1.
% \]
% This implies the slow root has not yet experienced significant exponential decay for the leading modes, giving $G_k^{\nesterov}(j) \eqsim \frac{1}{1-\rho}$. For the tail modes, we trivially have $G_k^{\nesterov}(j) \lesssim \frac{1}{1-\rho}$. Therefore, the kernel is dominated by these non-decayed leading modes:
% \[
%     \mathcal K_{\nesterov}(k) = \sum_{j\ge1}\lambda_j^2 G_k^{\nesterov}(j)^2 \eqsim \sum_{j=1}^{\mathcal{O}(1)}\lambda_j^2 \frac{1}{(1-\rho)^2} \eqsim \frac{1}{(1-\rho)^2}.
% \]

% \textbf{Stage 3: $k \gtrsim \frac{1-\rho}{\eta}$.} \\
% In this long-time regime, the exponential decay of the slow root becomes dominant across the spectrum. Substituting the two-sided root bounds into $G_k^{\nesterov}(j) \simeq \frac{y_{j,+}^{k+1}}{y_{j,+}-y_{j,-}}$, there exist constants $c, C > 0$ such that:
% \begin{equation}
% \label{eq:nes_green_sgd_pointwise}
%     \frac{c}{1-\rho}
%     \exp\left(-C\frac{\eta\lambda_j}{1-\rho}k\right)
%     \le
%     |G_k^{\nesterov}(j)|
%     \le
%     \frac{C}{1-\rho}
%     \exp\left(-c\frac{\eta\lambda_j}{1-\rho}k\right).
% \end{equation}
% Squaring this and summing over the spectrum gives:
% \begin{equation}
% \label{eq:nes_kernel_sgd_spectral_comparison}
%     \frac{c}{(1-\rho)^2}
%     \sum_{j\ge1}\lambda_j^2
%     \exp\left(-C\frac{\eta\lambda_j}{1-\rho}k\right)
%     \le
%     \mathcal K_{\nesterov}(k)
%     \le
%     \frac{C}{(1-\rho)^2}
%     \sum_{j\ge1}\lambda_j^2
%     \exp\left(-c\frac{\eta\lambda_j}{1-\rho}k\right).
% \end{equation}
% Using $\lambda_j \eqsim j^{-\beta}$, we estimate the corresponding spectral sum:
% \[
%     \sum_{j\ge1}j^{-2\beta}
%     \exp\left(-c\frac{\eta k}{1-\rho}j^{-\beta}\right).
% \]
% For the lower bound, let $J:=\left\lceil \left(\frac{\eta k}{1-\rho}\right)^{1/\beta}\right\rceil$. For $j\ge J$, we have $\frac{\eta k}{1-\rho}j^{-\beta}\le1$, yielding $\exp\left(-C\frac{\eta k}{1-\rho}j^{-\beta}\right)\ge e^{-C}$. Thus:
% \[
%     \sum_{j\ge1}j^{-2\beta}
%     \exp\left(-C\frac{\eta k}{1-\rho}j^{-\beta}\right)
%     \ge
%     e^{-C}\sum_{j\ge J}j^{-2\beta}
%     \ge
%     c J^{1-2\beta}
%     \ge
%     c\left(\frac{\eta k}{1-\rho}\right)^{-\left(2-\frac1\beta\right)}.
% \]
% For the upper bound, an integral comparison with $u=c\frac{\eta k}{1-\rho}x^{-\beta}$ gives:
% \begin{align*}
%     \sum_{j\ge1}j^{-2\beta}
%     \exp\left(-c\frac{\eta k}{1-\rho}j^{-\beta}\right)
%     &\le
%     C\int_1^\infty x^{-2\beta}
%     \exp\left(-c\frac{\eta k}{1-\rho}x^{-\beta}\right)\,dx \\
%     &=
%     C\left(\frac{\eta k}{1-\rho}\right)^{-\left(2-\frac1\beta\right)}
%     \int_0^{c\frac{\eta k}{1-\rho}}u^{1-\frac1\beta}e^{-u}\,du \\
%     &\le
%     C\left(\frac{\eta k}{1-\rho}\right)^{-\left(2-\frac1\beta\right)}.
% \end{align*}
% Combining these bounds yields:
% \[
%     \mathcal K_{\nesterov}(k)
%     \eqsim
%     \frac{1}{(1-\rho)^2}
%     \left(\frac{1-\rho}{\eta k}\right)^{2-\frac1\beta}.
% \]
% Synthesizing the three stages completes the proof.
% \end{proof}

% \begin{proposition}[Multiplicative Noise Bound, Fully Overdamped Regime]
% \label{prop:nes_self_noise_bounds0}
% In the fully overdamped regime, under the stability condition \(\eta\le cB(1-\rho)\) with \(c>0\) sufficiently small, and for 
% \[
%     k\ge C\max\{(1-\rho)^{-1},(1-\rho)/\eta\},
% \]
% we have:
% \[
%     \frac{\eta^2}{B}\sum_{\tau=0}^{k-1}\mathcal K_{\nesterov}(\tau)\mathbb{E}\mathcal E_{k-\tau}
%     \eqsim
%     \frac{\eta}{B(1-\rho)}
%     \left(\frac{1-\rho}{\eta k}\right)^{(2-\frac1\beta)\wedge s} + \frac{\eta^2\sigma^2}{B^2(1-\rho)^2}.
% \]
% \end{proposition}

% \begin{proof}
% Write the self-noise convolution in the equivalent time-forward form
% \[
%     \frac{\eta^2}{B}
%     \sum_{\tau=0}^{k-1}\mathcal K_{\nesterov}(\tau)
%     \mathbb{E}\mathcal E_{k-\tau}
%     =
%     \frac{\eta^2}{B}
%     \sum_{m=1}^{k}\mathcal K_{\nesterov}(k-m)
%     \mathbb{E}\mathcal E_m.
% \]
% In the proof, we denote this quantity by \(\mathcal M_k^{\nesterov}\).  Lemma~\ref{lem:nes_replace_lookahead_by_next_loss} and Proposition~\ref{prop:nes_risk_recursion} give
% \[
%     \mathbb{E}\mathcal E_m
%     \le
%     C\left[
%         \mathcal D_{\nesterov}(m)
%         +
%         \frac{\eta^2\sigma^2}{B}
%         \sum_{r=0}^{m-1}\mathcal K_{\nesterov}(r)
%         +
%         \mathcal M_m^{\nesterov}
%     \right].
% \]
% Substituting this bound into the definition of $\mathcal M_k^{\nesterov}$ yields three convolution terms:
% \begin{align}
% \label{eq:nes_prop_bootstrap_start}
%     \mathcal M_k^{\nesterov}
%     \le{}&
%     C\frac{\eta^2}{B}\sum_{m=1}^{k}
%         \mathcal K_{\nesterov}(k-m)\mathcal D_{\nesterov}(m) \nonumber \\
%     & + C\frac{\eta^4\sigma^2}{B^2}\sum_{m=1}^{k}\mathcal K_{\nesterov}(k-m)\sum_{r=0}^{m-1}\mathcal K_{\nesterov}(r) \nonumber \\
%     & + C\frac{\eta^2}{B}\sum_{m=1}^{k}
%         \mathcal K_{\nesterov}(k-m)\mathcal M_m^{\nesterov} .
% \end{align}

% We first bound the label-noise convolution term. Changing the summation indices and exploiting the non-negativity of the kernel, we have:
% \begin{align*}
%     \frac{\eta^4\sigma^2}{B^2}\sum_{m=1}^{k}\mathcal K_{\nesterov}(k-m)\sum_{r=0}^{m-1}\mathcal K_{\nesterov}(r)
%     &= \frac{\eta^4\sigma^2}{B^2} \sum_{n=0}^{k-1} \mathcal K_{\nesterov}(n) \sum_{r=0}^{k-1-n} \mathcal K_{\nesterov}(r) \\
%     &\le \frac{\eta^4\sigma^2}{B^2} \left( \sum_{n=0}^{\infty} \mathcal K_{\nesterov}(n) \right)^2.
% \end{align*}
% From the saturated Green-energy estimate in Lemma~\ref{lem:nes_label_noise}, we have $\sum_{n=0}^{\infty} \mathcal K_{\nesterov}(n) \le \frac{C}{\eta(1-\rho)}$. Therefore, this term is bounded by $\mathcal{O}\left( \frac{\eta^2\sigma^2}{B^2(1-\rho)^2} \right)$. For the lower bound, since $k \ge C\frac{1-\rho}{\eta}$, the sum over the restricted range $n, r \le k/2$ captures a constant fraction of the infinite sum, yielding the matching lower bound $\eqsim \frac{\eta^2\sigma^2}{B^2(1-\rho)^2}$. This non-decaying constant term represents the fundamental noise floor amplified by momentum.

% Next, we rigorously evaluate the deterministic convolution term. For asymptotic simplicity, we write $I_{\mathrm{det}} := \sum_{m=1}^{k} \mathcal K_{\nesterov}(k-m)\mathcal D_{\nesterov}(m)$. Based on the piecewise asymptotic forms of $\mathcal K_{\nesterov}$ and $\mathcal D_{\nesterov}$ established in Proposition~\ref{prop:nes_deterministic_bias0} and Lemma~\ref{lem:nes_kernel0}, we divide the summation into five distinct stages:

% \textbf{Stage 1:} $1 \le m \lesssim \frac{1}{1-\rho}$. \\
% Here, $\mathcal D_{\nesterov}(m) \eqsim 1$ and $k-m \eqsim k \gtrsim \frac{1-\rho}{\eta}$, giving $\mathcal K_{\nesterov}(k-m) \eqsim \frac{1}{(1-\rho)^2} \left(\frac{1-\rho}{\eta k}\right)^{2-\frac1\beta}$.
% \[
%     \sum_{m=1}^{\frac{1}{1-\rho}} 1 \cdot \frac{1}{(1-\rho)^2} \left(\frac{1-\rho}{\eta k}\right)^{2-\frac1\beta} \eqsim \frac{1}{(1-\rho)^3} \left(\frac{1-\rho}{\eta k}\right)^{2-\frac1\beta}.
% \]

% \textbf{Stage 2:} $\frac{1}{1-\rho} \lesssim m \lesssim \frac{1-\rho}{\eta}$. \\
% Here, $\mathcal D_{\nesterov}(m) \eqsim 1$ and $\mathcal K_{\nesterov}(k-m) \eqsim \frac{1}{(1-\rho)^2} \left(\frac{1-\rho}{\eta k}\right)^{2-\frac1\beta}$. The length of this interval is $\mathcal{O}(\frac{1-\rho}{\eta})$.
% \[
%     \sum_{m=\frac{1}{1-\rho}}^{\frac{1-\rho}{\eta}} 1 \cdot \frac{1}{(1-\rho)^2} \left(\frac{1-\rho}{\eta k}\right)^{2-\frac1\beta} \eqsim \frac{1-\rho}{\eta} \cdot \frac{1}{(1-\rho)^2} \left(\frac{1-\rho}{\eta k}\right)^{2-\frac1\beta} = \frac{1}{\eta(1-\rho)} \left(\frac{1-\rho}{\eta k}\right)^{2-\frac1\beta}.
% \]

% \textbf{Stage 3:} $\frac{1-\rho}{\eta} \lesssim m \lesssim k - \frac{1-\rho}{\eta}$. \\
% Here, $\mathcal D_{\nesterov}(m) \eqsim \left(\frac{1-\rho}{\eta m}\right)^s$ and $\mathcal K_{\nesterov}(k-m) \eqsim \frac{1}{(1-\rho)^2} \left(\frac{1-\rho}{\eta (k-m)}\right)^{2-\frac1\beta}$.
% \begin{align*}
%     &\sum_{m=\frac{1-\rho}{\eta}}^{k - \frac{1-\rho}{\eta}} \left(\frac{1-\rho}{\eta m}\right)^s \frac{1}{(1-\rho)^2} \left(\frac{1-\rho}{\eta (k-m)}\right)^{2-\frac1\beta} \\
%     &\eqsim \frac{1}{(1-\rho)^2} \left(\frac{1-\rho}{\eta}\right)^{2-\frac1\beta+s} \left[ \sum_{m=\frac{1-\rho}{\eta}}^{k/2} \frac{1}{m^s} \frac{1}{k^{2-\frac1\beta}} + \sum_{m=k/2}^{k-\frac{1-\rho}{\eta}} \frac{1}{k^s} \frac{1}{(k-m)^{2-\frac1\beta}} \right] \\
%     &\eqsim \frac{1}{(1-\rho)^2} \left(\frac{1-\rho}{\eta}\right)^{2-\frac1\beta+s} \left[ k^{-\left(2-\frac1\beta\right)} \begin{cases} k^{1-s}, & s<1 \\ \left(\frac{1-\rho}{\eta}\right)^{1-s}, & s>1 \end{cases} + k^{-s} \left(\frac{1-\rho}{\eta}\right)^{-\left(1-\frac1\beta\right)} \right] \\
%     &\eqsim \frac{1}{\eta(1-\rho)} \left(\frac{1-\rho}{\eta k}\right)^{2-\frac1\beta} + \frac{1}{\eta(1-\rho)} \left(\frac{1-\rho}{\eta k}\right)^s \eqsim \frac{1}{\eta(1-\rho)} \left(\frac{1-\rho}{\eta k}\right)^{\left(2-\frac1\beta\right)\wedge s}.
% \end{align*}

% \textbf{Stage 4:} $k - \frac{1-\rho}{\eta} \lesssim m \lesssim k - \frac{1}{1-\rho}$. \\
% Here, $\mathcal D_{\nesterov}(m) \eqsim \left(\frac{1-\rho}{\eta k}\right)^s$ and $\mathcal K_{\nesterov}(k-m) \eqsim \frac{1}{(1-\rho)^2}$ since $k-m \in [\frac{1}{1-\rho}, \frac{1-\rho}{\eta}]$. The length is $\mathcal{O}(\frac{1-\rho}{\eta})$.
% \[
%     \sum_{m=k-\frac{1-\rho}{\eta}}^{k-\frac{1}{1-\rho}} \left(\frac{1-\rho}{\eta k}\right)^s \frac{1}{(1-\rho)^2} \eqsim \frac{1-\rho}{\eta} \cdot \frac{1}{(1-\rho)^2} \left(\frac{1-\rho}{\eta k}\right)^s = \frac{1}{\eta(1-\rho)} \left(\frac{1-\rho}{\eta k}\right)^s.
% \]

% \textbf{Stage 5:} $k - \frac{1}{1-\rho} \lesssim m \le k$. \\
% Here, $\mathcal D_{\nesterov}(m) \eqsim \left(\frac{1-\rho}{\eta k}\right)^s$ and $\mathcal K_{\nesterov}(k-m) \eqsim (k-m)^2$ since $k-m \in [0, \frac{1}{1-\rho}]$.
% \[
%     \sum_{m=k-\frac{1}{1-\rho}}^{k} \left(\frac{1-\rho}{\eta k}\right)^s (k-m)^2 \eqsim \left(\frac{1-\rho}{\eta k}\right)^s \sum_{r=0}^{\frac{1}{1-\rho}} r^2 \eqsim \frac{1}{(1-\rho)^3} \left(\frac{1-\rho}{\eta k}\right)^s.
% \]

% Since $\eta \lesssim (1-\rho)^2$ in the Fully Overdamped Regime, we have $\frac{1}{\eta(1-\rho)} \gtrsim \frac{1}{(1-\rho)^3}$. Thus, the sums from Stages 1 and 5 are strictly lower order. The dominant contribution to $I_{\mathrm{det}}$ comes from Stages 2, 3, and 4, yielding:
% \[
%     I_{\mathrm{det}} \eqsim \frac{1}{\eta(1-\rho)} \left(\frac{1-\rho}{\eta k}\right)^{\left(2-\frac1\beta\right)\wedge s}.
% \]
% Multiplying by $\frac{\eta^2}{B}$, the deterministic convolution term is bounded by $\mathcal{O}\left( \frac{\eta}{B(1-\rho)} \left(\frac{1-\rho}{\eta k}\right)^{\left(2-\frac1\beta\right)\wedge s} \right)$. The matching lower bound follows identically from the two-sided nature of the constituent bounds.

% It remains to absorb the convolution with \(\mathcal M_m^{\nesterov}\).  Expanding \(\mathcal M_m^{\nesterov}\) and exchanging the order of summation gives
% \[
% \begin{aligned}
%     &\frac{\eta^2}{B}
%     \sum_{m=1}^{k}
%     \mathcal K_{\nesterov}(k-m)
%     \mathcal M_m^{\nesterov} \\
%     &\qquad=
%     \frac{\eta^2}{B}
%     \sum_{\ell=1}^{k}
%     \left[
%         \frac{\eta^2}{B}
%         \sum_{m=\ell}^{k}
%         \mathcal K_{\nesterov}(k-m)
%         \mathcal K_{\nesterov}(m-\ell)
%     \right]
%     \mathbb{E}\mathcal E_\ell.
% \end{aligned}
% \]
% All terms are nonnegative.  Lemma~\ref{lem:nes_kernel0}, together with the saturated Green-energy estimate behind Lemma~\ref{lem:nes_label_noise}, gives
% \[
%     \sum_{r=0}^{n}
%     \mathcal K_{\nesterov}(n-r)
%     \mathcal K_{\nesterov}(r)
%     \le
%     C\left(\sum_{r=0}^{\infty}\mathcal K_{\nesterov}(r)\right)
%     \mathcal K_{\nesterov}(n)
%     \le
%     \frac{C}{\eta(1-\rho)}
%     \mathcal K_{\nesterov}(n).
% \]
% Here, in the fully overdamped regime, Lemma~\ref{lem:nes_label_noise} yields
% \[
%     \sum_{r=0}^{\infty}\mathcal K_{\nesterov}(r)
%     \le
%     \frac{C}{\eta(1-\rho)}.
% \]
% Consequently,
% \[
% \begin{aligned}
%     \frac{\eta^2}{B}
%     \sum_{m=1}^{k}
%     \mathcal K_{\nesterov}(k-m)
%     \mathcal M_m^{\nesterov}
%     &\le
%     C\frac{\eta}{B(1-\rho)}
%     \frac{\eta^2}{B}
%     \sum_{\ell=1}^{k}
%     \mathcal K_{\nesterov}(k-\ell)
%     \mathbb{E}\mathcal E_\ell \\
%     &\le
%     C\frac{\eta}{B(1-\rho)}
%     \mathcal M_k^{\nesterov}.
% \end{aligned}
% \]
% Choosing the constant in \(\eta\le cB(1-\rho)\) sufficiently small absorbs this term into the left-hand side.  We obtain
% \[
%     \mathcal M_k^{\nesterov}
%     \eqsim
%     \frac{\eta}{B(1-\rho)}
%     \left(\frac{1-\rho}{\eta k}\right)^{(2-\frac1\beta)\wedge s} + \frac{\eta^2\sigma^2}{B^2(1-\rho)^2}.
% \]
% For the lower bound, the risk recursion and the nonnegativity of the variance terms imply that \(\mathbb{E}\mathcal E_m\) is bounded below by a constant multiple of \(\mathcal D_{\nesterov}(m)\).  The deterministic-convolution lower bound above then gives
% \[
%     \mathcal M_k^{\nesterov}
%     \ge
%     c\frac{\eta}{B(1-\rho)}
%     \left(\frac{1-\rho}{\eta k}\right)^{(2-\frac1\beta)\wedge s} + c'\frac{\eta^2\sigma^2}{B^2(1-\rho)^2}.
% \]
% This proves the proposition.
% \end{proof}


% \begin{theorem}[Scaling Law, Fully Overdamped Regime]
% \label{thm:nes_scaling_sgd}
% Under Assumptions~\ref{ass:diagonal-covariance} and~\ref{ass:power-law}, the stability condition $\eta \lesssim B(1-\rho)$, and for the fully overdamped regime $\eta\lesssim(1-\rho)^2$ with $k \gtrsim \frac{1-\rho}{\eta}$, the expected excess risk satisfies:
% \[
%     \mathbb{E}\mathcal E_k
%     \eqsim
%     \left(\frac{1-\rho}{\eta k}\right)^s
%     +
%     \frac{\sigma^2\eta}{B(1-\rho)}.
% \]
% \end{theorem}

% \begin{proof}
% By the expected risk recursion established in Proposition~\ref{prop:nes_risk_recursion} and the look-ahead replacement from Lemma~\ref{lem:nes_replace_lookahead_by_next_loss}, the total expected excess risk decomposes into three additive components: the deterministic bias, the accumulated label noise, and the multiplicative noise feedback $\mathcal{M}_k^{\nesterov}$:
% \[
%     \mathbb{E}\mathcal E_k
%     \eqsim
%     \cS_{\nesterov}(k)
%     +
%     \frac{\eta^2\sigma^2}{B}\sum_{\tau=0}^{k-1}\mathcal{K}_{\nesterov}(\tau)
%     +
%     \mathcal M_k^{\nesterov}.
% \]

% For the Fully Overdamped Regime ($k \gtrsim \frac{1-\rho}{\eta}$), we substitute the corresponding asymptotic bounds derived in the preceding propositions:
% \begin{enumerate}
%     \item \textbf{Deterministic Bias} (Proposition~\ref{prop:nes_deterministic_bias0}):
%     \[
%         \cS_{\nesterov}(k) \eqsim \left(\frac{1-\rho}{\eta k}\right)^s.
%     \]
%     \item \textbf{Label Noise} (Lemma~\ref{lem:nes_label_noise}): 
%     In the fully overdamped regime $\eta \lesssim (1-\rho)^2$, we have $\frac{\eta}{1-\rho} \lesssim 1-\rho \ll 1$. Since $\beta > 1$, the minimum in Lemma~\ref{lem:nes_label_noise} is attained by the linear term. Thus:
%     \[
%         \frac{\eta^2\sigma^2}{B}\sum_{\tau=0}^{k-1}\mathcal{K}_{\nesterov}(\tau) 
%         \eqsim \frac{\sigma^2}{B} \min\left\{ \left(\frac{\eta}{1-\rho}\right)^{\frac1\beta}, \frac{\eta}{1-\rho} \right\} 
%         = \frac{\sigma^2\eta}{B(1-\rho)}.
%     \]
%     \item \textbf{Multiplicative Noise} (Proposition~\ref{prop:nes_self_noise_bounds0}):
%     \[
%         \mathcal M_k^{\nesterov} \eqsim \frac{\eta}{B(1-\rho)} \left(\frac{1-\rho}{\eta k}\right)^{\left(2-\frac1\beta\right)\wedge s} + \frac{\eta^2\sigma^2}{B^2(1-\rho)^2}.
%     \]
% \end{enumerate}

% Summing these three terms, we analyze the structural absorption of the multiplicative noise fragments. First, under the given stability condition $\eta \lesssim B(1-\rho)$, the constant noise floor originating from the multiplicative noise is strictly bounded by the primary label noise:
% \[
%     \frac{\eta^2\sigma^2}{B^2(1-\rho)^2} = \left(\frac{\eta}{B(1-\rho)}\right) \frac{\sigma^2\eta}{B(1-\rho)} \lesssim \frac{\sigma^2\eta}{B(1-\rho)}.
% \]

% Next, we bound the polynomial tail of the multiplicative noise, denoted $M_3(k) = \frac{\eta}{B(1-\rho)} \left(\frac{1-\rho}{\eta k}\right)^{\left(2-\frac1\beta\right)\wedge s}$:
% \begin{itemize}
%     \item If $s \le 2-\frac{1}{\beta}$, then $M_3(k) = \frac{\eta}{B(1-\rho)} \left(\frac{1-\rho}{\eta k}\right)^s \lesssim \left(\frac{1-\rho}{\eta k}\right)^s \eqsim \cS_{\nesterov}(k)$. The term is completely absorbed into the deterministic polynomial bias.
%     \item If $s > 2-\frac{1}{\beta}$, then $M_3(k) = \frac{\eta}{B(1-\rho)} \left(\frac{1-\rho}{\eta k}\right)^{2-\frac{1}{\beta}}$.
%     When the deterministic bias dominates the label noise $\left( \left(\frac{1-\rho}{\eta k}\right)^{-s} \le \frac{B(1-\rho)}{\eta} \right)$, reversing the scaling inequality guarantees explicit containment:
%     \[
%         M_3(k) \le \left(\frac{\eta}{B(1-\rho)}\right)^{\frac{2-1/\beta}{s}} \left(\frac{1-\rho}{\eta k}\right)^s \lesssim \cS_{\nesterov}(k).
%     \]
%     When the label noise dominates the deterministic bias, since $k \gtrsim \frac{1-\rho}{\eta}$, we mathematically maintain $\left(\frac{1-\rho}{\eta k}\right)^{2-\frac{1}{\beta}} \lesssim 1$. Consequently:
%     \[
%         M_3(k) \lesssim \frac{\eta}{B(1-\rho)} \cdot 1 \eqsim \text{LN}(k).
%     \]
% \end{itemize}
% In all distinct bounding configurations, the multiplicative structural element is unconditionally bounded by $\max(\cS_{\nesterov}(k), \text{LN}(k))$. Thus, it is systematically absorbed, leaving the identical final scaling law:
% \[
%     \mathbb{E}\mathcal E_k
%     \eqsim
%     \left(\frac{1-\rho}{\eta k}\right)^s
%     +
%     \frac{\sigma^2\eta}{B(1-\rho)}.
% \]
% This completes the proof.
% \end{proof}

% \begin{theorem}[Scaling Law, No Label Noise, Fully Overdamped Regime]
% \label{thm:nes_scaling_sgd_noiseless}
% Under Assumptions~\ref{ass:diagonal-covariance} and~\ref{ass:power-law}, with \(\sigma=0\), the stability condition \(\eta \lesssim B(1-\rho)\), and in the fully overdamped regime \(\eta\lesssim(1-\rho)^2\), for \(k \gtrsim \frac{1-\rho}{\eta}\), the expected excess risk satisfies:
% \begin{equation}
% \label{eq:nes_scaling_sgd_noiseless}
%     \EE[\cE(\btheta_k)]
%     \eqsim
%     \left(\frac{1-\rho}{\eta k}\right)^s
%     +
%     \frac{\eta}{B(1-\rho)}
%     \left(\frac{1-\rho}{\eta k}\right)^{s \wedge \left(2-\frac1\beta\right)}.
% \end{equation}
% \end{theorem}

% \begin{proof}
% From Proposition~\ref{prop:nes_risk_recursion} and Lemma~\ref{lem:nes_replace_lookahead_by_next_loss}, with \(\sigma=0\), the expected risk simplifies to:
% \[
%     \EE[\cE(\btheta_k)]
%     \eqsim
%     \cS_{\nesterov}(k)
%     +
%     \mathcal M_k^{\nesterov}.
% \]
% In the fully overdamped regime \(k \gtrsim \frac{1-\rho}{\eta}\), Proposition~\ref{prop:nes_deterministic_bias0} gives:
% \[
%     \cS_{\nesterov}(k) \eqsim \left(\frac{1-\rho}{\eta k}\right)^s,
% \]
% and Proposition~\ref{prop:nes_self_noise_bounds0} yields:
% \[
%     \mathcal M_k^{\nesterov}
%     \eqsim
%     \frac{\eta}{B(1-\rho)}
%     \left(\frac{1-\rho}{\eta k}\right)^{s \wedge \left(2-\frac1\beta\right)}.
% \]
% Substituting these bounds directly completes the proof:
% \[
%     \EE[\cE(\btheta_k)]
%     \eqsim
%     \left(\frac{1-\rho}{\eta k}\right)^s
%     +
%     \frac{\eta}{B(1-\rho)}
%     \left(\frac{1-\rho}{\eta k}\right)^{s \wedge \left(2-\frac1\beta\right)}.
% \]
% \end{proof}





% \subsection{Loss in the mixed overdamped--underdamped regime}
% \label{app:nes_mixed_loss}

% When the leading spectral modes are underdamped, the loss contains an exponentially damped edge contribution in addition to the polynomial overdamped tail. The crossover boundary between the two spectral blocks is $J_{\max}$ from~\eqref{eq:J_max_nes}.


% Define the Nesterov edge damping factor
% \[
%     \gamma_{\nesterov}
%     :=
%     \rho\left(1-\frac{(1-\rho)^2}{(1+\rho)^2}\right)
%     =
%     \frac{4\rho^2}{(1+\rho)^2}.
% \]
% It satisfies
% \[
%     \gamma_{\nesterov}=1-(1-\rho)+O((1-\rho)^2),
% \]
% so the underdamped edge contribution decays exponentially once $k(1-\rho)\to\infty$.

% \begin{proposition}[Deterministic Bias, Mixed Regime]
% \label{prop:nes_deterministic_bias}
% Under Assumptions~\ref{ass:diagonal-covariance} and~\ref{ass:power-law}, assume the stability condition $\eta\lesssim \min\{1,B(1-\rho)\}$. For the mixed regime \(\eta\gtrsim(1-\rho)^2\), the deterministic bias decomposes into $\cS_{\nesterov}(k) = \cS_{\mathrm{over}}(k) + \cS_{\mathrm{under}}(k)$, where:
% \begin{equation*}
%     \cS_{\mathrm{over}}(k)
%     \eqsim
%     \begin{cases}
%     \left(\frac{(1-\rho)^2}{\eta}\right)^s, &\qquad k\lesssim\frac{1}{1-\rho},\\
%     \left(\frac{1-\rho}{\eta k}\right)^s, &\qquad k\gtrsim\frac{1}{1-\rho},
%     \end{cases}
% \end{equation*}
% and
% \begin{equation*}
%     \cS_{\mathrm{under}}(k)
%     \eqsim
%     \begin{cases}
%     \rho^k, &\qquad k\lesssim\frac{1}{\eta},\\
%     \frac{\rho^k}{(\eta k)^s}, &\qquad \frac{1}{\eta} \lesssim k \lesssim \frac{1}{(1-\rho)^2},\\
%     \frac{1}{\eta k}\left(\frac{(1-\rho)^2}{\eta}\right)^{s-1}\gamma_{\nesterov}^{\,k}, &\qquad \frac{1}{(1-\rho)^2} \lesssim k \lesssim \eta^{1/\beta}(1-\rho)^{-2-\frac{2}{\beta}},\\
%     \left(\frac{(1-\rho)^2}{\eta}\right)^{s+\frac{1}{\beta}}\gamma_{\nesterov}^{\,k}, &\qquad k \gtrsim \eta^{1/\beta}(1-\rho)^{-2-\frac{2}{\beta}}.
%     \end{cases}
% \end{equation*}
% \end{proposition}

% \begin{proof}
% For each coordinate, write \(e_k^{\mathrm{det}}(j)=-h_k^{\nesterov}(j)\theta_j^*\). The multiplier satisfies
% \[
%     h_{k+1}^{\nesterov}(j)
%     =
%     (1+\rho)(1-\eta\lambda_j)h_k^{\nesterov}(j)
%     -\rho(1-\eta\lambda_j)h_{k-1}^{\nesterov}(j),
%     \qquad
%     h_0^{\nesterov}(j)=1,
%     \qquad
%     h_1^{\nesterov}(j)=1-\eta\lambda_j.
% \]
% The roots are
% \[
%     y_{j,+},y_{j,-}
%     =
%     \frac{(1+\rho)(1-\eta\lambda_j)
%     \pm
%     \sqrt{(1+\rho)^2(1-\eta\lambda_j)^2-4\rho(1-\eta\lambda_j)}}{2}.
% \]
% The discriminant is negative precisely when
% \[
%     (1-\sqrt\rho)^2<\eta\lambda_j<(1+\sqrt\rho)^2.
% \]
% Thus the comparison between $\eta$ and $(1-\rho)^2$ determines whether a coordinate is overdamped or underdamped. In the mixed regime $\eta \gtrsim (1-\rho)^2$, the crossover index is defined by $\eta\lambda_{J_{\max}} \eqsim (1-\rho)^2$. We decompose the deterministic bias into the overdamped tail ($j > J_{\max}$) and the underdamped leading block ($j \le J_{\max}$).

% \textbf{1. Overdamped Tail ($\cS_{\mathrm{over}}$):} \\
% For overdamped coordinates in the stable small-step range, the root comparison used in Proposition~\ref{prop:nes_deterministic_bias0} gives constants $c_1,C_1>0$ such that
% \begin{equation*}
%     |h_k^{\nesterov}(j)|
%     \le
%     C_1\exp\left(-c_1\frac{\eta\lambda_j}{1-\rho}k\right).
% \end{equation*}
% Moreover, if $\eta\lambda_j\le c_1(1-\rho)^2$ and $k\ge C_1(1-\rho)^{-1}$, the slow root dominates the fast root and
% \begin{equation*}
%     |h_k^{\nesterov}(j)|
%     \ge
%     c_1\exp\left(-C_1\frac{\eta\lambda_j}{1-\rho}k\right).
% \end{equation*}

% For $k \lesssim \frac{1}{1-\rho}$, the exponent satisfies $\frac{\eta\lambda_j}{1-\rho}k \lesssim \frac{(1-\rho)^2}{1-\rho}\frac{1}{1-\rho} = 1$. The exponential decay is bounded away from zero, so $|h_k^{\nesterov}(j)| \eqsim 1$. The sum is strictly dominated by its lower limit at the crossover:
% \[
%     \mathcal D_{\mathrm{over}}(k) = \sum_{j > J_{\max}} \lambda_j(\theta_j^*)^2 h_k^{\nesterov}(j)^2 \eqsim \sum_{j > J_{\max}} j^{-1-\beta s} \eqsim (J_{\max})^{-\beta s} \eqsim \left(\frac{(1-\rho)^2}{\eta}\right)^s.
% \]

% For $k \gtrsim \frac{1}{1-\rho}$, applying $\lambda_j\eqsim j^{-\beta}$, the integral comparison yields the strict upper bound:
% \begin{equation*}
%     \mathcal D_{\mathrm{over}}(k)
%     \le
%     C\sum_{j > J_{\max}}j^{-1-\beta s}
%     \exp\left(-2c_1\frac{\eta k}{1-\rho}j^{-\beta}\right)
%     \le
%     C\left(\frac{1-\rho}{\eta k}\right)^s .
% \end{equation*}
% By restricting the sum to $j \ge (C\frac{\eta k}{1-\rho})^{1/\beta}$, we obtain the matching lower bound $c\left(\frac{1-\rho}{\eta k}\right)^s$.

% \textbf{2. Underdamped Leading Block ($\cS_{\mathrm{under}}$):} \\
% For $j \le J_{\max}$, write the two roots as
% \[
%     \sqrt{\rho(1-\eta\lambda_j)}e^{i\omega_j},
%     \quad
%     \sqrt{\rho(1-\eta\lambda_j)}e^{-i\omega_j},
%     \qquad \text{where} \quad
%     \cos \omega_j=\frac{(1+\rho)\sqrt{1-\eta\lambda_j}}{2\sqrt \rho}.
% \]
% Solving the initial conditions gives the exact identity
% \begin{equation*}
%     h_k^{\nesterov}(j)
%     =\bigl(\rho(1-\eta\lambda_j)\bigr)^{k/2}
%     \left[
%         \cos(k\omega_j)
%         +\frac{(1-\rho)\sqrt{1-\eta\lambda_j}}{2\sqrt{\rho}\sin \omega_j}
%         \sin(k\omega_j)
%     \right].
% \end{equation*}
% Set $A_j:=\frac{(1-\rho)\sqrt{1-\eta\lambda_j}}{2\sqrt{\rho}\sin \omega_j}$ and define $\tan\psi_j = A_j$. The squared multiplier is bounded by the envelope $\bigl(\rho(1-\eta\lambda_j)\bigr)^k (1+A_j^2) \cos^2(k\omega_j - \psi_j)$. 

% To evaluate the underdamped contribution, we define the positive spectral weights $w_j(k) = \lambda_j (\theta_j^*)^2 \bigl(\rho(1-\eta\lambda_j)\bigr)^k (1+A_j^2)$ and decompose the cosine-squared terms:
% \[
%     \cS_{\mathrm{under}}(k) = \frac{1}{2} \sum_{j \le J_{\max}} w_j(k) + \frac{1}{2} \sum_{j \le J_{\max}} w_j(k) \cos(2k\omega_j - 2\psi_j).
% \]
% The first sum represents the stable static mass $M_w(k) = \frac{1}{2} \sum_{j \le J_{\max}} w_j(k)$. To analyze the oscillatory second sum, we view the indices as a continuous variable $x \in [1, J_{\max}]$. For the continuous oscillatory integral $I(k)$, we apply the substitution $u = \omega(x)$ mapping $[1, J_{\max}]$ to $[u_{\min}, u_{\max}]$:
% \[
%     I(k) = \int_{u_{\min}}^{u_{\max}} g(u, k) \cos(2ku - \phi(u)) \, du, \qquad \text{where } g(u, k) = w(x(u), k) \left| x'(u) \right|, \phi(u) = 2\psi(x(u)).
% \]
% By integration by parts with the differential $d\left(\sin(2ku - \phi(u))\right) = (2k - \phi'(u)) \cos(2ku - \phi(u)) \, du$:
% \[
%     I(k) = \left[ v_k(u) \sin(2ku - \phi(u)) \right]_{u_{\min}}^{u_{\max}} - \int_{u_{\min}}^{u_{\max}} v_k'(u) \sin(2ku - \phi(u)) \, du,
% \]
% where $v_k(u) = \frac{g(u, k)}{2k - \phi'(u)}$. This establishes the temporal decay $|I(k)| \le \frac{C_1}{k} \max_u |g(u, k)| = \mathcal{O}(1/k)$. Thus, for sufficiently large $k$, the oscillatory integral is strictly subsumed by the static mass: $|I(k)| \le \frac{1}{4} M_w(k)$. This guarantees the discrete sum is strictly bounded away from zero by the static part.

% We now rigorously evaluate $M_w(k)$. Asymptotically, $1+A_j^2 \eqsim 1$ everywhere except for a subdominant logarithmic singularity strictly at the continuous edge, which does not alter the polynomial/exponential macroscopic order. Thus:
% \[
%     M_w(k) \eqsim \sum_{j=1}^{J_{\max}} j^{-1-\beta s} \bigl(\rho(1-\eta\lambda_j)\bigr)^k \eqsim \sum_{j=1}^{J_{\max}} j^{-1-\beta s} \gamma_{\nesterov}^k e^{-ckx_j},
% \]
% where $\gamma_{\nesterov} = \rho(1-u_c)$, $u_c \eqsim (1-\rho)^2$, and the distance to the edge is $x_j = \eta\lambda_j - u_c \approx \eta\beta J_{\max}^{-\beta-1}(J_{\max} - j)$. We proceed in four temporal stages:

% \textit{Stage 1: $k \lesssim \frac{1}{\eta}$.} \\
% The exponential decay $e^{-ckx_j} \approx 1$. The sum is  dominated by the leading modes ($j \sim 1$):
% \[
%     \cS_{\mathrm{under}}(k) \eqsim \rho^k \sum_{j=1}^{J_{\max}} j^{-1-\beta s} \eqsim \rho^k.
% \]

% \textit{Stage 2: $\frac{1}{\eta} \lesssim k \lesssim \frac{1}{(1-\rho)^2}$.} \\
% The dominant bulk peak $j_k \eqsim (\eta k)^{1/\beta} \ll J_{\max}$ is fully captured inside the interval. The integral approximation perfectly applies:
% \[
%     \cS_{\mathrm{under}}(k) \eqsim \int_1^{\infty} x^{-1-\beta s} \rho^k e^{-c k \eta x^{-\beta}} dx \eqsim \rho^k (\eta k)^{-s}.
% \]

% \textit{Stage 3 \& 4: $k \gtrsim \frac{1}{(1-\rho)^2}$.} \\
% The bulk peak $j_k$ passes the boundary $J_{\max}$, so the summation is strictly dominated by the discrete tail items $j \to J_{\max}$. Let $m = J_{\max} - j \ge 0$. The distance to the edge is $x_j \approx \eta\beta J_{\max}^{-\beta-1} m \eqsim (1-\rho)^{2+2/\beta}\eta^{-1/\beta} m$. 
% Using the discrete geometric series sum for the dominant tail terms:
% \begin{align*}
%     \cS_{\mathrm{under}}(k) 
%     &\eqsim \sum_{m=0}^{J_{\max}} J_{\max}^{-1-\beta s} \gamma_{\nesterov}^k \exp\left(-c k (1-\rho)^{2+2/\beta}\eta^{-1/\beta} m\right) \\
%     &\eqsim J_{\max}^{-1-\beta s} \gamma_{\nesterov}^k \sum_{m=0}^{\infty} \left( \exp\left(-c k (1-\rho)^{2+2/\beta}\eta^{-1/\beta}\right) \right)^m \\
%     &\eqsim \left(\frac{(1-\rho)^2}{\eta}\right)^{s+\frac{1}{\beta}} \gamma_{\nesterov}^k \frac{1}{1 - \exp\left(-c k (1-\rho)^{2+2/\beta}\eta^{-1/\beta}\right)}.
% \end{align*}

% \textit{Stage 3: $\frac{1}{(1-\rho)^2} \lesssim k \lesssim \eta^{1/\beta}(1-\rho)^{-2-\frac{2}{\beta}}$.} \\
% In this stage, the exponent is small, $k (1-\rho)^{2+2/\beta}\eta^{-1/\beta} \lesssim 1$, meaning the decay window covers multiple discrete points near the edge. We apply the Taylor expansion $1 - e^{-x} \eqsim x$:
% \begin{align*}
%     \cS_{\mathrm{under}}(k) 
%     &\eqsim \left(\frac{(1-\rho)^2}{\eta}\right)^{s+\frac{1}{\beta}} \gamma_{\nesterov}^k \frac{1}{k (1-\rho)^{2+2/\beta}\eta^{-1/\beta}} \\
%     &= \left(\frac{(1-\rho)^2}{\eta}\right)^{s+\frac{1}{\beta}} \gamma_{\nesterov}^k \frac{1}{k(1-\rho)^2} \left(\frac{\eta}{(1-\rho)^2}\right)^{\frac{1}{\beta}} \\
%     &= \frac{1}{\eta k}\left(\frac{(1-\rho)^2}{\eta}\right)^{s-1}\gamma_{\nesterov}^{\,k}.
% \end{align*}

% \textit{Stage 4: $k \gtrsim \eta^{1/\beta}(1-\rho)^{-2-\frac{2}{\beta}}$.} \\
% In this extreme long-time regime, the exponent is large, $k (1-\rho)^{2+2/\beta}\eta^{-1/\beta} \gtrsim 1$. The exponential decay is so sharp that the single point $j = J_{\max}$ strictly dominates the entire sum. The denominator satisfies $1 - \exp\left(-c k (1-\rho)^{2+2/\beta}\eta^{-1/\beta}\right) \eqsim 1$:
% \[
%     \cS_{\mathrm{under}}(k) \eqsim \left(\frac{(1-\rho)^2}{\eta}\right)^{s+\frac{1}{\beta}}\gamma_{\nesterov}^{\,k} \cdot 1 = \left(\frac{(1-\rho)^2}{\eta}\right)^{s+\frac{1}{\beta}}\gamma_{\nesterov}^{\,k}.
% \]
% This concludes the proof.
% \end{proof}


% \begin{lemma}[Convolution Kernel, Mixed Regime]
% \label{lem:nes_kernel}
% Under Assumptions~\ref{ass:diagonal-covariance} and~\ref{ass:power-law}, assume the stability condition $\eta\lesssim \min\{1,B(1-\rho)\}$. For the mixed regime \(\eta\gtrsim(1-\rho)^2\), write \(\mathcal{K}_{\nesterov}(k)=\cK_{\mathrm{over}}(k)+\cK_{\mathrm{under}}(k)\), we have:
% \begin{equation*}
%     \mathcal{K}_{\mathrm{over}}(k)
%     \eqsim
%     \begin{cases}
%     k^2 \left(\frac{(1-\rho)^2}{\eta}\right)^{2-\frac1\beta},&\qquad k\lesssim \frac{1}{1-\rho},\\
%     \frac{1}{(1-\rho)^2}\left(\frac{1-\rho}{\eta k}\right)^{2-\frac1\beta},&\qquad k\gtrsim\frac{1}{1-\rho}.
%     \end{cases}
% \end{equation*}
% \begin{equation*}
%     \mathcal{K}_{\mathrm{under}}(k)
%     \eqsim
%     \begin{cases}
%     k^2\rho^k,&\qquad k\lesssim\frac{1}{\sqrt{\eta}},\\
%     \frac{\rho^k}{\eta},&\qquad \frac{1}{\sqrt{\eta}} \lesssim k \lesssim \frac{1}{\eta},\\
%     \frac{\rho^k}{\eta^{2-\frac1\beta}k^{1-\frac1\beta}},&\qquad \frac{1}{\eta} \lesssim k \lesssim \frac{1}{(1-\rho)^2},\\
%     \frac{1}{\eta^2 k}\left(\frac{(1-\rho)^2}{\eta}\right)^{-1/\beta}\gamma_{\nesterov}^{\,k}, &\qquad \frac{1}{(1-\rho)^2} \lesssim k \lesssim \eta^{1/\beta}(1-\rho)^{-2-\frac{2}{\beta}},\\
%     \frac{(1-\rho)^2}{\eta^2} \gamma_{\nesterov}^{\,k}, &\qquad k \gtrsim \eta^{1/\beta}(1-\rho)^{-2-\frac{2}{\beta}}.
%     \end{cases}
% \end{equation*}
% \end{lemma}

% \begin{proof}
% Recall that the convolution kernel is defined as
% \[
%     \mathcal K_{\nesterov}(k)=\sum_{j\ge1}\lambda_j^2G_k^{\nesterov}(j)^2,
%     \qquad
%     G_k^{\nesterov}(j)=\frac{y_{j,+}^{k+1}-y_{j,-}^{k+1}}{y_{j,+}-y_{j,-}}.
% \]
% We split the sum into the overdamped tail $\mathcal{K}_{\mathrm{over}}(k)$ where $\eta\lambda_j\lesssim(1-\rho)^2$, and the underdamped leading block $\mathcal{K}_{\mathrm{under}}(k)$ where $\eta\lambda_j\gtrsim(1-\rho)^2$. The crossover index is $J_{\max} \eqsim \left(\frac{\eta}{(1-\rho)^2}\right)^{1/\beta}$.

% \textbf{1. Overdamped Contribution $\mathcal{K}_{\mathrm{over}}(k)$:} \\
% We divide the evaluation into two temporal stages, mirroring the proof method of Lemma~\ref{lem:nes_kernel0}.

% \textit{Stage 1: $k \lesssim \frac{1}{1-\rho}$.} \\
% For the overdamped coordinates with $\eta\lambda_j\lesssim (1-\rho)^2$, the roots $y_{j,+}, y_{j,-}$ are real and close to 1. In this short-time regime, the exponential decay is negligible, and all terms in the algebraic expansion of $G_k^{\nesterov}(j)$ are strictly positive, yielding $G_k^{\nesterov}(j) \eqsim k$. Thus, the overdamped part is evaluated by summing over the tail $j \gtrsim J_{\max}$:
% \[
%     \mathcal{K}_{\mathrm{over}}(k) \eqsim \sum_{\eta\lambda_j\lesssim (1-\rho)^2} \lambda_j^2 k^2 = k^2 \sum_{j \gtrsim J_{\max}} j^{-2\beta} \eqsim k^2 (J_{\max})^{1-2\beta}.
% \]
% Substituting $J_{\max} \eqsim \left(\frac{\eta}{(1-\rho)^2}\right)^{1/\beta}$, we obtain $(J_{\max})^{1-2\beta} \eqsim \left(\frac{\eta}{(1-\rho)^2}\right)^{\frac{1-2\beta}{\beta}} = \left(\frac{(1-\rho)^2}{\eta}\right)^{2-\frac1\beta}$. Therefore:
% \[
%     \mathcal{K}_{\mathrm{over}}(k) \eqsim k^2 \left(\frac{(1-\rho)^2}{\eta}\right)^{2-\frac1\beta}.
% \]

% \textit{Stage 2: $k \gtrsim \frac{1}{1-\rho}$.} \\
% For $k \gtrsim \frac{1}{1-\rho}$, the root comparison gives
% \[
%     |G_k^{\nesterov}(j)|
%     \eqsim
%     \frac{1}{1-\rho}
%     \exp\left(-c\frac{\eta\lambda_j}{1-\rho}k\right),
% \]
% with constants independent of $j,\eta,\rho,k$. Hence the overdamped part is comparable to
% \[
%     \mathcal{K}_{\mathrm{over}}(k)
%     \eqsim
%     \frac{1}{(1-\rho)^2}
%     \sum_{\eta\lambda_j\lesssim (1-\rho)^2}
%     \lambda_j^2
%     \exp\left(-c\frac{\eta\lambda_j}{1-\rho}k\right).
% \]
% Using $\lambda_j\eqsim j^{-\beta}$, the weighted spectral counting rule is $\sum_j \lambda_j^2 F(\lambda_j) \eqsim \int u^{1-\frac1\beta}F(u)\,du$. Therefore the overdamped sum is bounded above and below by constant multiples of
% \[
%     \frac{1}{(1-\rho)^2}
%     \int_0^{(1-\rho)^2/\eta}
%     u^{1-\frac1\beta}
%     \exp\left(-c\frac{\eta k}{1-\rho}u\right)\,du .
% \]
% Since $k \gtrsim \frac{1}{1-\rho}$, the exponential decay activates well within the integration domain. The integral evaluates to $\mathcal{O}\left( \left(\frac{1-\rho}{\eta k}\right)^{2-\frac1\beta} \right)$, giving exactly
% \begin{equation*}
%     \mathcal K_{\mathrm{over}}(k)
%     \eqsim
%     \frac{1}{(1-\rho)^2}
%     \left(\frac{1-\rho}{\eta k}\right)^{2-\frac1\beta}.
% \end{equation*}

% \textbf{2. Underdamped Contribution $\mathcal{K}_{\mathrm{under}}(k)$:} \\
% For $j \le J_{\max}$, write
% \[
%     y_{j,+}=\sqrt{\rho(1-\eta\lambda_j)} e^{i\omega_j},
%     \quad
%     y_{j,-}=\sqrt{\rho(1-\eta\lambda_j)} e^{-i\omega_j},
%     \qquad \text{where} \quad
%     \cos \omega_j=\frac{(1+\rho)\sqrt{1-\eta\lambda_j}}{2\sqrt \rho}.
% \]
% Then
% \begin{equation*}
%     G_k^{\nesterov}(j)=\bigl(\rho(1-\eta\lambda_j)\bigr)^{k/2}\frac{\sin((k+1)\omega_j)}{\sin \omega_j}.
% \end{equation*}
% On the separated underdamped region $\eta\lambda_j\gtrsim(1-\rho)^2$, stability gives $\sin^2\omega_j \eqsim \eta\lambda_j$. 

% To evaluate the underdamped contribution, we define the weights $\tilde{w}_j(k) = \lambda_j^2 \frac{\bigl(\rho(1-\eta\lambda_j)\bigr)^k}{\sin^2\omega_j}$ and decompose the sine-squared terms:
% \[
%     \mathcal{K}_{\mathrm{under}}(k) = \sum_{j \le J_{\max}} \tilde{w}_j(k) \sin^2((k+1)\omega_j) = \frac{1}{2} \sum_{j \le J_{\max}} \tilde{w}_j(k) - \frac{1}{2} \sum_{j \le J_{\max}} \tilde{w}_j(k) \cos(2(k+1)\omega_j).
% \]
% The first sum represents the stable static mass $M_{\tilde{w}}(k) = \frac{1}{2} \sum_{j \le J_{\max}} \tilde{w}_j(k)$. For the oscillatory second sum, we view the indices as a continuous variable $x \in [1, J_{\max}]$. Defining the continuous oscillatory integral $I(k) = \int_{1}^{J_{\max}} \tilde{w}(x, k) \cos(2(k+1)\omega(x)) \, dx$ and applying the substitution $u = \omega(x)$, we have:
% \[
%     I(k) = \int_{u_{\min}}^{u_{\max}} h(u, k) \cos(2(k+1)u) \, du, \qquad \text{where } h(u, k) = \tilde{w}(x(u), k) \left| x'(u) \right|.
% \]
% By integration by parts with the differential $d\left(\sin(2(k+1)u)\right) = 2(k+1) \cos(2(k+1)u) \, du$:
% \[
%     I(k) = \left[ \frac{h(u, k)}{2(k+1)} \sin(2(k+1)u) \right]_{u_{\min}}^{u_{\max}} - \int_{u_{\min}}^{u_{\max}} \frac{h'(u, k)}{2(k+1)} \sin(2(k+1)u) \, du.
% \]
% This establishes the temporal decay $|I(k)| \le \frac{C_1}{k} \max_u |h(u, k)| = \mathcal{O}(1/k)$. Thus, for $k \gtrsim \frac{1}{\sqrt{\eta}}$, the oscillatory integral is strictly subsumed by the static mass: $|I(k)| \le \frac{1}{4} M_{\tilde{w}}(k)$. This guarantees the discrete sum is strictly bounded away from zero by the static part.

% We now rigorously evaluate the static mass $M_{\tilde{w}}(k)$. Using the global approximation $\sin^2\omega_j \eqsim \eta\lambda_j$, the summand simplifies to:
% \[
%     M_{\tilde{w}}(k) \eqsim \sum_{j=1}^{J_{\max}} \frac{\lambda_j^2}{\eta\lambda_j} \gamma_j^k = \frac{1}{\eta} \sum_{j=1}^{J_{\max}} \lambda_j \gamma_j^k,
% \]
% where $\gamma_j = \rho(1-\eta\lambda_j) \approx \gamma_{\nesterov} e^{-c x_j}$, $\gamma_{\nesterov} = \rho(1-u_c)$, $u_c \eqsim (1-\rho)^2$, and the distance to the edge is $x_j = \eta\lambda_j - u_c \approx \eta\beta J_{\max}^{-\beta-1}(J_{\max} - j)$. We proceed in five temporal stages:

% \textit{Stage 1: $k \lesssim \frac{1}{\sqrt{\eta}}$.} \\
% In this short-time regime, the angle $(k+1)\omega_j \lesssim 1$, so the small-angle approximation applies: $\sin((k+1)\omega_j) \eqsim k\sin\omega_j$.
% Substituting this directly into the exact sum gives:
% \[
%     \mathcal K_{\mathrm{under}}(k) \eqsim \sum_{j \le J_{\max}} \lambda_j^2 k^2 \rho^k = k^2 \rho^k \sum_{j \le J_{\max}} \lambda_j^2 \eqsim k^2 \rho^k.
% \]

% \textit{Stage 2: $\frac{1}{\sqrt{\eta}} \lesssim k \lesssim \frac{1}{\eta}$.} \\
% The oscillation averages out, leaving $M_{\tilde{w}}(k)$. Since $k \eta \lesssim 1$, the exponential decay $e^{-ckx_j} \approx 1$ throughout the bulk spectrum $j \ge 1$. Since $\beta > 1$, the sum $\sum \lambda_j$ converges to $\mathcal{O}(1)$:
% \[
%     \mathcal K_{\mathrm{under}}(k) \eqsim \frac{1}{\eta} \sum_{j=1}^{J_{\max}} \lambda_j \rho^k \eqsim \frac{\rho^k}{\eta}.
% \]

% \textit{Stage 3: $\frac{1}{\eta} \lesssim k \lesssim \frac{1}{(1-\rho)^2}$.} \\
% In this intermediate regime, the exponential factor $e^{-k\eta \lambda_j}$ suppresses the leading modes, while the bulk peak at $j_k \eqsim (\eta k)^{1/\beta}$ remains strictly inside $[1, J_{\max}]$. Evaluating the continuous integral over $v = k u = k \eta x^{-\beta}$:
% \[
%     \mathcal K_{\mathrm{under}}(k) \eqsim \frac{\rho^k}{\beta \eta^{2-\frac{1}{\beta}} k^{1-\frac{1}{\beta}}} \int_{k u_c}^{k\eta} v^{-\frac{1}{\beta}} e^{-v} dv \eqsim \frac{\rho^k}{\eta^{2-\frac1\beta} k^{1-\frac1\beta}}.
% \]

% \textit{Stage 4 \& 5: $k \gtrsim \frac{1}{(1-\rho)^2}$.} \\
% The bulk peak passes the boundary $J_{\max}$, and the summation is strictly dominated by the discrete tail items $j \to J_{\max}$. Let $m = J_{\max} - j \ge 0$. The distance to the edge is $x_j \approx m \eta\beta J_{\max}^{-\beta-1} \eqsim m (1-\rho)^{2+2/\beta}\eta^{-1/\beta}$. 
% At the boundary, $\lambda_j \approx \lambda_{J_{\max}} \eqsim \frac{(1-\rho)^2}{\eta}$. Using the discrete geometric series sum for the dominant tail terms:
% \begin{align*}
%     \mathcal K_{\mathrm{under}}(k) 
%     &\eqsim \frac{1}{\eta} \sum_{m=0}^{J_{\max}} \frac{(1-\rho)^2}{\eta} \gamma_{\nesterov}^k \exp\left(-c k (1-\rho)^{2+2/\beta}\eta^{-1/\beta} m\right) \\
%     &\eqsim \frac{(1-\rho)^2}{\eta^2} \gamma_{\nesterov}^k \sum_{m=0}^{\infty} \left( \exp\left(-c k (1-\rho)^{2+2/\beta}\eta^{-1/\beta}\right) \right)^m \\
%     &\eqsim \frac{(1-\rho)^2}{\eta^2} \gamma_{\nesterov}^k \frac{1}{1 - \exp\left(-c k (1-\rho)^{2+2/\beta}\eta^{-1/\beta}\right)}.
% \end{align*}

% \textit{Stage 4: $\frac{1}{(1-\rho)^2} \lesssim k \lesssim \eta^{1/\beta}(1-\rho)^{-2-\frac{2}{\beta}}$.} \\
% In this stage, the exponent is small, $k (1-\rho)^{2+2/\beta}\eta^{-1/\beta} \lesssim 1$, meaning the decay window covers multiple discrete points. We apply the Taylor expansion $1 - e^{-x} \eqsim x$:
% \begin{align*}
%     \mathcal K_{\mathrm{under}}(k) 
%     &\eqsim \frac{(1-\rho)^2}{\eta^2} \gamma_{\nesterov}^k \frac{1}{k (1-\rho)^{2+2/\beta}\eta^{-1/\beta}} \\
%     &= \frac{1}{\eta^2 k} \frac{(1-\rho)^2}{(1-\rho)^{2+2/\beta}} \eta^{1/\beta} \gamma_{\nesterov}^k \\
%     &= \frac{1}{\eta^2 k} (1-\rho)^{-2/\beta} \eta^{1/\beta} \gamma_{\nesterov}^k \\
%     &= \frac{1}{\eta^2 k}\left(\frac{(1-\rho)^2}{\eta}\right)^{-1/\beta}\gamma_{\nesterov}^{\,k}.
% \end{align*}

% \textit{Stage 5: $k \gtrsim \eta^{1/\beta}(1-\rho)^{-2-\frac{2}{\beta}}$.} \\
% In this extreme long-time regime, $k (1-\rho)^{2+2/\beta}\eta^{-1/\beta} \gtrsim 1$. The exponential decay is so sharp that the single point $j = J_{\max}$ strictly dominates the entire sum, giving $1 - \exp\left(-c k (1-\rho)^{2+2/\beta}\eta^{-1/\beta}\right) \eqsim 1$:
% \[
%     \mathcal K_{\mathrm{under}}(k) \eqsim \frac{(1-\rho)^2}{\eta^2} \gamma_{\nesterov}^k.
% \]
% This concludes the proof.
% \end{proof}

% \begin{lemma}[Properties of $\rho^m$ and $\gamma_{\nesterov}^m$]
% \label{lem:nes_rho_gamma}
% Let $\gamma_{\nesterov} = \rho\left(1-\frac{(1-\rho)^2}{(1+\rho)^2}\right)$. In the mixed regime $\eta \gtrsim (1-\rho)^2$, for any time index $m$ satisfying $m \lesssim \frac{1}{(1-\rho)^2}$, we have 
% \[
% \rho^m \eqsim \gamma_{\nesterov}^m.
% \]
% Moreover, for $m \lesssim \frac{1}{1-\rho}$, we have $\rho^m \eqsim \gamma_{\nesterov}^m \eqsim 1$. Consequently, since $\frac{1}{\sqrt{\eta}} \lesssim \frac{1}{1-\rho}$, we obtain $\rho^{\frac{1}{\sqrt{\eta}}} \eqsim \gamma_{\nesterov}^{\frac{1}{\sqrt{\eta}}} \eqsim 1$.
% \end{lemma}
% \begin{proof}
% By definition, we can explicitly expand $\gamma_{\nesterov}^m$:


% \begin{equation*}
%     \gamma_{\nesterov}^m = \rho^m \left(1-\frac{(1-\rho)^2}{(1+\rho)^2}\right)^m.
% \end{equation*}
% For $m \lesssim \frac{1}{(1-\rho)^2}$, the exponent associated with the second factor is bounded, meaning $m \frac{(1-\rho)^2}{(1+\rho)^2} \lesssim 1$. Using the asymptotic approximation $(1-x)^m \eqsim e^{-mx}$ for $x \in (0,1)$, we obtain $\left(1-\frac{(1-\rho)^2}{(1+\rho)^2}\right)^m \eqsim \exp\left(-m \frac{(1-\rho)^2}{(1+\rho)^2}\right) \eqsim 1$, which strictly implies $\rho^m \eqsim \gamma_{\nesterov}^m$.

% Furthermore, when the condition tightens to $m \lesssim \frac{1}{1-\rho}$, we have:
% \begin{equation*}
%     \rho^m = (1-(1-\rho))^m \eqsim e^{-m(1-\rho)} \eqsim 1.
% \end{equation*}

% Thus, $\rho^m \eqsim \gamma_{\nesterov}^m \eqsim 1$. Since $\eta \gtrsim (1-\rho)^2$, the relation $\frac{1}{\sqrt{\eta}} \lesssim \frac{1}{1-\rho}$ naturally holds, completing the proof for $\rho^{\frac{1}{\sqrt{\eta}}} \eqsim \gamma_{\nesterov}^{\frac{1}{\sqrt{\eta}}} \eqsim 1$.

% \end{proof}


% \begin{lemma}[Convolution $\cS_{\mathrm{over}} * \mathcal{K}_{\mathrm{over}}$]
% \label{lem:nes_conv_over_over}
% Under the assumptions of Proposition~\ref{prop:nes_deterministic_bias} and Lemma~\ref{lem:nes_kernel}, and under the stability condition $\eta\lesssim B(1-\rho)$, for $k \gtrsim \eta^{1/\beta}(1-\rho)^{-2-\frac{2}{\beta}}$, the convolution of the overdamped deterministic bias with the overdamped kernel satisfies:
% \begin{equation*}
%     \sum_{m=0}^{k-1} \cS_{\mathrm{over}}(m) \mathcal{K}_{\mathrm{over}}(k-1-m) \eqsim \frac{1}{(1-\rho)^3} \left(\frac{(1-\rho)^2}{\eta}\right)^{s \vee \left(2-\frac{1}{\beta}\right)} \left(\frac{1-\rho}{\eta k}\right)^{s \wedge \left(2-\frac{1}{\beta}\right)}.
% \end{equation*}
% \end{lemma}
% \begin{proof}
% We rigorously evaluate the convolution by dividing the discrete integral into three distinct temporal stages.

% \textbf{Stage 1:} $0 \le m \lesssim \frac{1}{1-\rho}$. Here $\cS_{\mathrm{over}}(m) \eqsim \left(\frac{(1-\rho)^2}{\eta}\right)^s$ and $\mathcal{K}_{\mathrm{over}}(k-1-m) \eqsim \frac{1}{(1-\rho)^2} \left(\frac{1-\rho}{\eta k}\right)^{2-\frac{1}{\beta}}$.
% The summation evaluates to:
% \begin{align*}
%     &\sum_{m=0}^{\frac{1}{1-\rho}} \left(\frac{(1-\rho)^2}{\eta}\right)^s \frac{1}{(1-\rho)^2} \left(\frac{1-\rho}{\eta k}\right)^{2-\frac{1}{\beta}} \\
%     &\eqsim \frac{1}{1-\rho} \left(\frac{(1-\rho)^2}{\eta}\right)^s \frac{1}{(1-\rho)^2} \left(\frac{1-\rho}{\eta k}\right)^{2-\frac{1}{\beta}} \\
%     &= \frac{1}{(1-\rho)^3} \left(\frac{(1-\rho)^2}{\eta}\right)^s \left(\frac{1-\rho}{\eta k}\right)^{2-\frac{1}{\beta}}.
% \end{align*}

% \textbf{Stage 2:} $\frac{1}{1-\rho} \lesssim m \le k - \frac{1}{1-\rho}$. We substitute the bulk forms $\cS_{\mathrm{over}}(m) \eqsim \left(\frac{1-\rho}{\eta m}\right)^s$ and $\mathcal{K}_{\mathrm{over}}(k-1-m) \eqsim \frac{1}{(1-\rho)^2} \left(\frac{1-\rho}{\eta (k-m)}\right)^{2-\frac{1}{\beta}}$, and explicitly split the sum at $m = k/2$:
% \begin{align*}
%     &\sum_{m=\frac{1}{1-\rho}}^{k-\frac{1}{1-\rho}} \left(\frac{1-\rho}{\eta m}\right)^s \frac{1}{(1-\rho)^2} \left(\frac{1-\rho}{\eta (k-m)}\right)^{2-\frac{1}{\beta}} \\
%     &\eqsim \frac{1}{(1-\rho)^2} \sum_{m=\frac{1}{1-\rho}}^{\frac{k}{2}} \left(\frac{1-\rho}{\eta m}\right)^s \left(\frac{1-\rho}{\eta k}\right)^{2-\frac{1}{\beta}} + \frac{1}{(1-\rho)^2} \sum_{m=\frac{k}{2}}^{k-\frac{1}{1-\rho}} \left(\frac{1-\rho}{\eta k}\right)^s \left(\frac{1-\rho}{\eta (k-m)}\right)^{2-\frac{1}{\beta}} \\
%     &= \frac{1}{(1-\rho)^2} \left( \frac{1-\rho}{\eta} \right)^{s+2-\frac{1}{\beta}} \left[ \frac{1}{k^{2-\frac{1}{\beta}}} \sum_{m=\frac{1}{1-\rho}}^{\frac{k}{2}} m^{-s} + \frac{1}{k^s} \sum_{m=\frac{k}{2}}^{k-\frac{1}{1-\rho}} (k-m)^{-\left(2-\frac{1}{\beta}\right)} \right].
% \end{align*}
% Evaluating the polynomial sums over $m$: for $\sum m^{-s}$, it yields $k^{1-s}$ if $s < 1$, and is strictly dominated by its lower limit $(\frac{1}{1-\rho})^{1-s}$ if $s > 1$. For $\sum (k-m)^{-(2-\frac{1}{\beta})}$, since $\beta > 1 \implies 2-\frac{1}{\beta} > 1$, it is rigorously bounded by its lower limit $(\frac{1}{1-\rho})^{1-(2-\frac{1}{\beta})}$. Guided by the polynomial convolution logic, we continue the derivation:
% \begin{align*}
%     &\eqsim \frac{1}{(1-\rho)^2} \left( \frac{1-\rho}{\eta} \right)^{s+2-\frac{1}{\beta}} \left[ \begin{cases} k^{1-s} \cdot k^{-\left(2-\frac{1}{\beta}\right)} & \text{if } s < 1 \\ \left(\frac{1}{1-\rho}\right)^{1-s} \cdot k^{-\left(2-\frac{1}{\beta}\right)} & \text{if } s > 1 \end{cases} + \left(\frac{1}{1-\rho}\right)^{1-\left(2-\frac{1}{\beta}\right)} \cdot k^{-s} \right] \\
%     &= \frac{1}{(1-\rho)^3} \left( \frac{1-\rho}{\eta} \right)^{s+2-\frac{1}{\beta}} (1-\rho) \left[ \begin{cases} k^{1-s} \cdot k^{-\left(2-\frac{1}{\beta}\right)} & \text{if } s < 1 \\ (1-\rho)^{s-1} \cdot k^{-\left(2-\frac{1}{\beta}\right)} & \text{if } s > 1 \end{cases} + (1-\rho)^{1-\frac{1}{\beta}} \cdot k^{-s} \right] \\
%     &= \frac{1}{(1-\rho)^3} \left( \frac{1-\rho}{\eta} \right)^{s+2-\frac{1}{\beta}} \left[ \begin{cases} k^{1-s} (1-\rho) \cdot k^{-\left(2-\frac{1}{\beta}\right)} & \text{if } s < 1 \\ (1-\rho)^s \cdot k^{-\left(2-\frac{1}{\beta}\right)} & \text{if } s > 1 \end{cases} + (1-\rho)^{2-\frac{1}{\beta}} \cdot k^{-s} \right].
% \end{align*}
% Since $k \gtrsim \frac{1}{1-\rho}$, the asymptotically larger term dominates. For $s < 2-\frac{1}{\beta}$, the term $(1-\rho)^{2-\frac{1}{\beta}} k^{-s}$ is strictly larger. For $s > 2-\frac{1}{\beta}$, the term $(1-\rho)^s k^{-(2-\frac{1}{\beta})}$ is strictly larger. We can elegantly unify these regimes using the $\wedge$ (min) and $\vee$ (max) operators:
% \begin{align*}
%     &\eqsim \frac{1}{(1-\rho)^3} \left( \frac{1-\rho}{\eta} \right)^{s+2-\frac{1}{\beta}} (1-\rho)^{s+2-\frac{1}{\beta}} \left( \frac{1}{(1-\rho)k} \right)^{s \wedge \left(2-\frac{1}{\beta}\right)} \\
%     &= \frac{1}{(1-\rho)^3} \left( \frac{(1-\rho)^2}{\eta} \right)^{s+2-\frac{1}{\beta}} \left( \frac{1}{(1-\rho)k} \right)^{s \wedge \left(2-\frac{1}{\beta}\right)} \\
%     &= \frac{1}{(1-\rho)^3} \left( \frac{(1-\rho)^2}{\eta} \right)^{s \vee \left(2-\frac{1}{\beta}\right)} \left(\frac{(1-\rho)^2}{\eta}\right)^{s \wedge \left(2-\frac{1}{\beta}\right)} \left( \frac{1}{(1-\rho)k} \right)^{s \wedge \left(2-\frac{1}{\beta}\right)} \\
%     &= \frac{1}{(1-\rho)^3} \left( \frac{(1-\rho)^2}{\eta} \right)^{s \vee \left(2-\frac{1}{\beta}\right)} \left( \frac{1-\rho}{\eta k} \right)^{s \wedge \left(2-\frac{1}{\beta}\right)}.
% \end{align*}

% \textbf{Stage 3:} $k - \frac{1}{1-\rho} \lesssim m \le k-1$. Here $\cS_{\mathrm{over}}(m) \eqsim \left(\frac{1-\rho}{\eta k}\right)^s$ and $\mathcal{K}_{\mathrm{over}}(k-1-m) \eqsim (k-m)^2 \left(\frac{(1-\rho)^2}{\eta}\right)^{2-\frac{1}{\beta}}$.
% \begin{align*}
%     \sum_{m=k-\frac{1}{1-\rho}}^{k-1} \left(\frac{1-\rho}{\eta k}\right)^s (k-m)^2 \left(\frac{(1-\rho)^2}{\eta}\right)^{2-\frac{1}{\beta}} &\eqsim \left(\frac{1-\rho}{\eta k}\right)^s \left(\frac{(1-\rho)^2}{\eta}\right)^{2-\frac{1}{\beta}} \sum_{n=1}^{\frac{1}{1-\rho}} n^2 \\
%     &\eqsim \frac{1}{(1-\rho)^3} \left(\frac{1-\rho}{\eta k}\right)^s \left(\frac{(1-\rho)^2}{\eta}\right)^{2-\frac{1}{\beta}}.
% \end{align*}

% Aggregating the three stages yields the total bound. Let us verify that the results from Stage 1 and Stage 3 are strictly absorbed into the dominant term from Stage 2.

% If $s \ge 2-\frac{1}{\beta}$, we have $T_2 = T_1$. For $T_3$, its ratio to $T_2$ is:
% \begin{equation*}
%     \frac{T_3}{T_2} = \frac{\left(\frac{1-\rho}{\eta k}\right)^s \left(\frac{(1-\rho)^2}{\eta}\right)^{2-\frac{1}{\beta}}}{\left(\frac{(1-\rho)^2}{\eta}\right)^s \left(\frac{1-\rho}{\eta k}\right)^{2-\frac{1}{\beta}}} = \left(\frac{\frac{1-\rho}{\eta k}}{\frac{(1-\rho)^2}{\eta}}\right)^{s-\left(2-\frac{1}{\beta}\right)} = \left(\frac{1}{(1-\rho)k}\right)^{s-\left(2-\frac{1}{\beta}\right)}.
% \end{equation*}
% Since $k \gtrsim \frac{1}{1-\rho}$, we have $\frac{1}{(1-\rho)k} \lesssim 1$. Because $s - (2-\frac{1}{\beta}) \ge 0$, it follows that $T_3 \lesssim T_2$.

% If $s < 2-\frac{1}{\beta}$, we have $T_2 = T_3$. For $T_1$, its ratio to $T_2$ is:
% \begin{equation*}
%     \frac{T_1}{T_2} = \frac{\left(\frac{(1-\rho)^2}{\eta}\right)^s \left(\frac{1-\rho}{\eta k}\right)^{2-\frac{1}{\beta}}}{\left(\frac{(1-\rho)^2}{\eta}\right)^{2-\frac{1}{\beta}} \left(\frac{1-\rho}{\eta k}\right)^s} = \left(\frac{\frac{1-\rho}{\eta k}}{\frac{(1-\rho)^2}{\eta}}\right)^{\left(2-\frac{1}{\beta}\right)-s} = \left(\frac{1}{(1-\rho)k}\right)^{\left(2-\frac{1}{\beta}\right)-s} \lesssim 1.
% \end{equation*}
% Therefore, in all cases, the terms from Stage 1 and Stage 3 are rigorously subsumed by the unified expression derived in Stage 2. This concludes the proof.
% \end{proof}

% \begin{lemma}[Convolution $\cS_{\mathrm{over}} * \mathcal{K}_{\mathrm{under}}$]
% \label{lem:nes_conv_over_under}
% Under the assumptions of Proposition~\ref{prop:nes_deterministic_bias} and Lemma~\ref{lem:nes_kernel}, and under the stability condition $\eta\lesssim B(1-\rho)$, for $k \gtrsim \eta^{1/\beta}(1-\rho)^{-2-\frac{2}{\beta}}$, we have:
% \begin{equation*}
%     \sum_{m=0}^{k-1} \cS_{\mathrm{over}}(m) \mathcal{K}_{\mathrm{under}}(k-1-m) \eqsim \frac{1}{\eta(1-\rho)} \left(\frac{(1-\rho)^2}{\eta}\right)^{s+1} \gamma_{\nesterov}^k + \frac{1}{\eta(1-\rho)} \left(\frac{1-\rho}{\eta k}\right)^s.
% \end{equation*}
% \end{lemma}
% \begin{proof}
% We evaluate this convolution by transforming the index to $n = k-1-m \eqsim k-m$, and rigorously dividing the summation into six distinct stages based on the piecewise definitions of $\cS_{\mathrm{over}}$ and $\mathcal{K}_{\mathrm{under}}$.

% \textbf{Stage 1:} $0 \le m \lesssim \frac{1}{1-\rho}$. 
% Here $\cS_{\mathrm{over}}(m) \eqsim \left(\frac{(1-\rho)^2}{\eta}\right)^s$. Since $k$ is extremely large, $n \eqsim k \gtrsim \eta^{1/\beta}(1-\rho)^{-2-\frac{2}{\beta}}$, giving $\mathcal{K}_{\mathrm{under}}(n) \eqsim \frac{(1-\rho)^2}{\eta^2} \gamma_{\nesterov}^n$. For $m \lesssim \frac{1}{1-\rho}$, we have $\gamma_{\nesterov}^{-m} \eqsim 1$. Thus, summing over $m$ yields:
% \begin{align*}
%     \sum_{m=0}^{\frac{1}{1-\rho}} \left(\frac{(1-\rho)^2}{\eta}\right)^s \frac{(1-\rho)^2}{\eta^2} \gamma_{\nesterov}^{k-m} 
%     &\eqsim \frac{1}{1-\rho} \left(\frac{(1-\rho)^2}{\eta}\right)^s \frac{(1-\rho)^2}{\eta^2} \gamma_{\nesterov}^k \\
%     &= \frac{1}{\eta(1-\rho)} \left(\frac{(1-\rho)^2}{\eta}\right)^{s+1} \gamma_{\nesterov}^k.
% \end{align*}

% \textbf{Stage 2:} $\frac{1}{1-\rho} \lesssim m \lesssim k - \eta^{1/\beta}(1-\rho)^{-2-\frac{2}{\beta}}$. 
% Here $\cS_{\mathrm{over}}(m) \eqsim \left(\frac{1-\rho}{\eta m}\right)^s$. The kernel index satisfies $n \gtrsim \eta^{1/\beta}(1-\rho)^{-2-\frac{2}{\beta}}$, so $\mathcal{K}_{\mathrm{under}}(n) \eqsim \frac{(1-\rho)^2}{\eta^2} \gamma_{\nesterov}^n$. The sum is dominated by the upper bound of $m$ (i.e., the lower bound of $n$) due to the exponential decay $\gamma_{\nesterov}^n$, with an effective decay width of $\frac{1}{1-\gamma_{\nesterov}} \eqsim \frac{1}{1-\rho}$. Let the boundary be $n_0 = \eta^{1/\beta}(1-\rho)^{-2-\frac{2}{\beta}}$, we have:
% \begin{align*}
%     \sum_{n=n_0}^{k - \frac{1}{1-\rho}} \left(\frac{1-\rho}{\eta (k-n)}\right)^s \frac{(1-\rho)^2}{\eta^2} \gamma_{\nesterov}^n 
%     &\eqsim \left(\frac{1-\rho}{\eta k}\right)^s \frac{(1-\rho)^2}{\eta^2} \sum_{n=n_0}^{k} \gamma_{\nesterov}^n \\
%     &\eqsim \frac{1}{\eta(1-\rho)} \frac{(1-\rho)^2}{\eta} \left(\frac{1-\rho}{\eta k}\right)^s \gamma_{\nesterov}^{\eta^{1/\beta}(1-\rho)^{-2-\frac{2}{\beta}}}.
% \end{align*}

% \textbf{Stage 3:} $k - \eta^{1/\beta}(1-\rho)^{-2-\frac{2}{\beta}} \lesssim m \lesssim k - \frac{1}{(1-\rho)^2}$. 
% Here $m \eqsim k$, so $\cS_{\mathrm{over}}(m) \eqsim \left(\frac{1-\rho}{\eta k}\right)^s$. The kernel index is $n \in \left[ \frac{1}{(1-\rho)^2}, \eta^{1/\beta}(1-\rho)^{-2-\frac{2}{\beta}} \right]$, where $\mathcal{K}_{\mathrm{under}}(n) \eqsim \frac{1}{\eta^2 n}\left(\frac{(1-\rho)^2}{\eta}\right)^{-1/\beta} \gamma_{\nesterov}^n$. The sum is again dominated by the lower bound of $n$, say $n_1 = \frac{1}{(1-\rho)^2}$:
% \begin{align*}
%     \sum_{n=n_1}^{\eta^{1/\beta}(1-\rho)^{-2-\frac{2}{\beta}}} \left(\frac{1-\rho}{\eta k}\right)^s \frac{1}{\eta^2 n}\left(\frac{(1-\rho)^2}{\eta}\right)^{-\frac{1}{\beta}} \gamma_{\nesterov}^n 
%     &\eqsim \left(\frac{1-\rho}{\eta k}\right)^s \frac{1}{\eta^2 n_1} \left(\frac{(1-\rho)^2}{\eta}\right)^{-\frac{1}{\beta}} \gamma_{\nesterov}^{n_1} \frac{1}{1-\rho} \\
%     &= \frac{1}{\eta(1-\rho)} \left(\frac{(1-\rho)^2}{\eta}\right)^{1-\frac{1}{\beta}} \left(\frac{1-\rho}{\eta k}\right)^s \gamma_{\nesterov}^{\frac{1}{(1-\rho)^2}}.
% \end{align*}

% \textbf{Stage 4:} $k - \frac{1}{(1-\rho)^2} \lesssim m \lesssim k - \frac{1}{\eta}$. 
% Here $\cS_{\mathrm{over}}(m) \eqsim \left(\frac{1-\rho}{\eta k}\right)^s$. The kernel index is $n \in \left[ \frac{1}{\eta}, \frac{1}{(1-\rho)^2} \right]$, where $\mathcal{K}_{\mathrm{under}}(n) \eqsim \frac{\rho^n}{\eta^{2-1/\beta} n^{1-1/\beta}}$. The sum is dominated by the lower limit $n_2 = \frac{1}{\eta}$:
% \begin{align*}
%     \sum_{n=\frac{1}{\eta}}^{\frac{1}{(1-\rho)^2}} \left(\frac{1-\rho}{\eta k}\right)^s \frac{\rho^n}{\eta^{2-\frac{1}{\beta}} n^{1-\frac{1}{\beta}}} 
%     &\eqsim \left(\frac{1-\rho}{\eta k}\right)^s \frac{1}{\eta^{2-\frac{1}{\beta}}} \frac{\rho^{\frac{1}{\eta}}}{(1/\eta)^{1-\frac{1}{\beta}}} \frac{1}{1-\rho} \\
%     &= \frac{1}{\eta(1-\rho)} \left(\frac{1-\rho}{\eta k}\right)^s \rho^{\frac{1}{\eta}}.
% \end{align*}

% \textbf{Stage 5:} $k - \frac{1}{\eta} \lesssim m \lesssim k - \frac{1}{\sqrt{\eta}}$. 
% Here $\cS_{\mathrm{over}}(m) \eqsim \left(\frac{1-\rho}{\eta k}\right)^s$. The kernel index is $n \in \left[ \frac{1}{\sqrt{\eta}}, \frac{1}{\eta} \right]$, where $\mathcal{K}_{\mathrm{under}}(n) \eqsim \frac{\rho^n}{\eta}$. Using the properties of $\rho^m$ from Lemma~\ref{lem:nes_rho_gamma}, the sum is dominated by its lower limit $n_3 = \frac{1}{\sqrt{\eta}}$, where $\rho^{1/\sqrt{\eta}} \eqsim 1$:
% \begin{align*}
%     \sum_{n=\frac{1}{\sqrt{\eta}}}^{\frac{1}{\eta}} \left(\frac{1-\rho}{\eta k}\right)^s \frac{\rho^n}{\eta} 
%     &\eqsim \left(\frac{1-\rho}{\eta k}\right)^s \frac{1}{\eta} \rho^{\frac{1}{\sqrt{\eta}}} \frac{1}{1-\rho} \\
%     &\eqsim \frac{1}{\eta(1-\rho)} \left(\frac{1-\rho}{\eta k}\right)^s.
% \end{align*}

% \textbf{Stage 6:} $k - \frac{1}{\sqrt{\eta}} \lesssim m \le k - 1$. 
% Here $\cS_{\mathrm{over}}(m) \eqsim \left(\frac{1-\rho}{\eta k}\right)^s$. The kernel index is $n \in \left[ 1, \frac{1}{\sqrt{\eta}} \right]$, where $\mathcal{K}_{\mathrm{under}}(n) \eqsim n^2 \rho^n$. Since $n \le \frac{1}{\sqrt{\eta}}$, $\rho^n \eqsim 1$. Thus, the polynomial sum evaluates to:
% \begin{align*}
%     \sum_{n=1}^{\frac{1}{\sqrt{\eta}}} \left(\frac{1-\rho}{\eta k}\right)^s n^2 \rho^n 
%     &\eqsim \left(\frac{1-\rho}{\eta k}\right)^s \sum_{n=1}^{\frac{1}{\sqrt{\eta}}} n^2 \\
%     &\eqsim \left(\frac{1-\rho}{\eta k}\right)^s \left(\frac{1}{\sqrt{\eta}}\right)^3 = \eta^{-\frac{3}{2}} \left(\frac{1-\rho}{\eta k}\right)^s.
% \end{align*}

% \textbf{Final Magnitude Comparison:}
% We now compare the magnitudes of all six stages to derive the final conclusion.
% Stage 1 produces the dominant exponential transient $\frac{1}{\eta(1-\rho)} \left(\frac{(1-\rho)^2}{\eta}\right)^{s+1} \gamma_{\nesterov}^k$. 
% Stage 5 produces the dominant polynomial tail $\frac{1}{\eta(1-\rho)} \left(\frac{1-\rho}{\eta k}\right)^s$.
% We rigorously verify that Stages 2, 3, 4, and 6 are strictly bounded by Stage 5:
% \begin{itemize}
%     \item Stage 2 contains an exponentially small decay factor $\gamma_{\nesterov}^{\eta^{1/\beta}(1-\rho)^{-2-\frac{2}{\beta}}} \ll 1$.
%     \item Stage 3 contains an exponentially small decay factor $\gamma_{\nesterov}^{\frac{1}{(1-\rho)^2}} \ll 1$.
%     \item Stage 4 contains the decay factor $\rho^{\frac{1}{\eta}} \lesssim 1$, since $\eta \gtrsim (1-\rho)^2 \implies \frac{1}{\eta} \lesssim \frac{1}{(1-\rho)^2}$. It is at most comparable to Stage 5.
%     \item For Stage 6, the stability condition $\eta \gtrsim (1-\rho)^2$ strictly implies $\sqrt{\eta} \gtrsim 1-\rho$. Therefore, the prefactor satisfies $\eta^{-\frac{3}{2}} = \frac{1}{\eta\sqrt{\eta}} \lesssim \frac{1}{\eta(1-\rho)}$. Thus, Stage 6 is strictly subsumed by Stage 5.
% \end{itemize}
% Consequently, aggregating all stages yields exactly the sum of Stage 1 and Stage 5, completing the proof.
% \end{proof}


% \begin{lemma}[Convolution $\cS_{\mathrm{under}} * \mathcal{K}_{\mathrm{over}}$]
% \label{lem:nes_conv_under_over}
% Under the assumptions of Proposition~\ref{prop:nes_deterministic_bias} and Lemma~\ref{lem:nes_kernel}, and under the stability condition $\eta\lesssim B(1-\rho)$, for $k \gtrsim \eta^{1/\beta}(1-\rho)^{-2-\frac{2}{\beta}}$, we have:
% \begin{equation*}
%     \sum_{m=0}^{k-1} \cS_{\mathrm{under}}(m) \mathcal{K}_{\mathrm{over}}(k-1-m) \eqsim \frac{1}{(1-\rho)^3} \left(\frac{1-\rho}{\eta k}\right)^{2-\frac{1}{\beta}} + \frac{1}{(1-\rho)^3} \left(\frac{(1-\rho)^2}{\eta}\right)^{s+2} \gamma_{\nesterov}^k.
% \end{equation*}
% \end{lemma}
% \begin{proof}
% We evaluate this convolution by transforming the kernel index to $n = k-1-m \eqsim k-m$, and rigorously dividing the summation over $m$ into five distinct stages corresponding to the piecewise bounds of $\cS_{\mathrm{under}}$.

% \textbf{Stage 1:} $0 \le m \lesssim \frac{1}{\eta}$. 
% Here $\cS_{\mathrm{under}}(m) \eqsim \rho^m$. Since $k$ is extremely large, the kernel index $n \eqsim k \gtrsim \frac{1}{1-\rho}$, giving $\mathcal{K}_{\mathrm{over}}(n) \eqsim \frac{1}{(1-\rho)^2} \left(\frac{1-\rho}{\eta k}\right)^{2-\frac{1}{\beta}}$. 
% Summing the geometric series $\rho^m$ up to $\frac{1}{\eta}$ yields a factor of $\frac{1}{1-\rho}$:
% \begin{align*}
%     \sum_{m=0}^{\frac{1}{\eta}} \rho^m \frac{1}{(1-\rho)^2} \left(\frac{1-\rho}{\eta k}\right)^{2-\frac{1}{\beta}} 
%     &\eqsim \frac{1}{(1-\rho)^2} \left(\frac{1-\rho}{\eta k}\right)^{2-\frac{1}{\beta}} \sum_{m=0}^{\frac{1}{\eta}} \rho^m \\
%     &\eqsim \frac{1}{1-\rho} \frac{1}{(1-\rho)^2} \left(\frac{1-\rho}{\eta k}\right)^{2-\frac{1}{\beta}}\\
%     &= \frac{1}{(1-\rho)^3} \left(\frac{1-\rho}{\eta k}\right)^{2-\frac{1}{\beta}}.
% \end{align*}

% \textbf{Stage 2:} $\frac{1}{\eta} \lesssim m \lesssim \frac{1}{(1-\rho)^2}$. 
% Here $\cS_{\mathrm{under}}(m) \eqsim \frac{\rho^m}{(\eta m)^s}$. The kernel index remains $n \eqsim k$, so $\mathcal{K}_{\mathrm{over}}(n) \eqsim \frac{1}{(1-\rho)^2} \left(\frac{1-\rho}{\eta k}\right)^{2-\frac{1}{\beta}}$. 
% Since $m \gtrsim \frac{1}{\eta}$, we have $\eta m \gtrsim 1$, which strictly implies $(\eta m)^{-s} \lesssim 1$ for any $s > 0$. Therefore, $\cS_{\mathrm{under}}(m) \lesssim \rho^m$. The sum is bounded by:
% \begin{align*}
%     \sum_{m=\frac{1}{\eta}}^{\frac{1}{(1-\rho)^2}} \frac{\rho^m}{(\eta m)^s} \frac{1}{(1-\rho)^2} \left(\frac{1-\rho}{\eta k}\right)^{2-\frac{1}{\beta}} 
%     &\lesssim \frac{1}{(1-\rho)^2} \left(\frac{1-\rho}{\eta k}\right)^{2-\frac{1}{\beta}} \sum_{m=\frac{1}{\eta}}^{\frac{1}{(1-\rho)^2}} \rho^m \\
%     &\lesssim \frac{1}{(1-\rho)^2} \left(\frac{1-\rho}{\eta k}\right)^{2-\frac{1}{\beta}} \frac{\rho^{\frac{1}{\eta}}}{1-\rho} \\
%     &\lesssim \frac{1}{(1-\rho)^3} \left(\frac{1-\rho}{\eta k}\right)^{2-\frac{1}{\beta}}.
% \end{align*}

% \textbf{Stage 3:} $\frac{1}{(1-\rho)^2} \lesssim m \lesssim \eta^{1/\beta}(1-\rho)^{-2-\frac{2}{\beta}}$. 
% Here $\cS_{\mathrm{under}}(m) \eqsim \frac{1}{\eta m} \left(\frac{(1-\rho)^2}{\eta}\right)^{s-1} \gamma_{\nesterov}^m$. The kernel index is still $n \eqsim k$, so $\mathcal{K}_{\mathrm{over}}(n) \eqsim \frac{1}{(1-\rho)^2} \left(\frac{1-\rho}{\eta k}\right)^{2-\frac{1}{\beta}}$. 
% Due to the exponential decay of $\gamma_{\nesterov}^m$, the sum is strictly dominated by its lower limit $m_1 = \frac{1}{(1-\rho)^2}$. At this limit, $\cS_{\mathrm{under}}(m_1) \eqsim \frac{(1-\rho)^2}{\eta} \left(\frac{(1-\rho)^2}{\eta}\right)^{s-1} \gamma_{\nesterov}^{m_1} = \left(\frac{(1-\rho)^2}{\eta}\right)^s \gamma_{\nesterov}^{\frac{1}{(1-\rho)^2}}$. 
% Summing the geometric decay introduces a factor of $\frac{1}{1-\gamma_{\nesterov}} \eqsim \frac{1}{1-\rho}$:
% \begin{align*}
%     \sum_{m=m_1}^{\eta^{1/\beta}(1-\rho)^{-2-\frac{2}{\beta}}} \cS_{\mathrm{under}}(m) \mathcal{K}_{\mathrm{over}}(k) 
%     &\eqsim \frac{1}{(1-\rho)^2} \left(\frac{1-\rho}{\eta k}\right)^{2-\frac{1}{\beta}} \left(\frac{(1-\rho)^2}{\eta}\right)^s \gamma_{\nesterov}^{\frac{1}{(1-\rho)^2}} \sum_{m=0}^{\infty} \gamma_{\nesterov}^m \\
%     &\eqsim \frac{1}{(1-\rho)^3} \left(\frac{1-\rho}{\eta k}\right)^{2-\frac{1}{\beta}} \left(\frac{(1-\rho)^2}{\eta}\right)^s \gamma_{\nesterov}^{\frac{1}{(1-\rho)^2}}.
% \end{align*}

% \textbf{Stage 4:} $\eta^{1/\beta}(1-\rho)^{-2-\frac{2}{\beta}} \lesssim m \lesssim k - \frac{1}{1-\rho}$. 
% Here $\cS_{\mathrm{under}}(m) \eqsim \left(\frac{(1-\rho)^2}{\eta}\right)^{s+\frac{1}{\beta}} \gamma_{\nesterov}^m$. The kernel index $n$ ranges from $\frac{1}{1-\rho}$ to $k$, giving $\mathcal{K}_{\mathrm{over}}(n) \eqsim \frac{1}{(1-\rho)^2} \left(\frac{1-\rho}{\eta n}\right)^{2-\frac{1}{\beta}}$. 
% Since $\gamma_{\nesterov}^m$ decays exponentially while $\mathcal{K}_{\mathrm{over}}(n)$ grows polynomially as $m$ approaches $k$, the sum over this interval is bounded by the maximum of its two endpoints. \\
% At the lower endpoint $m_2 = \eta^{1/\beta}(1-\rho)^{-2-\frac{2}{\beta}}$, the contribution is \[\eqsim \frac{1}{(1-\rho)^3} \left(\frac{1-\rho}{\eta k}\right)^{2-\frac{1}{\beta}} \left(\frac{(1-\rho)^2}{\eta}\right)^{s+\frac{1}{\beta}} \gamma_{\nesterov}^{m_2}.\]\\ 
% At the upper endpoint $m_3 = k - \frac{1}{1-\rho}$, where $n \eqsim \frac{1}{1-\rho}$, the contribution smoothly connects to Stage 5, yielding \[\eqsim \frac{1}{(1-\rho)^3} \left(\frac{(1-\rho)^2}{\eta}\right)^{s+2} \gamma_{\nesterov}^k.\]

% \textbf{Stage 5:} $k - \frac{1}{1-\rho} \lesssim m \le k - 1$. 
% Here $m \eqsim k$, so $\cS_{\mathrm{under}}(m) \eqsim \left(\frac{(1-\rho)^2}{\eta}\right)^{s+\frac{1}{\beta}} \gamma_{\nesterov}^k$. The kernel index $n = k-1-m \in \left[0, \frac{1}{1-\rho}\right]$, giving the short-time form $\mathcal{K}_{\mathrm{over}}(n) \eqsim n^2 \left(\frac{(1-\rho)^2}{\eta}\right)^{2-\frac{1}{\beta}}$. 
% Summing the polynomial $n^2$ over this short interval yields $\sum_{n=0}^{\frac{1}{1-\rho}} n^2 \eqsim \left(\frac{1}{1-\rho}\right)^3 = \frac{1}{(1-\rho)^3}$. Thus, the sum strictly evaluates to:
% \begin{align*}
%     \sum_{n=0}^{\frac{1}{1-\rho}} \left(\frac{(1-\rho)^2}{\eta}\right)^{s+\frac{1}{\beta}} \gamma_{\nesterov}^k \cdot n^2 \left(\frac{(1-\rho)^2}{\eta}\right)^{2-\frac{1}{\beta}} 
%     &\eqsim \left(\frac{(1-\rho)^2}{\eta}\right)^{s+\frac{1}{\beta}} \gamma_{\nesterov}^k \left(\frac{(1-\rho)^2}{\eta}\right)^{2-\frac{1}{\beta}} \sum_{n=0}^{\frac{1}{1-\rho}} n^2 \\
%     &\eqsim \left(\frac{(1-\rho)^2}{\eta}\right)^{s+\frac{1}{\beta} + 2 - \frac{1}{\beta}} \gamma_{\nesterov}^k \frac{1}{(1-\rho)^3} \\
%     &= \frac{1}{(1-\rho)^3} \left(\frac{(1-\rho)^2}{\eta}\right)^{s+2} \gamma_{\nesterov}^k.
% \end{align*}

% \textbf{Final Magnitude Comparison:}
% We now compare the magnitudes of all five stages to derive the final conclusion.
% Stage 1 yields the dominant polynomial tail $\frac{1}{(1-\rho)^3} \left(\frac{1-\rho}{\eta k}\right)^{2-\frac{1}{\beta}}$.
% Stage 5 yields the dominant exponential transient $\frac{1}{(1-\rho)^3} \left(\frac{(1-\rho)^2}{\eta}\right)^{s+2} \gamma_{\nesterov}^k$.
% We rigorously verify that Stages 2, 3, and 4 are strictly bounded by the sum of Stage 1 and Stage 5:
% \begin{itemize}
%     \item Stage 2 is explicitly bounded by $\lesssim$ Stage 1, as shown in its derivation.
%     \item Stage 3 contains the exponential decay factor $\gamma_{\nesterov}^{\frac{1}{(1-\rho)^2}} \ll 1$, making it exponentially suppressed compared to Stage 1.
%     \item Stage 4 is bounded by the maximum of its endpoints. Its lower endpoint is exponentially suppressed by $\gamma_{\nesterov}^{\eta^{1/\beta}(1-\rho)^{-2-2/\beta}} \ll 1$ (much smaller than Stage 1), and its upper endpoint directly matches the magnitude of Stage 5.
% \end{itemize}
% Consequently, aggregating all stages yields exactly the sum of Stage 1 and Stage 5, completing the proof.
% \end{proof}

% \begin{lemma}[Convolution $\cS_{\mathrm{under}} * \mathcal{K}_{\mathrm{under}}$]
% \label{lem:nes_conv_under_under}
% Under the assumptions of Proposition~\ref{prop:nes_deterministic_bias} and Lemma~\ref{lem:nes_kernel}, and under the stability condition $\eta\lesssim B(1-\rho)$, for $k \gtrsim \eta^{1/\beta}(1-\rho)^{-2-\frac{2}{\beta}}$, we have:
% \begin{equation*}
%     \sum_{m=0}^{k-1} \cS_{\mathrm{under}}(m) \mathcal{K}_{\mathrm{under}}(k-1-m) \eqsim \frac{k}{\eta} \left(\frac{(1-\rho)^2}{\eta}\right)^{s+1+\frac{1}{\beta}} \gamma_{\nesterov}^k.
% \end{equation*}
% \end{lemma}
% \begin{proof}
% We evaluate this convolution by transforming the kernel index to $n = k-1-m \eqsim k-m$, and rigorously dividing the summation over $m$ into eight distinct stages based on the precise piecewise definitions of $\cS_{\mathrm{under}}$ and $\mathcal{K}_{\mathrm{under}}$.

% \textbf{Stage 1:} $0 \le m \lesssim \frac{1}{\eta}$. 
% Here $\cS_{\mathrm{under}}(m) \eqsim \rho^m$. Since $k$ is extremely large, the kernel index $n \eqsim k \gtrsim \eta^{1/\beta}(1-\rho)^{-2-\frac{2}{\beta}}$, giving $\mathcal{K}_{\mathrm{under}}(n) \eqsim \frac{(1-\rho)^2}{\eta^2} \gamma_{\nesterov}^n$. 
% For this range, $\rho^m \gamma_{\nesterov}^{k-m} \eqsim \gamma_{\nesterov}^k (\rho/\gamma_{\nesterov})^m$. Applying Lemma~\ref{lem:nes_rho_gamma}, we effectively bound the geometric sum $\sum \rho^m \lesssim \frac{1}{1-\rho}$, yielding:
% \begin{align*}
%     \sum_{m=0}^{\frac{1}{\eta}} \rho^m \frac{(1-\rho)^2}{\eta^2} \gamma_{\nesterov}^{k-m} 
%     &\lesssim \frac{1}{1-\rho} \frac{(1-\rho)^2}{\eta^2} \gamma_{\nesterov}^k = \frac{1-\rho}{\eta^2} \gamma_{\nesterov}^k.
% \end{align*}

% \textbf{Stage 2:} $\frac{1}{\eta} \lesssim m \lesssim \frac{1}{(1-\rho)^2}$. 
% Here $\cS_{\mathrm{under}}(m) \eqsim \frac{\rho^m}{(\eta m)^s}$. The kernel remains $\mathcal{K}_{\mathrm{under}}(n) \eqsim \frac{(1-\rho)^2}{\eta^2} \gamma_{\nesterov}^n$. 
% The summation is bounded by the polynomial limits of $m^{-s}$:
% \begin{align*}
%     \sum_{m=\frac{1}{\eta}}^{\frac{1}{(1-\rho)^2}} \frac{\rho^m}{(\eta m)^s} \frac{(1-\rho)^2}{\eta^2} \gamma_{\nesterov}^{k-m} 
%     &\eqsim \frac{(1-\rho)^2}{\eta^{s+2}} \gamma_{\nesterov}^k \sum_{m=\frac{1}{\eta}}^{\frac{1}{(1-\rho)^2}} m^{-s} \\
%     &\lesssim \frac{(1-\rho)^2}{\eta^{s+2}} \gamma_{\nesterov}^k \max\left\{ \left(\frac{1}{\eta}\right)^{1-s}, \left(\frac{1}{(1-\rho)^2}\right)^{1-s} \right\}.
% \end{align*}
% In either case, it is merely a polynomial in $\eta$ and $(1-\rho)$ multiplied by $\gamma_{\nesterov}^k$, lacking the global $\mathcal{O}(k)$ multiplier.

% \textbf{Stage 3:} $\frac{1}{(1-\rho)^2} \lesssim m \lesssim \eta^{1/\beta}(1-\rho)^{-2-\frac{2}{\beta}}$. 
% Here $\cS_{\mathrm{under}}(m) \eqsim \frac{1}{\eta m} \left(\frac{(1-\rho)^2}{\eta}\right)^{s-1} \gamma_{\nesterov}^m$. The kernel is $\mathcal{K}_{\mathrm{under}}(n) \eqsim \frac{(1-\rho)^2}{\eta^2} \gamma_{\nesterov}^{k-m}$. 
% The product simplifies elegantly:
% \begin{equation*}
%     \frac{1}{\eta m} \left(\frac{(1-\rho)^2}{\eta}\right)^{s-1} \gamma_{\nesterov}^m \frac{(1-\rho)^2}{\eta^2} \gamma_{\nesterov}^{k-m} = \frac{1}{\eta^2 m} \left(\frac{(1-\rho)^2}{\eta}\right)^s \gamma_{\nesterov}^k.
% \end{equation*}
% Summing $m^{-1}$ over this specific interval yields a logarithmic factor $\ln\left( \frac{\eta^{1/\beta}(1-\rho)^{-2-2/\beta}}{(1-\rho)^{-2}} \right) \eqsim \ln\left(\frac{1}{1-\rho}\right)$:
% \begin{align*}
%     \sum_{m=\frac{1}{(1-\rho)^2}}^{\eta^{1/\beta}(1-\rho)^{-2-\frac{2}{\beta}}} \frac{1}{\eta^2 m} \left(\frac{(1-\rho)^2}{\eta}\right)^s \gamma_{\nesterov}^k 
%     &\eqsim \frac{1}{\eta^2} \left(\frac{(1-\rho)^2}{\eta}\right)^s \ln\left(\frac{1}{1-\rho}\right) \gamma_{\nesterov}^k.
% \end{align*}

% \textbf{Stage 4 (The Bulk):} $\eta^{1/\beta}(1-\rho)^{-2-\frac{2}{\beta}} \lesssim m \lesssim k - \eta^{1/\beta}(1-\rho)^{-2-\frac{2}{\beta}}$. 
% Here $\cS_{\mathrm{under}}(m) \eqsim \left(\frac{(1-\rho)^2}{\eta}\right)^{s+\frac{1}{\beta}} \gamma_{\nesterov}^m$. The kernel is $\mathcal{K}_{\mathrm{under}}(n) \eqsim \frac{(1-\rho)^2}{\eta^2} \gamma_{\nesterov}^{k-m}$. 
% The product is strictly constant with respect to $m$:
% \begin{equation*}
%     \left(\frac{(1-\rho)^2}{\eta}\right)^{s+\frac{1}{\beta}} \gamma_{\nesterov}^m \frac{(1-\rho)^2}{\eta^2} \gamma_{\nesterov}^{k-m} = \frac{1}{\eta} \left(\frac{(1-\rho)^2}{\eta}\right)^{s+1+\frac{1}{\beta}} \gamma_{\nesterov}^k.
% \end{equation*}
% Since the length of this integration region is asymptotically $k - 2\eta^{1/\beta}(1-\rho)^{-2-\frac{2}{\beta}} \eqsim k$, summing this constant term directly yields:
% \begin{align*}
%     \sum_{m} \cS_{\mathrm{under}}(m) \mathcal{K}_{\mathrm{under}}(k-1-m) \eqsim \frac{k}{\eta} \left(\frac{(1-\rho)^2}{\eta}\right)^{s+1+\frac{1}{\beta}} \gamma_{\nesterov}^k.
% \end{align*}

% \textbf{Stage 5:} $k - \eta^{1/\beta}(1-\rho)^{-2-\frac{2}{\beta}} \lesssim m \lesssim k - \frac{1}{(1-\rho)^2}$. 
% This stage is perfectly symmetric to Stage 3. Here $n = k-m \in \left[\frac{1}{(1-\rho)^2}, \eta^{1/\beta}(1-\rho)^{-2-\frac{2}{\beta}}\right]$. $\cS_{\mathrm{under}}(m) \eqsim \left(\frac{(1-\rho)^2}{\eta}\right)^{s+\frac{1}{\beta}} \gamma_{\nesterov}^m$, and $\mathcal{K}_{\mathrm{under}}(n) \eqsim \frac{1}{\eta^2 n} \left(\frac{(1-\rho)^2}{\eta}\right)^{-\frac{1}{\beta}} \gamma_{\nesterov}^n$. 
% The product evaluates to:
% \begin{equation*}
%     \left(\frac{(1-\rho)^2}{\eta}\right)^{s+\frac{1}{\beta}} \gamma_{\nesterov}^m \frac{1}{\eta^2 n} \left(\frac{(1-\rho)^2}{\eta}\right)^{-\frac{1}{\beta}} \gamma_{\nesterov}^n = \frac{1}{\eta^2 n} \left(\frac{(1-\rho)^2}{\eta}\right)^s \gamma_{\nesterov}^k.
% \end{equation*}
% Summing $n^{-1}$ over its interval again produces the logarithmic factor:
% \begin{align*}
%     \sum_{n} \frac{1}{\eta^2 n} \left(\frac{(1-\rho)^2}{\eta}\right)^s \gamma_{\nesterov}^k \eqsim \frac{1}{\eta^2} \left(\frac{(1-\rho)^2}{\eta}\right)^s \ln\left(\frac{1}{1-\rho}\right) \gamma_{\nesterov}^k.
% \end{align*}

% \textbf{Stage 6:} $k - \frac{1}{(1-\rho)^2} \lesssim m \lesssim k - \frac{1}{\eta}$. 
% Here $n \in \left[\frac{1}{\eta}, \frac{1}{(1-\rho)^2}\right]$. $\cS_{\mathrm{under}}(m) \eqsim \left(\frac{(1-\rho)^2}{\eta}\right)^{s+\frac{1}{\beta}} \gamma_{\nesterov}^m$, and $\mathcal{K}_{\mathrm{under}}(n) \eqsim \frac{\rho^n}{\eta^{2-\frac{1}{\beta}} n^{1-\frac{1}{\beta}}}$. 
% Since $\rho^n \gamma_{\nesterov}^m \eqsim \gamma_{\nesterov}^k$ for bounded short stages, we sum over $n$:
% \begin{align*}
%     \sum_{n=\frac{1}{\eta}}^{\frac{1}{(1-\rho)^2}} \left(\frac{(1-\rho)^2}{\eta}\right)^{s+\frac{1}{\beta}} \frac{1}{\eta^{2-\frac{1}{\beta}} n^{1-\frac{1}{\beta}}} \gamma_{\nesterov}^k 
%     &\eqsim \left(\frac{(1-\rho)^2}{\eta}\right)^{s+\frac{1}{\beta}} \frac{1}{\eta^{2-\frac{1}{\beta}}} \gamma_{\nesterov}^k \left( n^{\frac{1}{\beta}} \Big|_{\frac{1}{\eta}}^{\frac{1}{(1-\rho)^2}} \right) \\
%     &\eqsim \left(\frac{(1-\rho)^2}{\eta}\right)^{s+\frac{1}{\beta}} \frac{1}{\eta^{2-\frac{1}{\beta}}} (1-\rho)^{-\frac{2}{\beta}} \gamma_{\nesterov}^k = \frac{1}{\eta^2} \left(\frac{(1-\rho)^2}{\eta}\right)^s \gamma_{\nesterov}^k.
% \end{align*}

% \textbf{Stage 7:} $k - \frac{1}{\eta} \lesssim m \lesssim k - \frac{1}{\sqrt{\eta}}$. 
% Here $n \in \left[\frac{1}{\sqrt{\eta}}, \frac{1}{\eta}\right]$. $\cS_{\mathrm{under}}(m) \eqsim \left(\frac{(1-\rho)^2}{\eta}\right)^{s+\frac{1}{\beta}} \gamma_{\nesterov}^m$, and $\mathcal{K}_{\mathrm{under}}(n) \eqsim \frac{\rho^n}{\eta}$. 
% The sum over $n$ has interval length $\mathcal{O}(\frac{1}{\eta})$:
% \begin{align*}
%     \sum_{n=\frac{1}{\sqrt{\eta}}}^{\frac{1}{\eta}} \left(\frac{(1-\rho)^2}{\eta}\right)^{s+\frac{1}{\beta}} \frac{\gamma_{\nesterov}^k}{\eta} \eqsim \frac{1}{\eta} \left(\frac{(1-\rho)^2}{\eta}\right)^{s+\frac{1}{\beta}} \frac{\gamma_{\nesterov}^k}{\eta} = \frac{1}{\eta^2} \left(\frac{(1-\rho)^2}{\eta}\right)^{s+\frac{1}{\beta}} \gamma_{\nesterov}^k.
% \end{align*}

% \textbf{Stage 8:} $k - \frac{1}{\sqrt{\eta}} \lesssim m \le k-1$. 
% Here $n \in \left[1, \frac{1}{\sqrt{\eta}}\right]$. $\cS_{\mathrm{under}}(m) \eqsim \left(\frac{(1-\rho)^2}{\eta}\right)^{s+\frac{1}{\beta}} \gamma_{\nesterov}^m$, and $\mathcal{K}_{\mathrm{under}}(n) \eqsim n^2 \rho^n$. 
% Summing the polynomial $n^2$ up to $1/\sqrt{\eta}$ yields:
% \begin{align*}
%     \sum_{n=1}^{\frac{1}{\sqrt{\eta}}} \left(\frac{(1-\rho)^2}{\eta}\right)^{s+\frac{1}{\beta}} \gamma_{\nesterov}^k n^2 \eqsim \left(\frac{(1-\rho)^2}{\eta}\right)^{s+\frac{1}{\beta}} \gamma_{\nesterov}^k \left(\frac{1}{\sqrt{\eta}}\right)^3 = \eta^{-\frac{3}{2}} \left(\frac{(1-\rho)^2}{\eta}\right)^{s+\frac{1}{\beta}} \gamma_{\nesterov}^k.
% \end{align*}

% \textbf{Final Magnitude Comparison:}
% We meticulously compare all 8 stages to the global bulk extracted in Stage 4: $\frac{k}{\eta} \left(\frac{(1-\rho)^2}{\eta}\right)^{s+1+\frac{1}{\beta}} \gamma_{\nesterov}^k$.
% Stage 4 possesses a decisive global multiplier $k$. We verify that Stage 4 strictly dominates the largest symmetric logarithmic transients (Stage 3 and Stage 5). We take their ratio:
% \begin{equation*}
%     \frac{\frac{1}{\eta^2} \left(\frac{(1-\rho)^2}{\eta}\right)^s \ln\left(\frac{1}{1-\rho}\right)}{\frac{k}{\eta} \left(\frac{(1-\rho)^2}{\eta}\right)^{s+1+\frac{1}{\beta}}} = \frac{\ln\left(\frac{1}{1-\rho}\right)}{\eta k \left(\frac{(1-\rho)^2}{\eta}\right)^{1+\frac{1}{\beta}}} = \frac{\ln\left(\frac{1}{1-\rho}\right)}{k \eta^{-\frac{1}{\beta}} (1-\rho)^{2+\frac{2}{\beta}}}.
% \end{equation*}
% Given the established condition $k \gtrsim \eta^{1/\beta}(1-\rho)^{-2-\frac{2}{\beta}}$, the denominator inherently scales as $\gtrsim 1 \cdot k_0$, where $k_0 \to \infty$. Thus, the linear multiplier $k$ strictly overwhelms the logarithmic factor $\ln(1-\rho)^{-1}$.
% Consequently, the constant and logarithmic boundary stages (Stages 1, 2, 3, 5, 6, 7, 8) are rigorously subsumed by the primary bulk integration of Stage 4. Aggregating all segments completes the proof.
% \end{proof}

% \begin{lemma}[Total Deterministic Convolution]
% \label{lem:nes_conv_total}
% Under the assumptions of Proposition~\ref{prop:nes_deterministic_bias} and Lemma~\ref{lem:nes_kernel}, and under the stability condition $\eta\lesssim B(1-\rho)$, for $k \gtrsim \eta^{1/\beta}(1-\rho)^{-2-\frac{2}{\beta}}$, the total convolution between the deterministic bias and the Nesterov momentum kernel satisfies:
% \begin{equation*}
%     \sum_{m=0}^{k-1} \cS_{\nesterov}(m) \mathcal{K}_{\nesterov}(k-1-m) \eqsim \frac{k}{\eta} \left(\frac{(1-\rho)^2}{\eta}\right)^{s+1+\frac{1}{\beta}} \gamma_{\nesterov}^k + \frac{1}{\eta(1-\rho)} \left(\frac{1-\rho}{\eta k}\right)^s + \frac{1}{(1-\rho)^3} \left(\frac{1-\rho}{\eta k}\right)^{2-\frac{1}{\beta}}.
% \end{equation*}
% \end{lemma}
% \begin{proof}
% Synthesizing the results from Lemmas~\ref{lem:nes_conv_over_over},~\ref{lem:nes_conv_over_under},~\ref{lem:nes_conv_under_over}, and~\ref{lem:nes_conv_under_under}, we aggregate all individual convolution bounds.
% Because $\eta \gtrsim (1-\rho)^2$, the relation $\left(\frac{(1-\rho)^2}{\eta}\right)^{1-\frac{1}{\beta}} \lesssim 1$ mathematically holds. This explicitly guarantees that cross-mixed terms strictly relax into standard polynomial boundaries:
% \begin{equation*}
%     \frac{1}{(1-\rho)^3} \left(\frac{1-\rho}{\eta k}\right)^s \left(\frac{(1-\rho)^2}{\eta}\right)^{2-\frac{1}{\beta}} = \frac{1}{\eta(1-\rho)} \left(\frac{1-\rho}{\eta k}\right)^s \left(\frac{(1-\rho)^2}{\eta}\right)^{1-\frac{1}{\beta}} \lesssim \frac{1}{\eta(1-\rho)} \left(\frac{1-\rho}{\eta k}\right)^s.
% \end{equation*}
% Similarly, the main polynomial cross term gracefully decomposes:
% \begin{equation*}
%     \frac{1}{(1-\rho)^3} \left(\frac{(1-\rho)^2}{\eta}\right)^{s \vee (2-\frac{1}{\beta})} \left(\frac{1-\rho}{\eta k}\right)^{s \wedge (2-\frac{1}{\beta})} \lesssim \frac{1}{(1-\rho)^3} \left(\frac{1-\rho}{\eta k}\right)^{2-\frac{1}{\beta}} + \frac{1}{\eta(1-\rho)} \left(\frac{1-\rho}{\eta k}\right)^s.
% \end{equation*}
% For the exponential transient branches associated with $\gamma_{\nesterov}^k$, we isolate the bulk contribution from Lemma~\ref{lem:nes_conv_under_under}: $\frac{k}{\eta} \left(\frac{(1-\rho)^2}{\eta}\right)^{s+1+\frac{1}{\beta}} \gamma_{\nesterov}^k$. We verify its dominance by comparing it to the largest alternative bound from Lemma~\ref{lem:nes_conv_under_over}, $\frac{1}{(1-\rho)^3} \left(\frac{(1-\rho)^2}{\eta}\right)^{s+2} \gamma_{\nesterov}^k$. Taking their ratio gives:
% \begin{equation*}
%     \frac{\frac{1}{(1-\rho)^3} \left(\frac{(1-\rho)^2}{\eta}\right)^{s+2}}{\frac{k}{\eta} \left(\frac{(1-\rho)^2}{\eta}\right)^{s+1+\frac{1}{\beta}}} = \frac{\eta}{k(1-\rho)^3} \left(\frac{(1-\rho)^2}{\eta}\right)^{1-\frac{1}{\beta}}.
% \end{equation*}
% Since $k \gtrsim \eta^{1/\beta}(1-\rho)^{-2-\frac{2}{\beta}}$, $k$ contains a factor asymptotically larger than $(1-\rho)^{-2}$, ensuring this ratio is strictly $\lesssim 1$. Therefore, the total convolution dynamically resolves to the summation of the three unabsorbable primary scaling forms.
% \end{proof}


% \begin{proposition}[Multiplicative Noise Bounds, Mixed Regime]
% \label{prop:nes_self_noise_bounds}
% Under the assumptions of Proposition~\ref{prop:nes_deterministic_bias} and Lemma~\ref{lem:nes_kernel}, and under the stability condition $\eta\le cB(1-\rho)$ with $c>0$ sufficiently small, for $k \gtrsim \eta^{1/\beta}(1-\rho)^{-2-\frac{2}{\beta}}$, we have:
% \begin{align*}
%     \frac{\eta^2}{B}\sum_{\tau=0}^{k-1}\mathcal K_{\nesterov}(\tau)\mathbb{E}\mathcal E_{k-\tau}
%     \eqsim\quad
%     &\frac{\eta^2}{B}\bigg[
%         \frac{k}{\eta} \left(\frac{(1-\rho)^2}{\eta}\right)^{s+1+\frac{1}{\beta}} \gamma_{\nesterov}^k
%         + \frac1{(1-\rho)^3}\left(\frac{1-\rho}{\eta k}\right)^{2-\frac1\beta} 
%         + \frac1{\eta(1-\rho)}\left(\frac{1-\rho}{\eta k}\right)^s
%     \bigg] \\
%     &+ \frac{\eta^2\sigma^2}{B^2} \left( \min\left\{\left(\frac{\eta}{1-\rho}\right)^{\frac1\beta}, \frac{\eta}{1-\rho}\right\} \right)^2.
% \end{align*}
% \end{proposition}

% \begin{proof}
% To rigorously evaluate the multiplicative noise accumulation, we denote the convolution sequence as $\mathcal M_k^{\nesterov}$ and rewrite it in the time-forward form:
% \begin{equation*}
%     \mathcal M_k^{\nesterov} := \frac{\eta^2}{B} \sum_{\tau=0}^{k-1} \mathcal{K}_{\nesterov}(\tau) \mathbb{E}\mathcal{E}_{k-\tau} = \frac{\eta^2}{B} \sum_{m=1}^{k} \mathcal{K}_{\nesterov}(k-m) \mathbb{E}\mathcal{E}_m.
% \end{equation*}
% By applying the Nesterov expected risk recursion established in Proposition~\ref{prop:nes_risk_recursion} and substituting the look-ahead point variance replacement rigorously verified in Lemma~\ref{lem:nes_replace_lookahead_by_next_loss}, the next-step excess risk satisfies the bounded sequence:
% \begin{equation*}
%     \mathbb{E}\mathcal E_m \le C\left[ \mathcal D_{\nesterov}(m) + \frac{\eta^2\sigma^2}{B} \sum_{r=0}^{m-1}\mathcal K_{\nesterov}(r) + \mathcal M_m^{\nesterov} \right].
% \end{equation*}
% Substituting this uniform bound back into the definition of $\mathcal M_k^{\nesterov}$ explicitly isolates the noise accumulation into three distinct linear convolution operators:
% \begin{align}
% \label{eq:nes_mixed_bootstrap_start}
%     \mathcal M_k^{\nesterov} \le\quad 
%     &C \frac{\eta^2}{B}\sum_{m=1}^{k} \mathcal K_{\nesterov}(k-m)\mathcal D_{\nesterov}(m) \nonumber \\
%     &+ C \frac{\eta^4\sigma^2}{B^2}\sum_{m=1}^{k}\mathcal K_{\nesterov}(k-m)\sum_{r=0}^{m-1}\mathcal K_{\nesterov}(r) \nonumber \\
%     &+ C \frac{\eta^2}{B}\sum_{m=1}^{k} \mathcal K_{\nesterov}(k-m)\mathcal M_m^{\nesterov}.
% \end{align}

% We proceed to bound each of these three convolutions independently.

% \textbf{1. Label-Noise Convolution Term:} \\
% By exchanging the summation indices and exploiting the unconditional non-negativity of the Nesterov kernel sequence, the double sum representing the label-noise feedback is bounded by the square of the saturated kernel integral:
% \begin{align*}
%     \frac{\eta^4\sigma^2}{B^2}\sum_{m=1}^{k}\mathcal K_{\nesterov}(k-m)\sum_{r=0}^{m-1}\mathcal K_{\nesterov}(r) 
%     &= \frac{\eta^4\sigma^2}{B^2} \sum_{n=0}^{k-1} \mathcal K_{\nesterov}(n) \sum_{r=0}^{k-1-n} \mathcal K_{\nesterov}(r) \\
%     &\le \frac{\eta^4\sigma^2}{B^2} \left( \sum_{n=0}^{\infty} \mathcal K_{\nesterov}(n) \right)^2.
% \end{align*}
% By the exact saturated Green-energy identity established in Lemma~\ref{lem:nes_label_noise}, the infinite sum over the mixed-regime kernel converges uniformly to $\sum_{n=0}^{\infty} \mathcal K_{\nesterov}(n) \eqsim \frac{1}{\eta} \min\left\{\left(\frac{\eta}{1-\rho}\right)^{\frac1\beta}, \frac{\eta}{1-\rho}\right\}$. Substituting this resolves the exact label noise floor:
% \begin{align*}
%     \frac{\eta^4\sigma^2}{B^2} \left( \sum_{n=0}^{\infty} \mathcal K_{\nesterov}(n) \right)^2 
%     &\eqsim \frac{\eta^4\sigma^2}{B^2} \frac{1}{\eta^2} \left( \min\left\{\left(\frac{\eta}{1-\rho}\right)^{\frac1\beta}, \frac{\eta}{1-\rho}\right\} \right)^2 \\
%     &= \frac{\eta^2\sigma^2}{B^2} \left( \min\left\{\left(\frac{\eta}{1-\rho}\right)^{\frac1\beta}, \frac{\eta}{1-\rho}\right\} \right)^2.
% \end{align*}
% Since $k$ is sufficiently large ($k \gtrsim \eta^{1/\beta}(1-\rho)^{-2-\frac{2}{\beta}}$), the sum over the restricted finite range captures a constant fraction of the infinite sum limit, resulting in the rigorous lower bound of matching.

% \textbf{2. Deterministic Bias Convolution Term:} \\
% The convolution between the kernel and the deterministic bias is evaluated precisely by substituting the complete asymptotic summation formulated in Lemma~\ref{lem:nes_conv_total}:
% \begin{align*}
%     \frac{\eta^2}{B}\sum_{m=1}^{k}\mathcal K_{\nesterov}(k-m)\mathcal D_{\nesterov}(m) 
%     &\eqsim \frac{\eta^2}{B} \left[ \frac{k}{\eta} \left(\frac{(1-\rho)^2}{\eta}\right)^{s+1+\frac{1}{\beta}} \gamma_{\nesterov}^k 
%     + \frac{1}{(1-\rho)^3} \left(\frac{1-\rho}{\eta k}\right)^{2-\frac{1}{\beta}} \right.\\ 
%      &\left. + \frac{1}{\eta(1-\rho)} \left(\frac{1-\rho}{\eta k}\right)^s \right].
% \end{align*}

% \textbf{3. Self-Feedback Convolution Term:} \\
% We evaluate the recursive operator associated with $\mathcal M_m^{\nesterov}$. By explicitly expanding the inner sequence and carefully applying Fubini's theorem to exchange the finite discrete summations:
% \begin{align*}
%     \frac{\eta^2}{B}\sum_{m=1}^{k-1}\mathcal K_{\nesterov}(k-m)\mathcal M_m^{\nesterov} 
%     &= \frac{\eta^2}{B}\sum_{m=1}^{k-1} \mathcal K_{\nesterov}(k-m) \left[ \frac{\eta^2}{B}\sum_{\ell=1}^{m} \mathcal K_{\nesterov}(m-\ell) \mathbb{E}\mathcal E_\ell \right] \\
%     &= \frac{\eta^2}{B}\sum_{\ell=1}^{k-1} \left[ \frac{\eta^2}{B}\sum_{m=\ell}^{k-1} \mathcal K_{\nesterov}(k-m)\mathcal K_{\nesterov}(m-\ell) \right] \mathbb{E}\mathcal E_\ell.
% \end{align*}
% To analyze the bracketed inner kernel-kernel convolution $\sum_{r=0}^{n} \mathcal K_{\nesterov}(n-r)\mathcal K_{\nesterov}(r)$, we leverage the spectral decoupling property. Since the convolution of the discrete Green's function squared decays at most at the rate of the sequence itself, we apply the standard uniform sub-convolution bound:
% \begin{equation*}
%     \sum_{r=0}^n \mathcal K_{\nesterov}(n-r)\mathcal K_{\nesterov}(r) \le C \left( \sum_{r=0}^{\infty} \mathcal K_{\nesterov}(r) \right) \mathcal K_{\nesterov}(n).
% \end{equation*}
% Substituting the saturated mass sum $\sum_{r=0}^{\infty} \mathcal{K}_{\nesterov}(r) \le \frac{C}{\eta(1-\rho)}$ (which acts as a rigorous strict upper bound across all limits in the mixed regime), the bracketed operator reduces to:
% \begin{equation*}
%     \frac{\eta^2}{B}\sum_{\ell=1}^{k-1} \left[ \frac{\eta^2}{B} \frac{C}{\eta(1-\rho)} \mathcal K_{\nesterov}(k-\ell) \right] \mathbb{E}\mathcal E_\ell = C \frac{\eta}{B(1-\rho)} \left( \frac{\eta^2}{B} \sum_{\ell=1}^{k-1} \mathcal K_{\nesterov}(k-\ell) \mathbb{E}\mathcal E_\ell \right) = C \frac{\eta}{B(1-\rho)} \mathcal M_k^{\nesterov}.
% \end{equation*}
% Under the stated stability criterion $\eta \le cB(1-\rho)$, the global fractional multiplier mathematically satisfies:
% \begin{equation*}
%     C \frac{\eta}{B(1-\rho)} \le C \cdot c \le \frac{1}{2},
% \end{equation*}
% provided the tuning constant $c>0$ is chosen to be sufficiently small. Consequently, this self-feedback accumulation is structurally smaller than $\mathcal M_k^{\nesterov}$ itself and is  subtracted and permanently absorbed into the left-hand side of equation \eqref{eq:nes_mixed_bootstrap_start}.

% Synthesizing the exact boundaries formulated from the unabsorbable Label-Noise Convolution and Deterministic Bias Convolution directly yields the complete magnitude behavior stated in the proposition. The matching lower bound follows symmetrically since all individual variances and component matrices evaluate exclusively non-negatively.
% \end{proof}

% \begin{theorem}[Two-Sided Scaling Law, Mixed Regime]
% \label{thm:nes_two_sided_scaling}
% Under Assumption~\ref{ass:power-law} and Proposition~\ref{prop:noise-geometry}, the conditions of Proposition~\ref{prop:nes_deterministic_bias} and Lemma~\ref{lem:nes_kernel}, and the stability condition \(\eta\le cB(1-\rho)\) with \(c>0\) sufficiently small, for the extreme long-time mixed regime $k \gtrsim \eta^{1/\beta}(1-\rho)^{-2-\frac{2}{\beta}}$, the expected excess risk strictly satisfies:
% \begin{align*}
%     \mathbb{E}\mathcal E_k \eqsim\quad
%     &\left( 1 + \frac{k(1-\rho)^2}{B} \right) \left(\frac{(1-\rho)^2}{\eta}\right)^{s+\frac{1}{\beta}} \gamma_{\nesterov}^{\,k} 
%     + \left(\frac{1-\rho}{\eta k}\right)^s \\
%     &+ \frac{\sigma^2}{B} \min\left\{ \left(\frac{\eta}{1-\rho}\right)^{\frac1\beta}, \frac{\eta}{1-\rho} \right\}.
% \end{align*}
% Equivalently, up to the exponentially small Nesterov transient, the polynomial self-noise part adopts an identical analytical scaling form to the Polyak mixed overdamped--underdamped regime, while maintaining the distinctive Nesterov label-noise floor.
% \end{theorem}

% \begin{proof}
% To establish the final scaling law, we aggregate the constituent bounds derived in the previous propositions through the expected risk recursion (Proposition~\ref{prop:nes_risk_recursion}). The total expected excess risk decomposes into three macroscopic components: the deterministic bias $\cS_{\nesterov}(k)$, the accumulated label noise $\text{LN}(k)$, and the multiplicative noise feedback $\mathcal{M}_k^{\nesterov}$:
% \begin{equation*}
%     \mathbb{E}\mathcal E_k \eqsim \cS_{\nesterov}(k) + \text{LN}(k) + \mathcal{M}_k^{\nesterov}.
% \end{equation*}

% For the long-time regime $k \gtrsim \eta^{1/\beta}(1-\rho)^{-2-\frac{2}{\beta}}$, we systematically substitute the precise asymptotic sequences:
% \begin{enumerate}
%     \item \textbf{Deterministic Bias} (Proposition~\ref{prop:nes_deterministic_bias}):
%     \begin{equation*}
%         \cS_{\nesterov}(k) \eqsim \left(\frac{(1-\rho)^2}{\eta}\right)^{s+\frac{1}{\beta}} \gamma_{\nesterov}^k + \left(\frac{1-\rho}{\eta k}\right)^s.
%     \end{equation*}
%     \item \textbf{Label Noise} (Lemma~\ref{lem:nes_label_noise}):
%     \begin{equation*}
%         \text{LN}(k) = \frac{\eta^2\sigma^2}{B}\sum_{\tau=0}^{k-1}\mathcal{K}_{\nesterov}(\tau) \eqsim \frac{\sigma^2}{B} \min\left\{ \left(\frac{\eta}{1-\rho}\right)^{\frac1\beta}, \frac{\eta}{1-\rho} \right\}.
%     \end{equation*}
%     \item \textbf{Multiplicative Noise} (Proposition~\ref{prop:nes_self_noise_bounds}):
%     \begin{align*}
%         \mathcal M_k^{\nesterov} \eqsim\quad &\frac{\eta k}{B} \left(\frac{(1-\rho)^2}{\eta}\right)^{s+1+\frac{1}{\beta}} \gamma_{\nesterov}^k 
%         + \frac{\eta^2}{B(1-\rho)^3}\left(\frac{1-\rho}{\eta k}\right)^{2-\frac1\beta} 
%         + \frac{\eta}{B(1-\rho)}\left(\frac{1-\rho}{\eta k}\right)^s \\
%         &+ \frac{\eta^2\sigma^2}{B^2} \left( \min\left\{\left(\frac{\eta}{1-\rho}\right)^{\frac1\beta}, \frac{\eta}{1-\rho}\right\} \right)^2.
%     \end{align*}
% \end{enumerate}

% \textbf{Absorption Analysis:}
% We now meticulously evaluate the summation of these three components and demonstrate how lower-order fragments generated by the multiplicative noise $\mathcal M_k^{\nesterov}$ are strictly absorbed by the primary terms.

% \textit{1. Absorption of the Noise Floor:} \\
% We compare the multiplicative noise floor $\text{NF}(k) = \frac{\eta^2\sigma^2}{B^2} \left(\min\{\dots\}\right)^2$ with the primary label noise $\text{LN}(k)$. Their ratio is:
% \begin{equation*}
%     \frac{\text{NF}(k)}{\text{LN}(k)} = \frac{\frac{\eta^2\sigma^2}{B^2} \left( \min\left\{\left(\frac{\eta}{1-\rho}\right)^{\frac1\beta}, \frac{\eta}{1-\rho}\right\} \right)^2}{\frac{\sigma^2}{B} \min\left\{\left(\frac{\eta}{1-\rho}\right)^{\frac1\beta}, \frac{\eta}{1-\rho}\right\}} = \frac{\eta^2}{B} \min\left\{\left(\frac{\eta}{1-\rho}\right)^{\frac1\beta}, \frac{\eta}{1-\rho}\right\}.
% \end{equation*}
% Since the minimum function is trivially bounded by its linear argument, $\min\{\dots\} \le \frac{\eta}{1-\rho}$. Thus, the ratio satisfies:
% \begin{equation*}
%     \frac{\text{NF}(k)}{\text{LN}(k)} \le \frac{\eta^3}{B(1-\rho)} = \eta^2 \left(\frac{\eta}{B(1-\rho)}\right).
% \end{equation*}
% By the explicitly imposed stability condition $\eta \le cB(1-\rho)$, the factor $\frac{\eta}{B(1-\rho)} \le c \ll 1$. Furthermore, in the small learning rate regime, $\eta^2 \ll 1$. Therefore, the ratio is strictly $\ll 1$, mathematically guaranteeing that the multiplicative noise floor is  dominated by, and thus fully absorbed into, the label noise $\text{LN}(k)$.

% \textit{2. Absorption of the Multiplicative Polynomial Bias:} \\
% We examine the polynomial tail originating from the deterministic convolution in $\mathcal M_k^{\nesterov}$ that shares the source exponent $s$, denoted $P_2(k) = \frac{\eta}{B(1-\rho)}\left(\frac{1-\rho}{\eta k}\right)^s$. We compare it to the fundamental deterministic polynomial bias $P_1(k) = \left(\frac{1-\rho}{\eta k}\right)^s$ from $\cS_{\nesterov}(k)$. Their exact ratio is:
% \begin{equation*}
%     \frac{P_2(k)}{P_1(k)} = \frac{\eta}{B(1-\rho)}.
% \end{equation*}
% Under the identical stability condition $\eta \le cB(1-\rho)$ where $c$ is a sufficiently small constant, this ratio is $\le c \ll 1$. Thus, this polynomial cross-term is strictly absorbed into the primary deterministic polynomial bias.

% \textit{3. Absorption of the Multiplicative Polynomial Bias (exponent $2-\frac{1}{\beta}$):} \\
% We examine the polynomial tail originating from the multiplicative noise with exponent $2-\frac{1}{\beta}$, denoted $P_3(k) = \frac{\eta^2}{B(1-\rho)^3}\left(\frac{1-\rho}{\eta k}\right)^{2-\frac1\beta}$. We show that this term is universally absorbed by the primary label noise $\text{LN}(k) \eqsim \frac{\sigma^2}{B} \min\left\{ \left(\frac{\eta}{1-\rho}\right)^{\frac1\beta}, \frac{\eta}{1-\rho} \right\}$, treating $\sigma^2$ as a constant $\Theta(1)$. We separate the analysis into two distinct regimes:

% \textbf{Case 1: $\eta \le 1-\rho$.} \\
% Here, $\text{LN}(k) \eqsim \frac{\sigma^2\eta}{B(1-\rho)}$. Since $k \gtrsim \eta^{1/\beta}(1-\rho)^{-2-\frac{2}{\beta}}$ and $\eta \gtrsim (1-\rho)^2$, we have $k \gtrsim \frac{1}{1-\rho}$, which implies $\frac{1-\rho}{\eta k} \lesssim \frac{(1-\rho)^2}{\eta}$. Substituting this upper bound into $P_3(k)$ gives:
% \begin{align*}
%     P_3(k) &\lesssim \frac{\eta^2}{B(1-\rho)^3} \left(\frac{(1-\rho)^2}{\eta}\right)^{2-\frac{1}{\beta}} \\
%     &= \frac{\eta}{B(1-\rho)} \left(\frac{(1-\rho)^2}{\eta}\right)^{1-\frac{1}{\beta}}.
% \end{align*}
% Since $\eta \gtrsim (1-\rho)^2$, we have $\frac{(1-\rho)^2}{\eta} \lesssim 1$. Because $\beta > 1$, it follows that $P_3(k) \lesssim \frac{\eta}{B(1-\rho)} \eqsim \text{LN}(k)$. Thus, the term is  absorbed.

% \textbf{Case 2: $\eta > 1-\rho$.} \\
% Here, $\text{LN}(k) \eqsim \frac{\sigma^2}{B} \left(\frac{\eta}{1-\rho}\right)^{\frac1\beta}$. The extreme long-time condition $k \gtrsim \eta^{1/\beta}(1-\rho)^{-2-\frac{2}{\beta}}$ yields the strict upper bound on the time-dependent component:
% \begin{equation*}
%     \frac{1-\rho}{\eta k} \lesssim \frac{1-\rho}{\eta \cdot \eta^{1/\beta}(1-\rho)^{-2-\frac{2}{\beta}}} = \eta^{-1-\frac{1}{\beta}} (1-\rho)^{3+\frac{2}{\beta}}.
% \end{equation*}
% Substituting this directly into $P_3(k)$:
% \begin{align*}
%     P_3(k) &\lesssim \frac{\eta^2}{B(1-\rho)^3} \left( \eta^{-1-\frac{1}{\beta}} (1-\rho)^{3+\frac{2}{\beta}} \right)^{2-\frac{1}{\beta}} \\
%     &= \frac{1}{B} \eta^{2 - \left(1+\frac{1}{\beta}\right)\left(2-\frac{1}{\beta}\right)} (1-\rho)^{-3 + \left(3+\frac{2}{\beta}\right)\left(2-\frac{1}{\beta}\right)} \\
%     &= \frac{1}{B} \eta^{-\frac{1}{\beta} + \frac{1}{\beta^2}} (1-\rho)^{3 + \frac{1}{\beta} - \frac{2}{\beta^2}}.
% \end{align*}
% We now take the ratio of this bound to the label noise:
% \begin{align*}
%     \frac{P_3(k)}{\text{LN}(k)} &\lesssim \frac{\frac{1}{B} \eta^{-\frac{1}{\beta} + \frac{1}{\beta^2}} (1-\rho)^{3 + \frac{1}{\beta} - \frac{2}{\beta^2}}}{\frac{\sigma^2}{B} \eta^{\frac{1}{\beta}} (1-\rho)^{-\frac{1}{\beta}}} \\
%     &\eqsim \eta^{-\frac{2}{\beta} + \frac{1}{\beta^2}} (1-\rho)^{3 + \frac{2}{\beta} - \frac{2}{\beta^2}} \\
%     &= (1-\rho)^3 \left( \frac{1-\rho}{\eta} \right)^{\frac{1}{\beta}\left(2-\frac{1}{\beta}\right)}.
% \end{align*}
% Since $\eta > 1-\rho$, we have $\frac{1-\rho}{\eta} < 1$. Furthermore, since $\beta > 1$, the exponent $\frac{1}{\beta}\left(2-\frac{1}{\beta}\right) > 0$. Together with $(1-\rho)^3 < 1$, the entire ratio is rigorously bounded by $\lesssim 1$. Thus, $P_3(k) \lesssim \text{LN}(k)$ is completely absorbed by the label noise floor in all cases.

% \textbf{Final Scaling Law Construction:} \\
% After successfully absorbing the aforementioned terms, we gather all surviving, unabsorbable components. 
% For the exponential transient associated with $\gamma_{\nesterov}^k$, we algebraically combine the fundamental deterministic edge and the multiplicative edge:
% \begin{align*}
%     \left(\frac{(1-\rho)^2}{\eta}\right)^{s+\frac{1}{\beta}} \gamma_{\nesterov}^k + \frac{\eta k}{B} \left(\frac{(1-\rho)^2}{\eta}\right)^{s+1+\frac{1}{\beta}} \gamma_{\nesterov}^k 
%     &= \left( 1 + \frac{\eta k}{B} \frac{(1-\rho)^2}{\eta} \right) \left(\frac{(1-\rho)^2}{\eta}\right)^{s+\frac{1}{\beta}} \gamma_{\nesterov}^k \\
%     &= \left( 1 + \frac{k(1-\rho)^2}{B} \right) \left(\frac{(1-\rho)^2}{\eta}\right)^{s+\frac{1}{\beta}} \gamma_{\nesterov}^k.
% \end{align*}
% Neither sub-term can be universally discarded since $k$ can grow arbitrarily large with respect to $B(1-\rho)^{-2}$.

% Appending the surviving polynomial bias and the primary label noise floor yields the exact stated theorem:
% \begin{align*}
%     \mathbb{E}\mathcal E_k \eqsim \left( 1 + \frac{k(1-\rho)^2}{B} \right) \left(\frac{(1-\rho)^2}{\eta}\right)^{s+\frac{1}{\beta}} \gamma_{\nesterov}^{\,k} + \left(\frac{1-\rho}{\eta k}\right)^s + \frac{\sigma^2}{B} \min\left\{ \left(\frac{\eta}{1-\rho}\right)^{\frac1\beta}, \frac{\eta}{1-\rho} \right\}.
% \end{align*}
% This explicitly concludes the proof.
% \end{proof}

% \begin{theorem}[Two-Sided Scaling Law, No Label Noise, Mixed Regime]
% \label{thm:nes_two_sided_scaling_noiseless}
% Under Assumptions~\ref{ass:diagonal-covariance} and~\ref{ass:power-law}, with \(\sigma=0\), the stability condition \(\eta \lesssim B(1-\rho)\), and in the mixed regime \(k \gtrsim \eta^{1/\beta}(1-\rho)^{-2-\frac{2}{\beta}}\), the expected excess risk satisfies:
% \begin{equation}
% \label{eq:nes_two_sided_scaling_noiseless}
% \begin{aligned}
%     \EE[\cE(\btheta_k)]
%     \eqsim\quad
%     &\left(
%         1+\frac{k(1-\rho)^2}{B}
%     \right)
%     \left(\frac{(1-\rho)^2}{\eta}\right)^{s+\frac1\beta}
%     \gamma_{\nesterov}^{\,k} \\
%     &+
%     \left(\frac{1-\rho}{\eta k}\right)^s
%     +
%     \frac{\eta^2}{B(1-\rho)^3}
%     \left(\frac{1-\rho}{\eta k}\right)^{2-\frac1\beta}.
% \end{aligned}
% \end{equation}
% \end{theorem}

% \begin{proof}
% By Proposition~\ref{prop:nes_risk_recursion} and Lemma~\ref{lem:nes_replace_lookahead_by_next_loss}, when \(\sigma=0\) the expected excess risk reduces to:
% \[
%     \EE[\cE(\btheta_k)]
%     \eqsim
%     \cS_{\nesterov}(k)
%     +
%     \mathcal M_k^{\nesterov}.
% \]
% For the extreme long-time mixed regime \(k \gtrsim \eta^{1/\beta}(1-\rho)^{-2-\frac{2}{\beta}}\), Proposition~\ref{prop:nes_deterministic_bias} gives:
% \[
%     \cS_{\nesterov}(k)
%     \eqsim
%     \left(\frac{(1-\rho)^2}{\eta}\right)^{s+\frac1\beta}
%     \gamma_{\nesterov}^k
%     +
%     \left(\frac{1-\rho}{\eta k}\right)^s.
% \]
% Proposition~\ref{prop:nes_self_noise_bounds} provides the multiplicative noise:
% \[
%     \mathcal M_k^{\nesterov}
%     \eqsim
%     \frac{\eta k}{B}
%     \left(\frac{(1-\rho)^2}{\eta}\right)^{s+1+\frac1\beta}
%     \gamma_{\nesterov}^k
%     +
%     \frac{\eta^2}{B(1-\rho)^3}
%     \left(\frac{1-\rho}{\eta k}\right)^{2-\frac1\beta}
%     +
%     \frac{\eta}{B(1-\rho)}
%     \left(\frac{1-\rho}{\eta k}\right)^s.
% \]
% Grouping the exponential transient terms yields:
% \[
%     \left(
%         1+\frac{\eta k}{B}\frac{(1-\rho)^2}{\eta}
%     \right)
%     \left(\frac{(1-\rho)^2}{\eta}\right)^{s+\frac1\beta}
%     \gamma_{\nesterov}^k
%     =
%     \left(
%         1+\frac{k(1-\rho)^2}{B}
%     \right)
%     \left(\frac{(1-\rho)^2}{\eta}\right)^{s+\frac1\beta}
%     \gamma_{\nesterov}^k.
% \]
% For the polynomial term sharing the source exponent \(s\), we have:
% \[
%     \left(
%         1+\frac{\eta}{B(1-\rho)}
%     \right)
%     \left(\frac{1-\rho}{\eta k}\right)^s.
% \]
% By the stability condition \(\eta \lesssim B(1-\rho)\), we have \(\frac{\eta}{B(1-\rho)} \lesssim 1\). Hence the factor \(1+\frac{\eta}{B(1-\rho)}\) is bounded by a universal constant and can be absorbed into the asymptotic equivalence, giving:
% \[
%     1+\frac{\eta}{B(1-\rho)} \eqsim 1.
% \]
% Appending the remaining self-noise term concludes the proof.
% \end{proof}


\subsection{Uniform boundedness of expected risk}
\label{app:nes_uniform_boundedness}

As long as the additive label noise is strictly positive ($\sigma > 0$), we can prove that the expected risk is uniformly bounded, which allows the multiplicative gradient noise to be completely absorbed by the additive label noise floor.

% \begin{lemma}[Uniform boundedness of expected risk]
% \label{lem:nes_loss_bounded}
% Suppose $\sigma > 0$. Under the stability condition $\eta \le c B^\beta (1-\rho)$ for a sufficiently small universal constant $c > 0$, the expected excess risk of Nesterov is uniformly bounded for all $k \ge 0$:
% \[
%     \EE[\cE(\btheta_k)] \le M,
% \]
% where $M > 0$ is a constant independent of $k$, $\eta$, $B$, and $\rho$.
% \end{lemma}

% \begin{proof}
% We proceed by induction on $k$. By Proposition~\ref{prop:nes_risk_recursion}, there exist universal constants $C_1, C_2, C_3 > 0$ such that for any $k \ge 1$:
% \begin{equation} \label{eq:nes_risk_recursion_bound}
%     \EE[\cE(\btheta_k)] \le C_1 \cS_{\nesterov}(k) + C_2 \frac{\eta^2\sigma^2}{B}\sum_{\tau=0}^{k-1}\mathcal{K}_{\nesterov}(\tau) + C_3 \frac{\eta^2}{B}\sum_{\tau=0}^{k-1}\mathcal{K}_{\nesterov}(\tau)\EE[\cE(\btheta_{k-\tau})].
% \end{equation}

% First, Propositions~\ref{prop:nes_deterministic_bias0} and \ref{prop:nes_deterministic_bias} establish that the deterministic bias monotonically decays or is uniformly bounded by its initialization across all regimes. Thus:
% \[
%     \cS_{\nesterov}(k) \le C_h^2 \cE(\btheta_0) =: M_1.
% \]

% Second, applying Lemma~\ref{lem:nes_label_noise}, the infinite sum of the convolution kernel satisfies:
% \[
%     \frac{\eta^2}{B} \sum_{\tau=0}^\infty \mathcal{K}_{\nesterov}(\tau) \le \frac{C_K}{B} \min\left\{ \left(\frac{\eta}{1-\rho}\right)^{\frac1\beta}, \frac{\eta}{1-\rho} \right\}.
% \]
% Therefore, the label noise contribution is universally bounded by:
% \[
%     C_2 \frac{\eta^2\sigma^2}{B}\sum_{\tau=0}^{k-1}\mathcal{K}_{\nesterov}(\tau) \le C_2 \sigma^2 \frac{C_K}{B} \left(\frac{\eta}{1-\rho}\right)^{\frac1\beta}.
% \]
% Under the stability condition $\eta \le c B^\beta (1-\rho)$, we have $\left(\frac{\eta}{1-\rho}\right)^{1/\beta} \le c^{1/\beta} B$. Thus, the label noise is bounded by 
% \[
% C_2 \sigma^2 C_K c^{1/\beta} =: M_2.
% \]

% Assuming inductively $\EE[\cE(\btheta_m)] \le M$ for $m < k$, the multiplicative noise term yields:
% \[
%     C_3 \frac{\eta^2}{B}\sum_{\tau=0}^{k-1}\mathcal{K}_{\nesterov}(\tau)\EE[\cE(\btheta_{k-\tau})] \le C_3 \left( \frac{\eta^2}{B} \sum_{\tau=0}^{\infty} \mathcal{K}_{\nesterov}(\tau) \right) M \le C_3 C_K \frac{1}{B} \left(\frac{\eta}{1-\rho}\right)^{\frac1\beta} M \le C_3 C_K c^{1/\beta} M.
% \]
% Choosing $c$ sufficiently small such that $c^{1/\beta} \le \frac{1}{2 C_3 C_K}$ limits this term to $M/2$. Substituting into \eqref{eq:nes_risk_recursion_bound} yields:
% \[
%     \EE[\cE(\btheta_k)] \le M_1 + M_2 + \frac{M}{2}.
% \]
% Setting $M = 2(M_1 + M_2)$ guarantees $\EE[\cE(\btheta_k)] \le M$. The base case $k=0$ holds trivially since $\EE[\cE(\btheta_0)] = \cS_{\nesterov}(0) \le M_1 \le M$. This completes the induction.
% \end{proof}


\begin{lemma}[Uniform boundedness of expected risk]
\label{lem:nes_loss_bounded}
Suppose $\sigma > 0$. Under the stability condition $\eta \lesssim 1$ and $\eta \le c B^\beta (1-\rho)$ for a sufficiently small universal constant $c > 0$, the expected excess risk of Nesterov is uniformly bounded for all $k \ge 0$:
\[
    \EE[\cE(\btheta_k)] \le M,
\]
where $M > 0$ is a constant independent of $k$, $\eta$, $B$, and $\rho$.
\end{lemma}

\begin{proof}
We proceed by induction on $k$. By Proposition~\ref{prop:nes_risk_recursion}, there exist universal constants $C_1, C_2, C_3 > 0$ such that for any $k \ge 1$:
\begin{equation} \label{eq:nes_risk_recursion_bound}
    \EE[\cE(\btheta_k)] \le C_1 \cS_{\nesterov}(k) + C_2 \frac{\eta^2\sigma^2}{B}\sum_{\tau=0}^{k-1}\mathcal{K}_{\nesterov}(\tau) + C_3 \frac{\eta^2}{B}\sum_{\tau=0}^{k-1}\mathcal{K}_{\nesterov}(\tau)\EE[\cE(\btheta_{k-\tau})].
\end{equation}

First, Propositions~\ref{prop:nes_deterministic_bias0} and \ref{prop:nes_deterministic_bias} establish that the deterministic bias monotonically decays or is uniformly bounded by its initialization across all regimes. Thus:
\[
    \cS_{\nesterov}(k) \le C_h^2 \cE(\btheta_0) =: M_1.
\]

Second, applying Lemma~\ref{lem:nes_label_noise}, the infinite sum of the convolution kernel satisfies:
\[
    \frac{\eta^2}{B} \sum_{\tau=0}^\infty \mathcal{K}_{\nesterov}(\tau) \le \frac{C_K}{B} \min\left\{ \left(\frac{\eta}{1-\rho}\right)^{\frac1\beta}, \frac{\eta}{1-\rho} \right\}.
\]
Therefore, the label noise contribution is universally bounded by:
\[
    C_2 \frac{\eta^2\sigma^2}{B}\sum_{\tau=0}^{k-1}\mathcal{K}_{\nesterov}(\tau) \le C_2 \sigma^2 \frac{C_K}{B} \left(\frac{\eta}{1-\rho}\right)^{\frac1\beta}.
\]
Under the stability condition $\eta \le c B^\beta (1-\rho)$, we have $\left(\frac{\eta}{1-\rho}\right)^{1/\beta} \le c^{1/\beta} B$. Thus, the label noise is  bounded by
\[
    C_2 \sigma^2 C_K c^{1/\beta} =: M_2.
\]

To handle the multiplicative noise term, we must first isolate the $\tau=0$ component, which contains the current step $\EE[\cE(\btheta_k)]$. Note that
\[
    \mathcal{K}_{\nesterov}(0) = \sum_{j\ge1} \lambda_j^2 \bigl(G_0^{\nesterov}(j)\bigr)^2 = \sum_{j\ge1} \lambda_j^2 \le \kappa_0 < \infty.
\]
The $\tau=0$ term is therefore bounded by $C_3 \frac{\eta^2}{B} \kappa_0 \EE[\cE(\btheta_k)]$. Since the local stability condition implies $\eta \lesssim 1$, we can choose the learning rate prefactor sufficiently small such that $C_3 \frac{\eta^2}{B} \kappa_0 \le \frac{1}{2}$. Subtracting this term from both sides of \eqref{eq:nes_risk_recursion_bound} gives:
\begin{equation} \label{eq:nes_risk_recursion_bound_absorbed}
    \frac{1}{2} \EE[\cE(\btheta_k)] \le C_1 \cS_{\nesterov}(k) + C_2 \frac{\eta^2\sigma^2}{B}\sum_{\tau=0}^{k-1}\mathcal{K}_{\nesterov}(\tau) + C_3 \frac{\eta^2}{B}\sum_{\tau=1}^{k-1}\mathcal{K}_{\nesterov}(\tau)\EE[\cE(\btheta_{k-\tau})].
\end{equation}

Now, we assume inductively that $\EE[\cE(\btheta_m)] \le M$ for all $m < k$. For the remaining multiplicative noise summation ($\tau \ge 1 \implies k-\tau < k$), we can apply the induction hypothesis:
\begin{align*}
    C_3 \frac{\eta^2}{B}\sum_{\tau=1}^{k-1}\mathcal{K}_{\nesterov}(\tau)\EE[\cE(\btheta_{k-\tau})] 
    &\le C_3 \left( \frac{\eta^2}{B} \sum_{\tau=0}^{\infty} \mathcal{K}_{\nesterov}(\tau) \right) M \\
    &\le C_3 C_K \frac{1}{B} \left(\frac{\eta}{1-\rho}\right)^{\frac1\beta} M \\
    &\le C_3 C_K c^{1/\beta} M.
\end{align*}
By choosing $c$ sufficiently small such that $c^{1/\beta} \le \frac{1}{4 C_3 C_K}$, this term is limited to $M/4$. Substituting the bounds $M_1$, $M_2$, and $M/4$ into \eqref{eq:nes_risk_recursion_bound_absorbed} yields:
\[
    \frac{1}{2} \EE[\cE(\btheta_k)] \le M_1 + M_2 + \frac{M}{4}.
\]
Setting $M = 4(M_1 + M_2)$ guarantees that
\[
    \frac{1}{2} \EE[\cE(\btheta_k)] \le \frac{M}{4} + \frac{M}{4} = \frac{M}{2} 
    \implies 
    \EE[\cE(\btheta_k)] \le M.
\]
The base case $k=0$ holds trivially since $\EE[\cE(\btheta_0)] = \cS_{\nesterov}(0) \le M_1 < M$. This completes the induction.
\end{proof}

\begin{corollary}[Absorption of multiplicative noise]
\label{cor:nes_noise_absorption}
For $\sigma > 0$, under the stability condition $\eta \lesssim \min\{1, B^\beta (1-\rho)\}$, the multiplicative gradient noise is strictly bounded by the additive label noise up to a constant factor:
\[
    \frac{\eta^2}{B}\sum_{\tau=0}^{k-1}\mathcal{K}_{\nesterov}(\tau)\EE[\cE(\btheta_{k-\tau})] \le \frac{M}{\sigma^2} \left( \frac{\eta^2\sigma^2}{B}\sum_{\tau=0}^{k-1}\mathcal{K}_{\nesterov}(\tau) \right).
\]
Consequently, the multiplicative noise can be absorbed, and the expected excess risk is asymptotically equivalent to the sum of the deterministic bias and the label noise:
\[
    \EE[\cE(\btheta_k)] \eqsim \cS_{\nesterov}(k) + \frac{\eta^2\sigma^2}{B}\sum_{\tau=0}^{k-1}\mathcal{K}_{\nesterov}(\tau).
\]
\end{corollary}

% ==============================================================================

\subsection{Risk dynamics in the fully overdamped regime}
\label{app:nes_fully_overdamped_loss}

In the fully overdamped regime, every characteristic root is real and the deterministic bias has no oscillatory component.

\begin{proposition}[Deterministic bias, fully overdamped regime]
\label{prop:nes_deterministic_bias0}
Under Assumptions~\ref{ass:diagonal-covariance} and~\ref{ass:power-law}, in the fully overdamped regime $\eta\lambda_1\le c_0(1-\rho)^2$ with a sufficiently small universal constant $c_0>0$, we have:
\begin{equation*}
    \cS_{\nesterov}(k)
    \eqsim
    \begin{cases}
    1,&\qquad k\lesssim \frac{1-\rho}{\eta},\\
    \left(\frac{1-\rho}{\eta k}\right)^s, &\qquad k\gtrsim\frac{1-\rho}{\eta}.\\
    \end{cases}
\end{equation*}
\end{proposition}
\begin{proof}
For each coordinate, write the deterministic multiplier as
\[
    e_k^{\mathrm{det}}(j)=-h_k^{\nesterov}(j)\theta_j^*.
\]
The deterministic Nesterov coordinate recursion gives
\[
    h_{k+1}^{\nesterov}(j)
    =
    (1+\rho)(1-\eta\lambda_j)h_k^{\nesterov}(j)
    -\rho(1-\eta\lambda_j)h_{k-1}^{\nesterov}(j),
    \qquad
    h_0^{\nesterov}(j)=1,
    \qquad
    h_1^{\nesterov}(j)=1-\eta\lambda_j.
\]
Its characteristic roots are
\[
    y_{j,+},y_{j,-}
    =
    \frac{(1+\rho)(1-\eta\lambda_j)
    \pm
    \sqrt{(1+\rho)^2(1-\eta\lambda_j)^2-4\rho(1-\eta\lambda_j)}}{2}.
\]
In the fully overdamped regime, \(\eta\lambda_j\le c_0(1-\rho)^2\) for every \(j\).  To locate the two roots, evaluate the characteristic polynomial at \(1-u\).  A direct expansion gives
\[
\begin{aligned}
    &(1-u)^2
    -(1+\rho)(1-\eta\lambda_j)(1-u)
    +\rho(1-\eta\lambda_j) \\
    &\qquad=
    \eta\lambda_j
    -\bigl(1-\rho+(1+\rho)\eta\lambda_j\bigr)u
    +u^2.
\end{aligned}
\]
Taking \(c_0>0\) sufficiently small, the two evaluations
\[
    \eta\lambda_j
    -\bigl(1-\rho+(1+\rho)\eta\lambda_j\bigr)
    \frac{\eta\lambda_j}{2(1-\rho)}
    +
    \left(\frac{\eta\lambda_j}{2(1-\rho)}\right)^2
    >0
\]
and
\[
    \eta\lambda_j
    -\bigl(1-\rho+(1+\rho)\eta\lambda_j\bigr)
    \frac{2\eta\lambda_j}{1-\rho}
    +
    \left(\frac{2\eta\lambda_j}{1-\rho}\right)^2
    <0
\]
hold uniformly for all \(0<\eta\lambda_j\le c_0(1-\rho)^2\).  Hence the slow root satisfies
\[
    1-\frac{2\eta\lambda_j}{1-\rho}
    \le
    y_{j,+}
    \le
    1-\frac{\eta\lambda_j}{2(1-\rho)}.
\]
Moreover,
\[
    y_{j,+}y_{j,-}=\rho(1-\eta\lambda_j).
\]
Using the lower bound on \(y_{j,+}\), and decreasing \(c_0\) if necessary, the fast root obeys
\[
    0\le y_{j,-}
    \le
    1-\frac{1-\rho}{2}.
\]

The solution of the two-point initial value problem is
\[
    h_k^{\nesterov}(j)
    =
    \frac{1-\eta\lambda_j-y_{j,-}}{y_{j,+}-y_{j,-}}y_{j,+}^{k}
    +
    \frac{y_{j,+}-(1-\eta\lambda_j)}{y_{j,+}-y_{j,-}}y_{j,-}^{k}.
\]
The root bounds above imply
\[
    \frac12
    \le
    \frac{1-\eta\lambda_j-y_{j,-}}{y_{j,+}-y_{j,-}}
    \le
    2,
    \qquad
    \left|
    \frac{y_{j,+}-(1-\eta\lambda_j)}{y_{j,+}-y_{j,-}}
    \right|
    \le
    Cc_0.
\]
Indeed, \(y_{j,+}-y_{j,-}\eqsim1-\rho\), while \(1-\eta\lambda_j-y_{j,-}\eqsim1-\rho\), and
\[
    \left|y_{j,+}-(1-\eta\lambda_j)\right|
    \lesssim
    \frac{\eta\lambda_j}{1-\rho}
    \le
    c_0(1-\rho).
\]
Thus the slow component has a uniformly positive coefficient and the fast component has a coefficient of order \(c_0\).

To evaluate the deterministic bias $\cS_{\nesterov}(k) = \frac12\sum_{j\ge1}\lambda_j(\theta_j^*)^2 \left(h_k^{\nesterov}(j)\right)^2$, we divide the time $k$ into three successive stages.

\textbf{Stage 1: $k \lesssim \frac{1}{1-\rho}$.} \\
In this extremely short-time regime, for the leading modes ($j=\mathcal{O}(1)$), since $\eta\lambda_j \le c_0(1-\rho)^2$, the slow root satisfies $y_{j,+} \ge 1 - 2c_0(1-\rho)$. For $k \lesssim \frac{1}{1-\rho}$, we have
\[
    y_{j,+}^k \ge \left(1 - 2c_0(1-\rho)\right)^{\frac{C}{1-\rho}} \eqsim 1.
\]
Since the dynamics are unconditionally stable ($h_k^{\nesterov}(j) \le 1$), it follows that $h_k^{\nesterov}(j) \eqsim 1$ for the leading modes. Consequently, the bias is dominated by these non-decayed modes:
\[
    \cS_{\nesterov}(k) \eqsim \sum_{j\ge1}\lambda_j(\theta_j^*)^2 \eqsim 1.
\]

\textbf{Stage 2: $\frac{1}{1-\rho} \lesssim k \lesssim \frac{1-\rho}{\eta}$.} \\
For $k \gtrsim \frac{1}{1-\rho}$, the fast root decays exponentially since $y_{j,-}^k \le \left(1-\frac{1-\rho}{2}\right)^k \ll 1$. The multiplier is  dominated by the slow root:
\[
    h_k^{\nesterov}(j) \simeq \frac{1-\eta\lambda_j - y_{j,-}}{y_{j,+} - y_{j,-}} y_{j,+}^k \simeq y_{j,+}^k.
\]
For the leading modes with $\lambda_j = \Theta(1)$, within the time window $k \lesssim \frac{1-\rho}{\eta}$, the exponent is bounded by:
\[
    y_{j,+}^k \eqsim \left(1 - c\frac{\eta\lambda_j}{1-\rho}\right)^k \eqsim \exp\left(-c\frac{\eta\lambda_j}{1-\rho}k\right) \eqsim 1,
\]
where we used the condition $\frac{\eta\lambda_j}{1-\rho}k \lesssim \lambda_j \le \lambda_1 = \mathcal{O}(1)$. Thus, the leading modes have not yet experienced significant exponential decay, preserving $h_k^{\nesterov}(j) \eqsim 1$. As a result, the deterministic bias remains unchanged:
\[
    \cS_{\nesterov}(k) \eqsim 1.
\]

\textbf{Stage 3: $k \gtrsim \frac{1-\rho}{\eta}$.} \\
In this long-time regime, the exponential decay of the slow root becomes significant across the spectrum. For the slow root, the preceding bounds give
\[
    \exp\left(-C\frac{\eta\lambda_j}{1-\rho}k\right)
    \le
    y_{j,+}^{k}
    \le
    \exp\left(-c\frac{\eta\lambda_j}{1-\rho}k\right).
\]
For the fast root, since \(\eta\lambda_j/(1-\rho)\le c_0(1-\rho)\), the ratio of the fast component to the slow component is bounded by
\[
    Cc_0
    \exp\left(-c(1-\rho)k+C\frac{\eta\lambda_j}{1-\rho}k\right)
    \le
    Cc_0\exp\left(-c(1-\rho)k\right),
\]
after decreasing \(c_0\).  Therefore, for all \(k\ge C(1-\rho)^{-1}\), the fast component is at most a fixed fraction of the slow component.  Consequently there exist constants \(c,C>0\) such that
\[
    c\exp\left(-C\frac{\eta\lambda_j}{1-\rho}k\right)
    \le
    \left|h_k^{\nesterov}(j)\right|
    \le
    C\exp\left(-c\frac{\eta\lambda_j}{1-\rho}k\right).
\]

Substituting this coordinate estimate into the deterministic bias gives
\[
    c\sum_{j\ge1}\lambda_j(\theta_j^*)^2
    \exp\left(-C\frac{\eta\lambda_j}{1-\rho}k\right)
    \le
    \cS_{\nesterov}(k)
    \le
    C\sum_{j\ge1}\lambda_j(\theta_j^*)^2
    \exp\left(-c\frac{\eta\lambda_j}{1-\rho}k\right).
\]
By Assumption~\ref{ass:power-law},
\[
    \lambda_j(\theta_j^*)^2\eqsim j^{-1-\beta s},
    \qquad
    \lambda_j\eqsim j^{-\beta}.
\]
It remains to estimate
\[
    \sum_{j\ge1}j^{-1-\beta s}
    \exp\left(-c\frac{\eta k}{1-\rho}j^{-\beta}\right).
\]
The condition $k \gtrsim \frac{1-\rho}{\eta}$ gives \(\frac{\eta k}{1-\rho}\ge C\). \\
For the lower bound, take the tail
$j\ge
    \left\lceil
    \left(\frac{\eta k}{1-\rho}\right)^{1/\beta}
    \right\rceil.$
On this tail the exponential factor is bounded below by a positive constant, hence
\[
    \sum_{j\ge1}j^{-1-\beta s}
    \exp\left(-C\frac{\eta k}{1-\rho}j^{-\beta}\right)
    \ge
    c\left(\frac{\eta k}{1-\rho}\right)^{-s}.
\]
For the upper bound, an integral comparison gives
\[
\begin{aligned}
    \sum_{j\ge1}j^{-1-\beta s}
    \exp\left(-c\frac{\eta k}{1-\rho}j^{-\beta}\right)
    &\le
    C\int_1^\infty x^{-1-\beta s}
    \exp\left(-c\frac{\eta k}{1-\rho}x^{-\beta}\right)\,dx \\
    &\le
    C\left(\frac{\eta k}{1-\rho}\right)^{-s}.
\end{aligned}
\]
Combining the spectral upper and lower bounds yields
\[
    \cS_{\nesterov}(k)
    \eqsim
    \left(\frac{1-\rho}{\eta k}\right)^s.
\]
Synthesizing the three stages completes the proof.
\end{proof}

\begin{theorem}[Scaling law, fully overdamped regime]
\label{thm:nes_scaling_sgd}
Under Assumptions~\ref{ass:diagonal-covariance} and~\ref{ass:power-law},  assume the stability condition $\eta \lesssim \min\{1,B^\beta(1-\rho)\}$ holds. For the fully overdamped regime $\eta\lesssim(1-\rho)^2$ with $k \gtrsim \frac{1-\rho}{\eta}$, the expected excess risk satisfies:
\[
    \mathbb{E}\mathcal E_k
    \eqsim
    \left(\frac{1-\rho}{\eta k}\right)^s
    +
    \frac{\sigma^2\eta}{B(1-\rho)}.
\]
\end{theorem}
\begin{proof}
By Corollary~\ref{cor:nes_noise_absorption}, the expected risk is completely determined by the deterministic bias and the label noise:
\[
    \EE[\cE(\btheta_k)] \eqsim \cS_{\nesterov}(k) + \frac{\eta^2\sigma^2}{B}\sum_{\tau=0}^{k-1}\mathcal{K}_{\nesterov}(\tau).
\]
For $k \gtrsim \frac{1-\rho}{\eta}$, we substitute the deterministic bias from Proposition~\ref{prop:nes_deterministic_bias0}:
\[
    \cS_{\nesterov}(k) \eqsim \left(\frac{1-\rho}{\eta k}\right)^s.
\]
We substitute the label noise floor from Lemma~\ref{lem:nes_label_noise}. In the fully overdamped regime $\eta \lesssim (1-\rho)^2$, we have $\frac{\eta}{1-\rho} \lesssim 1-\rho \ll 1$. Since $\beta > 1$, the minimum is attained by the linear term:
\[
    \frac{\eta^2\sigma^2}{B}\sum_{\tau=0}^{k-1}\mathcal{K}_{\nesterov}(\tau) \eqsim \frac{\sigma^2}{B} \min\left\{ \left(\frac{\eta}{1-\rho}\right)^{\frac1\beta}, \frac{\eta}{1-\rho} \right\} = \frac{\sigma^2\eta}{B(1-\rho)}.
\]
Summing these two terms directly yields the exact scaling law.
\end{proof}



\subsection{Risk dynamics in the mixed-damping regime}
\label{app:nes_mixed_loss}


In the mixed-damping regime, the spectrum is partitioned into underdamped and overdamped modes by the crossover index $J_{\max} \eqsim \left(\frac{\eta}{(1-\rho)^2}\right)^{1/\beta}$. Consequently, the deterministic bias decomposes into an overdamped polynomial tail and an underdamped oscillatory transient, denoted by $\cS_{\mathrm{over}}(k)$ and $\cS_{\mathrm{under}}(k)$. Due to destructive interference among the underdamped modes, $\cS_{\mathrm{under}}(k)$ lacks a strictly positive pointwise lower bound. 



\begin{proposition}[Deterministic bias, mixed-damping regime]
\label{prop:nes_deterministic_bias}
Under Assumptions~\ref{ass:diagonal-covariance} and~\ref{ass:power-law}, for the mixed-damping regime \(\eta\gtrsim(1-\rho)^2\), the deterministic bias decomposes into $\cS_{\nesterov}(k) = \cS_{\mathrm{over}}(k) + \cS_{\mathrm{under}}(k)$, where the overdamped tail satisfies the two-sided bounds:
\begin{equation*}
    \cS_{\mathrm{over}}(k)
    \eqsim
    \begin{cases}
    \left(\frac{(1-\rho)^2}{\eta}\right)^s, &\qquad k\lesssim\frac{1}{1-\rho},\\
    \left(\frac{1-\rho}{\eta k}\right)^s, &\qquad k\gtrsim\frac{1}{1-\rho},
    \end{cases}
\end{equation*}
and $\cS_{\mathrm{under}}(k)$ represents the underdamped transient whose properties are formally established in Lemma~\ref{lem:nes_underdamped_transient_properties}.
\end{proposition}

\begin{proof}
For each coordinate, write \(e_k^{\mathrm{det}}(j)=-h_k^{\nesterov}(j)\theta_j^*\). The multiplier satisfies
\[
    h_{k+1}^{\nesterov}(j)
    =
    (1+\rho)(1-\eta\lambda_j)h_k^{\nesterov}(j)
    -\rho(1-\eta\lambda_j)h_{k-1}^{\nesterov}(j),
    \quad
    h_0^{\nesterov}(j)=1,\ 
    h_1^{\nesterov}(j)=1-\eta\lambda_j.
\]
The roots are
\[
    y_{j,+},y_{j,-}
    =
    \frac{(1+\rho)(1-\eta\lambda_j)
    \pm
    \sqrt{(1+\rho)^2(1-\eta\lambda_j)^2-4\rho(1-\eta\lambda_j)}}{2}.
\]
The discriminant is negative precisely when
\[
    \eta\lambda_j>\frac{(1-\rho)^2}{(1+\rho)^2}.
\]
Thus the comparison between $\eta$ and $(1-\rho)^2$ determines whether a coordinate is overdamped or underdamped. In the mixed-damping regime $\eta \gtrsim (1-\rho)^2$, the crossover index is defined by $\eta\lambda_{J_{\max}} \eqsim (1-\rho)^2$. We decompose the deterministic bias into the overdamped tail ($j > J_{\max}$) and the underdamped leading block ($j \le J_{\max}$).

For overdamped coordinates in the stable small-step range, the root comparison used in Proposition~\ref{prop:nes_deterministic_bias0} gives constants $c_1,C_1>0$ such that
\begin{equation*}
    |h_k^{\nesterov}(j)|
    \le
    C_1\exp\left(-c_1\frac{\eta\lambda_j}{1-\rho}k\right).
\end{equation*}
Moreover, if $\eta\lambda_j\le c_1(1-\rho)^2$ and $k\ge C_1(1-\rho)^{-1}$, the slow root dominates the fast root and
\begin{equation*}
    |h_k^{\nesterov}(j)|
    \ge
    c_1\exp\left(-C_1\frac{\eta\lambda_j}{1-\rho}k\right).
\end{equation*}

For $k \lesssim \frac{1}{1-\rho}$, the exponent satisfies $\frac{\eta\lambda_j}{1-\rho}k \lesssim \frac{(1-\rho)^2}{1-\rho}\frac{1}{1-\rho} = 1$. The exponential decay is bounded away from zero, so $|h_k^{\nesterov}(j)| \eqsim 1$. The sum is dominated by its lower limit at the crossover:
\[
    \cS_{\mathrm{over}}(k) = \sum_{j > J_{\max}} \lambda_j(\theta_j^*)^2 h_k^{\nesterov}(j)^2 \eqsim \sum_{j > J_{\max}} j^{-1-\beta s} \eqsim (J_{\max})^{-\beta s} \eqsim \left(\frac{(1-\rho)^2}{\eta}\right)^s.
\]

For $k \gtrsim \frac{1}{1-\rho}$, applying $\lambda_j\eqsim j^{-\beta}$, the integral comparison yields the strict upper bound:
\begin{equation*}
    \cS_{\mathrm{over}}(k)
    \le
    C\sum_{j > J_{\max}}j^{-1-\beta s}
    \exp\left(-2c_1\frac{\eta k}{1-\rho}j^{-\beta}\right)
    \le
    C\left(\frac{1-\rho}{\eta k}\right)^s .
\end{equation*}
By restricting the sum to $j \ge (C\frac{\eta k}{1-\rho})^{1/\beta}$, we obtain the matching lower bound $c\left(\frac{1-\rho}{\eta k}\right)^s$.

The remaining deterministic bias comes from the underdamped modes, defined as $\cS_{\mathrm{under}}(k) := \frac{1}{2} \sum_{j \le J_{\max}} \lambda_j (\theta_j^*)^2 h_k^{\nesterov}(j)^2$, which is analyzed in Lemma~\ref{lem:nes_underdamped_transient_properties}.
\end{proof}


\begin{lemma}[Properties of the underdamped transient]
\label{lem:nes_underdamped_transient_properties}
Let $\cS_{\mathrm{under}}(k)$ denote the deterministic bias contributed by the underdamped modes ($j \le J_{\max}$). Under the conditions of Proposition~\ref{prop:nes_deterministic_bias}, define the piecewise envelope $\mathcal{E}_{\mathrm{env}}(k)$ by:
\begin{equation*}
    \mathcal{E}_{\mathrm{env}}(k) := 
    \begin{cases}
    \rho^k, &\qquad k\lesssim\frac{1}{\eta},\\[2mm]
    \dfrac{\rho^k}{(\eta k)^s}, &\qquad \dfrac{1}{\eta} \lesssim k \lesssim \dfrac{1}{(1-\rho)^2},\\[3mm]
    \dfrac{1}{\eta k}\left(\dfrac{(1-\rho)^2}{\eta}\right)^{s-1}p_{\nesterov}^{\,k}, &\qquad \dfrac{1}{(1-\rho)^2} \lesssim k \lesssim \eta^{1/\beta}(1-\rho)^{-2-\frac{2}{\beta}},\\[3mm]
    \left(\dfrac{(1-\rho)^2}{\eta}\right)^{s+\frac{1}{\beta}}p_{\nesterov}^{\,k}, &\qquad k \gtrsim \eta^{1/\beta}(1-\rho)^{-2-\frac{2}{\beta}},
    \end{cases}
\end{equation*}
where $p_{\nesterov} := \rho\left(1 - \frac{(1-\rho)^2}{(1+\rho)^2}\right) = \frac{4\rho^2}{(1+\rho)^2} < \rho$. Then for every step $k \ge 0$, $\cS_{\mathrm{under}}(k)$ satisfies:
\begin{equation*}
    0 \le \cS_{\mathrm{under}}(k) \le C \mathcal{E}_{\mathrm{env}}(k).
\end{equation*}
\end{lemma}

\begin{proof}
For $j \le J_{\max}$, write the two characteristic roots as:
\[
    \sqrt{\rho(1-\eta\lambda_j)}e^{i\omega_j},
    \quad
    \sqrt{\rho(1-\eta\lambda_j)}e^{-i\omega_j},
    \qquad \text{where} \quad
    \cos \omega_j=\frac{(1+\rho)\sqrt{1-\eta\lambda_j}}{2\sqrt \rho}.
\]
Solving the initial conditions gives the exact identity for the coordinate multiplier:
\begin{equation*}
    h_k^{\nesterov}(j)
    =\bigl(\rho(1-\eta\lambda_j)\bigr)^{k/2}
    \left[
        \cos(k\omega_j)
        +\frac{(1-\rho)\sqrt{1-\eta\lambda_j}}{2\sqrt{\rho}\sin \omega_j}
        \sin(k\omega_j)
    \right].
\end{equation*}
Set $A_j:=\frac{(1-\rho)\sqrt{1-\eta\lambda_j}}{2\sqrt{\rho}\sin \omega_j}$ and define $\tan\psi_j = A_j$. The squared multiplier is bounded by the envelope:
\[
    \bigl(h_k^{\nesterov}(j)\bigr)^2 = \bigl(\rho(1-\eta\lambda_j)\bigr)^k (1+A_j^2) \cos^2(k\omega_j - \psi_j).
\]
On the separated underdamped band $\eta\lambda_j \ge \frac{(1-\rho)^2}{(1+\rho)^2}$, we have $4\rho\sin^2\omega_j \eqsim \eta\lambda_j$, which guarantees $|A_j| \le C$.

Defining the positive spectral weights:
\[
    w_j(k) := \frac{1}{2}\lambda_j (\theta_j^*)^2 \bigl(\rho(1-\eta\lambda_j)\bigr)^k (1+A_j^2),
\]
the total underdamped bias decomposes into:
\[
    \cS_{\mathrm{under}}(k) = \sum_{j \le J_{\max}} w_j(k) \cos^2(k\omega_j - \psi_j).
\]

\medskip
\noindent\textbf{Evaluation of the static mass envelope $M_w(k)$.}\\

We define the static mass:
\[
    M_w(k) := \frac{1}{2}\sum_{j=1}^{J_{\max}} \lambda_j (\theta_j^*)^2 \bigl(\rho(1-\eta\lambda_j)\bigr)^k (1+A_j^2).
\]
Since $|A_j| \le C$, we have $1+A_j^2 \eqsim 1$. Substituting $\lambda_j (\theta_j^*)^2 \eqsim j^{-1-\beta s}$ and $\lambda_j \eqsim j^{-\beta}$, the static mass is asymptotically equivalent to:
\[
    M_w(k) \eqsim \rho^k \sum_{j=1}^{J_{\max}} j^{-1-\beta s} (1-\eta\lambda_j)^k.
\]
We evaluate $M_w(k)$ across four distinct temporal stages:

\medskip
\noindent\textbf{Stage 1 ($k \lesssim \frac{1}{\eta}$):}
For all modes $j \in [1, J_{\max}]$, we have $\eta\lambda_j \le \eta\lambda_1 \lesssim \eta$. Since $k \lesssim \frac{1}{\eta}$, the product satisfies $k \eta \lambda_j \lesssim 1$, which implies:
\[
    (1-\eta\lambda_j)^k \eqsim 1, \qquad \forall\, j \in [1, J_{\max}].
\]
Consequently, the summation is dominated by the leading indices ($j \sim 1$):
\[
    M_w(k) \eqsim \rho^k \sum_{j=1}^{J_{\max}} j^{-1-\beta s} \eqsim \rho^k \sum_{j=1}^{\infty} j^{-1-\beta s} \eqsim \rho^k.
\]

\medskip
\noindent\textbf{Stage 2 ($\frac{1}{\eta} \lesssim k \lesssim \frac{1}{(1-\rho)^2}$):}
In this regime, $k\eta \gtrsim 1$, and $(1-\eta\lambda_j)^k \eqsim \exp\bigl(-c k \eta j^{-\beta}\bigr)$ suppresses the leading modes. The dominant spectral mass is concentrated at the continuous peak $j_k \eqsim (\eta k)^{1/\beta}$. Since $k \lesssim \frac{1}{(1-\rho)^2}$, this peak lies strictly within the summation range:
\[
    j_k \eqsim (\eta k)^{1/\beta} \lesssim \left(\frac{\eta}{(1-\rho)^2}\right)^{1/\beta} \eqsim J_{\max}.
\]
Applying Euler--Maclaurin integral comparison with the change of variables $u = c k \eta x^{-\beta}$:
\begin{align*}
    M_w(k) 
    &\eqsim \rho^k \int_1^{J_{\max}} x^{-1-\beta s} \exp\bigl(-c k \eta x^{-\beta}\bigr) \, dx \\
    &= \frac{\rho^k}{\beta (c k \eta)^s} \int_{c k \eta J_{\max}^{-\beta}}^{c k \eta} u^{s-1} e^{-u} \, du.
\end{align*}
Since the integration limits satisfy $c k \eta J_{\max}^{-\beta} \eqsim k(1-\rho)^2 \lesssim 1$ and $c k \eta \gtrsim 1$, the integral is bounded above and below by universal positive constants:
\[
    \int_{c k \eta J_{\max}^{-\beta}}^{c k \eta} u^{s-1} e^{-u} \, du \eqsim \int_0^{\infty} u^{s-1} e^{-u} \, du = \Gamma(s) \eqsim 1.
\]
Therefore, we obtain:
\[
    M_w(k) \eqsim \frac{\rho^k}{(\eta k)^s}.
\]

\medskip
\noindent\textbf{Stage 3 \& 4 ($k \gtrsim \frac{1}{(1-\rho)^2}$):}
When $k \gtrsim \frac{1}{(1-\rho)^2}$, the continuous peak $(\eta k)^{1/\beta}$ moves past $J_{\max}$. Inside the summation interval $j \le J_{\max}$, the decay factor $(1-\eta\lambda_j)^k$ is slowest at the boundary $j = J_{\max}$, where $\eta\lambda_{J_{\max}} = \frac{(1-\rho)^2}{(1+\rho)^2} =: u_c$. We define the boundary decay factor:
\[
    p_{\nesterov} := \rho(1-u_c) = \rho\left(1 - \frac{(1-\rho)^2}{(1+\rho)^2}\right) = \frac{4\rho^2}{(1+\rho)^2}.
\]
Let $m := J_{\max} - j \ge 0$. Linearizing the eigenvalues near the boundary gives:
\[
    \eta\lambda_j - u_c = \eta(\lambda_{J_{\max}-m} - \lambda_{J_{\max}}) = \Delta_x m + \mathcal{O}\bigl(\eta J_{\max}^{-\beta-2} m^2\bigr),
\]
where the discrete frequency spacing is:
\[
    \Delta_x := \eta \beta J_{\max}^{-\beta-1} \eqsim \eta \left(\frac{(1-\rho)^2}{\eta}\right)^{1+\frac{1}{\beta}} = (1-\rho)^{2+\frac{2}{\beta}} \eta^{-\frac{1}{\beta}}.
\]
Factoring out the boundary coefficient $J_{\max}^{-1-\beta s} \eqsim \left(\frac{(1-\rho)^2}{\eta}\right)^{s+\frac{1}{\beta}}$, the sum is governed by the geometric series over $m$:
\begin{align*}
    M_w(k) 
    &\eqsim \left(\frac{(1-\rho)^2}{\eta}\right)^{s+\frac{1}{\beta}} p_{\nesterov}^k \sum_{m=0}^{J_{\max}-1} \exp\bigl(-c k \Delta_x m\bigr) \\
    &= \left(\frac{(1-\rho)^2}{\eta}\right)^{s+\frac{1}{\beta}} p_{\nesterov}^k \left( \frac{1 - \exp\bigl(-c k \Delta_x J_{\max}\bigr)}{1 - \exp\bigl(-c k \Delta_x\bigr)} \right).
\end{align*}
Since $k \Delta_x J_{\max} \eqsim k u_c \eqsim k(1-\rho)^2 \gtrsim 1$, the numerator satisfies $1 - \exp\bigl(-c k \Delta_x J_{\max}\bigr) \eqsim 1$, yielding:
\begin{equation}
\label{eq:nes_Mw_geometric_form}
    M_w(k) \eqsim \left(\frac{(1-\rho)^2}{\eta}\right)^{s+\frac{1}{\beta}} p_{\nesterov}^k \left( \frac{1}{1 - \exp\bigl(-c k \Delta_x\bigr)} \right).
\end{equation}

\medskip
\noindent\textbf{Stage 3 ($\frac{1}{(1-\rho)^2} \lesssim k \lesssim \eta^{1/\beta}(1-\rho)^{-2-\frac{2}{\beta}}$):}
In this regime, the exponent in the denominator of \eqref{eq:nes_Mw_geometric_form} is small: $k \Delta_x = k (1-\rho)^{2+\frac{2}{\beta}} \eta^{-\frac{1}{\beta}} \lesssim 1$. Using the Taylor expansion $1 - e^{-u} \eqsim u$ for $u \in (0, 1]$:
\begin{align*}
    M_w(k) 
    &\eqsim \left(\frac{(1-\rho)^2}{\eta}\right)^{s+\frac{1}{\beta}} p_{\nesterov}^k \left( \frac{1}{c k \Delta_x} \right) \\
    &\eqsim \left(\frac{(1-\rho)^2}{\eta}\right)^{s+\frac{1}{\beta}} p_{\nesterov}^k \left( \frac{\eta^{\frac{1}{\beta}}}{k (1-\rho)^{2+\frac{2}{\beta}}} \right) \\
    &= \frac{1}{\eta k} \left(\frac{(1-\rho)^2}{\eta}\right)^{s-1} p_{\nesterov}^{\,k}.
\end{align*}

\medskip
\noindent\textbf{Stage 4 ($k \gtrsim \eta^{1/\beta}(1-\rho)^{-2-\frac{2}{\beta}}$):}
Here, the exponent is large: $k \Delta_x = k (1-\rho)^{2+\frac{2}{\beta}} \eta^{-\frac{1}{\beta}} \gtrsim 1$. Thus, the denominator in \eqref{eq:nes_Mw_geometric_form} is bounded away from zero: $1 - \exp\bigl(-c k \Delta_x\bigr) \eqsim 1$. The discrete mode $j = J_{\max}$ ($m=0$) dominates the sum:
\[
    M_w(k) \eqsim \left(\frac{(1-\rho)^2}{\eta}\right)^{s+\frac{1}{\beta}} p_{\nesterov}^{\,k}.
\]

\medskip
\noindent Combining the four stages establishes that $M_w(k) \eqsim \mathcal{E}_{\mathrm{env}}(k)$ uniformly across all $k \ge 0$.

Since $\cos^2(\cdot) \le 1$, we immediately obtain the pointwise upper bound:
\[
    \cS_{\mathrm{under}}(k) \le \sum_{j \le J_{\max}} w_j(k) = M_w(k) \le C \mathcal{E}_{\mathrm{env}}(k).
\]
The lower bound is trivially zero due to the potential for destructive interference of the oscillatory terms at specific iterations.
\end{proof}


\begin{theorem}[Scaling law, mixed-damping regime]
\label{thm:nes_two_sided_scaling}
Under Assumptions~\ref{ass:diagonal-covariance} and~\ref{ass:power-law}, assuming the stability condition \(\eta\lesssim\min\{1, B^\beta(1-\rho)\}\), for the mixed-damping regime $\eta\gtrsim(1-\rho)^2$ with $k \gtrsim \max\left\{\frac{1}{1-\rho}, \frac{1-\rho}{\eta}\right\}$, the expected excess risk satisfies:
\begin{equation}
\label{eq:nes_mixed_scaling_law}
    \EE[\cE(\btheta_k)] \eqsim \cS_{\mathrm{under}}(k) + \left(\frac{1-\rho}{\eta k}\right)^s + \frac{\sigma^2}{B} \min\left\{ \left(\frac{\eta}{1-\rho}\right)^{\frac1\beta}, \frac{\eta}{1-\rho} \right\},
\end{equation}
where the properties of the underdamped transient $\cS_{\mathrm{under}}(k)$ are established in Lemma~\ref{lem:nes_underdamped_transient_properties}.
\end{theorem}

\begin{proof}
By Corollary~\ref{cor:nes_noise_absorption}, the multiplicative noise is absorbed by the additive label noise floor, yielding:
\begin{equation}
\label{eq:nes_total_risk_mixed}
    \EE[\cE(\btheta_k)] 
    \eqsim 
    \cS_{\mathrm{under}}(k) + \cS_{\mathrm{over}}(k) + \frac{\eta^2\sigma^2}{B}\sum_{\tau=0}^{k-1}\mathcal{K}_{\nesterov}(\tau).
\end{equation}

By Lemma~\ref{lem:nes_label_noise}, for $k \gtrsim \max\left\{\frac{1}{1-\rho}, \frac{1-\rho}{\eta}\right\}$, the label noise summation captures the saturated floor:
\begin{equation*}
    \frac{\eta^2\sigma^2}{B}\sum_{\tau=0}^{k-1}\mathcal{K}_{\nesterov}(\tau) 
    \eqsim 
    \frac{\sigma^2}{B} \min\left\{ \left(\frac{\eta}{1-\rho}\right)^{\frac1\beta}, \frac{\eta}{1-\rho} \right\}.
\end{equation*}

By Proposition~\ref{prop:nes_deterministic_bias}, since $k \gtrsim \frac{1}{1-\rho}$, the overdamped tail reaches its terminal polynomial decay phase:
\begin{equation*}
    \cS_{\mathrm{over}}(k) \eqsim \left(\frac{1-\rho}{\eta k}\right)^s.
\end{equation*}

Substituting these components explicitly into \eqref{eq:nes_total_risk_mixed} establishes the equivalence \eqref{eq:nes_mixed_scaling_law}.
\end{proof}





% % ==============================================================================

% \subsection{Risk dynamics in the mixed regime}
% \label{app:nes_mixed_loss}

% When the leading spectral modes are underdamped, the loss contains an exponentially damped edge contribution in addition to the polynomial overdamped tail. 

% Define the Nesterov edge damping factor:
% \[
%     p_{\nesterov}
%     :=
%     \rho\left(1-\frac{(1-\rho)^2}{(1+\rho)^2}\right)
%     =
%     \frac{4\rho^2}{(1+\rho)^2}.
% \]
% It strictly satisfies $p_{\nesterov} = 1-(1-\rho)+\mathcal{O}((1-\rho)^2) < \rho$, ensuring that the underdamped edge contribution decays exponentially.


% \begin{proposition}[Deterministic bias, mixed regime]
% \label{prop:nes_deterministic_bias}
% Under Assumptions~\ref{ass:diagonal-covariance} and~\ref{ass:power-law}, assume the local stability condition $\eta\lesssim 1$. For the mixed regime \(\eta\gtrsim(1-\rho)^2\), the deterministic bias decomposes into $\cS_{\nesterov}(k) = \cS_{\mathrm{over}}(k) + \cS_{\mathrm{under}}(k)$, where:
% \begin{equation*}
%     \cS_{\mathrm{over}}(k)
%     \eqsim
%     \begin{cases}
%     \left(\frac{(1-\rho)^2}{\eta}\right)^s, &\qquad k\lesssim\frac{1}{1-\rho},\\
%     \left(\frac{1-\rho}{\eta k}\right)^s, &\qquad k\gtrsim\frac{1}{1-\rho},
%     \end{cases}
% \end{equation*}
% and
% \begin{equation*}
%     \cS_{\mathrm{under}}(k)
%     \eqsim
%     \begin{cases}
%     \rho^k, &\qquad k\lesssim\frac{1}{\eta},\\
%     \frac{\rho^k}{(\eta k)^s}, &\qquad \frac{1}{\eta} \lesssim k \lesssim \frac{1}{(1-\rho)^2},\\
%     \frac{1}{\eta k}\left(\frac{(1-\rho)^2}{\eta}\right)^{s-1}p_{\nesterov}^{\,k}, &\qquad \frac{1}{(1-\rho)^2} \lesssim k \lesssim \eta^{1/\beta}(1-\rho)^{-2-\frac{2}{\beta}},\\
%     \left(\frac{(1-\rho)^2}{\eta}\right)^{s+\frac{1}{\beta}}p_{\nesterov}^{\,k}, &\qquad k \gtrsim \eta^{1/\beta}(1-\rho)^{-2-\frac{2}{\beta}}.
%     \end{cases}
% \end{equation*}
% \end{proposition}

% \begin{proof}
% For each coordinate, write \(e_k^{\mathrm{det}}(j)=-h_k^{\nesterov}(j)\theta_j^*\). The multiplier satisfies
% \[
%     h_{k+1}^{\nesterov}(j)
%     =
%     (1+\rho)(1-\eta\lambda_j)h_k^{\nesterov}(j)
%     -\rho(1-\eta\lambda_j)h_{k-1}^{\nesterov}(j),
%     \qquad
%     h_0^{\nesterov}(j)=1,
%     \qquad
%     h_1^{\nesterov}(j)=1-\eta\lambda_j.
% \]
% The roots are
% \[
%     y_{j,+},y_{j,-}
%     =
%     \frac{(1+\rho)(1-\eta\lambda_j)
%     \pm
%     \sqrt{(1+\rho)^2(1-\eta\lambda_j)^2-4\rho(1-\eta\lambda_j)}}{2}.
% \]
% The discriminant is negative precisely when
% \[
%     (1-\sqrt\rho)^2<\eta\lambda_j<(1+\sqrt\rho)^2.
% \]
% Thus the comparison between $\eta$ and $(1-\rho)^2$ determines whether a coordinate is overdamped or underdamped. In the mixed regime $\eta \gtrsim (1-\rho)^2$, the crossover index is defined by $\eta\lambda_{J_{\max}} \eqsim (1-\rho)^2$. We decompose the deterministic bias into the overdamped tail ($j > J_{\max}$) and the underdamped leading block ($j \le J_{\max}$).

% \textbf{1. Overdamped Tail ($\cS_{\mathrm{over}}$):} \\
% For overdamped coordinates in the stable small-step range, the root comparison used in Proposition~\ref{prop:nes_deterministic_bias0} gives constants $c_1,C_1>0$ such that
% \begin{equation*}
%     |h_k^{\nesterov}(j)|
%     \le
%     C_1\exp\left(-c_1\frac{\eta\lambda_j}{1-\rho}k\right).
% \end{equation*}
% Moreover, if $\eta\lambda_j\le c_1(1-\rho)^2$ and $k\ge C_1(1-\rho)^{-1}$, the slow root dominates the fast root and
% \begin{equation*}
%     |h_k^{\nesterov}(j)|
%     \ge
%     c_1\exp\left(-C_1\frac{\eta\lambda_j}{1-\rho}k\right).
% \end{equation*}

% For $k \lesssim \frac{1}{1-\rho}$, the exponent satisfies $\frac{\eta\lambda_j}{1-\rho}k \lesssim \frac{(1-\rho)^2}{1-\rho}\frac{1}{1-\rho} = 1$. The exponential decay is bounded away from zero, so $|h_k^{\nesterov}(j)| \eqsim 1$. The sum is strictly dominated by its lower limit at the crossover:
% \[
%     \cS_{\mathrm{over}}(k) = \sum_{j > J_{\max}} \lambda_j(\theta_j^*)^2 h_k^{\nesterov}(j)^2 \eqsim \sum_{j > J_{\max}} j^{-1-\beta s} \eqsim (J_{\max})^{-\beta s} \eqsim \left(\frac{(1-\rho)^2}{\eta}\right)^s.
% \]

% For $k \gtrsim \frac{1}{1-\rho}$, applying $\lambda_j\eqsim j^{-\beta}$, the integral comparison yields the strict upper bound:
% \begin{equation*}
%     \cS_{\mathrm{over}}(k)
%     \le
%     C\sum_{j > J_{\max}}j^{-1-\beta s}
%     \exp\left(-2c_1\frac{\eta k}{1-\rho}j^{-\beta}\right)
%     \le
%     C\left(\frac{1-\rho}{\eta k}\right)^s .
% \end{equation*}
% By restricting the sum to $j \ge (C\frac{\eta k}{1-\rho})^{1/\beta}$, we obtain the matching lower bound $c\left(\frac{1-\rho}{\eta k}\right)^s$.

% \textbf{2. Underdamped Leading Block ($\cS_{\mathrm{under}}$):} \\
% For $j \le J_{\max}$, write the two roots as
% \[
%     \sqrt{\rho(1-\eta\lambda_j)}e^{i\omega_j},
%     \quad
%     \sqrt{\rho(1-\eta\lambda_j)}e^{-i\omega_j},
%     \qquad \text{where} \quad
%     \cos \omega_j=\frac{(1+\rho)\sqrt{1-\eta\lambda_j}}{2\sqrt \rho}.
% \]
% Solving the initial conditions gives the exact identity
% \begin{equation*}
%     h_k^{\nesterov}(j)
%     =\bigl(\rho(1-\eta\lambda_j)\bigr)^{k/2}
%     \left[
%         \cos(k\omega_j)
%         +\frac{(1-\rho)\sqrt{1-\eta\lambda_j}}{2\sqrt{\rho}\sin \omega_j}
%         \sin(k\omega_j)
%     \right].
% \end{equation*}
% Set $A_j:=\frac{(1-\rho)\sqrt{1-\eta\lambda_j}}{2\sqrt{\rho}\sin \omega_j}$ and define $\tan\psi_j = A_j$. The squared multiplier is bounded by the envelope $\bigl(\rho(1-\eta\lambda_j)\bigr)^k (1+A_j^2) \cos^2(k\omega_j - \psi_j)$. 

% To evaluate the underdamped contribution, we define the positive spectral weights $w_j(k) = \lambda_j (\theta_j^*)^2 \bigl(\rho(1-\eta\lambda_j)\bigr)^k (1+A_j^2)$ and decompose the cosine-squared terms:
% \[
%     \cS_{\mathrm{under}}(k) = \frac{1}{2} \sum_{j \le J_{\max}} w_j(k) + \frac{1}{2} \sum_{j \le J_{\max}} w_j(k) \cos(2k\omega_j - 2\psi_j).
% \]
% The first sum represents the stable static mass $M_w(k) = \frac{1}{2} \sum_{j \le J_{\max}} w_j(k)$. To analyze the oscillatory second sum, we view the indices as a continuous variable $x \in [1, J_{\max}]$. For the continuous oscillatory integral $I(k)$, we apply the substitution $u = \omega(x)$ mapping $[1, J_{\max}]$ to $[u_{\min}, u_{\max}]$:
% \[
%     I(k) = \int_{u_{\min}}^{u_{\max}} g(u, k) \cos(2ku - \phi(u)) \, du, \qquad \text{where } g(u, k) = w(x(u), k) \left| x'(u) \right|, \phi(u) = 2\psi(x(u)).
% \]
% By integration by parts with the differential $d\left(\sin(2ku - \phi(u))\right) = (2k - \phi'(u)) \cos(2ku - \phi(u)) \, du$:
% \[
%     I(k) = \left[ v_k(u) \sin(2ku - \phi(u)) \right]_{u_{\min}}^{u_{\max}} - \int_{u_{\min}}^{u_{\max}} v_k'(u) \sin(2ku - \phi(u)) \, du,
% \]
% where $v_k(u) = \frac{g(u, k)}{2k - \phi'(u)}$. This establishes the temporal decay $|I(k)| \le \frac{C_1}{k} \max_u |g(u, k)| = \mathcal{O}(1/k)$. Thus, for sufficiently large $k$, the oscillatory integral is strictly subsumed by the static mass: $|I(k)| \le \frac{1}{4} M_w(k)$. This guarantees the discrete sum is strictly bounded away from zero by the static part.

% We now rigorously evaluate $M_w(k)$. Asymptotically, $1+A_j^2 \eqsim 1$ everywhere except for a subdominant logarithmic singularity strictly at the continuous edge, which does not alter the polynomial/exponential macroscopic order. Thus:
% \[
%     M_w(k) \eqsim \sum_{j=1}^{J_{\max}} j^{-1-\beta s} \bigl(\rho(1-\eta\lambda_j)\bigr)^k \eqsim \sum_{j=1}^{J_{\max}} j^{-1-\beta s} p_{\nesterov}^k e^{-ckx_j},
% \]
% where $p_{\nesterov} = \rho(1-u_c)$, $u_c \eqsim (1-\rho)^2$, and the distance to the edge is $x_j = \eta\lambda_j - u_c \approx \eta\beta J_{\max}^{-\beta-1}(J_{\max} - j)$. We proceed in four temporal stages:

% \textit{Stage 1: $k \lesssim \frac{1}{\eta}$.} \\
% The exponential decay $e^{-ckx_j} \approx 1$. The sum is  dominated by the leading modes ($j \sim 1$):
% \[
%     \cS_{\mathrm{under}}(k) \eqsim \rho^k \sum_{j=1}^{J_{\max}} j^{-1-\beta s} \eqsim \rho^k.
% \]

% \textit{Stage 2: $\frac{1}{\eta} \lesssim k \lesssim \frac{1}{(1-\rho)^2}$.} \\
% The dominant bulk peak $j_k \eqsim (\eta k)^{1/\beta} \ll J_{\max}$ is fully captured inside the interval. The integral approximation perfectly applies:
% \[
%     \cS_{\mathrm{under}}(k) \eqsim \int_1^{\infty} x^{-1-\beta s} \rho^k e^{-c k \eta x^{-\beta}} dx \eqsim \rho^k (\eta k)^{-s}.
% \]

% \textit{Stage 3 \& 4: $k \gtrsim \frac{1}{(1-\rho)^2}$.} \\
% The bulk peak $j_k$ passes the boundary $J_{\max}$, so the summation is strictly dominated by the discrete tail items $j \to J_{\max}$. Let $m = J_{\max} - j \ge 0$. The distance to the edge is $x_j \approx \eta\beta J_{\max}^{-\beta-1} m \eqsim (1-\rho)^{2+2/\beta}\eta^{-1/\beta} m$. 
% Using the discrete geometric series sum for the dominant tail terms:
% \begin{align*}
%     \cS_{\mathrm{under}}(k) 
%     &\eqsim \sum_{m=0}^{J_{\max}} J_{\max}^{-1-\beta s} p_{\nesterov}^k \exp\left(-c k (1-\rho)^{2+2/\beta}\eta^{-1/\beta} m\right) \\
%     &\eqsim J_{\max}^{-1-\beta s} p_{\nesterov}^k \sum_{m=0}^{\infty} \left( \exp\left(-c k (1-\rho)^{2+2/\beta}\eta^{-1/\beta}\right) \right)^m \\
%     &\eqsim \left(\frac{(1-\rho)^2}{\eta}\right)^{s+\frac{1}{\beta}} p_{\nesterov}^k \frac{1}{1 - \exp\left(-c k (1-\rho)^{2+2/\beta}\eta^{-1/\beta}\right)}.
% \end{align*}

% \textit{Stage 3: $\frac{1}{(1-\rho)^2} \lesssim k \lesssim \eta^{1/\beta}(1-\rho)^{-2-\frac{2}{\beta}}$.} \\
% In this stage, the exponent is small, $k (1-\rho)^{2+2/\beta}\eta^{-1/\beta} \lesssim 1$, meaning the decay window covers multiple discrete points near the edge. We apply the Taylor expansion $1 - e^{-x} \eqsim x$:
% \begin{align*}
%     \cS_{\mathrm{under}}(k) 
%     &\eqsim \left(\frac{(1-\rho)^2}{\eta}\right)^{s+\frac{1}{\beta}} p_{\nesterov}^k \frac{1}{k (1-\rho)^{2+2/\beta}\eta^{-1/\beta}} \\
%     &= \left(\frac{(1-\rho)^2}{\eta}\right)^{s+\frac{1}{\beta}} p_{\nesterov}^k \frac{1}{k(1-\rho)^2} \left(\frac{\eta}{(1-\rho)^2}\right)^{\frac{1}{\beta}} \\
%     &= \frac{1}{\eta k}\left(\frac{(1-\rho)^2}{\eta}\right)^{s-1}p_{\nesterov}^{\,k}.
% \end{align*}

% \textit{Stage 4: $k \gtrsim \eta^{1/\beta}(1-\rho)^{-2-\frac{2}{\beta}}$.} \\
% In this extreme long-time regime, the exponent is large, $k (1-\rho)^{2+2/\beta}\eta^{-1/\beta} \gtrsim 1$. The exponential decay is so sharp that the single point $j = J_{\max}$ strictly dominates the entire sum. The denominator satisfies $1 - \exp\left(-c k (1-\rho)^{2+2/\beta}\eta^{-1/\beta}\right) \eqsim 1$:
% \[
%     \cS_{\mathrm{under}}(k) \eqsim \left(\frac{(1-\rho)^2}{\eta}\right)^{s+\frac{1}{\beta}}p_{\nesterov}^{\,k} \cdot 1 = \left(\frac{(1-\rho)^2}{\eta}\right)^{s+\frac{1}{\beta}}p_{\nesterov}^{\,k}.
% \]
% This concludes the proof.
% \end{proof}


% \begin{theorem}[Two-sided scaling law, mixed regime]
% \label{thm:nes_two_sided_scaling}
% Under Assumption~\ref{ass:power-law} and Proposition~\ref{prop:noise-geometry}, the conditions of Proposition~\ref{prop:nes_deterministic_bias}, and the stability condition \(\eta\le cB^\beta(1-\rho)\) with \(c>0\) sufficiently small, for the mixed regime $\eta\gtrsim(1-\rho)^2$ with $k \gtrsim \frac{1-\rho}{\eta}$, the expected excess risk strictly satisfies:
% \begin{align}
% \label{eq:nes_final_scaling_mixed}
%     \mathbb{E}\mathcal E_k \quad\eqsim\quad
%     &\min\left\{1, (\eta k)^{-s}\right\} \rho^k 
%     + \left(\frac{1-\rho}{\eta k}\right)^s 
%     + \frac{\sigma^2}{B} \min\left\{ \left(\frac{\eta}{1-\rho}\right)^{\frac1\beta}, \frac{\eta}{1-\rho} \right\}.
% \end{align}
% \end{theorem}

% \begin{proof}
% By Corollary~\ref{cor:nes_noise_absorption}, the multiplicative gradient noise is absorbed by the additive label noise floor:
% \[
%     \EE[\cE(\btheta_k)] 
%     \eqsim 
%     \cS_{\nesterov}(k) + \frac{\eta^2\sigma^2}{B}\sum_{\tau=0}^{k-1}\mathcal{K}_{\nesterov}(\tau).
% \]
% By Lemma~\ref{lem:nes_label_noise}, the domain condition $k \gtrsim \frac{1-\rho}{\eta}$ guarantees the label noise summation is saturated:
% \[
%     \frac{\eta^2\sigma^2}{B}\sum_{\tau=0}^{k-1}\mathcal{K}_{\nesterov}(\tau) 
%     \eqsim 
%     \frac{\sigma^2}{B} \min\left\{ \left(\frac{\eta}{1-\rho}\right)^{\frac1\beta}, \frac{\eta}{1-\rho} \right\}.
% \]

% It remains to establish that the true deterministic bias $\cS_{\nesterov}(k) = \cS_{\mathrm{under}}(k) + \cS_{\mathrm{over}}(k)$ from Proposition~\ref{prop:nes_deterministic_bias} matches the target analytical formula:
% \[
%     \widehat{\cS}_{\nesterov}(k) := \min\left\{1, (\eta k)^{-s}\right\} \rho^k + \left(\frac{1-\rho}{\eta k}\right)^s.
% \]
% We partition the domain $k \gtrsim \frac{1-\rho}{\eta}$ into four sequential stages. Since the relative order of $\frac{1}{\eta}$ and $\frac{1}{1-\rho}$ is unfixed, we explicitly discuss their relation in the intermediate regime.

% \textbf{Stage 1: $\frac{1-\rho}{\eta} \lesssim k \lesssim \min\left\{\frac{1}{\eta}, \frac{1}{1-\rho}\right\}$.}

% Here $\eta k \lesssim 1 \implies \min\{1, (\eta k)^{-s}\} = 1$, and $k(1-\rho) \lesssim 1 \implies \rho^k = (1-(1-\rho))^k \eqsim 1$. Thus:
% \[
%     \widehat{\cS}_{\nesterov}(k) \eqsim 1 + \left(\frac{1-\rho}{\eta k}\right)^s.
% \]
% Since $k \gtrsim \frac{1-\rho}{\eta}$, the polynomial term satisfies $\left(\frac{1-\rho}{\eta k}\right)^s \lesssim 1$, yielding $\widehat{\cS}_{\nesterov}(k) \eqsim 1$.

% Concurrently, Proposition~\ref{prop:nes_deterministic_bias} asserts $\cS_{\mathrm{under}}(k) \eqsim \rho^k \eqsim 1$ and $\cS_{\mathrm{over}}(k) \eqsim \left(\frac{(1-\rho)^2}{\eta}\right)^s$. Because $\eta \gtrsim (1-\rho)^2$, we have $\cS_{\mathrm{over}}(k) \lesssim 1$. Hence, the true bias $\cS_{\nesterov}(k) \eqsim 1 + \mathcal{O}(1) \eqsim 1 \eqsim \widehat{\cS}_{\nesterov}(k)$.

% \textbf{Stage 2: The intermediate regime.}

% We verify the equivalence across the intermediate window bounded by $\frac{1}{\eta}$ and $\frac{1}{1-\rho}$, examining two exhaustive cases:
% \begin{itemize}
%     \item \textit{Case A: $\eta \gtrsim 1-\rho \implies \frac{1}{\eta} \lesssim \frac{1}{1-\rho}$.} The interval is $\frac{1}{\eta} \lesssim k \lesssim \frac{1}{1-\rho}$.\\
    
%     Here $\eta k \gtrsim 1 \implies \min\{1, (\eta k)^{-s}\} = (\eta k)^{-s}$, and $k(1-\rho) \lesssim 1 \implies \rho^k \eqsim 1$. The target evaluates to:
%     \[
%         \widehat{\cS}_{\nesterov}(k) \eqsim (\eta k)^{-s} + \left(\frac{1-\rho}{\eta k}\right)^s \eqsim (\eta k)^{-s},
%     \]
%     where the second term is absorbed because $1-\rho \lesssim \eta \lesssim 1$.\\
    
%     By Proposition~\ref{prop:nes_deterministic_bias}, $\cS_{\mathrm{under}}(k) \eqsim \frac{\rho^k}{(\eta k)^s} \eqsim (\eta k)^{-s}$ and $\cS_{\mathrm{over}}(k) \eqsim \left(\frac{(1-\rho)^2}{\eta}\right)^s$. Since $k \lesssim \frac{1}{1-\rho}$, we have $(\eta k)^{-s} \gtrsim \left(\frac{\eta}{1-\rho}\right)^{-s} = \left(\frac{1-\rho}{\eta}\right)^s \ge \left(\frac{(1-\rho)^2}{\eta}\right)^s$. Thus, $\cS_{\nesterov}(k) \eqsim (\eta k)^{-s} \eqsim \widehat{\cS}_{\nesterov}(k)$.

%     \item \textit{Case B: $\eta \lesssim 1-\rho \implies \frac{1}{1-\rho} \lesssim \frac{1}{\eta}$.} The interval is $\frac{1}{1-\rho} \lesssim k \lesssim \frac{1}{\eta}$.\\
    
%     Here $\eta k \lesssim 1 \implies \min\{1, (\eta k)^{-s}\} = 1$. The target directly yields:
%     \[
%         \widehat{\cS}_{\nesterov}(k) = \rho^k + \left(\frac{1-\rho}{\eta k}\right)^s.
%     \]
%     By Proposition~\ref{prop:nes_deterministic_bias}, $\cS_{\mathrm{under}}(k) \eqsim \rho^k$ and $\cS_{\mathrm{over}}(k) \eqsim \left(\frac{1-\rho}{\eta k}\right)^s$. Summing them exactly matches $\widehat{\cS}_{\nesterov}(k)$.
% \end{itemize}

% \textbf{Stage 3: $\max\left\{\frac{1}{\eta}, \frac{1}{1-\rho}\right\} \lesssim k \lesssim \frac{1}{(1-\rho)^2}$.}

% Here $\eta k \gtrsim 1 \implies \min\{1, (\eta k)^{-s}\} = (\eta k)^{-s}$. The target yields:
% \[
%     \widehat{\cS}_{\nesterov}(k) = (\eta k)^{-s} \rho^k + \left(\frac{1-\rho}{\eta k}\right)^s.
% \]
% By Proposition~\ref{prop:nes_deterministic_bias}, the components are exactly $\cS_{\mathrm{under}}(k) \eqsim \frac{\rho^k}{(\eta k)^s}$ and $\cS_{\mathrm{over}}(k) \eqsim \left(\frac{1-\rho}{\eta k}\right)^s$, which precisely reconstructs $\widehat{\cS}_{\nesterov}(k)$.

% \textbf{Stage 4: $k \gtrsim \frac{1}{(1-\rho)^2}$.}

% The mixed regime condition $\eta \gtrsim (1-\rho)^2$ dictates $\eta k \gtrsim \frac{\eta}{(1-\rho)^2} \gtrsim 1 \implies \min\{1, (\eta k)^{-s}\} = (\eta k)^{-s}$. Thus:
% \[
%     \widehat{\cS}_{\nesterov}(k) = (\eta k)^{-s} \rho^k + \left(\frac{1-\rho}{\eta k}\right)^s.
% \]
% We rigorously demonstrate that the exponential transient is absorbed. The ratio satisfies:
% \[
%     \frac{(\eta k)^{-s} \rho^k}{\left(\frac{1-\rho}{\eta k}\right)^s} 
%     = (1-\rho)^{-s} \rho^k 
%     \le (1-\rho)^{-s} e^{-k(1-\rho)}.
% \]
% For $k \gtrsim \frac{1}{(1-\rho)^2}$, the exponent $k(1-\rho) \gtrsim \frac{1}{1-\rho} \gg 1$. Factoring the decay gives:
% \[
%     (1-\rho)^{-s} e^{-k(1-\rho)} \le (1-\rho)^{-s} e^{-\frac{1}{2(1-\rho)}} e^{-\frac{k(1-\rho)}{2}}.
% \]
% Since $x^{-s} e^{-x/2} \to 0$ as $x \to \infty$, the term $(1-\rho)^{-s} e^{-\frac{1}{2(1-\rho)}} \ll 1$. Thus, the transient is fully absorbed:
% \[
%     \widehat{\cS}_{\nesterov}(k) \eqsim \left(\frac{1-\rho}{\eta k}\right)^s.
% \]

% Concurrently, we bound the true underdamped bias $\cS_{\mathrm{under}}(k)$ from Proposition~\ref{prop:nes_deterministic_bias}, where $p_{\nesterov} < \rho \le e^{-(1-\rho)}$. We split into its two sub-branches:
% \begin{itemize}
%     \item \textit{Sub-branch $\frac{1}{(1-\rho)^2} \lesssim k \lesssim \eta^{1/\beta}(1-\rho)^{-2-\frac{2}{\beta}}$:}
%     \begin{align*}
%         \frac{\cS_{\mathrm{under}}(k)}{\cS_{\mathrm{over}}(k)} 
%         &\eqsim \frac{\frac{1}{\eta k}\left(\frac{(1-\rho)^2}{\eta}\right)^{s-1} p_{\nesterov}^k}{\left(\frac{1-\rho}{\eta k}\right)^s} \\
%         &= k^{s-1} (1-\rho)^{s-2} p_{\nesterov}^k \\
%         &\le \big[ k^{s-1} e^{-\frac{k(1-\rho)}{2}} \big] (1-\rho)^{s-2} e^{-\frac{k(1-\rho)}{2}}.
%     \end{align*}
%     Since $k(1-\rho) \gtrsim \frac{1}{1-\rho}$, we have $e^{-\frac{k(1-\rho)}{2}} \le e^{-\frac{1}{2(1-\rho)}}$. Thus, $(1-\rho)^{s-2} e^{-\frac{1}{2(1-\rho)}} \to 0$, absorbing $\cS_{\mathrm{under}}(k)$.
    
%     \item \textit{Sub-branch $k \gtrsim \eta^{1/\beta}(1-\rho)^{-2-\frac{2}{\beta}}$:}
%     Using $\eta \gtrsim (1-\rho)^2 \implies \eta^{-\frac{1}{\beta}} \lesssim (1-\rho)^{-\frac{2}{\beta}}$, we have:
%     \begin{align*}
%         \frac{\cS_{\mathrm{under}}(k)}{\cS_{\mathrm{over}}(k)} 
%         &\eqsim \frac{\left(\frac{(1-\rho)^2}{\eta}\right)^{s+\frac{1}{\beta}} p_{\nesterov}^k}{\left(\frac{1-\rho}{\eta k}\right)^s} \\
%         &= (1-\rho)^{s+\frac{2}{\beta}} \eta^{-\frac{1}{\beta}} k^s p_{\nesterov}^k \\
%         &\lesssim \big[ \big(k(1-\rho)\big)^s e^{-\frac{k(1-\rho)}{2}} \big] e^{-\frac{k(1-\rho)}{2}} \to 0.
%     \end{align*}
% \end{itemize}
% In both cases, $\cS_{\mathrm{under}}(k) \ll \cS_{\mathrm{over}}(k) \eqsim \left(\frac{1-\rho}{\eta k}\right)^s$. The true bias simplifies to $\cS_{\nesterov}(k) \eqsim \left(\frac{1-\rho}{\eta k}\right)^s$, perfectly matching $\widehat{\cS}_{\nesterov}(k)$. 

% Synthesizing all stages, $\cS_{\nesterov}(k) \eqsim \widehat{\cS}_{\nesterov}(k)$ holds uniformly for $k \gtrsim \frac{1-\rho}{\eta}$. Substituting this and the saturated label noise into the total risk equivalence completes the proof.
% \end{proof}



% \newpage
\section{Optimal data scaling across batch-size regimes}
\label{app:sec:data-scaling}
This appendix proves the batch-size-dependent scaling laws stated in Theorem~\ref{thm:batch-size-scaling}. Throughout this appendix, we work in the noisy setting $\sigma>0$ considered in Section~\ref{sec:data-scaling} and impose the data-budget constraint
\[
B\eqsim D^b,\qquad T=\frac{D}{B}\eqsim D^{1-b},\qquad 0\le b<1.
\]
Thus, the batch exponent $b$ is prescribed, and only the remaining hyperparameters are optimized. We use the transition exponents defined in Section~\ref{sec:data-scaling}:
\[
b_1=\frac{s}{s+1},\qquad b_2=\frac{2s+1/\beta}{2s+1+1/\beta}=\frac{2s\beta+1}{2s\beta+\beta+1},\qquad b_3=\frac{2s+1}{2s+2}.
\]
For the momentum methods, we parameterize the momentum gap and learning rate by $1-\rho\eqsim D^{-m}$ and $\eta\eqsim D^{-c}$, where $m,c\ge0$.
\subsection{SGD}
For SGD, write $\eta\eqsim D^{-c}$, where $c\ge0$. By the SGD risk scaling law in Section~\ref{sec:scaling law},
\[
\EE[\cE(\btheta_T)]\eqsim(\eta T)^{-s}+\frac{\sigma^2\eta}{B}\eqsim D^{-s(1-b-c)}+D^{-(b+c)}.
\]
Therefore, for a fixed batch exponent $b$, the achievable decay exponent is $\min\{s(1-b-c),b+c\}$. If $0\le b<b_1$, the two terms can be balanced by taking $c=b_1-b\ge0$. Since $b_1=s/(s+1)$, this gives
\[
s(1-b-c)=b+c=b_1=\frac{s}{s+1}.
\]
Hence $r_{\sgd}(b)=s/(s+1)$ for $0\le b<b_1$. 

If $b\ge b_1$, the unconstrained balancing choice would require $c=b_1-b\le0$. Because the learning rate is not allowed to grow polynomially with $D$, we must have $c\ge0$. Moreover, $b+c\ge b\ge b_1$, so the signal term is the slower-decaying term. It is optimized by taking $c=0$, which yields $r_{\sgd}(b)=s(1-b)$ for $b_1\le b<1$. 

Consequently,
\[
r_{\sgd}(b)=
\begin{cases}
\dfrac{s}{s+1},&0\le b<b_1,\\[4pt]
s(1-b),&b_1\le b<1.
\end{cases}
\]
Thus, $b_1$ is precisely the largest batch exponent for which SGD preserves the data-optimal rate $s/(s+1)$.

\subsection{Polyak}
We next prove the Polyak part of Theorem~\ref{thm:batch-size-scaling}. We retain $T\eqsim D^{1-b}$, $1-\rho\eqsim D^{-m}$, and $\eta\eqsim D^{-c}$, where $m,c\ge0$. We require $b+m<1$, so that $\rho^T\leq\exp\bigl(-(1-\rho)T\bigr)=\exp\bigl(-\Theta(D^{1-b-m})\bigr)$ is smaller than any polynomial power of $D$. The stability condition $\eta\lesssim B(1-\rho)$ becomes $b+c-m\ge0$. Moreover,
\[
\eta_\rho T=\frac{\eta T}{1-\rho}\eqsim D^{1-b-c+m},\qquad \frac{\eta_\rho}{B}=\frac{\eta}{B(1-\rho)}\eqsim D^{-(b+c-m)}.
\]
The damping threshold $\eta\eqsim(1-\rho)^2$ corresponds to $c=2m$. We consider the fully overdamped and mixed-damping regimes separately.

\paragraph{Fully overdamped regime.} If $c\ge2m$, Theorem~\ref{thm:hbm scaling law} gives
\[
\EE[\cE(\btheta_T)]\eqsim(\eta_\rho T)^{-s}+\frac{\sigma^2\eta_\rho}{B}\eqsim D^{-s(1-b-c+m)}+D^{-(b+c-m)}.
\]
Therefore, the polynomial decay exponent in the fully overdamped regime is $\min\{s(1-b-c+m),b+c-m\}$.

\paragraph{Mixed-damping regime.} If $c\le2m$, Theorem~\ref{thm:hbm scaling law} gives
\[
\EE[\cE(\btheta_T)]\eqsim P(T)+(\eta_\rho T)^{-s}+\frac{\sigma^2\eta_\rho}{B}.
\]
Since $b+m<1$, the transient $P(T)\lesssim\rho^T\leq
\exp\bigl(-(1-\rho)T\bigr)=\exp\bigl(-\Theta(D^{1-b-m})\bigr)$ is smaller than any polynomial power of $D$. Hence
\[
\EE[\cE(\btheta_T)]\eqsim D^{-s(1-b-c+m)}+D^{-(b+c-m)},
\]
and the polynomial decay exponent is again $\min\{s(1-b-c+m),b+c-m\}$. Thus, although the two regimes have different transient dynamics, they lead to the same polynomial optimization problem. The signal and additive-noise powers are balanced when
\[
s(1-b-c+m)=b+c-m
\quad\Longleftrightarrow\quad
b+c-m=\frac{s}{s+1}=b_1.
\]
If $0\le b<b_1$, this balance is attained by taking $m=0$ and $c=b_1-b$. Since $c\ge2m$, this choice lies in the fully overdamped regime and yields $r_{\polyak}(b)=s/(s+1)$. 

If $b_1\le b<b_3$, the same balance is attained by taking $c=0$ and $m=b-b_1$. For $b>b_1$, this choice satisfies $c\le2m$ and therefore lies in the mixed-damping regime. Moreover, $b+c-m=b_1$, while the transient condition becomes $b+m=2b-b_1<1$. Since $(1+b_1)/2=(2s+1)/(2s+2)=b_3$, this condition holds precisely when $b<b_3$. Hence $r_{\polyak}(b)=s/(s+1)$ for $0\le b<b_3$.

Now suppose that $b\ge b_3$. Since $c\ge0$ and $m<1-b$, every admissible choice satisfies $b+c-m>2b-1\ge b_1$, so the signal term is the slower-decaying term and $s(1-b-c+m)<2s(1-b)$. This upper bound is approached by taking $c=0$ and $m\uparrow1-b$, which lies in the mixed-damping regime and gives
\[
s(1-b-c+m)\uparrow2s(1-b),\qquad b+c-m\downarrow2b-1.
\]
Moreover, $2b-1\ge2s(1-b)$ holds exactly when $b\ge b_3$, so the signal term is indeed dominant. Consequently, $r_{\polyak}(b)=2s(1-b)$ for $b_3\le b<1$, where the endpoint value is understood in the supremal sense represented by the $o(1)$ term in Definition~\ref{def:batch-scaling-exponent}. Combining the two regimes gives
\[
r_{\polyak}(b)=
\begin{cases}
\dfrac{s}{s+1},&0\le b<b_3,\\[4pt]
2s(1-b),&b_3\le b<1.
\end{cases}
\]
Thus, Polyak preserves the optimal data-scaling exponent $s/(s+1)$ up to the critical batch exponent $b_3$. Compared with SGD, it extends the critical batch size from $D^{b_1}$ to $D^{b_3}$, while leaving the optimal small-batch exponent unchanged.

\subsection{Nesterov}
We finally prove the Nesterov part of Theorem~\ref{thm:batch-size-scaling}. We retain $B\eqsim D^b$, $T\eqsim D^{1-b}$, $1-\rho\eqsim D^{-m}$, and $\eta\eqsim D^{-c}$, where $m,c\ge0$. We require $b+m<1$, so that the Nesterov transient $N(T)$ is smaller than any polynomial power of $D$. The stability condition $\eta\lesssim B^\beta(1-\rho)$ becomes $\beta b+c-m\ge0$, while the damping threshold $\eta\eqsim(1-\rho)^2$ corresponds to $c=2m$. We optimize the fully overdamped and mixed-damping regimes separately and then compare their optimal decay exponents.

\paragraph{Fully overdamped regime.} If $c\ge2m$, Theorem~\ref{thm:nesterov scaling law} gives
\[
\EE[\cE(\btheta_T)]\eqsim(\eta_\rho T)^{-s}+\frac{\sigma^2\eta_\rho}{B}\eqsim D^{-s(1-b-c+m)}+D^{-(b+c-m)}.
\]
Thus, the decay exponent is $\min\{s(1-b-c+m),b+c-m\}$. Since $c\ge2m$ implies $c-m\ge0$, we have $b+c-m\ge b$. If $0\le b<b_1$, the signal and additive-noise powers can be balanced by taking $m=0$ and $c=b_1-b$, which gives $b+c-m=b_1$ and decay exponent $s/(s+1)$. If $b\ge b_1$, then $b+c-m\ge b\ge b_1$, so the signal term is limiting. Its exponent is maximized by taking $m=c=0$, which gives $s(1-b)$. Therefore, the optimal exponent in the fully overdamped regime is
\[
r_{\nesterov}^{\rm full}(b)=
\begin{cases}
\dfrac{s}{s+1},&0\le b<b_1,\\[4pt]
s(1-b),&b_1\le b<1.
\end{cases}
\]
In particular, the fully overdamped Nesterov dynamics have the same batch-size-dependent scaling exponent as SGD.
\paragraph{Mixed-damping regime.} If $c\le2m$, Theorem~\ref{thm:nesterov scaling law} gives
\[
\EE[\cE(\btheta_T)]
\eqsim
N(T)
+(\eta_\rho T)^{-s}
+\frac{\sigma^2}{B}\min\left\{\eta_\rho,\eta_\rho^{1/\beta}\right\}.
\]
Since $b+m<1$, the first term \(N(T)\lesssim \min\{1,(\eta T)^{-s}\}\rho^T\leq\exp\bigl(-(1-\rho)T\bigr)=\exp\bigl(-\Theta(D^{1-b-m})\bigr)\) is smaller than any polynomial power of $D$. Hence
\[
\EE[\cE(\btheta_T)]\eqsim D^{-s(1-b-c+m)}+\frac{\sigma^2}{B}\min\left\{\frac{\eta}{1-\rho},\left(\frac{\eta}{1-\rho}\right)^{1/\beta}\right\}.
\]
The signal term has decay power $s(1-b-c+m)$. Since $\eta/(1-\rho)\eqsim D^{m-c}$, the additive-noise decay power is
\[
\begin{cases}
b+c-m,&c\ge m,\\[4pt]
b+\dfrac{c-m}{\beta},&c\le m.
\end{cases}
\]
We optimize these two branches separately.
\paragraph{The branch $c\ge m$.} In this branch, the decay exponent is
\[
\min\{s(1-b-c+m),b+c-m\}.
\]
Because $c-m\ge0$, the same argument as in the fully overdamped regime applies. For $0\le b<b_1$, the balance $b+c-m=b_1$ is feasible in the mixed-damping regime. For example, one may take $m=b_1-b$ and $c=2(b_1-b)$, for which $b+m=b_1<1$. For $b\ge b_1$, the optimal choice satisfies $c-m=0$ and gives exponent $s(1-b)$. Therefore,
\[
r_{\nesterov}^{\rm mixed,\,c\ge m}(b)=
\begin{cases}
\dfrac{s}{s+1},&0\le b<b_1,\\[4pt]
s(1-b),&b_1\le b<1.
\end{cases}
\]
\paragraph{The branch $c\le m$.} In this branch, the decay exponent is
\[
\min\left\{s(1-b-c+m),\,b+\frac{c-m}{\beta}\right\}.
\]
If $0\le b<b_1$, then at $m-c=0$ we have $b<s(1-b)$. Increasing $m-c$ increases the signal power but decreases the additive-noise power, so the optimal choice is $m-c=0$, which gives decay exponent $b$.
If $b\ge b_1$, balancing the signal and additive-noise powers gives
\[
s(1-b-c+m)=b+\frac{c-m}{\beta},
\qquad
m-c=\frac{\beta\bigl((s+1)b-s\bigr)}{s\beta+1}.
\]
For a fixed value of $m-c$, the least restrictive choice for the transient condition is $c=0$, so we take
\[
c=0,\qquad m=\frac{\beta\bigl((s+1)b-s\bigr)}{s\beta+1}.
\]
This choice satisfies the stability condition because
\[
\beta b+c-m=\frac{\beta s\bigl(1+(\beta-1)b\bigr)}{s\beta+1}>0.
\]
The signal and additive-noise powers are then equal to
\[
s(1-b+m)=b-\frac{m}{\beta}=\frac{s\bigl(1+(\beta-1)b\bigr)}{s\beta+1}.
\]
The transient condition holds precisely when
\[
b+m<1
\quad\Longleftrightarrow\quad
b<\frac{2s\beta+1}{2s\beta+\beta+1}
=\frac{2s+1/\beta}{2s+1+1/\beta}
=b_2.
\]
Therefore, for $b_1\le b<b_2$, the optimal exponent on the branch $c\le m$ is
\[
\frac{s\bigl(1+(\beta-1)b\bigr)}{s\beta+1}.
\]
If $b\ge b_2$, the preceding balance is no longer compatible with $b+m<1$. Since $c\ge0$ and $m<1-b$, every admissible choice satisfies
\[
s(1-b-c+m)<2s(1-b).
\]
This upper bound is approached by taking $c=0$ and $m\uparrow1-b$. Under this choice,
\[
s(1-b+m)\uparrow2s(1-b),\qquad b-\frac{m}{\beta}\to\frac{(\beta+1)b-1}{\beta}.
\]
At $b=b_2$, the two limiting powers agree, while for $b>b_2$,
\[
\frac{(\beta+1)b-1}{\beta}>2s(1-b).
\]
Thus, the signal term is limiting and the supremal decay exponent is $2s(1-b)$. Consequently,
\[
r_{\nesterov}^{\rm mixed,\,c\le m}(b)=
\begin{cases}
b,&0\le b<b_1,\\[4pt]
\dfrac{s\bigl(1+(\beta-1)b\bigr)}{s\beta+1},&b_1\le b<b_2,\\[8pt]
2s(1-b),&b_2\le b<1.
\end{cases}
\]
Taking the larger exponent between the two additive-noise branches gives the optimal exponent in the mixed-damping regime:
\[
r_{\nesterov}^{\rm mixed}(b)=
\begin{cases}
\dfrac{s}{s+1},&0\le b<b_1,\\[6pt]
\dfrac{s\bigl(1+(\beta-1)b\bigr)}{s\beta+1},&b_1\le b<b_2,\\[8pt]
2s(1-b),&b_2\le b<1.
\end{cases}
\]
Indeed, the branch $c\ge m$ is optimal for $b<b_1$, whereas the branch $c\le m$ is optimal for $b>b_1$. Finally, optimizing over the fully overdamped and mixed-damping regimes gives
\[
r_{\nesterov}(b)=\max\left\{r_{\nesterov}^{\rm full}(b),r_{\nesterov}^{\rm mixed}(b)\right\}.
\]
For $0\le b<b_1$, the two regimes attain the same exponent $s/(s+1)$. For $b_1<b<b_2$, the mixed-damping regime exponent $s(1+(\beta-1)b)/(s\beta+1)$ is strictly larger than the fully overdamped exponent $s(1-b)$. For $b\ge b_2$, the mixed-damping regime exponent $2s(1-b)$ is also larger than the fully overdamped exponent $s(1-b)$. We therefore obtain
\[
r_{\nesterov}(b)=
\begin{cases}
\dfrac{s}{s+1},&0\le b<b_1,\\[6pt]
\dfrac{s\bigl(1+(\beta-1)b\bigr)}{s\beta+1},&b_1\le b<b_2,\\[8pt]
2s(1-b),&b_2\le b<1.
\end{cases}
\]
Thus, Nesterov agrees with SGD and Polyak in the small-batch regime, improves its data-scaling exponent once $b$ exceeds $b_1$, and reaches its maximal exponent at $b=b_2$. For $b\ge b_2$, its exponent becomes $2s(1-b)$, which coincides with the Polyak exponent in the ultra-large-batch regime.



% \section{Data-optimal scaling derivations}
% \label{app:data_optimal_scaling}

% Throughout this section, we assume $\sigma>0$ and impose the fixed data budget
% \[
%     BK\eqsim D.
% \]
% We derive the optimal polynomial decay rate directly from the loss bounds proved above.  To keep the dependence on the algorithmic parameters explicit, all constraints and decay powers are stated in terms of the scaling exponents of $B$, $K$, $1-\rho$, and $\eta$, without introducing auxiliary optimization variables.

% \subsection{SGD}
% \label{app:sgd_scaling_analysis}

% To establish the data-optimal scaling for SGD presented in the SGD part of Theorem~\ref{thm:batch-size-scaling}, we parameterize the hyperparameters in terms of the data budget \(D\) and balance the polynomial decay rates.

% \begin{proof}[Proof of the SGD part of Theorem~\ref{thm:batch-size-scaling}]
% Let \(B\eqsim D^a\), \(K\eqsim D^{1-a}\), and \(\eta\eqsim D^{-c}\), where \(a,c\ge0\). Since \(c\ge0\), the stability condition \(\eta\lesssim1\) is satisfied at the level of polynomial powers. Moreover, \(\eta K\eqsim D^{1-a-c}\), so decay of the deterministic bias requires \(a+c\le1\).

% By the SGD scaling law,
% \[
%     \EE[\cE(\btheta_K)]
%     \eqsim
%     \left(\frac{1}{\eta K}\right)^s
%     +
%     \frac{\sigma^2\eta}{B}.
% \]
% Under the data-budget parametrization, the deterministic-bias and additive-noise terms scale respectively as \(D^{-s(1-a-c)}\) and \(D^{-(a+c)}\). Therefore, the achievable decay power is
% \[
%     \min\{s(1-a-c),a+c\}
%     \leq
%     \frac{s}{s+1},
% \]
% with equality if and only if \(a+c=s/(s+1)\). Hence every choice satisfying \(0\le a\le s/(s+1)\) and \(c=s/(s+1)-a\) achieves
% \[
%     \EE[\cE(\btheta_K)]
%     \eqsim
%     D^{-\frac{s}{s+1}}.
% \]
% Equivalently, the optimal learning rate satisfies
% \[
%     \eta
%     \eqsim
%     D^{a-\frac{s}{s+1}}
%     \eqsim
%     B D^{-\frac{s}{s+1}},
% \]
% which gives the linear scaling rule. Finally, the largest admissible batch exponent is obtained by setting \(c=0\), yielding
% \[
%    B_{\sgd}^{\crit}(D)
%     \eqsim
%     D^{\frac{s}{s+1}}.
% \]
% \end{proof}

% % \subsection{SGD}
% % \label{app:sgd_scaling_analysis}

% % To establish the data-optimal scaling for SGD presented in Proposition \ref{prop:sgd data scaling} of Section 6, we parameterize the hyperparameters in terms of the data budget $D$ and balance the polynomial decay rates. 

% % \begin{proof}[Proof of Proposition \ref{prop:sgd data scaling}]
% % Let
% % \[
% %     B\eqsim D^a,
% %     \qquad
% %     K\eqsim D^{1-a},
% %     \qquad
% %     \eta\eqsim D^{-c},
% %     \qquad a,c\ge0.
% % \]
% % The stability condition $\eta\lesssim B$ is equivalent, at the level of polynomial powers, to $a+c\ge0$.  Moreover,
% % \[
% %     \eta K\eqsim D^{1-a-c},
% % \]
% % so a non-increasing deterministic bias requires $a+c\le1$.

% % By Theorem~\ref{thm:pol_scaling_sgd} with $\rho=0$,
% % \[
% %     \EE[\cE(\btheta_K)]
% %     \eqsim
% %     \left(\frac{1}{\eta K}\right)^s
% %     +\frac{\sigma^2\eta}{B}
% %     +\frac{\eta^2}{B}
% %         \left(\frac{1}{\eta K}\right)^{s\wedge(2-\frac1\beta)}.
% % \]
% % Substituting the data-budget parametrization yields
% % \[
% % \begin{aligned}
% %     \left(\frac{1}{\eta K}\right)^s
% %     &\eqsim D^{-s(1-a-c)},\\
% %     \frac{\sigma^2\eta}{B}
% %     &\eqsim D^{-(a+c)},\\
% %     \frac{\eta^2}{B}
% %         \left(\frac{1}{\eta K}\right)^{s\wedge(2-\frac1\beta)}
% %     &\eqsim
% %     D^{-\left[
% %         s\wedge\left(2-\frac1\beta\right)
% %         +\left(2-s\wedge\left(2-\frac1\beta\right)\right)(a+c)-a
% %     \right]}.
% % \end{aligned}
% % \]
% % Consequently, the achievable decay power cannot exceed
% % \[
% %     \min\{s(1-a-c),a+c\}.
% % \]
% % For every $a+c\in[0,1]$,
% % \[
% %     \min\{s(1-a-c),a+c\}
% %     \le \frac{s}{s+1},
% % \]
% % with equality if and only if
% % \[
% %     a+c=\frac{s}{s+1}.
% % \]
% % It remains to verify that the multiplicative-noise term is not active at this balance.  If $a+c=s/(s+1)$, its decay power minus $s/(s+1)$ equals
% % \[
% %     s\wedge\left(2-\frac1\beta\right)
% %     +\left(1-s\wedge\left(2-\frac1\beta\right)\right)\frac{s}{s+1}-a.
% % \]
% % Since $c\ge0$ implies $a\le s/(s+1)$, this quantity is bounded below by
% % \[
% %     \left(s\wedge\left(2-\frac1\beta\right)\right)
% %     \left(1-\frac{s}{s+1}\right)>0.
% % \]
% % Thus the multiplicative-noise contribution is strictly lower order.  We conclude that every choice satisfying
% % \[
% %     0\le a\le\frac{s}{s+1},
% %     \qquad
% %     c=\frac{s}{s+1}-a
% % \]
% % achieves
% % \[
% %     \EE[\cE(\btheta_K)]\eqsim D^{-\frac{s}{s+1}}.
% % \]
% % The largest admissible batch exponent is obtained by setting $c=0$:
% % \[
% %     B\eqsim D^{\frac{s}{s+1}}.
% % \]
% % \end{proof}


% \subsection{Polyak momentum}
% \label{app:hb_scaling_analysis}

% We now derive the data-optimal scaling for Polyak momentum stated in
% the Polyak part of Theorem~\ref{thm:batch-size-scaling}. We consider the fully overdamped and
% mixed overdamped--underdamped regimes separately.

% \begin{proof}[Proof of the Polyak part of Theorem~\ref{thm:batch-size-scaling}]
% Let
% \[
% B\eqsim D^a,
% \qquad
% K\eqsim D^{1-a},
% \qquad
% 1-\rho\eqsim D^{-b},
% \qquad
% \eta\eqsim D^{-c},
% \qquad
% a,b,c\geq0,
% \]
% and assume \(a+b<1\), so that
% \[
% \rho^K
% \eqsim
% \exp\!\left(-D^{1-a-b}\right)
% \]
% is smaller than any polynomial power of \(D\). We have
% \[
% \frac{\eta}{B(1-\rho)}
% \eqsim
% D^{-(a+c-b)},
% \qquad
% \frac{1-\rho}{\eta K}
% \eqsim
% D^{-(1-a-c+b)}.
% \]
% The stability condition \(\eta\lesssim B(1-\rho)\) becomes
% \[
% a+c-b\geq0,
% \]
% while decay of the polynomial bias requires \(a+c-b\leq1\). The damping
% threshold \(\eta\eqsim(1-\rho)^2\) corresponds to \(c=2b\).

% In the fully overdamped regime \(c\geq2b\), Theorem~\ref{thm:hbm scaling law}
% gives
% \[
% \EE[\cE(\btheta_K)]
% \eqsim
% \left(\frac{1-\rho}{\eta K}\right)^s
% +
% \frac{\eta\sigma^2}{B(1-\rho)}.
% \]
% The two terms have decay powers
% \[
% s(1-a-c+b)
% \qquad\text{and}\qquad
% a+c-b.
% \]
% Therefore,
% \[
% \min\left\{
% s(1-a-c+b),\,a+c-b
% \right\}
% \leq
% \frac{s}{s+1},
% \]
% with equality if and only if
% \[
% a+c-b=\frac{s}{s+1}.
% \]

% In the mixed overdamped--underdamped regime \(c\leq2b\),
% Theorem~\ref{thm:hbm scaling law} gives
% \[
% \begin{aligned}
% \EE[\cE(\btheta_K)]
% \eqsim{}&
% \rho^K
% +
% \left(\frac{1-\rho}{\eta K}\right)^s
% +
% \frac{\eta\sigma^2}{B(1-\rho)}\\
% &+
% \frac{\eta K}{B}\rho^K
% +
% \frac{\eta^2}{B(1-\rho)^3}
% \left(\frac{1-\rho}{\eta K}\right)^{2-\frac1\beta}.
% \end{aligned}
% \]
% Because \(a+b<1\), both terms containing \(\rho^K\) are exponentially small.
% The remaining multiplicative-noise term has decay power
% \[
% 2-\frac1\beta-a-b+\frac{a+c-b}{\beta}.
% \]
% At the bias--variance balance
% \[
% a+c-b=\frac{s}{s+1},
% \]
% the difference between this decay power and \(s/(s+1)\) is
% \[
% \begin{aligned}
% &2-\frac1\beta-a-b
% -\left(1-\frac1\beta\right)\frac{s}{s+1}\\
% &>
% 1-\frac1\beta
% -\left(1-\frac1\beta\right)\frac{s}{s+1}
% =
% \left(1-\frac1\beta\right)\frac{1}{s+1}
% >0,
% \end{aligned}
% \]
% where we used \(a+b<1\). Thus the multiplicative-noise term is strictly
% lower order and does not impose any additional constraint.

% Consequently, every admissible choice satisfying
% \[
% a+c-b=\frac{s}{s+1}
% \]
% achieves
% \[
% \EE[\cE(\btheta_K)]
% \eqsim
% D^{-\frac{s}{s+1}}.
% \]
% Moreover,
% \[
% \frac{\eta}{1-\rho}
% \eqsim
% D^{b-c}
% =
% D^{a-(a+c-b)}
% \eqsim
% B D^{-\frac{s}{s+1}},
% \]
% which proves the linear scaling rule.

% It remains to determine the largest admissible batch exponent. From
% \[
% c=\frac{s}{s+1}+b-a\geq0,
% \]
% we obtain
% \[
% b\geq a-\frac{s}{s+1}.
% \]
% If \(a\leq s/(s+1)\), then the batch exponent is already below the claimed
% maximum. If \(a>s/(s+1)\), combining the preceding inequality with
% \(a+b<1\) gives
% \[
% 2a-\frac{s}{s+1}<1,
% \]
% and hence
% \[
% a<
% \frac{2s+1}{2s+2}.
% \]
% For every sufficiently small \(\varepsilon>0\), take
% \[
% a=\frac{2s+1}{2s+2}-\varepsilon,
% \qquad
% b=\frac{1}{2s+2}-\varepsilon,
% \qquad
% c=0.
% \]
% Then
% \[
% a+b=1-2\varepsilon<1,
% \qquad
% a+c-b=\frac{s}{s+1},
% \qquad
% c\leq2b,
% \]
% so this choice lies in the mixed regime, satisfies the stability condition,
% and achieves the optimal loss rate. Therefore,
% \[
% B_{\polyak}^{\crit}(D)
% =
% D^{\frac{2s+1}{2s+2}-o(1)}.
% \]
% \end{proof}

% % \subsection{Polyak momentum}
% % \label{app:hb_scaling_analysis}

% % Next, we derive the data-optimal scaling for Polyak momentum as stated in Theorem \ref{thm:polyak data scaling} of Section 6. We evaluate the polynomial decay powers under the data budget constraint across both the Fully Overdamped Regime and the Mixed Overdamped--Underdamped Regime.

% % \begin{proof}[Proof of Theorem \ref{thm:polyak data scaling}]
% % Let
% % \[
% %     B\eqsim D^a,
% %     \qquad
% %     K\eqsim D^{1-a},
% %     \qquad
% %     1-\rho\eqsim D^{-b},
% %     \qquad
% %     \eta\eqsim D^{-c},
% %     \qquad a,b,c\ge0,
% % \]
% % and assume $a+b<1$.  Then
% % \[
% %     \frac{\eta}{B(1-\rho)}\eqsim D^{-(a+c-b)},
% %     \qquad
% %     \frac{1-\rho}{\eta K}\eqsim D^{-(1-a-c+b)}.
% % \]
% % Hence the stability condition $\eta\lesssim B(1-\rho)$ becomes
% % \[
% %     a+c-b\ge0,
% % \]
% % while decay of the polynomial bias requires $a+c-b\le1$.  The damping threshold $\eta\eqsim(1-\rho)^2$ is equivalent to $c=2b$; therefore $c\ge2b$ is the Fully Overdamped Regime and $c\le2b$ is the Mixed Overdamped--Underdamped Regime.

% % In the Fully Overdamped Regime, Theorem~\ref{thm:pol_scaling_sgd} gives
% % \[
% %     \EE[\cE(\btheta_K)]
% %     \eqsim
% %     \left(\frac{1-\rho}{\eta K}\right)^s
% %     +\frac{\sigma^2\eta}{B(1-\rho)}
% %     +\frac{\eta^2}{B(1-\rho)^2}
% %         \left(\frac{1-\rho}{\eta K}\right)^{s\wedge(2-\frac1\beta)}.
% % \]
% % The corresponding polynomial decay powers are
% % \[
% % \begin{aligned}
% %     &s(1-a-c+b),\\
% %     &a+c-b,\\
% %     &s\wedge\left(2-\frac1\beta\right)
% %       +\left(2-s\wedge\left(2-\frac1\beta\right)\right)(a+c-b)-a.
% % \end{aligned}
% % \]
% % In the Mixed Overdamped--Underdamped Regime, Theorem~\ref{thm:pol_two_sided_scaling} gives
% % \[
% % \begin{aligned}
% %     \EE[\cE(\btheta_K)]
% %     \eqsim{}&
% %     \rho^K
% %     +\left(\frac{1-\rho}{\eta K}\right)^s
% %     +\frac{\sigma^2\eta}{B(1-\rho)}\\
% %     &+\frac{\eta^2}{B(1-\rho)^3}
% %         \left(\frac{1-\rho}{\eta K}\right)^{2-\frac1\beta}
% %     +\frac{\eta^2}{B(1-\rho)^2}
% %         \left(\frac{1-\rho}{\eta K}\right)^{s\wedge(2-\frac1\beta)}.
% % \end{aligned}
% % \]
% % In addition to the three powers displayed above, the first self-noise term in the second line has decay power
% % \[
% %     2-\frac1\beta-a-b+\frac{a+c-b}{\beta}.
% % \]
% % The transient is exponentially small because
% % \[
% %     \rho^K\eqsim\exp\!\left(-D^{1-a-b}\right),
% %     \qquad a+b<1.
% % \]

% % The deterministic-bias and label-noise terms imply the universal upper bound
% % \[
% %     \min\{s(1-a-c+b),a+c-b\}
% %     \le\frac{s}{s+1}.
% % \]
% % Equality requires
% % \[
% %     a+c-b=\frac{s}{s+1}.
% % \]
% % We now verify that both self-noise contributions can be made no larger than the balanced bias and label-noise terms.  First, the decay power shared by both damping regimes is at least $s/(s+1)$ precisely when
% % \[
% %     a\le
% %     s\wedge\left(2-\frac1\beta\right)
% %     +\left(1-s\wedge\left(2-\frac1\beta\right)\right)
% %       \frac{s}{s+1}
% %     =\frac{s+s\wedge(2-\frac1\beta)}{s+1}.
% % \]
% % Second, in the Mixed Overdamped--Underdamped Regime, $a+c-b\le a+b<1$.  Therefore
% % \[
% % \begin{aligned}
% %     &2-\frac1\beta-a-b+\frac{a+c-b}{\beta}-(a+c-b)\\
% %     &\qquad=
% %     2-\frac1\beta-a-b
% %     -\left(1-\frac1\beta\right)(a+c-b)\\
% %     &\qquad>
% %     \left(1-\frac1\beta\right)
% %     \left[1-(a+c-b)\right]
% %     \ge0.
% % \end{aligned}
% % \]
% % Thus the Mixed Overdamped--Underdamped Regime-specific self-noise term is strictly lower order than the label-noise term.  It follows that every $a,b,c\ge0$ satisfying
% % \[
% % \begin{gathered}
% %     a+b<1,
% %     \qquad
% %     a+c-b=\frac{s}{s+1},\\
% %     a\le\frac{s+s\wedge(2-\frac1\beta)}{s+1}
% % \end{gathered}
% % \]
% % achieves
% % \[
% %     \EE[\cE(\btheta_K)]
% %     \eqsim
% %     D^{-\frac{s}{s+1}}
% %     +\exp\!\left(-D^{1-a-b}\right).
% % \]

% % It remains to maximize $a$ over this feasible set.  Since
% % \[
% %     c=\frac{s}{s+1}+b-a\ge0,
% % \]
% % we must have $b\ge a-s/(s+1)$.  Combining this with $a+b<1$ yields
% % \[
% %     a<\frac{1}{2}\left(1+\frac{s}{s+1}\right)
% %     =\frac{2s+1}{2s+2}
% % \]
% % whenever $a>s/(s+1)$; no stronger restriction is obtained when $a\le s/(s+1)$.  Together with the self-noise constraint, the supremal batch exponent is therefore
% % \[
% %     \min\left\{
% %         \frac{2s+1}{2s+2},
% %         \frac{s+s\wedge(2-\frac1\beta)}{s+1}
% %     \right\}.
% % \]
% % If $s\le2-1/\beta$, the second term is $2s/(s+1)$.  If $s>2-1/\beta$, the first term is smaller than the second because $2-1/\beta>1$.  Hence the same quantity may be written as
% % \[
% %     \min\left\{
% %         \frac{2s+1}{2s+2},
% %         \frac{2s}{s+1}
% %     \right\}.
% % \]
% % Accordingly,
% % \[
% %     B\eqsim
% %     D^{\min\left\{
% %         \frac{2s+1}{2s+2},
% %         \frac{2s}{s+1}
% %     \right\}},
% % \]
% % where the first exponent is a supremum when it is active, due to the strict constraint $a+b<1$.
% % \end{proof}

% \subsection{Nesterov momentum}
% \label{app:nag_scaling_analysis}

% We now derive the data-optimal scaling for Nesterov momentum stated in
% the Nesterov part of Theorem~\ref{thm:batch-size-scaling}.

% \begin{proof}[Proof of the Nesterov part of Theorem~\ref{thm:batch-size-scaling}]
% We retain the parametrization
% \[
% B\eqsim D^a,
% \qquad
% K\eqsim D^{1-a},
% \qquad
% 1-\rho\eqsim D^{-b},
% \qquad
% \eta\eqsim D^{-c},
% \qquad
% a,b,c\geq0,
% \]
% and assume \(a+b<1\). The stability condition
% \(\eta\lesssim B^\beta(1-\rho)\) becomes
% \[
% \beta a+c-b\geq0,
% \]
% while the damping threshold \(\eta\eqsim(1-\rho)^2\) corresponds to \(c=2b\).
% Moreover,
% \[
% \gamma_{\nesterov}^K
% \eqsim
% \exp\!\left(-D^{1-a-b}\right),
% \]
% so all terms containing \(\gamma_{\nesterov}^K\), including those
% multiplied by polynomial factors, are smaller than any polynomial power of
% \(D\).

% In the fully overdamped regime \(c\geq2b\),
% Theorem~\ref{thm:nesterov scaling law} gives
% \[
% \EE[\cE(\btheta_K)]
% \eqsim
% \left(\frac{1-\rho}{\eta K}\right)^s
% +
% \frac{\eta\sigma^2}{B(1-\rho)}.
% \]
% As in the Polyak case, balancing the two terms gives at most
% \[
% D^{-\frac{s}{s+1}}.
% \]

% We next consider the mixed regime \(c\leq2b\). After removing the
% exponentially small transient terms, Theorem~\ref{thm:nesterov scaling law}
% gives
% \[
% \begin{aligned}
% \EE[\cE(\btheta_K)]
% \eqsim{}&
% \left(\frac{1-\rho}{\eta K}\right)^s
% +
% \frac{\sigma^2}{B}
% \min\left\{
% \frac{\eta}{1-\rho},
% \left(\frac{\eta}{1-\rho}\right)^{1/\beta}
% \right\}\\
% &+
% \frac{\eta^2}{B(1-\rho)^3}
% \left(\frac{1-\rho}{\eta K}\right)^{2-\frac1\beta}.
% \end{aligned}
% \]
% The deterministic-bias term has decay power
% \[
% s(1-a-c+b),
% \]
% and the additive-noise term has decay power
% \[
% \begin{cases}
% a+c-b, & c\geq b,\\[0.4em]
% a+\dfrac{c-b}{\beta}, & c\leq b.
% \end{cases}
% \]
% The multiplicative-noise term has decay power
% \[
% 2-\frac1\beta-a-b+\frac{a+c-b}{\beta}.
% \]

% When \(c\geq b\), balancing the deterministic bias and additive noise gives
% at most
% \[
% \frac{s}{s+1}.
% \]
% Since
% \[
% \frac{2s\beta}{2s\beta+\beta+1}
% -
% \frac{s}{s+1}
% =
% \frac{s(\beta-1)}
% {(s+1)(2s\beta+\beta+1)}
% >0,
% \]
% the improved Nesterov rate can only arise on the branch \(c\leq b\).

% On this branch, the additive-noise decay power can be written as
% \[
% a+\frac{c-b}{\beta}
% =
% \left(1-\frac1\beta\right)a
% +
% \frac{a+c-b}{\beta}.
% \]
% For fixed \(a+c-b\), this power increases with \(a\). Since
% \[
% a+b
% =
% 2a+c-(a+c-b)<1
% \]
% and \(c\geq0\), we have
% \[
% a<
% \frac{1+a+c-b}{2}.
% \]
% Therefore, for fixed \(a+c-b\), the supremal additive-noise decay power is
% \[
% \frac{\beta-1+(\beta+1)(a+c-b)}{2\beta},
% \]
% and it is approached when \(c\to0\) and \(a+b\to1\). Balancing this quantity
% with the deterministic-bias decay power gives
% \[
% s(1-a-c+b)
% =
% \frac{\beta-1+(\beta+1)(a+c-b)}{2\beta}.
% \]
% Solving yields
% \[
% a+c-b
% =
% \frac{2s\beta-\beta+1}{2s\beta+\beta+1},
% \]
% and the common decay power is
% \[
% \frac{2s\beta}{2s\beta+\beta+1}.
% \]
% The corresponding limiting parameter exponents are
% \[
% a\to
% \frac{2s\beta+1}{2s\beta+\beta+1},
% \qquad
% b\to
% \frac{\beta}{2s\beta+\beta+1},
% \qquad
% c\to0.
% \]

% It remains to verify that the multiplicative-noise term is lower order.
% At these limiting exponents, its decay power is
% \[
% \frac{2s\beta+\beta-1}{2s\beta+\beta+1},
% \]
% which exceeds the optimal decay power by
% \[
% \frac{\beta-1}{2s\beta+\beta+1}>0.
% \]
% Thus the multiplicative-noise contribution is not active at the optimum.

% For every sufficiently small \(\varepsilon>0\), choose
% \[
% a=
% \frac{2s\beta+1}{2s\beta+\beta+1}-\varepsilon,
% \qquad
% b=
% \frac{\beta}{2s\beta+\beta+1}-\varepsilon,
% \qquad
% c=0.
% \]
% Then
% \[
% a+b=1-2\varepsilon<1,
% \]
% and the deterministic-bias, additive-noise, and multiplicative-noise decay
% powers are, respectively,
% \[
% \frac{2s\beta}{2s\beta+\beta+1},
% \qquad
% \frac{2s\beta}{2s\beta+\beta+1}
% -\left(1-\frac1\beta\right)\varepsilon,
% \qquad
% \frac{2s\beta+\beta-1}{2s\beta+\beta+1}
% +2\varepsilon.
% \]
% The stability condition is satisfied because
% \[
% \beta a+c-b
% =
% \frac{2s\beta^2}{2s\beta+\beta+1}
% -(\beta-1)\varepsilon
% >0
% \]
% for sufficiently small \(\varepsilon\). The choice also lies in the mixed
% regime since \(c=0\leq2b\). Hence
% \[
% \EE[\cE(\btheta_K)]
% \eqsim
% D^{-\frac{2s\beta}{2s\beta+\beta+1}
% +\left(1-\frac1\beta\right)\varepsilon}.
% \]
% Letting \(\varepsilon\downarrow0\) proves
% \[
% \cE_{\nesterov}^{\opt}(D)
% =
% D^{-\frac{2s\beta}{2s\beta+\beta+1}+o(1)},
% \qquad
% B_{\nesterov}^{\crit}(D)
% =
% D^{\frac{2s\beta+1}{2s\beta+\beta+1}-o(1)}.
% \]
% \end{proof}

% % \subsection{Nesterov momentum}
% % \label{app:nag_scaling_analysis}

% % Finally, we establish the data-optimal scaling for Nesterov momentum detailed in Theorem \ref{thm:nesterov data scaling} of Section 6. This requires solving a constrained polynomial optimization problem to balance the bias and noise terms across different regimes of the source exponent $s$.

% % \begin{proof}[Proof of Theorem \ref{thm:nesterov data scaling}]
% % We retain the parametrization
% % \[
% %     B\eqsim D^a,
% %     \qquad
% %     K\eqsim D^{1-a},
% %     \qquad
% %     1-\rho\eqsim D^{-b},
% %     \qquad
% %     \eta\eqsim D^{-c},
% %     \qquad a,b,c\ge0,
% % \]
% % with $a+b<1$.  The data-scaling stability condition $\eta\lesssim B^\beta(1-\rho)$ becomes
% % \[
% %     \beta a+c-b\ge0.
% % \]
% % The damping threshold remains $c=2b$.  As in the Polyak momentum case,
% % \[
% %     \gamma_{\nesterov}^K
% %     \eqsim\exp\!\left(-D^{1-a-b}\right),
% % \]
% % after imposing $a+b<1$.

% % Theorems~\ref{thm:nes_scaling_sgd} and~\ref{thm:nes_two_sided_scaling} show that the deterministic-bias contribution has decay power
% % \[
% %     s(1-a-c+b),
% % \]
% % and that the label-noise contribution is
% % \[
% %     \frac{\sigma^2}{B}
% %     \min\left\{
% %         \left(\frac{\eta}{1-\rho}\right)^{\frac1\beta},
% %         \frac{\eta}{1-\rho}
% %     \right\}.
% % \]
% % Since $\eta/(1-\rho)\eqsim D^{b-c}$, its decay power is
% % \[
% %     \begin{cases}
% %     a+c-b, & c\ge b,\\[4pt]
% %     a+\dfrac{c-b}{\beta}, & c\le b.
% %     \end{cases}
% % \]
% % The two self-noise contributions have decay powers
% % \[
% %     2-\frac1\beta-a-b+\frac{a+c-b}{\beta}
% % \]
% % and
% % \[
% %     s\wedge\left(2-\frac1\beta\right)
% %     +\left(2-s\wedge\left(2-\frac1\beta\right)\right)(a+c-b)-a,
% % \]
% % respectively.

% % When $c\ge b$, the label-noise power is $a+c-b$, and the same bias--variance balance as for SGD and Polyak momentum gives at most $s/(s+1)$.  The improved Nesterov rate can therefore only arise on the branch $c\le b$.  On this branch, the polynomial optimization is
% % \[
% % \begin{aligned}
% %     \sup_{\substack{a,b,c\ge0,\ a+b<1,\ c\le b,\\
% %                     \beta a+c-b\ge0,\ a+c-b\le1}}
% %     \min\bigg\{&
% %         s(1-a-c+b),
% %         a+\frac{c-b}{\beta},\\
% %         &2-\frac1\beta-a-b+\frac{a+c-b}{\beta},\\
% %         &s\wedge\left(2-\frac1\beta\right)
% %         +\left(2-s\wedge\left(2-\frac1\beta\right)\right)(a+c-b)-a
% %     \bigg\}.
% % \end{aligned}
% % \]
% % This formulation contains all polynomial terms in the Nesterov loss bound; hence solving it gives both an upper bound and an achievability condition.

% % For fixed $a$ and $a+c-b$, only the third displayed power depends on $a+b$, and it decreases as $a+b$ increases.  The constraints $b\ge0$ and $c\ge0$ imply
% % \[
% %     a+b\ge a,
% %     \qquad
% %     a+b\ge 2a-(a+c-b).
% % \]
% % On the branch $c\le b$, one has $a+c-b\le a$, so the smallest feasible value is
% % \[
% %     a+b=2a-(a+c-b),
% % \]
% % corresponding to $c=0$.  Larger values of $a+b$ remain admissible whenever the third power is slack; this observation accounts for the interval of optimal momentum scalings below.

% % \paragraph{Regime $0<s<(2\beta-1)/(2\beta)$.}
% % In this regime, $s<(2\beta-1)/(2\beta)<1<2-1/\beta$, and hence $s\wedge(2-1/\beta)=s$.  The deterministic bias, label noise, and second self-noise term are simultaneously active.  Solving
% % \[
% % \begin{aligned}
% %     s(1-a-c+b)
% %     &=a+\frac{c-b}{\beta},\\
% %     s(1-a-c+b)
% %     &=s+(2-s)(a+c-b)-a
% % \end{aligned}
% % \]
% % gives
% % \[
% %     a+c-b=\frac{s\beta}{s\beta+2\beta-1},
% %     \qquad
% %     a=\frac{2s\beta}{s\beta+2\beta-1}.
% % \]
% % Substitution into any of the three active powers gives
% % \[
% %     \frac{s(2\beta-1)}{s\beta+2\beta-1}.
% % \]
% % The remaining self-noise power is no smaller than this value provided
% % \[
% %     a+b
% %     \le
% %     2-\frac1\beta
% %     -2\left(1-\frac1\beta\right)
% %       \frac{s\beta}{s\beta+2\beta-1}.
% % \]
% % Using the two equalities above, the full family of attaining parameters can be written directly as
% % \[
% % \begin{aligned}
% %     a&=\frac{2s\beta}{s\beta+2\beta-1},\\
% %     c&=b-\frac{s\beta}{s\beta+2\beta-1},\\
% %     \frac{s\beta}{s\beta+2\beta-1}
% %     \le b
% %     &\le
% %     2-\frac1\beta
% %     -\left(4-\frac{2}{\beta}\right)
% %       \frac{s\beta}{s\beta+2\beta-1},
% % \end{aligned}
% % \]
% % together with $a+b<1$.  The lower bound on $b$ is equivalent to $c\ge0$.  At the smallest admissible $b$, one has
% % \[
% %     a+b=\frac{3s\beta}{s\beta+2\beta-1}<1
% % \]
% % if and only if $s<(2\beta-1)/(2\beta)$.  Thus the feasible set is nonempty precisely below the stated transition; at equality, the strict condition $a+b<1$ prevents attainment.  The stability condition is satisfied because
% % \[
% %     \beta a+c-b
% %     =\frac{s\beta(2\beta-1)}{s\beta+2\beta-1}>0.
% % \]
% % Therefore
% % \[
% %     \EE[\cE(\btheta_K)]
% %     \eqsim
% %     D^{-\frac{s(2\beta-1)}{s\beta+2\beta-1}}
% %     +\exp\!\left(-D^{1-a-b}\right).
% % \]

% % \paragraph{Regime $s\ge(2\beta-1)/(2\beta)$.}
% % For larger source exponent, the batch exponent is limited by the strict data-budget condition $a+b<1$.  The optimal polynomial power is approached as
% % \[
% %     c\to0,
% %     \qquad
% %     a+b\to1.
% % \]
% % Under these limiting relations, $b\to1-a$ and therefore $a+c-b\to2a-1$.  The deterministic-bias and label-noise powers converge respectively to
% % \[
% %     2s(1-a)
% %     \qquad\text{and}\qquad
% %     a-\frac{1-a}{\beta}.
% % \]
% % Balancing these two quantities gives
% % \[
% %     a\to
% %     \frac{2s\beta+1}{2s\beta+\beta+1},
% %     \qquad
% %     a+c-b\to
% %     \frac{2s\beta-\beta+1}{2s\beta+\beta+1}.
% % \]
% % The common limiting decay power is
% % \[
% %     \frac{2s\beta}{2s\beta+\beta+1}.
% % \]
% % More explicitly, for any sufficiently small $\varepsilon>0$, take
% % \[
% %     c=0,
% %     \qquad
% %     a=\frac{2s\beta+1}{2s\beta+\beta+1}-\varepsilon,
% %     \qquad
% %     b=a-\frac{2s\beta-\beta+1}{2s\beta+\beta+1}.
% % \]
% % Then $a+b=1-2\varepsilon<1$, the stability constraint is slack, and the polynomial decay power converges to $2s\beta/(2s\beta+\beta+1)$ as $\varepsilon\downarrow0$.  Hence this value is a supremum rather than an attained maximum.

% % Combining the two regimes, the optimal noisy Nesterov decay power is
% % \[
% %     \begin{cases}
% %     \dfrac{s(2\beta-1)}{s\beta+2\beta-1},
% %     &0<s\le\dfrac{2\beta-1}{2\beta},\\[1.2em]
% %     \dfrac{2s\beta}{2s\beta+\beta+1},
% %     &s\ge\dfrac{2\beta-1}{2\beta},
% %     \end{cases}
% % \]
% % where the value at the transition and the second branch are understood in the supremal sense under $a+b<1$.  The two expressions coincide at $s=(2\beta-1)/(2\beta)$.

% % Finally, the corresponding largest batch exponent is
% % \[
% %     \begin{cases}
% %     \dfrac{2s\beta}{s\beta+2\beta-1},
% %     &0<s\le\dfrac{2\beta-1}{2\beta},\\[1.2em]
% %     \dfrac{2s\beta+1}{2s\beta+\beta+1},
% %     &s\ge\dfrac{2\beta-1}{2\beta},
% %     \end{cases}
% % \]
% % with the same supremal interpretation at and above the transition.  Equivalently, the maximal batch scaling is
% % \[
% %     B\eqsim
% %     \begin{cases}
% %     D^{\frac{2s\beta}{s\beta+2\beta-1}},
% %     &0<s<\dfrac{2\beta-1}{2\beta},\\[1.2em]
% %     D^{\frac{2s\beta+1}{2s\beta+\beta+1}-o(1)},
% %     &s\ge\dfrac{2\beta-1}{2\beta}.
% %     \end{cases}
% % \]
% % \end{proof}


% \subsection{Batch size phase diagram for data-optimal scaling}
% \label{app:fix_bsz}

% \begin{theorem}[Data-optimal scaling for fixed batch size]
% Given a data budget \(D\), let \(B\eqsim D^a\), where \(a\in[0,1)\) is the batch size exponent. We optimally tune the remaining hyperparameters subject to the corresponding stability conditions. Define optimal data scaling rate \(r_A(a)\) by
% \[
% \inf_{\rho,\eta}
% \EE[\cE(\btheta_K)]
% =
% D^{-r_A(a)+o(1)},
% \qquad
% A\in\{\sgd,\polyak,\nesterov\}.
% \]
% Then the optimal decay exponents are
% \[
% r_{\sgd}(a)
% =
% \begin{cases}
% \dfrac{s}{s+1},
% & 0\leq a\leq \dfrac{s}{s+1},\\[1mm]
% s(1-a),
% & \dfrac{s}{s+1}<a<1,
% \end{cases}
% \]
% \[
% r_{\polyak}(a)
% =
% \begin{cases}
% \dfrac{s}{s+1},
% & 0\leq a\leq \dfrac{2s+1}{2s+2},\\[1mm]
% 2s(1-a),
% & \dfrac{2s+1}{2s+2}<a<1,
% \end{cases}
% \]
% and
% \[
% r_{\nesterov}(a)
% =
% \begin{cases}
% \dfrac{s}{s+1},
% & 0\leq a\leq \dfrac{s}{s+1},\\[1mm]
% \dfrac{s\bigl(1+(\beta-1)a\bigr)}{s\beta+1},
% & \dfrac{s}{s+1}<a\leq
% \dfrac{2s\beta+1}{2s\beta+\beta+1},\\[2mm]
% 2s(1-a),
% & \dfrac{2s\beta+1}{2s\beta+\beta+1}<a<1.
% \end{cases}
% \]
% \end{theorem}

% \begin{proof}
% We optimize the polynomial decay power separately for the three algorithms. At boundaries involving the strict condition \(a+b<1\), the stated value is understood in the supremal sense and is represented by the \(o(1)\) term.

% \paragraph{SGD.}
% Let \(\eta\eqsim D^{-c}\), with \(c\geq0\). The SGD scaling law gives
% \[
% \EE[\cE(\btheta_K)]
% \eqsim
% D^{-s(1-a-c)}+D^{-(a+c)}.
% \]
% Hence the decay power is
% \[
% \min\{s(1-a-c),a+c\}.
% \]
% The two powers are balanced when
% \[
% a+c=\frac{s}{s+1}.
% \]
% For \(a\leq s/(s+1)\), the choice \(c=s/(s+1)-a\) is admissible and yields \(r_{\sgd}(a)=s/(s+1)\). For \(a>s/(s+1)\), the balance would require \(c<0\). The optimal admissible choice is therefore \(c=0\), for which the signal term is dominant and
% \[
% r_{\sgd}(a)=s(1-a).
% \]

% \paragraph{Polyak momentum.}
% Let
% \[
% 1-\rho\eqsim D^{-b},
% \qquad
% \eta\eqsim D^{-c},
% \qquad
% b,c\geq0.
% \]
% The stability condition \(\eta\lesssim B(1-\rho)\) becomes
% \(a+c-b\geq0,\)
% and the signal and additive-noise powers are
% \(s(1-a-c+b)\) \text{and} \(a+c-b.\)
% The exponentially damped terms are negligible when \(a+b<1\). Balancing signal and additive noise requires
% \[
% a+c-b=\frac{s}{s+1}.
% \]

% If \(a\leq s/(s+1)\), this balance is attained by taking
% \(b=0, c=\frac{s}{s+1}-a.\)

% If \(s/(s+1)<a<(2s+1)/(2s+2)\), it is attained by taking
% \(c=0, b=a-\frac{s}{s+1},\)
% for which \(a+b<1\). The endpoint \(a=(2s+1)/(2s+2)\) is approached with an arbitrarily small slack in \(a+b<1\). Therefore,
% \[
% r_{\polyak}(a)=\frac{s}{s+1}
% \qquad\text{for}\qquad
% 0\leq a\leq\frac{2s+1}{2s+2}.
% \]

% Now suppose \(a>(2s+1)/(2s+2)\). Since \(c\geq0\) and \(a+b<1\),
% \[
% a+c-b>2a-1>\frac{s}{s+1}.
% \]
% Thus the signal power is the smaller of the two powers. Taking \(c=0\) and letting \(b\uparrow1-a\), the signal power approaches \(2s(1-a)\). Hence
% \[
% r_{\polyak}(a)=2s(1-a).
% \]

% \paragraph{Nesterov momentum.}
% Use the same parametrization. The signal power remains
% \[
% s(1-a-c+b),
% \]
% while the additive-noise power is
% \[
% \begin{cases}
% a+c-b, & c\geq b,\\[1mm]
% a+\dfrac{c-b}{\beta}, & c\leq b.
% \end{cases}
% \]
% The stability condition \(\eta\lesssim B^\beta(1-\rho)\) becomes
% \(\beta a+c-b\geq0.\)

% On the branch \(c\geq b\), balancing the signal and additive-noise powers gives
% \[
% a+c-b=\frac{s}{s+1}.
% \]
% This is admissible for \(a\leq s/(s+1)\), and therefore
% \[
% r_{\nesterov}(a)=\frac{s}{s+1}
% \qquad\text{for}\qquad
% 0\leq a\leq\frac{s}{s+1}.
% \]

% For \(a>s/(s+1)\), the improved rate is obtained on the branch \(c\leq b\). Balancing signal and additive noise gives
% \[
% s(1-a-c+b)
% =
% a+\frac{c-b}{\beta},
% \]
% or equivalently
% \[
% c-b
% =
% \frac{\beta\bigl(s-(s+1)a\bigr)}{s\beta+1}.
% \]
% Substitution into either power yields
% \[
% \frac{s\bigl(1+(\beta-1)a\bigr)}{s\beta+1}.
% \]
% This balance can be realized with \(c=0\) and
% \(b=\frac{\beta\bigl((s+1)a-s\bigr)}{s\beta+1},\)
% provided \(a+b<1\). This condition is equivalent to
% \[
% a<\frac{2s\beta+1}{2s\beta+\beta+1}.
% \]
% At equality, the same rate is approached with arbitrarily small slack. Hence
% \[
% r_{\nesterov}(a)
% =
% \frac{s\bigl(1+(\beta-1)a\bigr)}{s\beta+1}
% \]
% for
% \[
% \frac{s}{s+1}<a\leq
% \frac{2s\beta+1}{2s\beta+\beta+1}.
% \]

% Finally, suppose
% \[
% a>
% \frac{2s\beta+1}{2s\beta+\beta+1}.
% \]
% The balanced choice above violates \(a+b<1\). Taking \(c=0\) and letting \(b\uparrow1-a\), the signal and additive-noise powers approach
% \[
% 2s(1-a)
% \qquad\text{and}\qquad
% \frac{(\beta+1)a-1}{\beta},
% \]
% respectively. These powers coincide at
% \(a=(2s\beta+1)/(2s\beta+\beta+1)\); above this threshold, the signal power is smaller. Therefore,
% \[
% r_{\nesterov}(a)=2s(1-a).
% \]
% For all choices above, the terms containing \(\rho^K\) or \(\gamma_{\nesterov}^K\) decay faster than any polynomial because \(a+b<1\). This completes the proof.
% \end{proof}

% \newpage
% \section{Discussion for the noiseless case \texorpdfstring{\(\sigma=0\)}{sigma = 0}}
% \label{app:noiseless}

% \subsection{Scaling laws in the noiseless setting}

% In this section, we present the expected risk bounds for SGD, Polyak momentum, and Nesterov momentum in the strictly noiseless setting (\(\sigma = 0\)). In the absence of additive label noise, the optimization dynamics are entirely determined by the interplay between the deterministic signal decay and the multiplicative noise induced by mini-batch sampling. The exact scaling laws characterizing this behavior for SGD, Polyak momentum, and Nesterov momentum are formally stated in Theorem~\ref{thm:sgd scaling law noiseless}, Theorem~\ref{thm:hbm scaling law noiseless}, and Theorem~\ref{thm:nesterov scaling law noiseless}, respectively.

% \begin{theorem}[SGD scaling law, noiseless case]
% \label{thm:sgd scaling law noiseless}
% Assume that the SGD dynamics are risk stable, \(\eta\lesssim1\), and assume that \(\sigma=0\). When \(k \gtrsim 1/\eta\), the expected excess risk satisfies:
% \begin{equation}
% \label{eq:sgd_loss_noiseless}
% \EE[\cE(\btheta_k)]
% \eqsim
% \begin{cases}
% \underbrace{ \left(\eta k\right)^{-s} }_{\text{signal}}, & 0 < s \le 2-\frac{1}{\beta}, \\[12pt]
% \underbrace{ \left(\eta k\right)^{-s} }_{\text{signal}} + \underbrace{ \frac{\eta}{B}\left(\eta k\right)^{-\left(2-\frac{1}{\beta}\right)} }_{\text{multiplicative noise}}, & s > 2-\frac{1}{\beta}.
% \end{cases}
% \end{equation}
% \end{theorem}
% See Appendix~\ref{app:sgd_scaling_law} for the proof of Theorem~\ref{thm:sgd scaling law noiseless}.

% \begin{theorem}[Polyak momentum scaling law, noiseless case]
% \label{thm:hbm scaling law noiseless}
% Assume that the Polyak dynamics are risk stable, $\eta\lesssim\min\{1,B(1-\rho)\}$, and assume that \(\sigma=0\). Then, for $k\gtrsim\max\{\frac{1-\rho}{\eta},\frac{1}{1-\rho}\}$, we have:
% \begin{itemize}
% \item (Fully overdamped regime) When \(\eta\lesssim(1-\rho)^2\),
% \begin{equation}
% \label{eq:hb-loss-overdamped-noiseless}
% \EE[\cE(\btheta_k)]
% \eqsim
% \underbrace{
% \left(\frac{1-\rho}{\eta k}\right)^s
% }_{\text{signal}}
% +
% \underbrace{
% \frac{\eta}{B(1-\rho)}\left(\frac{1-\rho}{\eta k}\right)^{s\wedge (2-\frac 1\beta)}
% }_{\text{multiplicative noise}}.
% \end{equation}
% \vspace{-1em}
% \item (Mixed regime) When \(\eta\gtrsim(1-\rho)^2\),
% \begin{equation}
% \label{eq:hb-loss-mixed-noiseless}
% \EE[\cE(\btheta_k)]
% \eqsim
% \underbrace{\rho^k+\left(\frac{1-\rho}{\eta k}\right)^s}_{\text{signal}}
% +
% \underbrace{\frac{\eta k}{B}\rho^k+\frac{\eta^2}{B(1-\rho)^3}\left(\frac{1-\rho}{\eta k}\right)^{2-\frac{1}{\beta}}}_{\text{multiplicative noise}}.
% \vspace{-0.5em}
% \end{equation}
% \end{itemize}
% \end{theorem}
% See Appendix~\ref{app:hbm scaling law} for the proof of Theorem~\ref{thm:hbm scaling law noiseless}.

% \begin{theorem}[Nesterov momentum scaling law, noiseless case]
% \label{thm:nesterov scaling law noiseless}
% Assume that the Nesterov dynamics are risk stable, $\eta\lesssim\min\{1,B^\beta(1-\rho)\}$, and assume that \(\sigma=0\). Then, for $k\gtrsim\max\{\frac{1-\rho}{\eta},\frac{1}{1-\rho}\}$, we have:
% \begin{itemize}
% \item (Fully overdamped regime) When \(\eta\lesssim(1-\rho)^2\),
% \begin{equation}
% \label{eq:nag-loss-overdamped-noiseless}
% \EE[\cE(\btheta_k)]
% \eqsim
% \underbrace{
% \left(\frac{1-\rho}{\eta k}\right)^s
% }_{\text{signal}}
% +
% \underbrace{
% \frac{\eta}{B(1-\rho)}\left(\frac{1-\rho}{\eta k}\right)^{s\wedge(2-\frac{1}{\beta})}
% }_{\text{multiplicative noise}}.
% \end{equation}
% \vspace{-1em}
% \item (Mixed regime) When \(\eta\gtrsim(1-\rho)^2\),
% \begin{equation}
% \label{eq:nag-loss-mixed-noiseless}
% \begin{aligned}
% \EE[\cE(\btheta_k)]
% \eqsim
% &\underbrace{\left(\frac{(1-\rho)^2}{\eta}\right)^{s+\frac{1}{\beta}}\gamma_{\nesterov}^k+
% \left(\frac{1-\rho}{\eta k}\right)^s
% }_{\text{signal}}
% +\\
% &\underbrace{\frac{k(1-\rho)^2}{B} \left(\frac{(1-\rho)^2}{\eta}\right)^{s+\frac{1}{\beta}}\gamma_{\nesterov}^k + \frac{\eta^2}{B(1-\rho)^3}\left(\frac{1-\rho}{\eta k}\right)^{2-\frac 1\beta}}_{\text{multiplicative noise}}
% .
% \vspace{-0.5em}
% \end{aligned}
% \end{equation}
% where $\gamma_{\nesterov}:=\rho\left(1-\frac{(1-\rho)^2}{(1+\rho)^2}\right)$ denotes the damping factor.
% \end{itemize}
% \end{theorem}
% See Appendix~\ref{app:nesterov scaling law} for the proof of Theorem~\ref{thm:nesterov scaling law noiseless}.


% \subsection{Data-optimal scaling in the noiseless setting}
% \label{app:sigma0_data}

% In this section we determine, for each algorithm, the fastest polynomial decay of the final excess risk that is achievable under the fixed data budget \(D=BK\) when \(\sigma=0\), together with the hyperparameter scalings and the batch sizes that realize it. The results are stated in Theorem~\ref{thm:sgd data scaling noiseless}, Theorem~\ref{thm:polyak data scaling noiseless} and Theorem~\ref{thm:nesterov data scaling noiseless}. Similar to Appendix~\ref{app:data_optimal_scaling}, we use the parametrization~\eqref{eq:data-scaling-parameterization},
% \[
%     B\eqsim D^{a},
%     \qquad
%     K\eqsim D^{1-a},
%     \qquad
%     1-\rho\eqsim D^{-b},
%     \qquad
%     \eta\eqsim D^{-c},
%     \qquad
%     a,b,c\geq0,
% \]
% where \(b=0\) for SGD, and we abbreviate
% \[
%     \mu:=s\wedge\left(2-\frac1\beta\right),
%     \qquad
%     x:=a+c-b,
%     \qquad
%     \theta:=\min\left\{1,\left(2-\frac1\beta-s\right)_{+}\right\}.
% \]

% Because the additive label noise is absent, the only competition is between the deterministic signal and the multiplicative noise. The following lemma reduces this competition to two affine decay powers and treats the fully overdamped and the mixed regime simultaneously.

% \begin{lemma}[Decay powers in the noiseless setting]
% \label{lem:noiseless_decay_powers}
% Let \(a,b,c\geq0\) satisfy \(x\leq1\) and \(a+b<1\), and assume the stability condition of the corresponding algorithm. Then
% \[
%     \frac{1-\rho}{\eta K}\eqsim D^{-(1-x)},
%     \qquad
%     \frac{\eta}{B(1-\rho)}\eqsim D^{-x},
%     \qquad
%     \frac{\eta}{(1-\rho)^{2}}\eqsim D^{2b-c},
%     \qquad
%     (1-\rho)K\eqsim D^{1-a-b},
% \]
% and the final risk given by Theorem~\ref{thm:sgd scaling law noiseless}, Theorem~\ref{thm:hbm scaling law noiseless} and Theorem~\ref{thm:nesterov scaling law noiseless} satisfies
% \[
%     \EE[\cE(\btheta_K)]
%     \eqsim
%     D^{-P_{\rm sig}}+D^{-P_{\rm mult}},
% \]
% where
% \[
%     P_{\rm sig}=s(1-x),
%     \qquad
%     P_{\rm mult}
%     =
%     \left(2-\frac1\beta\right)
%     -\left(1-\frac1\beta\right)x
%     -\left(2b-c\right)_{+}.
% \]
% Moreover, eliminating \(b=a+c-x\),
% \[
%     \left(2b-c\right)_{+}=\left(2a+c-2x\right)_{+},
%     \qquad
%     2a+c=\log_{D}\frac{B^{2}}{\eta}.
% \]
% \end{lemma}

% \begin{proof}
% The four elementary identities follow by direct substitution:
% \[
%     \frac{1-\rho}{\eta K}
%     \eqsim
%     \frac{D^{-b}}{D^{-c}\,D^{1-a}}
%     =
%     D^{-(1-a-c+b)}
%     =
%     D^{-(1-x)},
%     \qquad
%     \frac{\eta}{B(1-\rho)}
%     \eqsim
%     \frac{D^{-c}}{D^{a}\,D^{-b}}
%     =
%     D^{-(a+c-b)}
%     =
%     D^{-x},
% \]
% \[
%     \frac{\eta}{(1-\rho)^{2}}
%     \eqsim
%     \frac{D^{-c}}{D^{-2b}}
%     =
%     D^{2b-c},
%     \qquad
%     (1-\rho)K
%     \eqsim
%     D^{-b}\,D^{1-a}
%     =
%     D^{1-a-b}.
% \]
% Substituting \(b=a+c-x\) gives \(2b-c=2(a+c-x)-c=2a+c-2x\).

% \textbf{Signal.} By the first identity,
% \[
%     \left(\frac{1-\rho}{\eta K}\right)^{s}
%     \eqsim
%     D^{-s(1-x)}
%     =
%     D^{-P_{\rm sig}}.
% \]

% \textbf{Transients.} For Polyak momentum,
% \[
%     \rho^{K}
%     \eqsim
%     \exp\left(-\Theta\left((1-\rho)K\right)\right)
%     =
%     \exp\left(-\Theta\left(D^{1-a-b}\right)\right),
%     \qquad
%     1+\frac{\eta K}{B}\eqsim1+D^{1-2a-c}.
% \]
% For Nesterov momentum,
% \[
%     1-\gamma_{\nesterov}
%     =
%     (1-\rho)+\rho\frac{(1-\rho)^{2}}{(1+\rho)^{2}}
%     \eqsim
%     1-\rho
%     \quad\Longrightarrow\quad
%     \gamma_{\nesterov}^{K}
%     \eqsim
%     \exp\left(-\Theta\left(D^{1-a-b}\right)\right),
% \]
% while its prefactor satisfies
% \[
%     \left(1+\frac{K(1-\rho)^{2}}{B}\right)
%     \left(\frac{(1-\rho)^{2}}{\eta}\right)^{s+\frac1\beta}
%     \eqsim
%     \left(1+D^{1-2a-2b}\right)D^{-\left(s+\frac1\beta\right)(2b-c)}.
% \]
% In both cases the transient is a fixed power of \(D\) multiplied by \(\exp\left(-\Theta\left(D^{1-a-b}\right)\right)\), hence smaller than any polynomial power of \(D\) because \(a+b<1\).

% \textbf{Multiplicative noise, fully overdamped regime \(c\geq2b\).} Here
% \[
%     x=a+c-b\geq a+2b-b=a+b\geq0
%     \quad\Longrightarrow\quad
%     \frac{\eta}{B(1-\rho)}\eqsim D^{-x}\lesssim1.
% \]
% Writing \(t:=\frac{1-\rho}{\eta K}\lesssim1\) and using \(t^{p\wedge q}\eqsim t^{p}+t^{q}\), the multiplicative term splits as
% \[
%     \frac{\eta}{B(1-\rho)}\,t^{\,s\wedge\left(2-\frac1\beta\right)}
%     \eqsim
%     \frac{\eta}{B(1-\rho)}\,t^{\,s}
%     +
%     \frac{\eta}{B(1-\rho)}\,t^{\,2-\frac1\beta}
%     \lesssim
%     t^{\,s}
%     +
%     \frac{\eta}{B(1-\rho)}\,t^{\,2-\frac1\beta},
% \]
% so the first summand is absorbed by the signal, while the second one equals
% \[
%     D^{-x}\,D^{-\left(2-\frac1\beta\right)(1-x)}
%     =
%     D^{-\left[\left(2-\frac1\beta\right)-\left(1-\frac1\beta\right)x\right]}
%     =
%     D^{-P_{\rm mult}},
% \]
% since \(c\geq2b\) gives \(\left(2b-c\right)_{+}=0\).

% \textbf{Multiplicative noise, mixed regime \(c\leq2b\).} Here the multiplicative term is
% \[
%     \frac{\eta^{2}}{B(1-\rho)^{3}}\left(\frac{1-\rho}{\eta K}\right)^{2-\frac1\beta},
%     \qquad
%     \frac{\eta^{2}}{B(1-\rho)^{3}}
%     =
%     \frac{\eta}{B(1-\rho)}\cdot\frac{\eta}{(1-\rho)^{2}}
%     \eqsim
%     D^{-x}\,D^{2b-c},
% \]
% so its decay power is
% \[
%     x-(2b-c)+\left(2-\frac1\beta\right)(1-x)
%     =
%     \left(2-\frac1\beta\right)-\left(1-\frac1\beta\right)x-(2b-c)
%     =
%     P_{\rm mult},
% \]
% because \(\left(2b-c\right)_{+}=2b-c\) in this regime. Collecting the signal, the transient and the multiplicative contributions completes the proof.
% \end{proof}

% \begin{theorem}[Data-optimal scaling for SGD, noiseless case]
% \label{thm:sgd data scaling noiseless}
% Given a data budget \(D=BK\), consider the final loss of SGD in Theorem~\ref{thm:sgd scaling law noiseless} under the stability condition \(\eta\lesssim1\), and assume \(\sigma=0\). Then the data-optimal hyperparameters saturate the stability boundary and obey the \textbf{linear scaling rule}
% \[
%     \eta\eqsim B\eqsim1,
%     \qquad
%     K\eqsim D,
% \]
% the \textbf{optimal data scaling rate} is
% \[
%     \cE_{\sgd}^{\opt}(D)
%     \eqsim
%     D^{-\left(s\wedge\left(2-\frac1\beta\right)\right)},
% \]
% and the \textbf{largest batch size} that preserves this optimal rate is
% \[
%    B_{\sgd}^{\opt}(D)\eqsim1.
% \]
% More precisely, for \(B\eqsim D^{a}\) with \(0\leq a<1\) the optimal decay exponent equals
% \[
%     r_{\sgd}^{\,0}(a)
%     =
%     \min\left\{
%         s(1-a),\;
%         \left(2-\frac1\beta\right)-\left(1-\frac1\beta\right)a
%     \right\},
% \]
% which is strictly decreasing in \(a\).
% \end{theorem}

% \begin{proof}[Proof of Theorem~\ref{thm:sgd data scaling noiseless}]
% For SGD we have \(\rho=0\), hence \(b=0\) and \(x=a+c\). The stability condition \(\eta\lesssim1\) reads \(c\geq0\), and $\left(2b-c\right)_{+}=\left(-c\right)_{+}=0.$
% By Lemma~\ref{lem:noiseless_decay_powers}, the decay exponent is \(\min\left\{P_{\rm sig},P_{\rm mult}\right\}\) with
% \[
%     P_{\rm sig}(x)=s(1-x),
%     \qquad
%     P_{\rm mult}(x)=\left(2-\frac1\beta\right)-\left(1-\frac1\beta\right)x,
%     \qquad
%     x=a+c\geq0.
% \]

% \textbf{Step 1: both powers decrease.} Since \(s>0\) and \(\beta>1\),
% \[
%     P_{\rm sig}(x)-P_{\rm sig}(0)=-sx,
%     \qquad
%     P_{\rm mult}(x)-P_{\rm mult}(0)=-\left(1-\frac1\beta\right)x,
% \]
% so both are strictly decreasing on \(x\geq0\), and hence so is their pointwise minimum.

% \textbf{Step 2: the optimum.} Consequently
% \[
%     \min\left\{P_{\rm sig}(x),P_{\rm mult}(x)\right\}
%     \leq
%     \min\left\{P_{\rm sig}(0),P_{\rm mult}(0)\right\}
%     =
%     \min\left\{s,\,2-\frac1\beta\right\}
%     =
%     \mu,
% \]
% with equality if and only if \(x=0\). Since \(x=a+c\) with \(a,c\geq0\), this forces
% \[
%     a=0,
%     \qquad
%     c=0,
%     \qquad\text{i.e.}\qquad
%     B\eqsim1,
%     \quad
%     \eta\eqsim1,
%     \quad
%     K\eqsim D,
% \]
% and therefore \(\cE_{\sgd}^{\opt}(D)\eqsim D^{-\mu}\).

% \textbf{Step 3: fixed batch size.} For fixed \(a\), Step 1 shows that \(\min\left\{P_{\rm sig},P_{\rm mult}\right\}\) is maximized over \(c\geq0\) at \(c=0\), i.e. \(x=a\), which gives
% \[
%     r_{\sgd}^{\,0}(a)
%     =
%     \min\left\{
%         s(1-a),\;
%         \left(2-\frac1\beta\right)-\left(1-\frac1\beta\right)a
%     \right\}
%     <
%     \mu
%     \qquad\text{for every }a>0 .
% \]
% Hence \(B_{\sgd}^{\opt}(D)\eqsim1\).
% \end{proof}

% \begin{theorem}[Data-optimal scaling for Polyak momentum, noiseless case]
% \label{thm:polyak data scaling noiseless}
% Given a data budget \(D=BK\), consider the final loss of Polyak momentum in Theorem~\ref{thm:hbm scaling law noiseless} under the stability condition \(\eta\lesssim\min\{1,B(1-\rho)\}\), and assume \(\sigma=0\). Then the data-optimal effective learning rate saturates the stability boundary and obeys the \textbf{linear scaling rule}
% \(\frac{\eta}{1-\rho}\eqsim B .\)
% Every choice satisfying this rule together with
% \[
%     \frac{B^{2}}{\eta}\lesssim D^{\theta},
%     \qquad
%     \theta=\min\left\{1,\left(2-\frac1\beta-s\right)_{+}\right\},
% \]
% achieves the same \textbf{optimal data scaling rate}
% \[
%     \cE_{\polyak}^{\opt}(D)
%     \eqsim
%     D^{-\left(s\wedge\left(2-\frac1\beta\right)\right)}
%     \eqsim
%     \cE_{\sgd}^{\opt}(D).
% \]
% Moreover, the \textbf{largest batch size} that preserves this optimal rate is
% \[
%    B_{\polyak}^{\opt}(D)
%     =
%     D^{\frac{\theta}{2}-o(1)}
%     =
%     \begin{cases}
%     D^{\frac12-o(1)},
%     & 0<s\leq1-\dfrac1\beta,\\[2mm]
%     D^{\frac12\left(2-\frac1\beta-s\right)},
%     & 1-\dfrac1\beta<s\leq2-\dfrac1\beta,\\[2mm]
%     1,
%     & s>2-\dfrac1\beta,
%     \end{cases}
% \]
% where the \(o(1)\) is needed only on the first branch, and where the maximal batch size is realized by \(\eta\eqsim1\) and \(1-\rho\eqsim B^{-1}\).
% \end{theorem}

% \begin{proof}[Proof of Theorem~\ref{thm:polyak data scaling noiseless}]
% The two stability constraints read
% \[
%     \eta\lesssim B(1-\rho)
%     \iff
%     D^{-c}\lesssim D^{a-b}
%     \iff
%     x=a+c-b\geq0,
%     \qquad
%     \eta\lesssim1
%     \iff
%     c\geq0 .
% \]
% We also impose \(a+b<1\): if \(a+b\geq1\), then \((1-\rho)K\eqsim D^{1-a-b}\lesssim1\), hence \(\rho^{K}\eqsim1\) and the risk does not decay. By Lemma~\ref{lem:noiseless_decay_powers}, the decay exponent equals \(\min\left\{P_{\rm sig},P_{\rm mult}\right\}\) with
% \[
%     P_{\rm sig}=s(1-x),
%     \qquad
%     P_{\rm mult}
%     =
%     \left(2-\frac1\beta\right)-\left(1-\frac1\beta\right)x-\left(2a+c-2x\right)_{+}.
% \]

% \textbf{Step 1: upper bound.} Since \(x\geq0\), \(\beta>1\) and \(\left(\cdot\right)_{+}\geq0\),
% \[
%     P_{\rm sig}\leq s,
%     \qquad
%     P_{\rm mult}\leq\left(2-\frac1\beta\right)-\left(1-\frac1\beta\right)x\leq2-\frac1\beta,
% \]
% and therefore
% \[
%     \min\left\{P_{\rm sig},P_{\rm mult}\right\}
%     \leq
%     s\wedge\left(2-\frac1\beta\right)
%     =
%     \mu .
% \]

% \textbf{Step 2: optimality forces \(x=0\).} Assume \(\min\left\{P_{\rm sig},P_{\rm mult}\right\}=\mu\).

% If \(s\leq2-\frac1\beta\), then \(\mu=s\) and
% \[
%     s(1-x)=P_{\rm sig}\geq\mu=s
%     \quad\Longrightarrow\quad
%     x\leq0
%     \quad\Longrightarrow\quad
%     x=0 .
% \]

% If \(s>2-\frac1\beta\), then \(\mu=2-\frac1\beta\) and
% \[
%     P_{\rm mult}\geq2-\frac1\beta
%     \iff
%     \left(1-\frac1\beta\right)x+\left(2a+c-2x\right)_{+}\leq0 .
% \]
% Both summands on the left are nonnegative, so \(x=0\) and \(2a+c\leq0\), i.e. \(a=c=b=0\).

% \textbf{Step 3: the constraint on \(\left(a,c\right)\) at \(x=0\).} Let \(x=0\), that is \(b=a+c\), which is exactly
% \[
%     \eta\eqsim B(1-\rho)
%     \qquad\Longleftrightarrow\qquad
%     \frac{\eta}{1-\rho}\eqsim B .
% \]
% Then \(2a+c-2x=2a+c\geq0\), so
% \[
%     P_{\rm sig}=s,
%     \qquad
%     P_{\rm mult}=\left(2-\frac1\beta\right)-\left(2a+c\right),
% \]
% and consequently
% \[
%     \min\left\{P_{\rm sig},P_{\rm mult}\right\}
%     =
%     \min\left\{s,\;\left(2-\frac1\beta\right)-\left(2a+c\right)\right\}
%     =
%     \mu
%     \iff
%     2a+c\leq2-\frac1\beta-s .
% \]
% Combining this with the transient condition \(a+b=2a+c<1\) yields the single constraint
% \[
%     2a+c\leq\theta=\min\left\{1,\left(2-\frac1\beta-s\right)_{+}\right\},
%     \qquad\text{strictly if }\theta=1,
% \]
% which, by Lemma~\ref{lem:noiseless_decay_powers}, is precisely \(\frac{B^{2}}{\eta}\lesssim D^{\theta}\).

% \textbf{Step 4: the largest batch size.} Maximizing \(a\) subject to \(2a+c\leq\theta\) and \(c\geq0\) gives
% \[
%     c=0,
%     \qquad
%     a\leq\frac{\theta}{2}.
% \]
% Conversely, take \(c=0\) and \(b=a=\frac{\theta}{2}\), replaced by \(a=b=\frac{\theta}{2}-\varepsilon\) with \(\varepsilon>0\) arbitrarily small when \(\theta=1\). Then
% \[
%     x=a+c-b=0,
%     \qquad
%     c=0\leq2b,
%     \qquad
%     a+b=2a+c=\theta\leq1,
% \]
% so that the choice is stable, lies in the mixed regime \(\eta\gtrsim(1-\rho)^{2}\), has a negligible transient, and satisfies
% \[
%     \min\left\{P_{\rm sig},P_{\rm mult}\right\}
%     =
%     \min\left\{s,\;\left(2-\frac1\beta\right)-\theta\right\}
%     =
%     \mu ,
% \]
% where the last equality uses \(\theta\leq\left(2-\frac1\beta-s\right)_{+}\). This choice corresponds to
% \[
%     \eta\eqsim1,
%     \qquad
%     1-\rho\eqsim D^{-\frac\theta2}\eqsim B^{-1},
%     \qquad
%     B\eqsim D^{\frac\theta2}.
% \]

% \textbf{Step 5: explicit form of \(\theta\).} Finally,
% \[
%     2-\frac1\beta-s\geq1
%     \iff
%     s\leq1-\frac1\beta,
%     \qquad
%     2-\frac1\beta-s\leq0
%     \iff
%     s\geq2-\frac1\beta,
% \]
% so that
% \[
%     \theta
%     =
%     \begin{cases}
%     1, & 0<s\leq1-\frac1\beta,\\[1mm]
%     2-\frac1\beta-s, & 1-\frac1\beta<s\leq2-\frac1\beta,\\[1mm]
%     0, & s>2-\frac1\beta,
%     \end{cases}
% \]
% which gives the stated three branches of \(B_{\polyak}^{\opt}(D)=D^{\theta/2-o(1)}\) and completes the proof.
% \end{proof}

% \begin{theorem}[Data-optimal scaling for Nesterov momentum, noiseless case]
% \label{thm:nesterov data scaling noiseless}
% Given a data budget \(D=BK\), consider the final loss of Nesterov momentum in Theorem~\ref{thm:nesterov scaling law noiseless} under the stability condition \(\eta\lesssim\min\{1,B^{\beta}(1-\rho)\}\), and assume \(\sigma=0\). Write
% \[
%     \Delta_{\beta}:=s\beta(\beta-1)+\beta^{2}+2\beta-1 .
% \]
% Then the data-optimal hyperparameters saturate the Nesterov stability boundary and obey the \textbf{\(\beta\)-th power scaling rule}
% \[
%     \eta\eqsim1,
%     \qquad
%     \frac{\eta}{1-\rho}\eqsim B^{\beta} .
% \]
% The corresponding \textbf{optimal batch size} is
% \[
%     B_{\nesterov}^{\opt}(D)
%     =
%     \begin{cases}
%     D^{\frac{1}{\beta+1}-o(1)},
%     & 0<s\leq\dfrac{\beta-1}{2\beta},\\[3mm]
%     D^{\frac{2\beta-1-s\beta}{\Delta_{\beta}}},
%     & \dfrac{\beta-1}{2\beta}<s\leq2-\dfrac1\beta,\\[3mm]
%     1,
%     & s>2-\dfrac1\beta,
%     \end{cases}
% \]
% and the \textbf{optimal data scaling rate} is
% \[
%     \cE_{\nesterov}^{\opt}(D)
%     =
%     \begin{cases}
%     D^{-\frac{2s\beta}{\beta+1}+o(1)},
%     & 0<s\leq\dfrac{\beta-1}{2\beta},\\[3mm]
%     D^{-\frac{s\beta(3\beta-1)}{\Delta_{\beta}}},
%     & \dfrac{\beta-1}{2\beta}<s\leq2-\dfrac1\beta,\\[3mm]
%     D^{-\left(2-\frac1\beta\right)},
%     & s>2-\dfrac1\beta .
%     \end{cases}
% \]
% Both expressions are continuous across the two thresholds, and
% \[
%     \cE_{\nesterov}^{\opt}(D)
%     \ll
%     \cE_{\polyak}^{\opt}(D)
%     \eqsim
%     \cE_{\sgd}^{\opt}(D)
%     \qquad\text{whenever } s<2-\frac1\beta .
% \]
% \end{theorem}

% \begin{proof}[Proof of Theorem~\ref{thm:nesterov data scaling noiseless}]
% The stability condition reads
% \[
%     \eta\lesssim B^{\beta}(1-\rho)
%     \iff
%     D^{-c}\lesssim D^{\beta a-b}
%     \iff
%     \beta a+c-b\geq0
%     \iff
%     x+(\beta-1)a\geq0 ,
% \]
% so that, in contrast with Polyak momentum, \(x\) may now be negative. We again impose \(a+b<1\), \(c\geq0\), and use Lemma~\ref{lem:noiseless_decay_powers} in the form
% \[
%     P_{\rm sig}=s(1-x),
%     \qquad
%     P_{\rm mult}
%     =
%     \left(2-\frac1\beta\right)-\left(1-\frac1\beta\right)x-\left(2a+c-2x\right)_{+}.
% \]

% \textbf{Step 1: reduction for fixed \(x\).} The power \(P_{\rm sig}\) depends only on \(x\), while \(P_{\rm mult}\) is nonincreasing in \(2a+c\). Hence, for each fixed \(x\), it suffices to minimize \(2a+c\) subject to
% \[
%     a\geq0,
%     \qquad
%     c\geq0,
%     \qquad
%     b=a+c-x\geq0,
%     \qquad
%     (\beta-1)a\geq-x .
% \]

% \textbf{Step 2: the branch \(x\geq0\).} Choosing \(a=0\) and \(c=x\) gives \(b=0\) and \(2a+c=x\leq2x\), hence \(\left(2a+c-2x\right)_{+}=0\) and
% \[
%     P_{\rm mult}=\left(2-\frac1\beta\right)-\left(1-\frac1\beta\right)x .
% \]
% Both \(P_{\rm sig}\) and \(P_{\rm mult}\) are then strictly decreasing in \(x\geq0\), so exactly as in Steps 1--2 of the proof of Theorem~\ref{thm:sgd data scaling noiseless},
% \[
%     \sup_{x\geq0}\min\left\{P_{\rm sig},P_{\rm mult}\right\}
%     =
%     \min\left\{s,\,2-\frac1\beta\right\}
%     =
%     \mu ,
% \]
% attained only at \(x=0\) with \(2a+c=0\), i.e. \(a=b=c=0\).

% \textbf{Step 3: the branch \(x<0\).} Put \(z:=-x>0\). Then \(b=a+c+z\geq0\) holds automatically, the stability condition becomes $ a\geq\frac{z}{\beta-1},$
% and \(2a+c-2x=2a+c+2z>0\), so the positive part is active and
% \[
%     P_{\rm mult}
%     =
%     \left(2-\frac1\beta\right)+\left(1-\frac1\beta\right)z-\left(2a+c+2z\right)
%     =
%     \left(2-\frac1\beta\right)-\left(1+\frac1\beta\right)z-\left(2a+c\right).
% \]
% Minimizing \(2a+c\) therefore gives
% \[
%     c=0,
%     \qquad
%     a=\frac{z}{\beta-1},
%     \qquad
%     b=a+z=\frac{\beta z}{\beta-1}=\beta a ,
% \]
% and hence, using
% \[
%     \frac{\beta+1}{\beta}+\frac{2}{\beta-1}
%     =
%     \frac{(\beta+1)(\beta-1)+2\beta}{\beta(\beta-1)}
%     =
%     \frac{\beta^{2}+2\beta-1}{\beta(\beta-1)}
%     =:C_{\beta}>0,
% \]
% we obtain
% \[
%     P_{\rm mult}
%     =
%     \left(2-\frac1\beta\right)-\left(1+\frac1\beta\right)z-\frac{2z}{\beta-1}
%     =
%     \left(2-\frac1\beta\right)-C_{\beta}z ,
%     \qquad
%     P_{\rm sig}=s(1+z).
% \]
% The transient condition becomes
% \[
%     a+b=\frac{z}{\beta-1}+\frac{\beta z}{\beta-1}=\frac{(\beta+1)z}{\beta-1}<1
%     \iff
%     0<z<z_{\max}:=\frac{\beta-1}{\beta+1}.
% \]
% Since \(P_{\rm sig}\) is strictly increasing and \(P_{\rm mult}\) is strictly decreasing in \(z\), the maximum of their minimum on \(\left[0,z_{\max}\right)\) is governed by the crossing point
% \[
%     s(1+z^{*})=\left(2-\frac1\beta\right)-C_{\beta}z^{*}
%     \qquad\Longleftrightarrow\qquad
%     z^{*}=\frac{2-\frac1\beta-s}{s+C_{\beta}} .
% \]

% \textbf{Step 4: the regime \(s>2-\frac1\beta\).} Here \(z^{*}<0\), so \(P_{\rm sig}(z)>P_{\rm mult}(z)\) for all \(z>0\), and \(\min\left\{P_{\rm sig},P_{\rm mult}\right\}=P_{\rm mult}\) is strictly decreasing. Its maximum over \(\left[0,z_{\max}\right)\) is attained at \(z=0\), where it equals \(2-\frac1\beta=\mu\), forcing \(a=b=c=0\). Together with Step 2,
% \[
%     \cE_{\nesterov}^{\opt}(D)\eqsim D^{-\left(2-\frac1\beta\right)},
%     \qquad
%     B_{\nesterov}^{\opt}(D)\eqsim1 .
% \]

% \textbf{Step 5: location of the crossing point.} Let \(s<2-\frac1\beta\), so that \(z^{*}>0\). Then
% \[
%     z^{*}\geq z_{\max}
%     \iff
%     (\beta+1)\left(2-\frac1\beta-s\right)\geq(\beta-1)\left(s+C_{\beta}\right).
% \]
% Multiplying both sides by \(\beta\) and using \(\beta(\beta-1)C_{\beta}=\beta^{2}+2\beta-1\),
% \[
%     (\beta+1)\left(2\beta-1-s\beta\right)
%     \geq
%     s\beta(\beta-1)+\beta^{2}+2\beta-1 ,
% \]
% that is,
% \[
%     2\beta^{2}+\beta-1-s\beta(\beta+1)
%     \geq
%     s\beta(\beta-1)+\beta^{2}+2\beta-1 ,
% \]
% which simplifies to
% \[
%     \beta^{2}-\beta\geq s\beta\left[(\beta+1)+(\beta-1)\right]=2s\beta^{2}
%     \qquad\Longleftrightarrow\qquad
%     s\leq\frac{\beta-1}{2\beta}.
% \]

% \textbf{Step 6: the regime \(\frac{\beta-1}{2\beta}<s<2-\frac1\beta\).} Here \(0<z^{*}<z_{\max}\), so the maximum is attained exactly at \(z=z^{*}\), where
% \[
%     1+z^{*}
%     =
%     \frac{s+C_{\beta}+2-\frac1\beta-s}{s+C_{\beta}}
%     =
%     \frac{C_{\beta}+2-\frac1\beta}{s+C_{\beta}} .
% \]
% We compute the two ingredients:
% \[
%     C_{\beta}+2-\frac1\beta
%     =
%     \frac{\beta^{2}+2\beta-1}{\beta(\beta-1)}+\frac{2\beta-1}{\beta}
%     =
%     \frac{\beta^{2}+2\beta-1+(2\beta-1)(\beta-1)}{\beta(\beta-1)}
%     =
%     \frac{3\beta^{2}-\beta}{\beta(\beta-1)}
%     =
%     \frac{3\beta-1}{\beta-1},
% \]
% \[
%     s+C_{\beta}
%     =
%     \frac{s\beta(\beta-1)+\beta^{2}+2\beta-1}{\beta(\beta-1)}
%     =
%     \frac{\Delta_{\beta}}{\beta(\beta-1)} .
% \]
% Therefore
% \[
%     1+z^{*}=\frac{3\beta-1}{\beta-1}\cdot\frac{\beta(\beta-1)}{\Delta_{\beta}}=\frac{\beta(3\beta-1)}{\Delta_{\beta}},
% \]
% \[
%     z^{*}
%     =
%     \frac{\beta(3\beta-1)-\Delta_{\beta}}{\Delta_{\beta}}
%     =
%     \frac{2\beta^{2}-3\beta+1-s\beta(\beta-1)}{\Delta_{\beta}}
%     =
%     \frac{(\beta-1)\left(2\beta-1-s\beta\right)}{\Delta_{\beta}} ,
% \]
% and the optimal decay power and batch exponent are
% \[
%     s(1+z^{*})=\frac{s\beta(3\beta-1)}{\Delta_{\beta}},
%     \qquad
%     a=\frac{z^{*}}{\beta-1}=\frac{2\beta-1-s\beta}{\Delta_{\beta}} .
% \]
% Since \(P_{\rm sig}\) and \(P_{\rm mult}\) are simultaneously active at \(z=z^{*}\), the minimizing choice \(c=0\), \(a=\frac{z^{*}}{\beta-1}\), \(b=\beta a\) is the unique optimum; it satisfies \(a+b=\frac{(\beta+1)z^{*}}{\beta-1}<1\), so the rate is attained. Because \(z^{*}>0\), this rate is strictly larger than \(\mu=s\).

% \textbf{Step 7: the regime \(0<s\leq\frac{\beta-1}{2\beta}\).} Here \(z^{*}\geq z_{\max}\), hence \(P_{\rm sig}\leq P_{\rm mult}\) on \(\left[0,z_{\max}\right)\) and \(\min\left\{P_{\rm sig},P_{\rm mult}\right\}=s(1+z)\) is strictly increasing, so the value
% \[
%     \lim_{z\uparrow z_{\max}}s(1+z)
%     =
%     s\left(1+\frac{\beta-1}{\beta+1}\right)
%     =
%     \frac{2s\beta}{\beta+1}
% \]
% is a supremum. Explicitly, for every sufficiently small \(\varepsilon>0\) take
% \[
%     c=0,
%     \qquad
%     z=z_{\max}-\varepsilon,
%     \qquad
%     a=\frac{z}{\beta-1}=\frac{1}{\beta+1}-\frac{\varepsilon}{\beta-1},
%     \qquad
%     b=\beta a .
% \]
% Then
% \[
%     a+b=\frac{(\beta+1)z}{\beta-1}=1-\frac{(\beta+1)\varepsilon}{\beta-1}<1,
%     \qquad
%     P_{\rm sig}=\frac{2s\beta}{\beta+1}-s\varepsilon
%     \leq
%     P_{\rm mult},
% \]
% so that
% \[
%     \EE[\cE(\btheta_K)]
%     \eqsim
%     D^{-\frac{2s\beta}{\beta+1}+s\varepsilon}.
% \]
% Letting \(\varepsilon\downarrow0\) gives
% \[
%     \cE_{\nesterov}^{\opt}(D)=D^{-\frac{2s\beta}{\beta+1}+o(1)},
%     \qquad
%     B_{\nesterov}^{\opt}(D)=D^{\frac{1}{\beta+1}-o(1)} .
% \]
% Because \(\frac{2\beta}{\beta+1}>1\), this rate is again strictly larger than \(\mu=s\).

% \textbf{Step 8: hyperparameter rule, continuity, and conclusion.} In Steps 4, 6 and 7 the optimal choice satisfies \(c=0\) and \(\beta a+c-b=0\), that is
% \[
%     \eta\eqsim1,
%     \qquad
%     \frac{\eta}{1-\rho}\eqsim B^{\beta}.
% \]
% At \(s=\frac{\beta-1}{2\beta}\) we have \(s\beta=\frac{\beta-1}{2}\), hence
% \[
%     \Delta_{\beta}
%     =
%     \frac{(\beta-1)^{2}}{2}+\beta^{2}+2\beta-1
%     =
%     \frac{3\beta^{2}+2\beta-1}{2}
%     =
%     \frac{(3\beta-1)(\beta+1)}{2},
% \]
% \[
%     \frac{s\beta(3\beta-1)}{\Delta_{\beta}}
%     =
%     \frac{\beta-1}{\beta+1}
%     =
%     \frac{2s\beta}{\beta+1},
%     \qquad
%     \frac{2\beta-1-s\beta}{\Delta_{\beta}}
%     =
%     \frac{\frac{3\beta-1}{2}}{\frac{(3\beta-1)(\beta+1)}{2}}
%     =
%     \frac{1}{\beta+1}.
% \]
% At \(s=2-\frac1\beta\) we have \(s\beta=2\beta-1\), hence
% \[
%     \Delta_{\beta}=(2\beta-1)(\beta-1)+\beta^{2}+2\beta-1=3\beta^{2}-\beta=\beta(3\beta-1),
% \]
% \[
%     \frac{s\beta(3\beta-1)}{\Delta_{\beta}}=\frac{2\beta-1}{\beta}=2-\frac1\beta,
%     \qquad
%     \frac{2\beta-1-s\beta}{\Delta_{\beta}}=0 .
% \]
% This proves the continuity of both expressions. Finally, comparing Steps 4, 6 and 7 with Step 2 shows that the branch \(x<0\) is strictly better than the branch \(x\geq0\) whenever \(s<2-\frac1\beta\), and equal to it otherwise, which completes the proof.
% \end{proof}


% \subsection{Data-optimal scaling for fixed batch size in the noiseless setting}
% \label{app:fix_bsz_noiseless}

% Given a data budget \(D\), let \(B\eqsim D^a\), where \(a\in[0,1)\), and assume \(\sigma=0\). We optimally tune the remaining hyperparameters subject to the corresponding stability conditions and define
% \[
% \inf_{\rho,\eta}
% \EE[\cE(\btheta_K)]
% =
% D^{-r_A(a)+o(1)},
% \qquad
% A\in\{\sgd,\polyak,\nesterov\}.
% \]
% We state the fixed-batch rates separately in the three source regimes appearing in Theorem~\ref{thm:nesterov data scaling noiseless}.

% \begin{theorem}[Fixed-batch scaling in the hard noiseless regime]
% \label{thm:fixed-batch-data-scaling-noiseless-hard}
% Suppose
% \(0<s\leq\frac{\beta-1}{2\beta}.\)
% Then
% \[
% r_{\sgd}(a)=s(1-a),
% \]
% \[
% r_{\polyak}(a)
% =
% \begin{cases}
% s,
% &0\leq a\leq\dfrac12,\\[1mm]
% 2s(1-a),
% &\dfrac12<a<1,
% \end{cases}
% \]
% and
% \[
% r_{\nesterov}(a)
% =
% \begin{cases}
% s\bigl(1+(\beta-1)a\bigr),
% &0\leq a\leq\dfrac1{\beta+1},\\[2mm]
% 2s(1-a),
% &\dfrac1{\beta+1}<a<1.
% \end{cases}
% \]
% At the boundaries involving the strict transient condition, the stated value is understood in the supremal sense represented by the \(o(1)\) term.
% \end{theorem}

% \begin{proof}
% By Lemma~\ref{lem:noiseless_decay_powers}, the signal and multiplicative-noise powers are
% \[
% s(1-a-c+b)
% \]
% and
% \[
% \left(2-\frac1\beta\right)
% -\left(1-\frac1\beta\right)(a+c-b)
% -(2b-c)_{+},
% \]
% respectively. For the momentum methods, decay of the transient further requires \(a+b<1\).

% \medskip

% \noindent\textbf{SGD.}
% Here \(b=0\). Both powers decrease with \(c\), so the optimal choice is \(c=0\). Since \(s<2-1/\beta\), the signal power is smaller for every \(a\in[0,1)\), and therefore
% \[
% r_{\sgd}(a)=s(1-a).
% \]

% \medskip

% \noindent\textbf{Polyak momentum.}
% The stability condition is \(a+c-b\geq0\). It suffices to take \(c=0\): if \(c\leq b\), replacing \((b,c)\) by \((b-c,0)\) preserves \(a+c-b\), weakens the transient constraint, and does not decrease either power; if \(c\geq b\), both powers are no larger than their values at \(b=c=0\). Hence
% \[
% 0\leq b\leq a,
% \qquad
% b<1-a.
% \]
% With \(c=0\), the signal power increases with \(b\), whereas the multiplicative-noise power decreases with \(b\). In the present regime the signal remains the smaller power throughout the admissible interval. \\
% If \(a\leq1/2\), the largest admissible choice is \(b=a\), giving
% \[
% r_{\polyak}(a)=s.
% \]
% If \(a>1/2\), take \(b\uparrow1-a\), which gives
% \[
% r_{\polyak}(a)=2s(1-a).
% \]

% \medskip

% \noindent\textbf{Nesterov momentum.}
% The stability condition is \(\beta a+c-b\geq0\). The same argument reduces the optimization to \(c=0\), so
% \[
% 0\leq b\leq\beta a,
% \qquad
% b<1-a.
% \]
% The signal remains the smaller power throughout the admissible interval because \(s\leq(\beta-1)/(2\beta)\). If \(a\leq1/(\beta+1)\), the stability boundary is active and \(b=\beta a\), which gives
% \[
% r_{\nesterov}(a)=s\bigl(1+(\beta-1)a\bigr).
% \]
% If \(a>1/(\beta+1)\), the transient boundary is active and taking \(b\uparrow1-a\) gives
% \[
% r_{\nesterov}(a)=2s(1-a).
% \]
% This proves the theorem.
% \end{proof}

% \begin{theorem}[Fixed-batch scaling in the intermediate noiseless regime]
% \label{thm:fixed-batch-data-scaling-noiseless-intermediate}
% Suppose
% \(\frac{\beta-1}{2\beta}<s<2-\frac1\beta.\)
% Then
% \[
% r_{\sgd}(a)=s(1-a).
% \]
% For Polyak momentum, when
% \(\frac{\beta-1}{2\beta}<s\leq1-\frac1\beta,\)
% we have
% \[
% r_{\polyak}(a)
% =
% \begin{cases}
% s,
% &0\leq a\leq\dfrac12,\\[1mm]
% 2s(1-a),
% &\dfrac12<a<1,
% \end{cases}
% \]
% whereas, when
% \(1-\frac1\beta<s<2-\frac1\beta,\)
% we have
% \[
% r_{\polyak}(a)
% =
% \begin{cases}
% s,
% &0\leq a\leq\dfrac{2-\frac1\beta-s}{2},\\[3mm]
% \dfrac{s\beta(3-2a)}{s\beta+\beta+1},
% &\dfrac{2-\frac1\beta-s}{2}<a
% \leq\dfrac{2s\beta+2-\beta}{2(s\beta+1)},\\[3mm]
% 2s(1-a),
% &\dfrac{2s\beta+2-\beta}{2(s\beta+1)}<a<1.
% \end{cases}
% \]
% For Nesterov momentum,
% \[
% r_{\nesterov}(a)
% =
% \begin{cases}
% s\bigl(1+(\beta-1)a\bigr),
% &0\leq a\leq
% \dfrac{2\beta-1-s\beta}
% {s\beta(\beta-1)+\beta^2+2\beta-1},\\[4mm]
% \dfrac{s\beta(3-2a)}{s\beta+\beta+1},
% &\dfrac{2\beta-1-s\beta}
% {s\beta(\beta-1)+\beta^2+2\beta-1}<a
% \leq\dfrac{2s\beta+2-\beta}{2(s\beta+1)},\\[4mm]
% 2s(1-a),
% &\dfrac{2s\beta+2-\beta}{2(s\beta+1)}<a<1.
% \end{cases}
% \]
% At the boundaries involving the strict transient condition, the stated value is understood in the supremal sense represented by the \(o(1)\) term.
% \end{theorem}

% \begin{proof}
% The signal and multiplicative-noise powers are again
% \[
% s(1-a-c+b)
% \]
% and
% \[
% \left(2-\frac1\beta\right)
% -\left(1-\frac1\beta\right)(a+c-b)
% -(2b-c)_{+}.
% \]

% \medskip

% \noindent\textbf{SGD.}
% The optimal choice is \(c=0\). Since \(s<2-1/\beta\), the signal power is smaller for all \(a\in[0,1)\), giving
% \[
% r_{\sgd}(a)=s(1-a).
% \]

% \medskip

% \noindent\textbf{Polyak momentum.}
% As in the proof of Theorem~\ref{thm:fixed-batch-data-scaling-noiseless-hard}, the optimum has \(c=0\) and
% $0\leq b\leq a,
% b<1-a.$
% The signal and multiplicative-noise powers are then
% \[
% s(1-a+b), \qquad\text{and} \quad
% \left(2-\frac1\beta\right)
% -\left(1-\frac1\beta\right)a
% -\left(1+\frac1\beta\right)b.
% \]
% The first is increasing and the second is decreasing in \(b\).

% If \(s\leq1-1/\beta\), the signal remains the smaller power throughout the admissible interval. Thus \(b=a\) is optimal for \(a\leq1/2\), while \(b\uparrow1-a\) is optimal for \(a>1/2\), giving the first Polyak formula.

% Suppose instead that \(s>1-1/\beta\). At the stability boundary \(b=a\), the signal remains the smaller power precisely when
% \(a\leq\frac{2-\frac1\beta-s}{2}.\)
% Beyond this point the two powers balance. Solving
% \[
% s(1-a+b)
% =
% \left(2-\frac1\beta\right)
% -\left(1-\frac1\beta\right)a
% -\left(1+\frac1\beta\right)b
% \]
% gives the common value
% \(\frac{s\beta(3-2a)}{s\beta+\beta+1}.\)
% The balancing choice reaches the transient boundary \(b=1-a\) when
% \[
% a=\frac{2s\beta+2-\beta}{2(s\beta+1)}.
% \]
% For larger \(a\), the signal power at \(b\uparrow1-a\) is dominant and equals \(2s(1-a)\).

% \medskip

% \noindent\textbf{Nesterov momentum.}
% The optimum has \(c=0\) and
% $0\leq b\leq\beta a, 
% b<1-a.$
% At the stability boundary \(b=\beta a\), the signal remains the smaller power precisely when
% \[
% a\leq
% \frac{2\beta-1-s\beta}
% {s\beta(\beta-1)+\beta^2+2\beta-1}.
% \]
% This gives the first Nesterov branch. Beyond this point the same two powers balance, so the middle branch is
% \(\frac{s\beta(3-2a)}{s\beta+\beta+1}.\)
% The balance reaches \(b=1-a\) at
% \[
% a=\frac{2s\beta+2-\beta}{2(s\beta+1)},
% \]
% after which the signal power at \(b\uparrow1-a\) gives \(2s(1-a)\). The adjacent expressions agree at both boundaries, which completes the proof.
% \end{proof}

% \begin{figure*}[t]
%       \centering

%       \begin{subfigure}[t]{0.49\textwidth}
%           \centering
%           \includegraphics[
%               width=\linewidth
%           ]{figures/noiseless_fixed_batch_hard.pdf}
%           \caption{Hard regime.}
%           \label{fig:noiseless-fixed-batch-hard}
%       \end{subfigure}
%       \hfill
%       \begin{subfigure}[t]{0.49\textwidth}
%           \centering
%           \includegraphics[
%               width=\linewidth
%           ]{figures/noiseless_fixed_batch_intermediate.pdf}
%           \caption{Intermediate regime.}
%           \label{fig:noiseless-fixed-batch-intermediate}
%       \end{subfigure}

%       \medskip

%       \begin{subfigure}[t]{0.49\textwidth}
%           \centering
%           \includegraphics[
%               width=\linewidth
%           ]{figures/noiseless_fixed_batch_easy.pdf}
%           \caption{Easy regime.}
%           \label{fig:noiseless-fixed-batch-easy}
%       \end{subfigure}

%       \caption{
%           Optimal fixed-batch data scaling rates in the noiseless setting
%           for $\beta=8$. Panels (a), (b), and (c) correspond to
%           Theorems~\ref{thm:fixed-batch-data-scaling-noiseless-hard},
%           \ref{thm:fixed-batch-data-scaling-noiseless-intermediate}, and
%           \ref{thm:fixed-batch-data-scaling-noiseless-easy}, and use
%           $s=0.3$, $s=1$, and $s=2.5$, respectively. The horizontal axis
%           is the batch-size exponent $a$ in $B\eqsim D^a$, and the vertical
%           axis is the optimal loss-decay exponent. Yellow, red, and blue
%           denote the rates attainable by SGD, Polyak momentum, and Nesterov
%           momentum, respectively.
%       }
%       \label{fig:noiseless-fixed-batch-phase-diagrams}
%   \end{figure*}
  
% \begin{theorem}[Fixed-batch scaling in the easy noiseless regime]
% \label{thm:fixed-batch-data-scaling-noiseless-easy}
% Suppose
% \(s\geq2-\frac1\beta.\)
% Then
% \[
% r_{\sgd}(a)
% =
% \begin{cases}
% \left(2-\dfrac1\beta\right)-\left(1-\dfrac1\beta\right)a,
% &0\leq a\leq\dfrac{s-2+\frac1\beta}{s-1+\frac1\beta},\\[3mm]
% s(1-a),
% &\dfrac{s-2+\frac1\beta}{s-1+\frac1\beta}<a<1,
% \end{cases}
% \]
% and Polyak and Nesterov momentum have the same optimal decay exponent,
% \[
% r_{\polyak}(a)
% =
% r_{\nesterov}(a)
% =
% \begin{cases}
% \left(2-\dfrac1\beta\right)-\left(1-\dfrac1\beta\right)a,
% &0\leq a\leq\dfrac{s-2+\frac1\beta}{s-1+\frac1\beta},\\[3mm]
% \dfrac{s\beta(3-2a)}{s\beta+\beta+1},
% &\dfrac{s-2+\frac1\beta}{s-1+\frac1\beta}<a
% \leq\dfrac{2s\beta+2-\beta}{2(s\beta+1)},\\[3mm]
% 2s(1-a),
% &\dfrac{2s\beta+2-\beta}{2(s\beta+1)}<a<1.
% \end{cases}
% \]
% At the boundaries involving the strict transient condition, the stated value is understood in the supremal sense represented by the \(o(1)\) term.
% \end{theorem}

% \begin{proof}
% \medskip

% \noindent\textbf{SGD.}
% The optimal choice is \(c=0\), so the signal and multiplicative-noise powers are
% \[
% s(1-a)
% \qquad\text{and}\qquad
% \left(2-\frac1\beta\right)-\left(1-\frac1\beta\right)a.
% \]
% They coincide when
% \(a=\frac{s-2+\frac1\beta}{s-1+\frac1\beta}.\)
% The multiplicative-noise power is smaller below this point and the signal power is smaller above it, giving the stated SGD formula.

% \medskip

% \noindent\textbf{Polyak momentum.}
% The optimum has \(c=0\), \(0\leq b\leq a\), and \(b<1-a\). At \(b=0\), the multiplicative-noise power is smaller precisely when
% \(a\leq\frac{s-2+\frac1\beta}{s-1+\frac1\beta},\)
% which gives the first branch. Beyond this point, increasing \(b\) balances the signal and multiplicative-noise powers, giving
% \(\frac{s\beta(3-2a)}{s\beta+\beta+1}.\)
% The balance reaches the transient boundary \(b=1-a\) when
% \[
% a=\frac{2s\beta+2-\beta}{2(s\beta+1)}.
% \]
% For larger \(a\), the signal power at \(b\uparrow1-a\) is smaller and equals \(2s(1-a)\).

% \medskip

% \noindent\textbf{Nesterov momentum.}
% The optimum has \(c=0\), \(0\leq b\leq\beta a\), and \(b<1-a\). In the easy regime, the optimum either occurs at \(b=0\), at the interior balance, or at the transient boundary. Each of these choices is admissible for both Polyak and Nesterov momentum, and the Nesterov stability boundary \(b=\beta a\) is never active. Consequently Nesterov has exactly the same three branches as Polyak momentum. The adjacent expressions agree at both boundaries, completing the proof.
% \end{proof}

% Figure~\ref{fig:noiseless-fixed-batch-phase-diagrams} summarizes the three fixed-batch phase diagrams in Theorems~\ref{thm:fixed-batch-data-scaling-noiseless-hard}, \ref{thm:fixed-batch-data-scaling-noiseless-intermediate}, and \ref{thm:fixed-batch-data-scaling-noiseless-easy}.


% \newpage
\section{Experimental details and additional results}
\label{app:experiment}

Our theoretical analysis is formulated in the infinite-dimensional PLK setting. For numerical experiments, we truncate the system to the first \(d=4096\) eigendirections, so that, for example, \(\bH=\diag(\lambda_1,\ldots,\lambda_d)\in\RR^{d\times d}\). Unless otherwise stated, all experiments use \(d=4096\).


\subsection{Risk stability boundaries}
\label{app:stability-experiment}

\paragraph{Simulation setup.} We run with \(\btheta_{-1}=\btheta_0=0\) for \(T=2000\) iterations. A run is classified as stable if its average risk over the final \(20\) iterations is below \(1\); it is  classified as unstable if the risk blows up or exceeds \(10^{12}\). For each setting and random seed, the empirical stability boundary is the largest stable learning rate on the search grid. We repeat each scan over five independent seeds and report the mean boundary.

\begin{itemize}
\item \textbf{Stability versus batch size.} For Figure~\ref{fig:numerical-summary} (left), we use \((s,\beta)=(3,2)\), fix \(1-\rho=2^{-24}\) for the momentum methods, and scan \(B\in\{2^0,2^1,\ldots,2^{12}\}\) and \(50\) logarithmically spaced learning rates between \(2^{-10}\) and \(8\). The dashed curves are post-hoc fits to the predicted scalings: a constant for SGD, \(B(1-\rho)\) for Polyak, and \(B^\beta(1-\rho)\) followed by a constant plateau for Nesterov. The dotted line marks the split used for the Nesterov fit.



\item \textbf{Stability versus momentum damping.} For Figure~\ref{fig:numerical-summary} (middle), we use \((s,\beta)=(3,2)\), \(\sigma=0\), and \(B=32\), vary \(1-\rho\in\{2^{-24},2^{-22},\ldots,2^{-2}\}\), and scan \(50\) logarithmically spaced learning rates between \(2^{-22}\) and \(8\). Since SGD is independent of \(\rho\), its boundary is constant across the damping grid. The dashed curves are post-hoc fits to the predicted form \(\eta^\crit=\min\{L,c(1-\rho)\}\) for Polyak and Nesterov, and to a constant for SGD.
\end{itemize}

\paragraph{Additional results.} Figure~\ref{fig:additional-stability-boundaries} shows additional results for different choices of \((s,\beta)\).



\begin{figure}[H]
    \centering
    \captionsetup[subfigure]{skip=0pt, font=small}
    % \begin{subfigure}[t]{0.325\textwidth}
    %     \centering
    %     \includegraphics[width=\linewidth]{figures/stability/beta1_5s3_plot.pdf}
    %     \caption{\(s=0.3,\beta=1.5\)}
    % \end{subfigure}
    % \hfill
    \begin{subfigure}[t]{0.325\textwidth}
        \centering
        \includegraphics[width=\linewidth]{figures/stability/beta1_5s5_plot.pdf}
        \caption{\(s=0.5,\beta=1.5\)}
    \end{subfigure}
    % \hfill
    \begin{subfigure}[t]{0.325\textwidth}
        \centering
        \includegraphics[width=\linewidth]{figures/stability/beta2s2_plot.pdf}
        \caption{\(s=0.2,\beta=2\)}
    \end{subfigure}
    % \medskip
    % \begin{subfigure}[t]{0.325\textwidth}
    %     \centering
    %     \includegraphics[width=\linewidth]{figures/stability/beta2s3_plot_main.pdf}
    %     \caption{\(s=3,\beta=2\)}
    % \end{subfigure}
    \begin{subfigure}[t]{0.32\textwidth}
        \centering
        \includegraphics[width=\linewidth]{figures/stability/beta2_5s5_plot.pdf}
        \caption{\(s=0.5,\beta=2.5\)}
    \end{subfigure}
    % \hfill
    % \begin{subfigure}[t]{0.325\textwidth}
    %     \centering
    %     \includegraphics[width=\linewidth]{figures/stability/beta2s5_plot.pdf}
    %     \caption{\(s=0.5,\beta=2\)}
    % \end{subfigure}
    % \hfill
    % \begin{subfigure}[t]{0.325\textwidth}
    %     \centering
    %     \includegraphics[width=\linewidth]{figures/stability/beta2_5s3_plot.pdf}
    %     \caption{\(s=0.3,\beta=2.5\)}
    % \end{subfigure}
    % \medskip

    % \vspace{-1em}
\caption{Additional stability boundaries across \((s,\beta)\).}
    \label{fig:additional-stability-boundaries}
\end{figure}

\subsection{Risk-dynamics scaling laws}
\label{app:scaling-law-experiments}

\paragraph{Simulation setup.} We simulate PLK regression initialized with \(\btheta_{-1}=\btheta_0=0\). For Figure~\ref{fig:numerical-summary} (right), we set \((s,\beta)=(0.3,4)\), \(\sigma=1\), \(B=32\), \(\eta=0.03\), and \(\rho=0.98\). Each method is run for \(2\times10^7\) iterations. Let \(\overline{\cE}(k)\) denote the population excess risk averaged over the \(100\) runs.

\paragraph{Scaling-law fits.} Theorems~\ref{thm:hbm scaling law} and~\ref{thm:nesterov scaling law} determine the functional forms of the risk dynamics up to multiplicative constants. We therefore introduce \(\Theta=(C_1,C_2,C_3)\in\RR_+^3\) and fit the following ansatz separately for each method. For Polyak,
\[
F_{\polyak}(k;\Theta)=C_1\rho^k+C_2(\eta_\rho k)^{-s}+C_3\frac{\sigma^2}{B}\eta_\rho,\qquad \eta_\rho:=\frac{\eta}{1-\rho},
\]
while for Nesterov,
\[
F_{\nesterov}(k;\Theta)=C_1\min\{1,(\eta k)^{-s}\}\rho^k+C_2(\eta_\rho k)^{-s}+C_3\frac{\sigma^2}{B}\min\{\eta_\rho,\eta_\rho^{1/\beta}\}.
\]
The first terms use the transient envelopes from Theorems~\ref{thm:hbm scaling law} and~\ref{thm:nesterov scaling law}, so the fits test the predicted scaling forms rather than their multiplicative constants. For each method, the coefficients are fitted by minimizing the mean squared relative error:
\[
\min_{\Theta\in\RR_+^3}\frac{1}{|\mathcal I_{\rm fit}|}\sum_{k\in \mathcal I_{\rm fit}}\left(\frac{F(k;\Theta)-\overline{\cE}(k)}{\overline{\cE}(k)}\right)^2,\qquad \mathcal I_{\rm fit}:=\{k\geq1:\eta_\rho k\geq1\},
\]
where \(F=F_{\polyak}\) or \(F_{\nesterov}\) for the corresponding method.

\paragraph{Additional results.} We repeat the same experiment for \((s,\beta)\in\{(0.3,6),(0.5,4),(0.5,6)\}\), keeping \(\sigma=1\), \(B=32\), \(\eta=0.03\), and \(\rho=0.98\). Each run uses \(5\times10^5\) iterations, and the reported risk is averaged over \(100\) independent runs. Figure~\ref{fig:additional-scaling-beta2} shows that the predicted scaling forms consistently track the empirical risk dynamics across these choices of \((s,\beta)\).

\begin{figure}[H]
    \centering
    \captionsetup[subfigure]{skip=0pt, font=small}
    \begin{subfigure}[t]{0.32\textwidth}
        \centering
        \includegraphics[width=\linewidth]{figures/scaling_law_new/scaling_law_s0_3_beta_6.pdf}
        \caption{\((s,\beta)=(0.3,6)\)}
    \end{subfigure}
    \hfill
    \begin{subfigure}[t]{0.32\textwidth}
        \centering
        \includegraphics[width=\linewidth]{figures/scaling_law_new/scaling_law_s0_5_beta_4.pdf}
        \caption{\((s,\beta)=(0.5,4)\)}
    \end{subfigure}
    \hfill
    \begin{subfigure}[t]{0.32\textwidth}
        \centering
        \includegraphics[width=\linewidth]{figures/scaling_law_new/scaling_law_s0_5_beta_6.pdf}
        \caption{\((s,\beta)=(0.5,6)\)}
    \end{subfigure}
    % \vspace{-0.5em}
    \caption{\textbf{Additional validation of the risk-dynamics scaling laws for Polyak and Nesterov.} Thin translucent lines show empirical  excess risks, while thick solid curves show the fitted scaling-law predictions. All experiments use \(\sigma=1\), \(B=32\), \(\eta=0.03\), and \(\rho=0.98\), run for \(5\times10^5\) iterations, and report averages over \(100\) independent runs.}
    \label{fig:additional-scaling-beta2}
\end{figure}



% \begin{figure}[H]
%     \centering
%     \begin{subfigure}[t]{0.32\textwidth}
%         \centering
%         \includegraphics[width=\linewidth]{figures/scaling_law/beta-2_s-0p3_sigma-0p5_dim-4096_runs-100_steps-500000.pdf}
%         \caption{\((s,\beta)=(0.3,2)\)}
%     \end{subfigure}
%     \hfill
%     \begin{subfigure}[t]{0.32\textwidth}
%         \centering
%         \includegraphics[width=\linewidth]{figures/scaling_law/beta-2_s-0p5_sigma-0p5_dim-4096_runs-100_steps-500000.pdf}
%         \caption{\((s,\beta)=(0.5,2)\)}
%     \end{subfigure}
%     \hfill
%         \begin{subfigure}[t]{0.32\textwidth}
%         \centering
%         \includegraphics[width=\linewidth]{figures/scaling_law/beta-4_s-0p3_sigma-0p5_dim-4096_runs-100_steps-500000.pdf}
%         \caption{\((s,\beta)=(0.3,4)\)}
%     \end{subfigure}
%     \medskip
%     \begin{subfigure}[t]{0.32\textwidth}
%         \centering
%         \includegraphics[width=\linewidth]{figures/scaling_law/beta-4_s-0p5_sigma-0p5_dim-4096_runs-100_steps-500000.pdf}
%         \caption{\((s,\beta)=(0.5,4)\)}
%     \end{subfigure}
%     \hfill
%     \begin{subfigure}[t]{0.32\textwidth}
%         \centering
%         \includegraphics[width=\linewidth]{figures/scaling_law/beta-6_s-0p3_sigma-0p5_dim-4096_runs-100_steps-500000.pdf}
%         \caption{\((s,\beta)=(0.3,6)\)}
%     \end{subfigure}
%     \hfill
%     \begin{subfigure}[t]{0.32\textwidth}
%         \centering
%         \includegraphics[width=\linewidth]{figures/scaling_law/beta-6_s-0p5_sigma-0p5_dim-4096_runs-100_steps-500000.pdf}
%         \caption{\((s,\beta)=(0.5,6)\)}
%     \end{subfigure}
%     \caption{Additional scaling-law trajectories for Polyak and Nesterov. All experiments use dimension \(d=4096\), run for \(5\times10^5\) iterations, and report the mean trajectory over 100 independent runs.}
%     \label{fig:additional-scaling-beta2}
% \end{figure}

\subsection{Batch-size phase diagram at a fixed data budget}
\label{app:critical-batch-size}

We simulate PLK regression with \((s,\beta)=(0.3,4)\) and \(\sigma=0.1\), with all methods initialized at \(\btheta_{-1}=\btheta_0=0\). We fix the total data budget at \(D=2^{14}\). 
% .
For each batch size, we search locally around the hyperparameter scalings predicted by Theorem~\ref{thm:batch-size-scaling}. Each candidate is evaluated over 50 independent runs. Let \(W=\max\{1,\lceil0.01 T\rceil\}\). We score each candidate by the mean excess risk over the final \(W\) iterations and across runs:
\[
\overline{\mathcal E}=\frac1{50}\sum_{r=1}^{50}\frac1 W\sum_{k=T-W+1}^{T}\mathcal E(\btheta_k^{(r)}).
\]
For each method and batch size, Figure~\ref{fig:intro-result} (right) reports the minimum \(\overline{\mathcal E}\) over the corresponding search grid.



Figure~\ref{fig:additional-cbs} shows additional  results on the batch-size phase diagram across \((s,\beta)\).



\begin{figure}[H]
    \centering
    \captionsetup[subfigure]{skip=0pt, font=small}
    \begin{subfigure}[t]{0.325\textwidth}
        \centering
        \includegraphics[width=\linewidth]{figures/cbs/s0.1beta2.pdf}
        \caption{\(s=0.1,\beta=2\)}
    \end{subfigure}
    \hfill
    \begin{subfigure}[t]{0.325\textwidth}
        \centering
        \includegraphics[width=\linewidth]{figures/cbs/s0.1beta4.pdf}
        \caption{\(s=0.1,\beta=4\)}
    \end{subfigure}
    % \hfill
    % \begin{subfigure}[t]{0.325\textwidth}
    %     \centering
    %     \includegraphics[width=\linewidth]{figures/cbs/s0.3beta2.pdf}
    %     \caption{\(s=0.3,\beta=2\)}
    % \end{subfigure}
    % \medskip
    \begin{subfigure}[t]{0.325\textwidth}
        \centering
        \includegraphics[width=\linewidth]{figures/cbs/s0.3beta6.pdf}
        \caption{\(s=0.3,\beta=6\)}
    \end{subfigure}
    % \hfill
    % \begin{subfigure}[t]{0.325\textwidth}
    %     \centering
    %     \includegraphics[width=\linewidth]{figures/cbs/s0.3beta8.pdf}
    %     \caption{\(s=0.3,\beta=8\)}
    % \end{subfigure}
    % \hfill
    % \begin{subfigure}[t]{0.32\textwidth}
    %     \centering
    %     \includegraphics[width=\linewidth]{figures/cbs/s0.5beta4.pdf}
    %     \caption{\(s=0.5,\beta=4\)}
    % \end{subfigure}
    % \vspace{-1em}
    \caption{Additional  experiments on the batch-size phase diagram under  different \((s,\beta)\).}
    \label{fig:additional-cbs}
\end{figure}


\subsection{Data-scaling exponents across batch-size regimes}
\label{app:data-scaling-experiments}

In Figure~\ref{fig:three-regime-plk}, we simulate PLK regression with \((s,\beta)=(0.3,4)\), \(\sigma=0.1\), and \(\btheta_{-1}=\btheta_0=0\), using data budgets \(D\in\{2000,4000,8000,16000,32000,64000,128000\}\). For a prescribed batch-size exponent \(b\), we set \(B=\operatorname{round}(C_BD^b)\) and \(T=\lfloor D/B\rfloor\), where \(C_B\) is an order-one constant. For each method and \(D\), we tune the multiplicative constants around the learning-rate and momentum scalings prescribed by Theorem~\ref{thm:batch-size-scaling}, and report the best stable configuration. Each run is evaluated by averaging the risk over the final \(1\%\) of optimization steps, and each point is averaged over 10 independent runs.

We evaluate one representative batch-size exponent from each of the three regimes.

\begin{itemize}
\item \textbf{Small batch.} We take \(b_{\rm small}=0\). For all three methods, we search around \(\eta\eqsim D^{-s/(s+1)}\); for Polyak and Nesterov, we additionally use \(1-\rho\eqsim1\).

\item \textbf{Large batch.} We take \(b_{\rm large}=b_2-\delta\), where \(b_2=\frac{2s\beta+1}{2s\beta+\beta+1}\) and \(\delta>0\) is small. SGD uses \(\eta\eqsim1\). For Polyak, we search around \(\eta\eqsim1\) and \(1-\rho\eqsim D^{-(b_{\rm large}-s/(s+1))}\). For Nesterov, we search around \(\eta\eqsim1\) and
\[
1-\rho\eqsim D^{-c_{\nesterov}},\qquad c_{\nesterov}=b_{\rm large}-\frac{s\beta-(\beta-1)b_{\rm large}}{s\beta+1}.
\]

\item \textbf{Ultra-large batch.} We choose \(b_3<b_{\rm ultra}<1\). SGD uses \(\eta\eqsim1\), while for both momentum methods we search around \(\eta\eqsim1\) and \(1-\rho\eqsim D^{-(1-b_{\rm ultra}-\delta)}\), with \(\delta>0\) chosen so that \(T(1-\rho)\to\infty\).
\end{itemize}

\paragraph{Additional results.} Figures~\ref{fig:app data s0.1beta4} and \ref{fig:app data s0.3beta2} show additional data-scaling results across \((s,\beta)\).


\begin{figure}[H]
    \centering
    \begin{subfigure}[t]{0.325\textwidth}
        \centering
        \includegraphics[width=\linewidth]{figures/data_s0.1_beta4/small_batch.pdf}
    \end{subfigure}
    \hfill
    \begin{subfigure}[t]{0.325\textwidth}
        \centering
        \includegraphics[width=\linewidth]{figures/data_s0.1_beta4/large_batch.pdf}
    \end{subfigure}
    \hfill
    \begin{subfigure}[t]{0.325\textwidth}
        \centering
        \includegraphics[width=\linewidth]{figures/data_s0.1_beta4/super_large_batch.pdf}
    \end{subfigure}
    \vspace{-1em}
    \caption{Data scaling with \((s,\beta)=(0.1,4)\).}
    \label{fig:app data s0.1beta4}
    \vspace{1em}
    \begin{subfigure}[t]{0.325\textwidth}
        \centering
        \includegraphics[width=\linewidth]{figures/data_s0.3_beta2/small_batch.pdf}
    \end{subfigure}
    \hfill
    \begin{subfigure}[t]{0.325\textwidth}
        \centering
        \includegraphics[width=\linewidth]{figures/data_s0.3_beta2/large_batch.pdf}
    \end{subfigure}
    \hfill
    \begin{subfigure}[t]{0.325\textwidth}
        \centering
        \includegraphics[width=\linewidth]{figures/data_s0.3_beta2/super_large_batch.pdf}
    \end{subfigure}
     \vspace{-1em}
    \caption{Data scaling with \((s,\beta)=(0.3,2)\).}
    \label{fig:app data s0.3beta2}
    % \vspace{1em}
    % \begin{subfigure}[t]{0.325\textwidth}
    %     \centering
    %     \includegraphics[width=\linewidth]{figures/data s0.8beta4/small_batch.pdf}
    % \end{subfigure}
    % \hfill
    % \begin{subfigure}[t]{0.325\textwidth}
    %     \centering
    %     \includegraphics[width=\linewidth]{figures/data s0.8beta4/large_batch.pdf}
    % \end{subfigure}
    % % \hfill
    % % \begin{subfigure}[t]{0.325\textwidth}
    % %     \centering
    % %     \includegraphics[width=\linewidth]{figures/data s0.8beta4/super_large_batch.pdf}
    % % \end{subfigure}
    %  \vspace{-1em}
    % \caption{Data scaling with \((s,\beta)=(0.8,4)\).}
    % \label{fig:app data s0.8beta4}
\end{figure}

% We simulate the power-law regression model with \(d=4096\), \(s=0.3\), \(\beta=4\), \(\sigma=0.1\), and \(D\in\{2000,4000,8000,16000,32000,64000,128000\}\). All algorithms are initialized at \(\btheta_{-1}=\btheta_0=\bm0\). At each iteration, we independently sample Gaussian noise according to the curvature-aligned covariance model in Proposition~\ref{prop:noise-geometry}.

% \paragraph{Polyak critical-batch experiment.}
% We compare SGD and Polyak at \(B\eqsim D^{\frac{s}{s+1}}\) and \(B\eqsim D^{\frac{2s}{s+1}}\), with \(K=\lfloor D/B\rfloor\). At the smaller batch size, SGD uses \(\eta=0.1\), while Polyak uses \(\eta=0.1D^{-\frac{s}{s+1}}\) and \(1-\rho=D^{-\frac{s}{s+1}}\). At the larger batch size, SGD uses \(\eta=0.5\), while Polyak uses \(\eta=0.1\) and \(1-\rho=D^{-\frac{s}{s+1}}\). Each point reports the mean final excess risk over \(100\) independent runs.

% \paragraph{Nesterov large-batch experiment.}
% All three algorithms use \(B=\operatorname{round}(0.15D^a)\) and \(K=\lfloor D/B\rfloor\), where
% \[
% a=\frac{2s\beta+1}{2s\beta+\beta+1}-0.05.
% \]
% SGD uses \(\eta=0.3\). Polyak uses \(\eta=0.1\) and \(1-\rho=D^{-(a-\frac{s}{s+1})}\). For Nesterov, defining \(x=\frac{s\beta-(\beta-1)a}{s\beta+1}\), we use \(\eta=0.1\) and \(1-\rho=2D^{-(a-x)}\). Each point reports the mean over \(100\) independent runs of the excess risk averaged over the final \(1\%\) of training.

\subsection{Delayed instability under increasing momentum}
\label{app:subsec:blow-up}

We simulate PLK regression  with \((s=0.3,\beta=4)\). We set the learning rate to \(\eta=0.05\) and the batch size to \(B=4\). To isolate the effect of multiplicative noise, we set the additive noise to \(\sigma=0\). For Polyak and Nesterov, we use the time-varying momentum schedule
\[
\rho_k=1-\frac{1}{k+1}.
\]
Each method is run for at most \(5\times10^5\) iterations over 50 independent runs, and the figure reports the median risk. 




\end{document}
